\documentclass[12pt,a4paper,oneside]{book}
\usepackage[margin=25mm]{geometry}
\usepackage{fontspec}
\setmainfont{Times New Roman}
\setsansfont{Times New Roman}
\setmonofont{DejaVu Sans Mono}
\usepackage{amsmath,amssymb,mathtools}
\usepackage{unicode-math}
\setmathfont{TeX Gyre Termes Math}
\usepackage{newunicodechar}
\newunicodechar{∅}{\ensuremath{\varnothing}}
\newunicodechar{ℰ}{\ensuremath{\mathcal{E}}}
\newunicodechar{ℜ}{\ensuremath{\mathfrak{R}}}
\newunicodechar{ℱ}{\ensuremath{\mathfrak{F}}}
\newunicodechar{ℒ}{\ensuremath{\mathfrak{L}}}
\newunicodechar{𝒜}{\ensuremath{\mathfrak{A}}}
\newunicodechar{↦}{\ensuremath{\mapsto}}
\newunicodechar{∈}{\ensuremath{\in}}
\newunicodechar{∉}{\ensuremath{\notin}}
\newunicodechar{∘}{\ensuremath{\circ}}
\newunicodechar{∧}{\ensuremath{\wedge}}
\newunicodechar{⊂}{\ensuremath{\subset}}
\newunicodechar{⊗}{\ensuremath{\otimes}}
\newunicodechar{⊥}{\ensuremath{\perp}}
\usepackage{graphicx}
\usepackage{longtable,booktabs,array}
\usepackage{xcolor,framed,fancyvrb}
\usepackage{setspace}
\usepackage{enumitem}
\usepackage{microtype}
\usepackage[hidelinks,unicode,pdfusetitle,bookmarksnumbered]{hyperref}
\usepackage[open,openlevel=2]{bookmark}
\doublespacing
\setlength{\parindent}{1.5em}
\setlength{\parskip}{0pt}
\setcounter{tocdepth}{2}
\setcounter{secnumdepth}{2}
\setlist{itemsep=0.25\baselineskip,topsep=0.25\baselineskip}
\providecommand{\tightlist}{\setlength{\itemsep}{0pt}\setlength{\parskip}{0pt}}
\providecommand{\pandocbounded}[1]{#1}
\definecolor{shadecolor}{RGB}{248,248,248}
\newenvironment{Shaded}{\begin{snugshade}}{\end{snugshade}}
\newenvironment{Highlighting}[1][]{}{}
\newcommand{\NormalTok}[1]{#1\par}
\makeatletter
\def\maxwidth{\ifdim\Gin@nat@width>\linewidth\linewidth\else\Gin@nat@width\fi}
\def\maxheight{\ifdim\Gin@nat@height>0.8\textheight 0.8\textheight\else\Gin@nat@height\fi}
\makeatother
\setkeys{Gin}{width=\maxwidth,height=\maxheight,keepaspectratio}
\graphicspath{{images/}}
\emergencystretch=3em
\pagestyle{plain}
\title{Algebra, Chapter 9}
\author{Nicolas Bourbaki}
\date{}
\begin{document}
\pagenumbering{gobble}
\maketitle
\clearpage
\frontmatter
\pagenumbering{roman}
\tableofcontents
\mainmatter
\pagenumbering{arabic}
\setcounter{chapter}{8}
\chapter{Sesquilinear forms and quadratic forms}

\section{Sesquilinear forms}

\subsection{Bilinear maps.}

In this no., A and B denote two rings, E a left A-module, F a right B-module, and G an (A, B)-bimodule, that is, a commutative group equipped with a left A-module structure and a right B-module structure such that one has \((ag)b = a(gb)\) for all \(a \in A\) , \(b \in B\) , \(g \in G\) .

DEFINITION 1. --- One says that a map \(\Phi\) of the product \(\mathbf{E} \times \mathbf{F}\) in \(\mathbf{G}\) is bilinear if it satisfies the following conditions:

\begin{enumerate}
\def\labelenumi{(\arabic{enumi})}
\tightlist
\item
  \(\Phi(x + x', y) = \Phi(x, y) + \Phi(x', y)\)
\end{enumerate}

for all \(x \in \mathbf{E}\) , \(x' \in \mathbf{E}\) , \(y \in \mathbf{F}\) ;

\begin{enumerate}
\def\labelenumi{(\arabic{enumi})}
\setcounter{enumi}{1}
\tightlist
\item
  \(\Phi(x, y + y') = \Phi(x, y) + \Phi(x, y')\)
\end{enumerate}

for all \(x \in \mathbf{E}, y \in \mathbf{F}, y' \in \mathbf{F}\) ;

\begin{enumerate}
\def\labelenumi{(\arabic{enumi})}
\setcounter{enumi}{2}
\item
  \(\Phi(ax, y) = a\Phi(x, y)\) for all \(a \in \mathbf{A}\) , \(x \in \mathbf{E}\) , \(y \in \mathbf{F}\) ;
\item
  \(\Phi(x, yb) = \Phi(x, y)b\) for all \(x \in \mathbf{E}, y \in \mathbf{F}, b \in \mathbf{B}\) .
\end{enumerate}

The tensor product \(E \otimes_{Z} F\) is canonically endowed with a structure of (A, B)-bimodule characterized by \(a(x \otimes y)b = ax \otimes yb\) (Chap. III, 2nd ed., App. II, no. 3), and the data of a bilinear map \(\Phi\) of \(E \times F\) into G is equivalent to that of a map \(\Psi\) of \(E \otimes_{Z} F\) into G that is a homomorphism for the structures of (A, B)-bimodules and that satisfies \(\Psi(x \otimes y) = \Phi(x, y)\) for all \(x \in E\) and \(y \in F\) .

The conditions imposed on \(\Phi\) by Definition 1 mean that the partial maps \(d_{\Phi}(y): x \to \Phi(x, y)\) and \(s_{\Phi}(x): y \to \Phi(x, y)\) are respectively an A-linear map from E to G and a B-linear map from F to G. Let us equip the commutative group \(\mathcal{L}_{\mathbf{A}}(\mathbf{E}, \mathbf{G})\) (resp. \(\mathcal{L}_{\mathbf{B}}(\mathbf{F}, \mathbf{G})\) ) with the structure of a right B-module (resp. left A-module) defined by \(ub(x) = u(x). b\) ( \(u \in \mathcal{L}_{\mathbf{A}}(\mathbf{E}, \mathbf{G}), x \in \mathbf{E}, b \in \mathbf{B}\) ) (respectively \(av(y) = a. v(y)\) ( \(a \in \mathbf{A}, v \in \mathcal{L}_{\mathbf{B}}(\mathbf{F}, \mathbf{G}), y \in \mathbf{F}\) )). Then conditions (1) through (4) are respectively equivalent to:

\[
s _ {\Phi} (x + x ^ {\prime}) = s _ {\Phi} (x) + s _ {\Phi} (x ^ {\prime})\tag{1'}
\]

\[
d _ {\Phi} (y + y ^ {\prime}) = d _ {\Phi} (y) + d _ {\Phi} (y ^ {\prime})\tag{2'}
\]

\[
s _ {\Phi} (a x) = a. s _ {\Phi} (x)\tag{3'}
\]

\[
d _ {\Phi} (y b) = d _ {\Phi} (y). b,\tag{4'}
\]

for all \(x, x'\) in E, \(y, y'\) in F, \(a \in A\) , \(b \in B\) ; in other words, the map \(d_{\Phi}\) of F into \(\mathcal{L}_{A}(E, G)\) is B-linear, and the map \(s_{\Phi}\) from E into \(\mathcal{L}_{B}(F, G)\) is A-linear. We have, by definition,

\begin{enumerate}
\def\labelenumi{(\arabic{enumi})}
\setcounter{enumi}{4}
\tightlist
\item
  \(\Phi(x, y) = d_{\Phi}(y)(x) = s_{\Phi}(x)(y)\) for all \(x \in \mathbf{E}, y \in \mathbf{F}\) .
\end{enumerate}

DEFINITION 2. --- Given a bilinear map \(\Phi\) from E × F to G, the map \(d_{\Phi}\) of F into \(\mathfrak{L}_{\mathrm{A}}(\mathrm{E}, \mathrm{G})\) (resp. the map \(s_{\Phi}\) from E into \(\mathfrak{L}_{\mathrm{B}}(\mathrm{F}, \mathrm{G})\) ) characterized by (5) is called the linear map associated on the right (resp. on the left) to \(\Phi\) .

Conversely, the data of a B-linear map \(d\) of F into \(\mathfrak{L}_{\mathbf{A}}(\mathbf{E},\mathbf{G})\) (resp. of an A-linear map \(s\) from E into \(\mathfrak{L}_{\mathbf{B}}(\mathbf{F},\mathbf{G})\) ) determines uniquely, by the formula

\[
\Phi (x, y) = d (y) (x) \quad (\text { resp. } \Phi (x, y) = s (x) (y))
\]

a bilinear map \(\Phi\) of \(\mathbf{E} \times \mathbf{F}\) into \(\mathbf{G}\) , of which \(d\) (resp. \(s\) ) is the linear map associated on the right (resp. on the left).

DEFINITION 3. --- A bilinear map \(\Phi\) from E × F to G is said to be degenerate on the right (resp. on the left) if there exists a nonzero element \(y_0\) of F (resp. \(x_0\) of E) such that \(\Phi(x, y_0) = 0\) for all \(x \in \mathbf{E}\) , (resp. \(\Phi(x_0, y) = 0\) for all \(y \in \mathbf{F}\) ). One says that \(\Phi\) is degenerate if it is degenerate on the right or if it is degenerate on the left.

For \(\Phi\) to be nondegenerate on the right (resp. on the left) it is necessary and sufficient that the linear map associated on the right (resp. on the left) to \(\Phi\) be injective; to say that \(\Phi\) is nondegenerate therefore means that the associated linear maps \(d_{\Phi}\) and \(s_{\Phi}\) are both injective.

Let \((e_i)_{i\in I}\) and \((f_k)_{k\in K}\) be two families of elements of E and F, and let \((a_i)_{i\in I}\) and \((b_k)_{k\in K}\) be two families of elements of A and B, all but a finite number of which are zero. It follows from equalities (1) to (4), by induction on the number of nonzero coefficients, that one has

\[
\Phi (\sum_ {i} a _ {i} e _ {i}, \sum_ {k} f _ {k} b _ {k}) = \sum_ {i, k} a _ {i} \Phi (e _ {i}, f _ {k}) b _ {k}.\tag{6}
\]

If \((e_{i})\) and \((f_{k})\) are systems of generators of the modules E and F, \(\Phi\) is therefore completely determined by the elements \(g_{ik} = \Phi(e_{i}, f_{k})\) . If \((e_{i})\) and \((f_{k})\) are bases of E and F and elements \(g_{ik}\) of G \((i \in I, k \in K)\) are given, then the formula

\[
\Phi (\sum_ {i} a _ {i} e _ {i}, \sum_ {k} f _ {k} b _ {k}) = \sum_ {i, k} a _ {i} g _ {i k} b _ {k}\tag{6'}
\]

defines a map from E × F to G, which is bilinear and satisfies \(\Phi(e_{i}, f_{k}) = g_{ik}\) . When \((e_{i})\) and \((f_{k})\) are finite bases, we say that \((\Phi(e_{i}, f_{k}))\) is the matrix of \(\Phi\) with respect to these bases.

The bilinear maps from E × F into G obviously form a subgroup of the additive group of maps from E × F into G. On the other hand, let a (resp. b) be an element of the center of A (resp. B); then the map \(a\Phi b\) from E × F to G defined by \((a\Phi b)(x, y) = a \cdot \Phi(x, y) \cdot b\) is bilinear. The set of bilinear maps from E × F to G is thus endowed with a bimodule structure over the centers of A and B.

Let E' (resp. F') be a left A-module (resp. a right B-module), u (resp. v) a homomorphism from E to E' (resp. from F to F'), and \(\Phi'\) a bilinear map from \(E' \times F'\) into G. The inverse image of \(\Phi'\) (relative to u and v) is the bilinear map \(\Phi\) from E × F to G defined by

\[
\Phi (x, y) = \Phi^ {\prime} (u (x), v (y)) \quad (x \in \mathrm{E}, y \in \mathrm{F}).
\]

One verifies easily that one has

\[
d _ {\Phi} (y) = d _ {\Phi^ {\prime}} (\nu (y)) \circ u \quad \text { and } \quad s _ {\Phi} (x) = s _ {\Phi^ {\prime}} (u (x)) \circ \nu
\]

for all \(x \in \mathbf{E}, y \in \mathbf{F}\) .

Let \(\Phi\) be a bilinear map from \(\mathbf{E} \times \mathbf{F}\) into \(\mathbf{G}\) , and \(h\) a homomorphism (for the structures of (A, B)-bimodules) of \(\mathbf{G}\) into another (A, B)-bimodule \(\mathbf{G}'\) . Then \(h \circ \Phi\) is a bilinear map from \(\mathbf{E} \times \mathbf{F}\) into \(\mathbf{G}'\) .

\subsection{Sesquilinear maps.}

In this section, unless expressly stated otherwise, we denote by A and B two rings, by E a left A-module and by F a left B-module; we denote by \(b \to b^{\mathrm{J}} (b \in \mathrm{B})\) an anti-automorphism of B, that is, a bijection from B onto itself satisfying \((b + c)^{\mathrm{J}} = b^{\mathrm{J}} + c^{\mathrm{J}}\) and \((bc)^{\mathrm{J}} = c^{\mathrm{J}}b^{\mathrm{J}}\) for all \(b, c\) in B; we write \(\mathbf{J}'\) instead of \(\mathbf{J}^{-1}\) . We denote by G a \((\mathbf{A}, \mathbf{B})\) -bimodule (no. 1).

DEFINITION 4. --- A map \(\Phi\) from E \(\times\) F into G is said to be right-sesquilinear for J if it satisfies conditions (1), (2), (3) (def. 1, no. 1) as well as

\begin{enumerate}
\def\labelenumi{(\arabic{enumi})}
\setcounter{enumi}{6}
\tightlist
\item
  \(\Phi(x, by) = \Phi(x, y). b^J\) for all \(x \in \mathbf{E}, y \in \mathbf{F}\) and \(b \in \mathbf{B}\) .
\end{enumerate}

If J is the identity (which requires that B be commutative), we recover the notion of a bilinear map.

Let \((e_i)_{i\in I}\) and \((f_k)_{k\in K}\) be two families of elements of E and F, and let \((a_i)_{i\in I}\) and \((b_k)_{k\in K}\) be elements of A and B that are zero except for a finite number of them. We then have

\[
\Phi (\sum_ {i} a _ {i} e _ {i}, \sum_ {k} b _ {k} f _ {k}) = \sum_ {i, k} a _ {i} \Phi (e _ {i}, f _ {k}) b _ {k} ^ {\mathrm{J}}.\tag{8}
\]

As in the case of a bilinear map, the elements \(\Phi(e_i, f_k)\) determine \(\Phi\) in a unique way when \((e_i)\) and \((f_k)\) are systems of generators, and can be chosen arbitrarily when \((e_i)\) and \((f_k)\) are bases of E and F; when \((e_i)\) and \((f_k)\) are finite bases, we say that \((\Phi(e_i, f_k))\) is the matrix of \(\Phi\) with respect to these bases.

As for bilinear maps, one defines on the set of right sesquilinear maps (for J) from E × F into G a bimodule structure over the centers of A and B. One defines the notion of inverse image of a sesquilinear map by the same formula as for a bilinear map. We shall moreover see that the study of sesquilinear maps can be reduced to that of bilinear maps.

DEFINITION 5. --- Let B be a ring, F a left (resp. right) B-module and J an antiautomorphism of B. One denotes by F \(^{j}\) the right (resp. left) B-module having the same underlying additive group as F and in which the external composition law is \((b, y) \to b^{J'}y\) (resp. \((b, y) \to yb^{J'}\) ) \((b \in B, y \in F, J' = J^{-1})\) .

With the notations of Definition 5, a linear map from \(\mathbf{F}^{\mathrm{j}}\) into a right (resp. left) B-module H is therefore identified with a Z-linear map \(u\) of F into H satisfying

\[
u (b y) = u (y) b ^ {J} \quad (\text { resp. } u (y b) = b ^ {J} u (y)) \quad (b \in B, y \in F).
\]

The map \(u\) from F into H is a semilinear map of F into H relative to J (chap. II, App. I, no. 1), if one regards J as an isomorphism of the opposite ring B \(^{0}\) of B onto B, and F as a B \(^{0}\) -module on the right (resp. on the left).

Similarly, a right-sesquilinear map \(\Phi\) (for J) from E × F into G, where F is a left B-module, identifies with a bilinear map from E × F \(^{J}\) into G; if the latter is degenerate on the right (resp. degenerate on the left, nondegenerate), one says that \(\Phi\) is degenerate on the right (resp. degenerate on the left, nondegenerate).

Remark. --- Let A and B be two rings, \(J_{1}\) an anti-automorphism of A, M a right A-module, N a right B-module, and G an (A, B)-bimodule. One says that a map \(\Phi\) from M \(\times\) N into G is left-sesquilinear for J \(_{1}\) if it is Z-bilinear and if it satisfies

\[
\Phi (x a, y b) = a ^ {\mathrm{J} _ {1}} \Phi (x, y) b \quad (x \in \mathrm{M}, y \in \mathrm{N}, a \in \mathrm{A}, b \in \mathrm{B}).\tag{9}
\]

Such a map is identified with a bilinear map from \(M^{J_{1}} \times N\) into G. We will often leave it to the reader to transpose to left sesquilinear maps the definitions and properties given for right sesquilinear maps; when we speak of a sesquilinear map (without specifying), it will be a right sesquilinear map.

\subsection{Orthogonality. Direct sums of bilinear or sesquilinear maps.}

In this section, A and B denote rings, E a left A-module, F a right (resp. left) B-module, G an (A, B)-bimodule, and Φ a bilinear (resp. sesquilinear for a given antiautomorphism J of B) map from E × F to G.

DEFINITION 6. --- Two elements \(x \in E\) and \(y \in F\) are said to be orthogonal with respect to \(\Phi\) if \(\Phi(x, y) = 0\) . Two subsets \(E' \subset E\) and \(F' \subset F\) are said to be orthogonal if, for any \(x \in E'\) and \(y \in F'\) , \(x\) and \(y\) are orthogonal. The set of elements of E (resp. F) orthogonal to a given submodule N of F (resp. M of E) is a submodule of E (resp. F), called the totally orthogonal (or simply orthogonal) submodule to N (resp. M), and denoted \(N^0\) (resp. \(M^0\) ).

Let H and H' be two submodules of E or of F. One has \(\mathrm{H}\subset(\mathrm{H}^{0})^{0}\) (denoted \(H^{00}\) ); if \(H\subset H'\) , one has \(H^{\prime0}\supset H^{0}\) . It follows that one has \(H^{0}\supset(H^{00})^{0}\) and \(H^{0}\subset(H^{0})^{00}\) ; by setting

\[
\mathrm{H} ^ {0 0 0} = (\mathrm{H} ^ {0 0}) ^ {0} = (\mathrm{H} ^ {0}) ^ {0 0} = ((\mathrm{H} ^ {0}) ^ {0}) ^ {0},
\]

we therefore have \(\mathbf{H}^0 = \mathbf{H}^{000}\) .

For the map \(\Phi\) to be degenerate (no. 1, def. 3) it is necessary and sufficient that at least one of the two submodules E \(^{0}\) , F \(^{0}\) be \(\neq \{0\}\) . It is clear that \(\Phi(x, y)\) does not change when one adds to \(x\) (resp. \(y\) )

an element of F \(^{0}\) (resp. E \(^{0}\) ), and Φ therefore defines, by passing to the quotient, a bilinear (or sesquilinear) map on (E/F \(^{0}\) ) × (F/E \(^{0}\) ); this is visibly nondegenerate; it is called the nondegenerate bilinear (or sesquilinear) map associated with Φ.

Let \((\mathrm{E}_{i})_{i\in\mathrm{I}}\) be a family of left A-modules, \((\mathrm{F}_{i})_{i\in\mathrm{I}}\) a family of right B-modules (resp. left B-modules), \(\Phi_{i}\) a bilinear map (resp. sesquilinear on the right for J) from \(E_{i}\times F_{i}\) into G. Let E (resp. F) denote the direct sum module of the \(E_{i}\) (resp. \(F_{i}\) ). We see immediately that the map \(\Phi\) of \(E\times F\) into G defined by

\[
\Phi ((x _ {i}), (y _ {i})) = \sum_ {i} \Phi_ {i} (x _ {i}, y _ {i}) \quad (x _ {i} \in \mathrm{E} _ {i}, y _ {i} \in \mathrm{F} _ {i})\tag{10}
\]

(a sum that makes sense since its terms are zero except for finitely many of them) is bilinear (resp. sesquilinear on the right for J). It is called the direct sum of the maps \(\Phi_{i}\) . It is clear that \(\mathbf{E}_{i}\) is orthogonal to \(\mathbf{F}_{j}\) with respect to \(\Phi\) for \(i \neq j\) . Conversely, let \(\Phi\) be a bilinear or sesquilinear map from \(\mathbf{E} \times \mathbf{F}\) into G, and suppose that E is the direct sum of submodules \((\mathbf{E}_{i})_{i \in I}\) and F the direct sum of submodules \((\mathbf{F}_{i})_{i \in I}\) such that \(\mathbf{E}_{i}\) is orthogonal to \(\mathbf{F}_{j}\) for \(i \neq j\) ; then \(\Phi\) is the direct sum of its restrictions to the products \(\mathbf{E}_{i} \times \mathbf{F}_{i} (i \in I)\) .

For \(\Phi\) to be nondegenerate, it is necessary and sufficient that each of the \(\Phi_{i}\) be so; under these conditions, the submodule orthogonal to \(\mathbf{E}_{i}\) is \(\sum_{j\neq i}\mathbf{F}_{j}\) .

\subsection{Change of base rings.}

In this section, we denote by A, B, A', B' four rings, by h and h' homomorphisms from A to A' and from B to B' respectively, by G an (A, B)-bimodule, by G' an (A', B')-bimodule, and by u a homomorphism from the underlying abelian group of G to the underlying abelian group of G', satisfying

\[
u (a g b) = h (a) u (g) h ^ {\prime} (b) \quad (a \in \mathrm{A}, g \in \mathrm{G}, b \in \mathrm{B}).\tag{11}
\]

Let E (resp. F) be a left A-module (resp. a right B-module). Recall (Chap. III, 2nd ed., App. II, no. 10) that, if one regards A′ (resp. B′) as a right A-module (resp. left B-module), the tensor product \(E' = A' \otimes_{A} E\) (resp. \(F' = F \otimes_{B} B'\) ) is endowed with a left A'-module (resp. right B'-module) structure defined by

\[
\begin{array}{c c} a _ {1} ^ {\prime} (a ^ {\prime} \otimes x) = (a _ {1} ^ {\prime} a ^ {\prime}) \otimes x & (a ^ {\prime}, a _ {1} ^ {\prime} \in A ^ {\prime}, x \in E) \\ (\text { resp. } (y \otimes b ^ {\prime}) b _ {1} ^ {\prime} = y \otimes (b ^ {\prime} b _ {1} ^ {\prime}) & (b ^ {\prime}, b _ {1} ^ {\prime} \in B ^ {\prime}, y \in F)). \end{array}\tag{12}
\]

PROPOSITION 1. --- Let E be a left A-module and F a right B-module; set \(\mathbf{E}^{\prime} = \mathbf{A}^{\prime}\otimes_{\mathbf{A}}\mathbf{E}\) and \(\mathbf{F}^{\prime} = \mathbf{F}\otimes_{\mathbf{B}}\mathbf{B}^{\prime}\) . For every bilinear map \(\Phi\) of \(\mathbf{E}\times \mathbf{F}\) into G, there exists one and only one bilinear map \(\Phi^{\prime}\) from \(\mathbf{E}^{\prime}\times \mathbf{F}^{\prime}\) into G' such that one has

\[
\Phi^ {\prime} (a ^ {\prime} \otimes x, y \otimes b ^ {\prime}) = a ^ {\prime}. u (\Phi (x, y)). b ^ {\prime}\tag{13}
\]

for all \(a' \in \mathrm{A}'\) , \(b' \in \mathrm{B}'\) , \(x \in \mathrm{E}\) , \(y \in \mathrm{F}\) .

The uniqueness of \(\Phi'\) follows from the fact that the elements \(a' \otimes x\) and \(y \otimes b'\) generate E' and F' respectively. To prove their existence, consider the map

\[
m: (a ^ {\prime}, x, y, b ^ {\prime}) \rightarrow a ^ {\prime}. u (\Phi (x, y)). b ^ {\prime}
\]

from A' × E × F × B' into G; it is evidently Z-multilinear, and it satisfies

\[
\begin{array}{r} m (a ^ {\prime}, a x, y, b ^ {\prime}) = m (a ^ {\prime} h (a), x, y, b ^ {\prime}) \\ m (a ^ {\prime}, x, y b, b ^ {\prime}) = m (a ^ {\prime}, x, y, h ^ {\prime} (b) b ^ {\prime}) \\ (a \in \mathrm{A}, b \in \mathrm{B}, a ^ {\prime} \in \mathrm{A} ^ {\prime}, b ^ {\prime} \in \mathrm{B} ^ {\prime}, x \in \mathrm{E}, y \in \mathrm{F}). \end{array}
\]

There therefore exists a Z-bilinear map \(\Phi'\) from E' × F' into G' satisfying (13) (Chap. III, 2nd ed., App. II, no. 1, prop. 2). This relation and the definition of the module structures of E' and F' by (12) show that \(\Phi'\) is bilinear, which completes the proof.

Given the hypotheses and notation of Proposition 1, let us now study the linear maps associated with \(\Phi\) and with \(\Phi'\) (no. 1, def. 2). To this end, we shall first define a canonical homomorphism from \(\mathfrak{L}_{A}(\mathrm{E}, \mathrm{G})\) into \(\mathfrak{L}_{A'}(\mathrm{E}', \mathrm{G}')\) . For every \(\nu \in \mathfrak{L}_{\Lambda}(\mathrm{E}, \mathrm{G})\) the map \((a', x) \to a'. u(\nu(x))\) of \(\mathrm{A}' \times \mathrm{E}\) into \(\mathrm{G}'\) is Z-bilinear, and, in view of (11), maps \((a'h(a), x)\) and \((a', ax)\) ( \(a \in \mathrm{A}\) ) onto the same element of \(\mathrm{G}'\) ; it therefore defines (Chap. III, 2nd ed.,

App. II, nos. 1 and 10) a map \(k(\nu)\) of \(\mathbf{E}' = \mathbf{A}' \otimes_{\mathbf{A}} \mathbf{E}\) into \(\mathbf{G}'\) such that \(k(\nu)(a' \otimes x) = a'. u(\nu(x))\) , and which, by (12), is A'-linear. Moreover, one deduces immediately from (12) that the map \(\nu \to k(\nu)\) of \(\mathfrak{L}_{\mathbf{A}}(\mathbf{E}, \mathbf{G})\) into \(\mathfrak{L}_{\mathbf{A}'}(\mathbf{E}', \mathbf{G}')\) satisfies \(k(\nu b) = k(\nu) h'(b)\) for all \(b \in B\) . Let \(i\) denote the canonical map \(y \to y \otimes 1\) of F into \(\mathbf{F}'\) . Then the diagram

\[
\begin{array}{c c c} \mathrm{F} & \xrightarrow {d _ {\Phi}} & \mathcal {L} _ {A} (\mathrm{E}, \mathrm{G}) \\ \downarrow i & & \downarrow k \\ \mathrm{F} ^ {\prime} & \xrightarrow {d _ {\Phi^ {\prime}}} & \mathcal {L} _ {A^{\prime}} (\mathrm{E} ^ {\prime}, \mathrm{G} ^ {\prime}) \end{array}\tag{14}
\]

(where \(d_{\Phi}\) and \(d_{\Phi'}\) denote the linear maps associated on the right to \(\Phi\) and \(\Phi'\) ) is commutative. Indeed, for \(x \in \mathrm{E}, y \in \mathrm{F}\) and \(a' \in \mathrm{A}'\) , one has \(d_{\Phi'}(i(y))(a' \otimes x) = \Phi'(a' \otimes x, y \otimes 1) = a'. u(\Phi(x, y)) = a'. u(d_{\Phi}(y)(x))\) , that is to say \(d_{\Phi'}(i(y))(a' \otimes x) = k(d_{\Phi}(y))(a' \otimes x)\) . One has an analogous commutation relation for the linear maps \(s_{\Phi}\) and \(s_{\Phi'}\) associated on the left to \(\Phi\) and \(\Phi'\) .

PROPOSITION 2. --- Suppose that B and B' are equipped with anti-automorphisms J and I such that

\[
h ^ {\prime} (b ^ {J}) = h ^ {\prime} (b) ^ {I} \quad \text { for all } b \in \mathrm{B}.\tag{15}
\]

Let E be a left A-module and F a left B-module; set \(E' = A' \otimes_{A} E\) and \(F' = B' \otimes_{B} F\) . For every sesquilinear map (with respect to J) \(\Phi\) of \(E \times F\) into G, there exists one and only one sesquilinear map (for I) \(\Phi'\) from \(E' \times F'\) into \(G'\) such that

\[
\Phi^ {\prime} (a ^ {\prime} \otimes x, b ^ {\prime} \otimes y) = a ^ {\prime}. u (\Phi (x, y)). b ^ {\prime I}\tag{16}
\]

for all \(a' \in \mathbf{A}'\) , \(b' \in \mathbf{B}'\) , \(x \in \mathbf{E}\) , \(y \in \mathbf{F}\) .

The uniqueness of \(\Phi'\) follows from the fact that the tensor products \(a' \otimes x\) and \(b' \otimes y\) generate E′ and F′ respectively. To establish their existence, consider the map

\[
m: (a, x, b ^ {\prime}, y) \rightarrow a ^ {\prime}. u (\Phi (x, y)). b ^ {\prime I}
\]

of \(\mathbf{A}' \times \mathbf{E} \times \mathbf{B}' \times \mathbf{F}\) into \(\mathbf{G}\). It is evidently \(\mathbf{Z}\)-multilinear, and, taking (11) and (15) into account, satisfies \(m(a', ax, b', y) = m(a'h(a), x, b', y)\) and \(m(a', x, b', by) = a'. u(\Phi(x, y)) \cdot h'(b)^{I} b'^{I} = m(a', x, b'h'(b), y)\) ( \(a \in \mathbf{A}, b \in \mathbf{B}, a' \in \mathbf{A}'\) , \(b' \in \mathbf{B}'\) , \(x \in \mathbf{E}\) , \(y \in \mathbf{F}\) ). There therefore exists a \(\mathbf{Z}\)-bilinear map \(\Phi'\) from \(\mathbf{E}' \times \mathbf{F}'\) into \(\mathbf{G}'\)

satisfying (16) (chap. III, 2nd ed., App. II, no. 1, prop. 2). This relation, together with the definition of the module structures of E′ and F′ by (12), shows, in view of (15), that Φ′ is sesquilinear for I, which completes the proof.

The most important examples of \((\mathbf{A}^{\prime}, \mathbf{B}^{\prime})\) -bimodules \(G^{\prime}\) equipped with Z-linear maps u from G into \(G^{\prime}\) satisfying (11), are the following:

\begin{enumerate}
\def\labelenumi{\arabic{enumi})}
\tightlist
\item
  We take for G′ the tensor product A′ ⊗\_A G ⊗\_B B′ (chap. III, 2nd ed., App. II, no. 9) and for u the map
\end{enumerate}

\[
g \to 1 \otimes g \otimes 1\tag{\(g\in G\)}
\]

from G to G'. The pair (G', u) thus defined is visibly universal in the following sense: for every (A', B')-bimodule G₁' and every Z-linear map u₁ from G to G₁' satisfying the analogue of (11), there exists a unique Z-linear map f from G' to G₁' such that f(a'g'b') = a'f(g')b' (a' ∈ A', g' ∈ G', b' ∈ B'; in other words, f is a homomorphism for the bimodule structures of G' and G₁') and such that u₁ = f ∘ u.

\begin{enumerate}
\def\labelenumi{\arabic{enumi})}
\setcounter{enumi}{1}
\item
  When A = B = G (the structure of (A, A)-bimodule of A being defined by left and right homotheties), \(A' = B'\) , and \(h = h'\) one can take for \(G'\) the ring \(A'\) and for u the homomorphism h from A into \(A'\) .
\item
  Suppose that we have A = B, \(A' = B'\) , \(h = h'\) , that the rings A and \(A'\) are commutative, and that the structure of left A-module of G coincides with its structure of right A-module. One can then take for \(G'\) the tensor product \(A' \otimes_{A} G\) (the structure of \(A'\) -module on the right of \(G'\) coinciding with its structure of \(A'\) -module on the left) and for u the map \(g \to 1 \otimes g\) ( \(g \in G\) ) from G into \(G'\) . We shall then say that the bilinear (resp. sesquilinear) map \(\Phi'\) defined by prop. 1 (resp. prop. 2) is obtained from \(\Phi\) by extension of the base ring, or by extension of scalars.
\end{enumerate}

What follows is valid equally for bilinear maps and for sesquilinear maps; the hypotheses and notations are those of prop. 1 (resp. prop. 2). Given a submodule M of E or F, we denote by M' the submodule of E' or F' generated by the canonical image of M.

PROPOSITION 3. --- With the hypotheses and notations of prop. 1 (resp. prop. 2), suppose moreover that A, B, A', B' are fields and that the maps \(\alpha\) and \(\beta\) from A'⊗A G and G⊗B B' into G' characterized by \(\alpha(a'\otimes g)=a'u(g)\) and \(\beta(g\otimes b')=u(g)b'\) ( \(a'\in A'\) , \(b'\in B'\) , \(g\in G\) ) are injective. Let M be a subspace of E and N a subspace of F. Then the subspace \((M')^{0}\) of F' orthogonal to M' with respect to \(\Phi'\) is equal to \((M^{0})^{\prime}\) , and, similarly, we have \((N')^{0} = (N^{0})^{\prime}\) .

Indeed, the inclusions \((\mathbf{M}^{0})^{\prime}\subset(\mathbf{M}^{\prime})^{0}\) and \((\mathbf{N}^{0})^{\prime}\subset(\mathbf{N}^{\prime})^{0}\) are evident (and moreover true without hypotheses on A, B, A', B', \(\alpha\) or \(\beta\) ). We shall prove the inclusion \((\mathbf{M}^{\prime})^{0}\subset(\mathbf{M}^{0})^{\prime}\) ; we leave to the reader the verification of the inclusion \((\mathbf{N}^{\prime})^{0}\subset(\mathbf{N}^{0})^{\prime}\) , which is entirely analogous. Let \(y^{\prime}\) be an element of \((\mathbf{M}^{\prime})^{0}\) . One can write

\[
y ^ {\prime} = \sum_ {i = 1} ^ {s} y _ {i} \otimes b _ {i} ^ {\prime} \quad (\text { resp. } y ^ {\prime} = \sum_ {i = 1} ^ {s} b _ {i} ^ {\prime} \otimes y _ {i})
\]

where \(y_i \in \mathbf{F} (1 \leqslant i \leqslant s)\) , and where the \(b_i'\) are elements of \(\mathbf{B}'\) which are linearly independent over \(\mathbf{B}\) for the structure of \(\mathbf{B}\) -module on the left (resp. on the right) of \(\mathbf{B}'\) . Let \(x \in \mathbf{M}\) and \(x' = 1 \otimes x \in \mathbf{M}'\) . We have

\[
0 = \Phi^ {\prime} (x ^ {\prime}, y ^ {\prime}) = \sum_ {i} u (\Phi (x, y _ {i})) b _ {i} ^ {\prime} = \beta (\sum_ {i} \Phi (x, y _ {i}) \otimes b _ {i} ^ {\prime})
\]

\[
(\text { resp. } 0 = \Phi^ {\prime} (x ^ {\prime}, y ^ {\prime}) = \sum_ {i} u (\Phi (x, y _ {i})) b _ {i} ^ {\prime \mathrm{I}} = \beta (\sum_ {i} \Phi (x, y _ {i}) \otimes b _ {i} ^ {\prime \mathrm{I}}).
\]

Since β is injective and since the \(b_i'\) (resp. the \(b _ {i} ^ {\prime \mathrm{I}}\) , taking (15) into account) are linearly independent over B for the structure of left B-module of B', this entails \(\Phi(x, y_i) = 0\) for \(i = 1, \ldots, s\) . Since this last relation holds for every \(x \in \mathbf{M}\) , one has \(y_i \in \mathbf{M}^0\) for \(i = 1, \ldots, s\) , whence \(y' \in (\mathbf{M}^0)'\) . QED.

COROLLARY. --- With the hypotheses and notations of prop. 3, for \(\Phi'\) to be nondegenerate, it is necessary and sufficient that \(\Phi\) be nondegenerate.

Indeed, by prop. 3, we have \((\mathbf{F}')^0 = (\mathbf{F}^0)'\) and \((\mathbf{E}')^0 = (\mathbf{E}^0)'\) . On the other hand, for \(\Phi\) (resp. \(\Phi'\) ) to be nondegenerate, it is necessary and sufficient that one have \(\mathbf{F}^0 = \mathbf{E}^0 = \{0\}\) (resp. \((\mathbf{F}')^0 = (\mathbf{E}')^0 = \{0\}\) ).

Remark. --- Suppose that A, B, A' and B' are fields. Then, for the bimodules G' defined in the three examples above, the maps α and β are injective, as follows immediately from chap. III, 2nd ed., App. II, no. 6.

\subsection{Some identities.}

In this section, we denote by A a ring equipped with an anti-automorphism J, by E a left A-module, by G an (A, A)-bimodule, and by Φ a sesquilinear map (on the right) for J from E × E into G. We set Q(x) = Φ(x, x) (x ∈ E). We evidently have

\[
\mathrm{Q} (x + y) = \mathrm{Q} (x) + \Phi (x, y) + \Phi (y, x) + \mathrm{Q} (y)\tag{17}
\]

\[
\mathrm{Q} (x - y) = \mathrm{Q} (x) - \Phi (x, y) - \Phi (y, x) + \mathrm{Q} (y)\tag{18}
\]

whatever x, y in E. Hence, by subtraction,

\[
2 (\Phi (x, y) + \Phi (y, x)) = Q (x + y) - Q (x - y).\tag{19}
\]

Let \(a\) be an element of A; replacing \(y\) by \(ay\) in (19), we obtain

\[
2 (\Phi (x, y) a ^ {j} + a \Phi (y, x)) = Q (x + a y) - Q (x - a y).\tag{20}
\]

From (19) and (20), multiplying (19) by \(a\) (on the left) and subtracting, we obtain:

\[
\begin{array}{r l} 2 (a \Phi (x, y) - \Phi (x, y) a ^ {J}) & \\ = a Q (x + y) - a Q (x - y) - Q (x + a y) + Q (x - a y). \end{array}\tag{21}
\]

Suppose in particular that A is a quadratic extension K(i) of a commutative ring K, with \(i^{2} = -1\) (chap. II, § 7, no. 7), that J is the K-automorphism \(u + iv \rightarrow u - iv\) (u, v in K) of A, and that the left A-module and right A-module structures of G coincide. Setting a = i in (21), we obtain

\[
4 \Phi (x, y) = Q (x + y) - Q (x - y) + i Q (x + i y) - i Q (x - i y). \tag {22}
\]

\subsection{Bilinear and sesquilinear forms. Rank.}

In this section, A denotes a ring (resp. a ring equipped with an antiautomorphism J), E a left A-module, and F a right A-module (resp. a left A-module). One endows A with the structure of an (A, A)-bimodule defined by left multiplications and right multiplications. In this case, a bilinear (resp. right sesquilinear for J) map from E × F into the bimodule A is called a bilinear (resp. right sesquilinear for J) form on E × F.

When E = F (which implies that it is a sesquilinear form), one often says that a sesquilinear form on E × F is a sesquilinear form on E.

Given two left A-modules E and E′, and two sesquilinear forms \(\Phi\) and \(\Phi'\) for J on E and E' respectively, one says that \(\Phi\) and \(\Phi'\) are equivalent if there exists an isomorphism \(u\) of the A-module E onto the A-module E′ such that \(\Phi'(u(x), u(y)) = \Phi(x, y)\) for all \(x\) , \(y\) in E; then \(\Phi\) is the inverse image of \(\Phi'\) with respect to \(u\) and \(u\) , and \(\Phi'\) is the inverse image of \(\Phi\) with respect to \(u^{-1}\) and \(u^{-1}\) (no. 2).

Let \(\Phi\) be a bilinear form on \(E \times F\) (where F denotes a right A-module). The linear maps \(s_{\Phi}\) and \(d_{\Phi}\) associated with \(\Phi\) (no. 1, def. 2) are then maps from E into the dual \(F^{*}\) of F, and from F into the dual \(E^{*}\) of E.

Thus, by definition,

\[
\Phi (x, y) = \left\langle x, d _ {\Phi} (y) \right\rangle = \left\langle y, s _ {\Phi} (x) \right\rangle .\tag{23}
\]

We shall now define the linear maps associated with a sesquilinear form. Let J be an antiautomorphism of A, and \(\Phi\) a sesquilinear form (on the right) for J on E × F (where F denotes a left A-module); set J' = J \(^{-1}\) . The map \(\Phi'\) from F × E into A defined by

\[
\Phi^ {\prime} (y, x) = \Phi (x, y) ^ {J ^ {\prime}}
\]

\[
(x \in \mathrm{E}, y \in \mathrm{F})
\]

is, as is easily seen, a sesquilinear form (on the right) for J' on F × E. By no. 2 (def. 5), the sesquilinear forms Φ and Φ' identify respectively with bilinear forms on E × F' and on F × E'. The maps dΦ and dΦ' associated with the latter are called the maps associated on the right and on the left with the sesquilinear form \(\Phi\) , and are denoted \(d_{\Phi}\) and \(s_{\Phi}\) . We therefore have, by definition:

\[
\Phi (x, y) = \left\langle x, d _ {\Phi} (y) \right\rangle = \left\langle y, s _ {\Phi} (x) \right\rangle^ {J} \quad (x \in \mathrm{E}, y \in \mathrm{F}).\tag{24}
\]

Thus \(d_{\Phi}\) (resp. \(s_{\Phi}\) ) is a linear map from \(\mathbf{F}^{\mathrm{j}}\) into \(\mathbf{E}^{*}\) (resp. of \(\mathbf{E}^{\mathrm{j}'}\) into \(\mathbf{F}^{*}\) ), or equivalently a semilinear map from \(\mathbf{F}\) into \(\mathbf{E}^{*}\) (resp. of \(\mathbf{E}\) into \(\mathbf{F}^{*}\) ) relative to \(\mathbf{J}\) (resp. \(\mathbf{J}'\) ) if one regards \(\mathbf{J}\) (resp. \(\mathbf{J}'\) ) as a ring isomorphism \(\mathbf{A}^{0}\) (the opposite of \(\mathbf{A}\) ) on \(\mathbf{A}\) , and \(\mathbf{F}\) (resp. \(\mathbf{E}\) ) as a \(\mathbf{A}^{0}\) -module on the right.

Formula (24) and def. 6 of no. 3 immediately imply that for every submodule N of F (resp. M of E), we have

\[
\mathbf {N ^ {0}} = \stackrel {- 1} {s _ {\Phi}} (\mathbf {N ^ {\prime}}) \quad (\text { resp. } \mathbf {M ^ {0}} = \stackrel {- 1} {d _ {\Phi}} (\mathbf {M ^ {\prime}}))\tag{25}
\]

where N′ (resp. M′) is the submodule of the dual F* of F (resp. of the dual E* of E) orthogonal to N (resp. M) (Chap. II, § 4, no. 2).

PROPOSITION 4. --- Suppose that A is a field, and let \(\Phi\) be a bilinear form (resp. sesquilinear form with respect to J) on \(E \times F\) ; for \(E/F^{0}\) to be of finite dimension, it is necessary and sufficient that \(F/E^{0}\) be finite-dimensional, and these dimensions are then equal.

Indeed, let \(\Phi_{1}\) be the nondegenerate form associated with \(\Phi\) , on \((\mathrm{E} / \mathrm{F}^{0}) \times (\mathrm{F} / \mathrm{E}^{0})\) (no. 3). Suppose that \(\mathrm{E} / \mathrm{F}^{0}\) is of finite dimension \(n\) ; since the linear map \(d_{\Phi_1}\) of \(\mathrm{F} / \mathrm{E}^{0}\) (resp. \((\mathrm{F} / \mathrm{E}^{0})^{\mathrm{j}})\) into \((\mathrm{E} / \mathrm{F}^{0})^{*}\) is injective, \(\mathrm{F} / \mathrm{E}^{0}\) is of finite dimension \(n' \leqslant n\) ; by considering \(s_{\Phi_1}\) one sees likewise that \(n \leqslant n'\) .

COROLLARY 1. --- Suppose that A is a field and that \(\Phi\) is nondegenerate. For a subspace M of E to be of finite dimension, it is necessary and sufficient that M \(^{0}\) be of finite codimension in F, and one then has codim M \(^{0}\) = dim M, and M \(^{00}\) = M.

As \(F^{0}=\{0\}\) , the first two assertions follow from Prop. 4 applied to the restriction of \(\Phi\) to M × F. Moreover, \(M^{0}\) is the orthogonal of \(M^{00}\) , hence \(M^{00}\) is of finite dimension equal to codim \(M^{0}=dim M\) ; but since \(M^{00}\supset M\) , one has \(M^{00}=M\) .

COROLLARY 2. --- Under the hypotheses of cor. 1, let M, N be two subspaces of E; we then have \((M + N)^{0} = M^{0} \cap N^{0}\) ; if moreover M and N are of finite dimension, we have \((M \cap N)^{0} = M^{0} + N^{0}\) .

The first assertion is trivial. Suppose M and N are finite-dimensional, and let \(G = M^{0} + N^{0}\) ; we have \(G^{0} = M^{00} \cap N^{00} = M \cap N\) by cor. 1; prop. 4 applied to the restriction of \(\Phi\) to \(M \times G\) then shows (since \(M^{0} \subset G\) and \(G^{0} \subset M\) ) that one has dim \(M/(M \cap N) = \dim G/M^{0} = \text{codim } M^{0} - \text{codim } G\) , and since codim \(M^{0} = \dim M\) , we deduce dim \((M \cap N) = \text{codim } G\) . But we also have dim \((M \cap N) = \text{codim } (M \cap N)^{0}\) by cor. 1, and since \(G \subset G^{00} = (M \cap N)^{0}\) we have \(G = (M \cap N)^{0}\) .

Prop. 4 allows us to make the following definition:

DEFINITION 7. --- Let A be a field (resp. a field equipped with an antiautomorphism J), E a left vector space over A, F a right (resp. left) vector space over A, and Φ a bilinear (resp. sesquilinear for J) form on E × F. Suppose that E/F⁰ and F/E⁰ are of finite dimension over A. The rank of Φ is the (finite) common dimension of the vector spaces E/F⁰ and F/E⁰.

When \(E/F^{0}\) and \(F/E^{0}\) are of infinite dimension, one says that \(\Phi\) is of infinite rank.

PROPOSITION 5. --- With the hypotheses and notation as in def. 7, the linear maps \(s_{\Phi}\) and \(d_{\Phi}\) associated with \(\Phi\) have the same rank, and this rank is equal to the rank of the form \(\Phi\) .

Indeed, the kernel of the map \(d_{\Phi}\) of F into E* is evidently E°, hence its rank equals the dimension of F/E°. Likewise, the rank of \(s_{\Phi}\) is equal to the dimension of E/F°.

PROPOSITION 6. --- With the hypotheses and notation of Definition 7, suppose further that E and F have the same finite dimension. Then the following conditions are equivalent:

\begin{enumerate}
\def\labelenumi{\alph{enumi})}
\item
  \(d_{\Phi}\) is injective;
\item
  \(d_{\Phi}\) is surjective;
\item
  \(s_{\Phi}\) is injective;
\end{enumerate}

\(d)\) \(s_{\Phi}\) is surjective;

\begin{enumerate}
\def\labelenumi{\alph{enumi})}
\setcounter{enumi}{4}
\tightlist
\item
  \(\Phi\) is nondegenerate.
\end{enumerate}

Indeed, since E, F, E* and F* have the same finite dimension, a) and b) are equivalent, as are c) and d) (chap. II, § 3, no. 4). Since \(s_{\Phi}\) and \(d_{\Phi}\) have the same rank (prop. 5), a) and c) are equivalent. Since e) is equivalent to the relation \(\mathrm{E}^{0} = \mathrm{F}^{0} = \{0\}\) , it is equivalent to the conjunction of a) and c), whence the equivalence of the stated conditions.

COROLLARY. --- With the hypotheses and notations being those of def. 7, suppose further that E is finite-dimensional and that \(\Phi\) is nondegenerate. Then we have dim E = dim F and, for every basis \((e_i)\) ( \(1 \leqslant i \leqslant \dim E\) ) of E, there exists a basis \((f_i)\) of F such that \(\Phi(e_i, f_k) = \delta_{ik}\) ( \(i, k = 1, \ldots, \dim E\) ).

Indeed, as \(\Phi\) is nondegenerate, we have \(E^{0} = F^{0} = \{0\}\) , whence dim \(E = \dim F\) (prop. 4). It follows (prop. 6) that \(d_{\Phi}\) is an isomorphism of \(F\) (resp. \(F^{j}\) ) onto \(E^{*}\) ; therefore, if \((e_{i}^{*})\) is the dual basis of \((e_{i})\) , the elements \(f_{i} = d_{\Phi}^{-1}(e_{i}^{*})\) form a basis of \(F\) which, in view of formula (23) (resp. formula (24)), satisfies \(\Phi(e_{i}, f_{k}) = \delta_{ik}\) .

It is immediate that, in this corollary, one may interchange the roles of E and F, replacing \(d_{\Phi}\) by \(s_{\Phi}\) in the proof.

Remark. --- Let A be a ring equipped with an antiautomorphism J, and let M and N be right A-modules, and \(\Phi\) a left sesquilinear form for J on M × N (no. 2, Remark); it therefore satisfies the equality

\[
\Phi (x a, y a ^ {\prime}) = a ^ {j} \Phi (x, y) a ^ {\prime} \quad (a, a ^ {\prime} \in A, x \in M, y \in N).
\]

The map \(\Phi'\) from N × M into A defined by \(\Phi'(y, x) = \Phi(x, y)^{J'}\) (where J' = J \(^{-1}\) ) is a left sesquilinear form for J', and \(\Phi\) and \(\Phi'\) are identified with bilinear forms on M \(^{J'}\) × N and N \(^{J'}\) × M, respectively. The maps \(s_{\Phi}\) and \(s_{\Phi'}\) associated with these bilinear forms are called the maps associated on the left and on the right with the sesquilinear form \(\Phi\) , and are denoted \(s_{\Phi}\) and \(d_{\Phi}\) . We therefore have, by definition

\[
\Phi (x, y) = \left\langle y, s _ {\Phi} (x) \right\rangle = \left\langle x, d _ {\Phi} (y) \right\rangle^ {J} \quad (x \in M, y \in N),\tag{26}
\]

and \(s_{\Phi}\) (resp. \(d_{\Phi}\) ) is a linear map from \(\mathbf{M}^{\mathbf{j}}\) into \(\mathbf{N}^{*}\) (resp.

of N \(^{j'}\) into M*). One could easily state and prove the analogues, for the case considered here, of def. 7 and Props. 4, 5, 6.

\subsection{Inverse form of a bilinear or sesquilinear form.}

Let A be a ring, E a left A-module, F a right A-module, and \(\Phi\) a bilinear form on \(\mathbf{E} \times \mathbf{F}\) . It is assumed here that the associated maps of \(\Phi\) , which will be denoted \(s\) and \(d\) , are bijective. Then the product map \((s, d)\) is a bijection from \(\mathbf{E} \times \mathbf{F}\) onto \(\mathbf{F}^* \times \mathbf{E}^*\) , and defines, by transport of structure, a bilinear form \(\hat{\Phi}\) on \(\mathbf{F}^* \times \mathbf{E}^*\) . This therefore satisfies

\[
\begin{array}{r l} \hat {\Phi} (y ^ {\prime}, x ^ {\prime}) = & \Phi (s ^ {- 1} (y ^ {\prime}), d ^ {- 1} (x _ {i} ^ {\prime})) \\ = & \left\langle s ^ {- 1} (y ^ {\prime}), x ^ {\prime} \right\rangle = \left\langle d ^ {- 1} (x ^ {\prime}), y ^ {\prime} \right\rangle \qquad (x ^ {\prime} \in \mathrm{E} ^ {*}, y ^ {\prime} \in \mathrm{F} ^ {*}). \end{array}\tag{27}
\]

DEFINITION 8. --- Let \(\Phi\) be a bilinear form on E × F whose associated maps s and d are bijective. The bilinear form \(\hat{\Phi}\) on F* × E* defined by (27) is called the inverse form of \(\Phi\) .

Now let \(\hat{s}\) and \(\hat{d}\) be the linear maps from F* to E** and from E* to F** associated on the left and on the right to \(\hat{\Phi}\) . Since, for \(x' \in \mathrm{E}^*\) and \(y' \in \mathrm{F}^*\) , we have by definition

\[
\hat {\Phi} (y ^ {\prime}, x ^ {\prime}) = \left\langle y ^ {\prime}, \hat {d} (x ^ {\prime}) \right\rangle = \left\langle x ^ {\prime}, \hat {s} (y ^ {\prime}) \right\rangle
\]

Comparing with (27), one sees that the linear form \(\hat{d}(x')\) on F* is equal to that defined by the element \(d^{-1}(x')\) of F. It follows that the composite map \(\hat{d} \circ d\) is the canonical map from F to its bidual F**, and that this map is bijective since \(d\) and \(\hat{d}\) (the latter by transport of structure) are bijective; hence, if one canonically identifies F with F**, one has \(\hat{d} = d^{-1}\) . Similarly, E is canonically identified with E**, and the canonical map from E to E** is \(\hat{s} \circ s\) and one has \(\hat{s} = s^{-1}\) . It follows from this that the inverse form of \(\hat{\Phi}\) is \(\Phi\) .

Now consider a ring A equipped with an antiautomorphism J, two left A-modules E and F, and a right sesquilinear form \(\Phi\) for J on E \(\times\) F such that the associated maps of \(\Phi\) , which will be denoted \(s\) and \(d\) , are bijective. Define a map \(\hat{\Phi}\) of \(\mathbf{F}^{*} \times \mathbf{E}^{*}\) into A by the first equation (27). This map satisfies, by (24) (no. 6), the relation

\[
\hat {\Phi} (y ^ {\prime}, x ^ {\prime}) = \left\langle s ^ {- 1} (y ^ {\prime}), x ^ {\prime} \right\rangle = \left\langle d ^ {- 1} (x ^ {\prime}), y ^ {\prime} \right\rangle^ {J} \quad (x ^ {\prime} \in E ^ {*}, y ^ {\prime} \in F ^ {*}).\tag{28}
\]

The map \(\hat{\Phi}\) is evidently \(\mathbf{Z}\) -bilinear; moreover, we have, for \(a\) , \(b\) in A, \(x' \in \mathrm{E}^*\) and \(y' \in \mathrm{F}^*\) and by virtue of the definitions of \(s\) and \(d\) ,

\[
\widehat {\Phi} (y ^ {\prime} a, x ^ {\prime} b) = \Phi (a ^ {j} s ^ {- 1} (y ^ {\prime}), b ^ {j ^ {\prime}} d ^ {- 1} (x ^ {\prime})) = a ^ {j} \widehat {\Phi} (y ^ {\prime}, x ^ {\prime}) b;
\]

therefore \(\hat{\Phi}\) is a left sesquilinear form for J (no. 2) on \(\mathbf{F}^{*} \times \mathbf{E}^{*}\) .

DEFINITION 9. --- Let \(\Phi\) be a right-sesquilinear form for J on E × F, whose associated maps s and d are bijective. The left-sesquilinear form \(\hat{\Phi}\) for J on F* × E* is called the inverse form of \(\Phi\) .

We leave it to the reader to define and study the inverse form of a left sesquilinear form. This inverse form is a right sesquilinear form.

Let \(\hat{s}\) and \(\hat{d}\) be the associated maps of \(\hat{\Phi}\) ; by (26) (no. 6) we have

\[
\hat {\Phi} (y ^ {\prime}, x ^ {\prime}) = \left\langle y ^ {\prime}, \hat {d} (x ^ {\prime}) \right\rangle^ {\mathbf {J}} = \left\langle x ^ {\prime}, \hat {s} (y ^ {\prime}) \right\rangle .\tag{29}
\]

Because of the fact that \(s\) is bijective, and from the equality \(\langle s^{-1}(y'), x'\rangle = \langle y', \hat{d}(x') \rangle^{\mathfrak{J}}\) which results from (28) and (29), we deduce that \(\hat{d}\) is bijective; therefore \(\hat{d} \circ d\) is bijective. But the equality \(\langle d^{-1}(x'), y' \rangle = \langle y', \hat{d}(x') \rangle\) which results from (28) and (29), shows that \(\hat{d} \circ d\) is the canonical map from F into its bidual F**. One sees similarly that \(\hat{s}\) is bijective and that \(\hat{s} \circ s\) is the canonical map from E into E**. Therefore, if one identifies E** with E and F** with F by means of these canonical maps, we have \(\hat{s} = s^{-1}\) , \(\hat{d} = d^{-1}\) , and \(\Phi\) is the inverse form of \(\hat{\Phi}\) .

With the same notations and hypotheses, let a be an invertible element of the center of A. Then the maps associated with the form \(a\Phi\) are, by (23) (resp. (24)), equal to \(a.d\) and \(a.s\) (resp. \(a^{j'}\) .s), hence are bijective. It thus follows from (27) that the inverse form of \(a\Phi\) is \(a^{-1}\widehat{\Phi}\) (resp. \((a^{j'})^{-1}\widehat{\Phi}\) ).

\subsection{Adjoint of a homomorphism.}

In this section, we denote by A a ring (resp. a ring equipped with an antiautomorphism J), by E and E' two left A-modules, by F and F' two right A-modules (resp. left A-modules), and by Φ and Φ' two bilinear forms (resp. sesquilinear forms for J) on E × F and E' × F' respectively. We assume that Φ is nondegenerate, in other words (no. 1) that the linear maps \(d_{\Phi}\) and \(s_{\Phi}\) associated with Φ are injective.

Given a homomorphism \(u\) from E into \(\mathbf{E}'\) , consider the set \(\mathbf{F}_{1}'\) of elements \(y'\) of \(\mathbf{F}'\) such that there exists \(y \in \mathbf{F}\) for which one has \(d_{\Phi'}(y') \circ u = d_{\Phi}(y)\) , that is to say \(\Phi'(u(x), y') = \Phi(x, y)\) for all \(x \in \mathbf{E}\) . It is clear that \(\mathbf{F}_{1}'\) is a submodule of \(\mathbf{F}'\) . Since \(d_{\Phi}\) is injective, there exists, for every \(y' \in \mathbf{F}_{1}'\) , one and only one element \(y\) of F such that \(\Phi'(u(x), y') = \Phi(x, y)\) . The map \(y' \to y\) of \(\mathbf{F}_{1}'\) into F thus defined is A-linear; denoting it by \(u^{*}\) , one has, for every \(x \in \mathbf{E}\) and every \(y' \in \mathbf{F}_{1}'\)

\[
\Phi^ {\prime} (u (x), y ^ {\prime}) = \Phi (x, u ^ {*} (y ^ {\prime})).\tag{30}
\]

DEFINITION 10. --- The hypotheses and notations being as before, one says that the homomorphism \(u^*\) of \(F_1'\) into F satisfying (30) is the left adjoint of \(u\) , and that \(F_1'\) is the defining submodule of \(u^*\) .

One defines similarly the right adjoint of a homomorphism \(\nu\) from F into F' by the formula

\[
\Phi^ {\prime} (x ^ {\prime}, \nu (y)) = \Phi (\nu^ {*} (x ^ {\prime}), y) \quad (x ^ {\prime} \in \mathrm{E} _ {1} ^ {\prime}, y \in \mathrm{F}),\tag{31}
\]

where \(E_{1}^{\prime}\) denotes the submodule of \(E^{\prime}\) defined analogously to \(F_{1}^{\prime}\) .

Remark. --- If the left adjoint \(u^{*}\) of \(u: E \to E'\) is everywhere defined, and if \(s_{\Phi'}\) and \(d_{\Phi'}\) are injective, formula (30) shows that \(u\) is the right adjoint of \(u^{*}\) .

From (30), we deduce that, if \(u_1\) and \(u_2\) are two homomorphisms from E into E′ admitting everywhere-defined adjoints, and if c is an element of the center of A, then we have

\[
\left\{ \begin{array}{l} (u _ {1} + u _ {2}) ^ {*} = u _ {1} ^ {*} + u _ {2} ^ {*}; 1 ^ {*} = 1; \\ (c u _ {1}) ^ {*} = c. u _ {1} ^ {*} \quad \text { when } \Phi \text { and } \Phi^ {\prime} \text { are bilinear }; \\ (c u _ {1}) ^ {*} = c ^ {J ^ {\prime}}. u _ {1} ^ {*} \quad \text { when } \Phi \text { and } \Phi^ {\prime} \text { are sesquilinear }. \end{array} \right.\tag{32}
\]

Additionally, if \(E^{\prime\prime}\) is a third left A-module, \(F^{\prime\prime}\) a third right (resp. left) A-module, \(\Phi^{\prime\prime}\) a bilinear form (resp. sesquilinear form with respect to J) on \(E^{\prime\prime} \times F^{\prime\prime}\) , and if \(u^{\prime}\) is a homomorphism of \(E^{\prime}\) into \(E^{\prime\prime}\) admitting a (left) adjoint defined everywhere, we have

\[
(u ^ {\prime} \circ u) ^ {*} = u ^ {*} \circ u ^ {\prime *}.\tag{33}
\]

In particular, if \(u\) is an isomorphism from E onto E', and if the adjoints \(u^{*}\) and \((u^{-1})^{*}\) are everywhere defined, \(u^{*}\) is an isomorphism from F' onto F, and we have \((u^{*})^{-1} = (u^{-1})^{*}\) . Analogous properties hold for right adjoints.

PROPOSITION 7. --- With the same notation as above, assume that \(d_{\Phi}\) is bijective. Then every homomorphism \(u\) from E to E' admits an everywhere-defined left adjoint, and we have \(u^{*} = (d_{\Phi})^{-1} \circ {}^{t}u \circ d_{\Phi'}\) .

Indeed, as \(d_{\Phi}\) is bijective, we have, with the notations from the beginning of the paragraph, \(\mathbf{F}_{1}^{\prime}=\mathbf{F}^{\prime}\) , and \(u^{*}\) is therefore everywhere defined. On the other hand, (30) is equivalent to

\[
\left\langle u (x), d _ {\Phi^ {\prime}} \left(y ^ {\prime}\right) \right\rangle = \left\langle x, \left(d _ {\Phi} \circ u ^ {*}\right) \left(y ^ {\prime}\right) \right\rangle \quad \left(x \in \mathrm{E}, y ^ {\prime} \in \mathrm{F} ^ {\prime}\right);
\]

or \(\langle d_{\Phi'}(y'), u(x) \rangle = \langle {}^t u(d_{\Phi'}(y')), x \rangle\) ; one therefore has \({}^t u(d_{\Phi'}(y')) = d_{\Phi}(u^*(y'))\) for all \(y' \in F'\) , whence \({}^t u \circ d_{\Phi'} = d_{\Phi} \circ u^*\) , and consequently the announced expression of \(u^*\) . QED.

Remark. --- When \(s_{\Phi}\) is bijective, every homomorphism \(\nu\) from F to F' admits a right adjoint that is everywhere defined, and we have

\[
v ^ {*} = (s _ {\Phi}) ^ {- 1} \circ {} ^ {t} v \circ s _ {\Phi^ {\prime}}.\tag{34}
\]

PROPOSITION 8. --- With the same notation as above, assume that \(s_{\Phi}\) and \(d_{\Phi}\) are bijective. Let \(u\) and \(v\) be isomorphisms from E onto \(\mathbf{E}'\) and from F onto \(\mathbf{F}'\) respectively. Then, for \(\Phi\) to be the inverse image of \(\Phi'\) with respect to \(u\) and \(\nu\) (that is, so that one has \(\Phi(x,y)=\Phi'(u(x),\nu(y))\) for all \(x\in\mathbf{E},y\in\mathbf{F}\) ), it is necessary and sufficient that we have \(u^{-1}=\nu^{*}\) and \(\nu^{-1}=u^{*}\) .

Indeed \(\Phi'(u(x), v(y)) = \Phi(x, y)\) is also written \(\Phi(x, u^*(v(y))) = \Phi(x, y)\) . If this holds for all \(x \in E\) and \(y \in F\) , one has \(u^* \circ v = 1\) since \(\Phi\) is nondegenerate. We therefore also have \(v^* \circ u = 1\) by (33). The converse is immediate.

COROLLARY. --- Let A be a ring equipped with an antiautomorphism J, and let E be a left A-module, \(\Phi\) a sesquilinear form for J on E × E whose associated maps are bijective, and u an automorphism of the A-module E. For u to leave \(\Phi\) invariant (that is, so that we have \(\Phi(u(x), u(y)) = \Phi(x, y)\) for all x, y in E), it is necessary and sufficient that the two adjoints of u be equal and that we have \(u^{*} = u^{-1}\) .

This follows immediately from prop. 8.

Remark. --- Under the hypotheses of the corollary to Proposition 8, suppose further that A is a field and that E is finite-dimensional over A. Let \(w\) be an endomorphism of E, \(w_{1}\) and \(w_{2}\) its right and left adjoints. Each of the conditions \(ww_{1}=1\) , \(ww_{2}=1\) , \(w_{1}w=1\) , \(w_{2}w=1\) implies that \(w\) is an automorphism of E leaving \(\Phi\) invariant, and that \(w_{1}=w_{2}\) .

\subsection{Tensor products and exterior powers of sesquilinear forms.}

In this section, A denotes a commutative ring. A bilinear form on a product of two A-modules is therefore a special case of a sesquilinear form. We denote by J an automorphism of A, and by J' its inverse.

Let \(\mathbf{E}_i\) ( \(i = 1, \ldots, m\) ) be A-modules. The map

\[
(x _ {1}, \dots , x _ {m}) \rightarrow x _ {1} \otimes \dots \otimes x _ {m}
\]

of \(\prod_{i=1}^{m}\mathrm{E}_{i}^{j}\) into \((\bigotimes_{i=1}^{m}\mathrm{E}_{i})^{j}\) ( \(x_{i}\in\mathrm{E}_{i}^{j}\) ) (cf. def. 5, no. 2) is evidently A-multilinear; it therefore defines (chap. III, § 1, no. 7) an A-linear map \(f\) of \(\bigotimes_{i}\mathrm{E}_{i}^{\mathrm{J}}\) into \((\bigotimes_{i}\mathrm{E}_{i})^{\mathrm{J}}\) ; this map sends \(x_{1}\otimes\cdots\otimes x_{m}\) (where the symbols \(\otimes\) denote tensor products in \(\bigotimes_{i}\mathrm{E}_{i}^{\mathrm{J}}\) ) to \(x_{1}\otimes\cdots\otimes x_{m}\) (where the symbols \(\otimes\) denote tensor products in \((\bigotimes_{i}\mathrm{E}_{i})^{\mathrm{J}}\) ). Hence \(f\) is an isomorphism of \(\bigotimes_{i}\mathrm{E}_{i}^{\mathrm{J}}\) onto \((\bigotimes_{i}\mathrm{E}_{i})^{\mathrm{J}}\) . We shall identify these two modules by means of this isomorphism.

Similarly, let E be an A-module. The map

\[
(x _ {1}, \dots , x _ {m}) \rightarrow x _ {1} \wedge \dots \wedge x _ {m}
\]

of \((\mathrm{E}^{\mathfrak{J}})^{m}\) into \((\wedge^{\mathfrak{m}}\mathrm{E})^{\mathfrak{J}}\) is obviously A-multilinear and alternating. It therefore defines an A-linear map \(f\) of \(\wedge^{\mathfrak{m}}\mathrm{E}^{\mathfrak{J}}\) into \((\wedge^{\mathfrak{m}}\mathrm{E})^{\mathfrak{J}}\) , which is evidently an isomorphism. We shall identify \(\wedge^{\mathfrak{m}}\mathrm{E}^{\mathfrak{J}}\) and \((\wedge^{\mathfrak{m}}\mathrm{E})^{\mathfrak{J}}\) by means of this isomorphism.

Let \(x'\) be an element of the dual \(\mathbf{E}^*\) of \(\mathbf{E}\) . The map \(x \to \langle x, x' \rangle^{\mathfrak{J}}\) ( \(x \in \mathbf{E}\) ) is an element \(x'^{\mathfrak{J}}\) of \((\mathrm{E}^{\mathfrak{J}})^*\) , and it is immediate that \(x' \to x'^{\mathfrak{J}}\) is a bijection \(g\) of \(\mathbf{E}^*\) onto \((\mathrm{E}^{\mathfrak{J}})^*\) satisfying \(g(ax') = a^{\mathfrak{J}}g(x')\) for all \(a \in \mathbf{A}\) . Consequently the composite map of \(g\) and of the identity map of \((\mathrm{E}^*)^{\mathfrak{J}}\) onto \(\mathbf{E}^*\) is an isomorphism of \((\mathrm{E}^*)^{\mathfrak{J}}\) onto \((\mathrm{E}^{\mathfrak{J}})^*\) . We shall identify these modules by means of this isomorphism, and we shall denote them \(\mathbf{E}_j^*\) .

Let \(E_{i}, F_{i} (i = 1, \ldots, m)\) be A-modules, and \(\Phi_{i} (i = 1, \ldots, m)\) a sesquilinear form for J on \(E_{i} \times F_{i}\) . The map

\[
(x _ {1}, \dots , x _ {m}, y _ {1}, \dots , y _ {m}) \rightarrow \Phi_ {1} (x _ {1}, y _ {1}) \Phi_ {2} (x _ {2}, y _ {2}) \dots \Phi_ {m} (x _ {m}, y _ {m})
\]

\((x_{i} \in \mathrm{E}_{i}, y_{i} \in \mathrm{F}_{i}, i = 1, \ldots, m)\) is an A-multilinear map from \(\mathrm{E}_{1} \times \cdots \times \mathrm{E}_{m} \times \mathrm{F}_{1}^{J} \times \cdots \times \mathrm{F}_{m}^{J}\) into A, and thus defines a bilinear form on \((\bigotimes_{i} \mathrm{E}_{i}) \times (\bigotimes_{i} \mathrm{F}_{i}^{J})\) (chap. III, § 1, no. 7). Since the second factor has been identified with \((\bigotimes_{i} \mathrm{F}_{i})^{J}\) , we have therefore defined a sesquilinear form \(\Phi\) for J on \((\bigotimes_{i} \mathrm{E}_{i}) \times (\bigotimes_{i} \mathrm{F}_{i})\) . It is characterized by

\[
\Phi (x _ {1} \otimes \dots \otimes x _ {m}, y _ {1} \otimes \dots \otimes y _ {m}) = \prod_ {i = 1} ^ {m} \Phi_ {i} (x _ {i}, y _ {i}) \quad (x _ {i} \in \mathrm{E} _ {i}, y _ {i} \in \mathrm{F} _ {i}).\tag{35}
\]

DEFINITION 11. --- Given A-modules \(\mathbf{E}_{i}, \mathbf{F}_{i}\) ( \(i=1,\ldots,m\) ) and, for each \(i\) , a sesquilinear form \(\Phi_{i}\) for J on \(\mathbf{E}_{i} \times \mathbf{F}_{i}\) , the sesquilinear form \(\Phi\) for J on \((\bigotimes_{i}\mathbf{E}_{i}) \times (\bigotimes_{i}\mathbf{F}_{i})\) characterized by (35) is called the tensor product of the sesquilinear forms \(\Phi_{i}\) .

In the case where the \(E_{i}\) and the \(F_{i}\) are equal to a single module E, and where the \(\Phi_{i}\) are equal to a single form \(\Psi\) , one says that \(\Phi\) is the extension of \(\Psi\) to \(\otimes^{m}E\) .

With the notations of def. 11, let us study the maps associated with \(\Phi\) . From formula (24) (no. 6) and (35) we derive the relation

\[
\Phi (x _ {1} \otimes \dots \otimes x _ {m}, y _ {1} \otimes \dots \otimes y _ {m}) = \prod_ {i = 1} ^ {m} \left\langle x _ {i}, d _ {\Phi_ {i}} (y _ {i}) \right\rangle = \prod_ {i = 1} ^ {m} \left\langle y _ {i}, s _ {\Phi_ {i}} (x _ {i}) \right\rangle^ {j}.
\]

We therefore have:

\[
s _ {\Phi} = j _ {s} \circ (s _ {\Phi_ {1}} \otimes \dots \otimes s _ {\Phi_ {m}}), \quad d _ {\Phi} = j _ {d} \circ (d _ {\Phi_ {1}} \otimes \dots \otimes d _ {\Phi_ {m}})\tag{36}
\]

where \(j_s\) (resp. \(j_d\) ) denotes the canonical map of \(\bigotimes_i F_i^*\) into \((\bigotimes_i F_i)^*\) (resp. of \(\bigotimes_i E_i^*\) into \((\bigotimes_i E_i)^*\) ) (chap. III, § 1, nos. 4 and 7).

PROPOSITION 9. — Let A be a commutative field equipped with an automorphism J, E \(_{i}\) , F \(_{i}\) two finite-dimensional vector spaces over A, and Φ \(_{i}\) a sesquilinear form for J on E \(_{i}\) × F \(_{i}\) (1 ≤ i ≤ m). If the forms Φ \(_{i}\) are nondegenerate, then so is their tensor product Φ. In this case the inverse form \(\hat{\Phi}\) of \(\Phi\) is the tensor product of the inverse forms Φ \(_{i}\) .

Indeed, since A is a field, it follows from Props. 6 and 7 of Chap. III, § 1, no. 3 that a tensor product of injective (resp. surjective) linear maps of A-modules is an injective (resp. surjective) linear map. Since the \(s_{\Phi_i}\) are bijective by hypothesis (prop. 6, no. 6), hence so is their tensor product. On the other hand, the canonical map \(j_s\) of \(\bigotimes_i F_i^*\) into \((\bigotimes_i F_i)^*\) is bijective (chap. III, § 1, no. 5, prop. 11). Therefore, by virtue of (36), \(s_\Phi\) is bijective, and this establishes our first assertion (prop. 6, no. 6). Likewise \(d_\Phi\) is bijective.

In the second assertion we have implicitly identified \(\bigotimes_{i} \mathrm{F}_{i}^{*}\) with \((\bigotimes_{i} \mathrm{F}_{i})^{*}\) and \(\bigotimes_{i} \mathrm{E}_{i}^{*}\) with \((\bigotimes_{i} \mathrm{E}_{i})^{*}\) by means of the maps \(j_{s}\) and \(j_{d}\) , which are here isomorphisms. The inverse forms cited in the statement exist since the \(s_{\Phi_{i}}\) , the \(d_{\Phi_{i}}\) , \(s_{\Phi}\) and \(d_{\Phi}\) are bijective (no. 7). Then set \(x' = x_{1}' \otimes \cdots \otimes x_{m}'\) , \(y' = y_{1}' \otimes \cdots \otimes y_{m}'\) ( \(x_{i}' \in \mathrm{E}_{i}^{*}\) , \(y_{i}' \in \mathrm{F}_{i}^{*}\) , \(i = 1, \ldots, m\) ). By definition of the inverse forms, and in view of (36), we have

\[
\begin{array}{l} \widehat {\Phi} (j _ {s} (y ^ {\prime}), j _ {d} (x ^ {\prime})) = \Phi (s _ {\Phi_ {1}} ^ {- 1} (y _ {1} ^ {\prime}) \otimes \dots \otimes s _ {\Phi_ {m}} ^ {- 1} (y _ {m} ^ {\prime}), d _ {\Phi_ {1}} ^ {- 1} (x _ {1} ^ {\prime}) \otimes \dots \otimes d _ {\Phi_ {m}} ^ {- 1} (x _ {m} ^ {\prime})) \\ = \prod_ {i = 1} ^ {m} \Phi_ {i} (s _ {\Phi_ {i}} ^ {- 1} (y _ {i} ^ {\prime}), d _ {\Phi_ {i}} ^ {- 1} (x _ {i} ^ {\prime})) = \prod_ {i = 1} ^ {m} \widehat {\Phi} _ {i} (y _ {i} ^ {\prime}, x _ {i} ^ {\prime}), \end{array}
\]

whence our second assertion.

QED.

Let E and F be two modules over the commutative ring A, and \(\Phi\) a sesquilinear form for J on E × F. The map

\[
(x _ {1}, \dots , x _ {m}, y _ {1}, \dots , y _ {m}) \rightarrow \det (\Phi (x _ {i}, y _ {k})) \quad (x _ {i} \in \mathrm{E}, y _ {i} \in \mathrm{F}, i = 1, \dots , m)
\]

of \(\mathbf{E}^{m}\times(\mathbf{F}^{\mathrm{j}})^{m}\) into A is A-multilinear. It therefore defines a bilinear form \(\Phi'\) on \((\bigotimes^{\overline{m}}\mathbf{E})\times(\bigotimes^{\overline{m}}\mathbf{F}^{\mathrm{j}})\) characterized by

\[
\Phi^ {\prime} (x _ {1} \otimes \dots \otimes x _ {m}, y _ {1} \otimes \dots \otimes y _ {m}) = \det (\Phi (x _ {i}, y _ {k})).
\]

Since the first member is zero when \(x_{i}=x_{k}\) or \(y_{i}=y_{k}\) ( \(i\neq k\) ), \(\Phi^{\prime}\) defines, by passing to quotients, a bilinear form on \((\bigwedge^{m}\mathrm{E})\times(\bigwedge^{m}\mathrm{F}^{\mathrm{J}})\) or again, since \(\bigwedge^{m}\mathrm{F}^{\mathrm{J}}\) is identified with \((\bigwedge^{m}\mathrm{F})^{\mathrm{J}}\) , a form \(\Phi_{(m)}\) sesquilinear for J on \((\bigwedge^{m}\mathrm{E})\times(\bigwedge^{m}\mathrm{F})\) . It is characterized by

\[
\left\{ \begin{array}{l} \Phi_ {(m)} (x _ {1} \wedge \dots \wedge x _ {m}, y _ {1} \wedge \dots \wedge y _ {m}) = \det (\Phi (x _ {i}, y _ {k})) \\ (x _ {i} \in \mathrm{E}, y _ {i} \in \mathrm{F}, i = 1, \dots , m). \end{array} \right.\tag{37}
\]

DEFINITION 12. --- Given two A-modules E, F and a form \(\Phi\) sesquilinear for J on \(E \times F\) , the form \(\Phi_{(m)}\) sesquilinear for J on \((\bigwedge^{m} E) \times (\bigwedge^{m} F)\) characterized by (37) is called the extension of \(\Phi\) to the m-th exterior powers.

With the notation of def. 12, let us study the maps associated with \(\Phi_{(m)}\) . From formula (24) (no. 6) and (37) we derive the relations

\[
\begin{array}{r l} \Phi_ {(m)} (x _ {1} \wedge \dots \wedge x _ {m}, y _ {1} \wedge \dots \wedge y _ {m}) & = \det (\langle x _ {i}, d _ {\Phi} (y _ {k}) \rangle) \\ & = \det (\langle y _ {i}, s _ {\Phi} (x _ {k}) \rangle^ {j}). \end{array}
\]

We therefore have

\[
s _ {\Phi (m)} = k _ {s} \circ (\stackrel {m} {\wedge} s _ {\Phi}), \quad d _ {\Phi_ {m}} = k _ {d} \circ (\stackrel {m} {\wedge} d _ {\Phi}),\tag{38}
\]

where \(k_s\) (resp. \(k_d\) ) denotes the canonical map of \(\bigwedge^m F^*\) into \((\bigwedge^m F)^*\) (resp. of \(\bigwedge^m E^*\) into \((\bigwedge^m E)^*\) ) (cf. Chap. III, § 8, no. 2).

PROPOSITION 10. --- Let A be a commutative field equipped with an automorphism J, E and F two finite-dimensional vector spaces over A, and \(\Phi\) a sesquilinear form for J on E \(\times\) F. If \(\Phi\) is nondegenerate, then its extension \(\Phi_{(m)}\) to the m-th exterior powers is nondegenerate, and the inverse form of \(\Phi_{(m)}\) is the extension to the m-th exterior powers of the inverse form \(\hat{\Phi}\) of \(\Phi\) .

Indeed, as \(s_{\Phi}\) and \(d_{\Phi}\) are bijective by hypothesis (prop. 6, no. 6), and the same holds for their exterior powers (chap. III, § 5, no. 7). On the other hand, the canonical maps \(k_s\) and \(k_d\) are bijective (chap. III, § 8, no. 2, th. 1). Therefore, by virtue of (38), \(s_{\Phi(m)}\) and \(d_{\Phi(m)}\) are bijective, which proves that \(\Phi_{(m)}\) is nondegenerate (prop. 6, no. 6). In the second assertion we have implicitly identified \(\bigwedge^m F^*\) with \((\bigwedge^m F)^*\) and \(\bigwedge^m E^*\) with \((\bigwedge^m E)^*\) by means of the maps \(k_s\) and \(k_d\) , which are here isomorphisms (loc. cit.). The inverse forms considered in the statement exist since \(s_{\Phi}\) , \(d_{\Phi}\) , \(s_{\Phi(m)}\) and \(d_{\Phi(m)}\) are bijective (no. 7). Then set \(x' = x_1' \wedge \ldots \wedge x_m'\) and \(y' = y_1' \wedge \ldots \wedge y_m'\) ( \(x_i' \in E^*\) , \(y_i' \in F^*\) , \(i = 1, \ldots, m\) ). By definition of inverse forms (no. 7) and in view of (38), we have

\[
\begin{array}{r l} \hat {\Phi} _ {(m)} (k _ {s} (y ^ {\prime}), k _ {d} (x ^ {\prime})) & = \Phi_ {(m)} (s _ {\Phi} ^ {- 1} (y _ {1} ^ {\prime}) \wedge \dots \wedge s _ {\Phi} ^ {- 1} (y _ {m} ^ {\prime}), d _ {\Phi} ^ {- 1} (x _ {1} ^ {\prime}) \wedge \dots \wedge d _ {\Phi} ^ {- 1} (x _ {m} ^ {\prime})) \\ & = \det (\Phi (s _ {\Phi} ^ {- 1} (y _ {i} ^ {\prime}), d _ {\Phi} ^ {- 1} (x _ {k} ^ {\prime}))) = \det (\hat {\Phi} (y _ {i} ^ {\prime}, x _ {k} ^ {\prime})) \end{array}
\]

whence our second assertion.

Remark. --- Let E be a free A-module, and let θ be the canonical isomorphism from ∧E onto the submodule of antisymmetrized tensors of order m (chap. III, § 5, no. 6, prop. 6). Let Φ be a sesquilinear form on E, let Φ(m) be the extension of Φ to ∧E, and let Θ be the sesquilinear form on ∧E which is the inverse image with respect to θ of the extension of Φ to ⊗E. By the definition of θ and of the antisymmetrization of a tensor, and by (35), we have

\[
\Theta (x _ {1} \wedge \dots \wedge x _ {m}, y _ {1} \wedge \dots \wedge y _ {m}) = \sum_ {\sigma , \tau} \varepsilon_ {\sigma} \varepsilon_ {\tau} \Phi (x _ {\sigma (1)}, y _ {\tau (1)}) \dots \Phi (x _ {\sigma (m)}, y _ {\tau (m)})
\]

where \(\sigma\) and \(\tau\) range over the symmetric group \(\mathfrak{S}_{m}\) . By the formula for computing determinants and formula (37), this expression can be written

\[
\sum_ {\tau \in \mathfrak {S} _ {m}} \varepsilon_ {\tau} \det (\Phi (x _ {i}, y _ {\tau (k)})) = m! \det (\Phi (x _ {i}, y _ {k}));
\]

In other words, we have \(\Theta = m!\Phi_{(m)}\) .

\subsection{Matrix calculations.}

In this no., we propose to relax the matrix calculus introduced in Chap. II, § 6, and to apply it to translate certain results proved in this section.

I. --- Let I and K be two finite index sets, H a nonempty set, and \(M = (m_{ik})_{(i,k) \in \mathbb{I} \times \mathbb{K}}\) a matrix over H (chap. II, § 6, no. 1, def. 1).

The transpose of \(M\) is called, and denoted \({}^t M\) , the matrix \((m_{ki}')_{(k,i) \in \mathbb{K} \times \mathbf{I}}\) satisfying \(m_{ki}' = m_{ik} ((i, k) \in \mathbf{I} \times \mathbf{K})\) . We obviously have

\[
^ {t} (^ {t} M) = M.\tag{39}
\]

This generalizes the notion introduced in Chap. II, § 6, no. 6.

Suppose that H is a commutative group (written additively). The set of matrices over H with I and K as index sets admits a structure of commutative group, since it is the set of maps from I × K to H. This group is written additively.

Let H', H'' be two nonempty sets, H an abelian group (written additively), and \(f:(h',h'')\to h'h''\) a map of \(\mathbf{H}'\times\mathbf{H}''\) into H. Given two matrices

\[
M ^ {\prime} = (m _ {i k} ^ {\prime}) _ {(i, k) \in \mathrm{I} \times \mathrm{K}}, \quad M ^ {\prime \prime} = (m _ {k l} ^ {\prime \prime}) _ {(k, l) \in \mathrm{K} \times \mathrm{L}}
\]

on H' and H'' respectively, such that the set K of column indices of M' is equal to the set of row indices of M'' , one calls the product of M' and M'' (following f) and denotes it M'M'' the matrix

\[
M ^ {\prime}. M ^ {\prime \prime} = (\sum_ {k \in \mathbf {K}} m _ {i k} ^ {\prime} m _ {k l} ^ {\prime \prime}) _ {(i, l) \in \mathbf {I} \times \mathbf {L}}\tag{40}
\]

on H. This generalizes the notion introduced in chap. II, § 6, no. 4. If \(\mathbf{H}' = \mathbf{H}'' = \mathbf{H}\) and if H is a ring, the product \(M'M''\) will, unless expressly stated otherwise, be computed "in H", that is, according to the map \((x, y) \to xy\) . When \(\mathbf{H}'\) and \(\mathbf{H}''\) are commutative groups (written additively) and \(f\) is bilinear, one has

\[
\left\{ \begin{array}{l} (M ^ {\prime} + M _ {1} ^ {\prime}) M ^ {\prime \prime} = M ^ {\prime} M ^ {\prime \prime} + M _ {1} ^ {\prime} M ^ {\prime \prime}, \\ M ^ {\prime} (M ^ {\prime \prime} + M _ {1} ^ {\prime \prime}) = M ^ {\prime} M ^ {\prime \prime} + M ^ {\prime} M _ {1} ^ {\prime \prime}, \end{array} \right.\tag{41}
\]

where \(M'\) , \(M_1'\) are matrices over \(H'\) , \(M''\) , \(M_1''\) matrices over \(H''\) , and where the sums and products written are assumed to be defined. Let \(M'\) , \(M''\) be matrices over the sets \(H'\) , \(H''\) , and \(f^0\) the map of \(H'' \times H'\) into H defined by \((h'', h') \to h' h''\) ; then one has

\[
{ } ^ { t } ( M ^ { \prime } M ^ { \prime \prime } ) = { } ^ { t } M ^ { \prime \prime } . { } ^ { t } M ^ { \prime }\tag{42}
\]

where the product in the first (resp. second) member is computed according to \(f\) (resp. \(f^{0}\) ).

In the case where \(H' = H'' = H\) is a ring, one recovers formula (12) of chap. II, §6, no. 6.

Let A be a ring, J an antiautomorphism of A. For every matrix \(M = (m_{ik})\) over A, we shall denote by \(M^{J}\) the matrix \((m_{ik}^{J})\) . Let \(M_{1}\) , \(M_{2}\) be two matrices over A such that \(M_{1}M_{2}\) is defined. Since J is an isomorphism of A onto the opposite ring \(A^{0}\) , one has \((M_{1}M_{2})^{J} = M_{1}^{J}\) . \(M_{2}^{J}\) where the first (resp. second) member is computed in A (resp. \(A^{0}\) ). In view of (42) and (39), this gives

\[
(M _ {1} M _ {2}) ^ {\mathrm{J}} = ^ {t} (^ {t} M _ {2} ^ {\mathrm{J}}. ^ {t} M _ {1} ^ {\mathrm{J}})\tag{43}
\]

where both members are computed in A.

Let \(H_{1}\) , \(H_{2}\) , \(H_{3}\) , \(H_{12}\) , \(H_{23}\) and H be commutative groups (written additively), \(f_{12}: H_{1} \times H_{2} \to H_{12}\) , \(f_{23}: H_{2} \times H_{3} \to H_{23}\) , \(f_{3}: H_{12} \times H_{3} \to H\) , \(f_{1}: H_{1} \times H_{23} \to H\) maps, and let \(M_{1}\) , \(M_{2}\) , \(M_{3}\) be matrices over \(H_{1}\) , \(H_{2}\) , \(H_{3}\) respectively. If \(f_{3}(f_{12}(x_{1}, x_{2}), x_{3}) = f_{1}(x_{1}, f_{23}(x_{2}, x_{3}))\) for all \(x_{i} \in H_{i}\) (i = 1, 2, 3), then the products \((M_{1}M_{2})M_{3}\) and \(M_{1}(M_{2}M_{3})\) (computed according to \(f_{12}\) , \(f_{3}\) , \(f_{23}\) and \(f_{1}\) ), if they are defined, are equal; they will be denoted \(M_{1}M_{2}M_{3}\) . When \(H_{1} = H_{2} = H_{3} = H_{12} = H_{23} = H\) , that H is a ring, and that \(f_{12}\) , \(f_{23}\) , \(f_{3}\) , \(f_{1}\) are equal to the map \((x, y) \to xy\) , the preceding condition expresses the associativity of the latter, and is therefore satisfied. Analogous conventions will be made for products of more than three factors.

Let A, B be two rings, \(M = (m_{ik})_{(i,k) \in I \times K}\) and \(M' = (m_{ik}')_{(i,k) \in I \times K}\) be two matrices over an (A, B)-bimodule G (no. 1). If, for every one-row matrix \(L = (a_i)_{i \in I}\) with elements in A and every one-column matrix \(C = (b_k)_{k \in K}\) with elements in B, one has \(L.M.C = L.M'.C\) (the products being computed according to the maps defining the structure of (A, B)-bimodule of G), then the matrices \(M\) and \(M'\) are equal. Indeed, if one takes \(a_i = 1\) , \(a_s = 0\) for \(s \neq i\) , \(b_k = 1\) , \(b_t = 0\) for \(t \neq k\) , the matrices \(L.M.C\) and \(L.M'.C\) , which are scalar matrices, are respectively equal to \(m_{ik}\) and \(m_{ik}'\) .

\begin{enumerate}
\def\labelenumi{\Roman{enumi}.}
\setcounter{enumi}{1}
\tightlist
\item
  Consider a ring A and a (right or left) A-module E, admitting a finite basis \((e_{i})_{i\in I}\) . For every element x of E, the one-column matrix formed by the components \(x_{i}\) ( \(i\in I\) ) of x with respect to \((e_{i})\) is called the matrix of x with respect to the basis \((e_{i})\) , and is denoted \(M(x)\) or x (cf. chap. II, § 6, no. 4); in the calculations it will be convenient, in order to recall that the index i is a row index, to adjoin to it a column index capable of taking only one value, and to write \((x_{i0})\) for the matrix \(M(x)\) .
\end{enumerate}

Now consider two A-modules (left or right) E and F, having finite bases \((e_{i})_{i\in I}\) and \((f_{k})_{k\in K}\) respectively; let \((f_{k}^{*})\) be the basis of \(F^{*}\) dual to \((f_{k})\) . We shall define the matrix, with respect to these bases, of a map u from E to F in the following four cases:

(D) E and F are right modules, and u is A-linear;

(G) E and F are left modules, and u is A-linear;

(GD) E is a left module, F a right module, A is equipped with an antiautomorphism J, u is Z-linear and satisfies \(u(ax) = u(x)a^{J}\) ( \(a \in A, x \in E\) ) (in other words, u is an A-linear map from \(E^{J}\) to F (no. 2, def. 5)).

(DG) E is a right A-module, F a left A-module, A is equipped with an antiautomorphism J, u is Z-linear and satisfies \(u(xa) = a^{J}u(x)\) ( \(x \in E, a \in A\) ) (in other words, u is an A-linear map from \(E^{J}\) into F).

In each of these four cases, the matrix of the map u is, by definition, the matrix \((u_{ki})_{(k,i)\in\mathbb{K}\times\mathbb{I}}\) such that

\[
u _ {k i} = \left\langle u (e _ {i}), f _ {k} ^ {*} \right\rangle .\tag{44}
\]

This definition coincides, in case (D), with that given in Chap. II, § 6, no. 3. Under these conditions, the matrix \(M(u(x))\) of the image of an element \(x\) of E is given by the following formulas:

\[
M (u (x)) = M (u). M (x)\tag{45 D}
\]

\[
{ } ^ { t } M ( u ( x ) ) = { } ^ { t } M ( x ) . { } ^ { t } M ( u )\tag{45 G}
\]

\[
M (u (x)) = M (u). M (x) ^ {\mathrm{J}}\tag{45 GD}
\]

\[
{ } ^ { t } M ( u ( x ) ) = { } ^ { t } M ( x ) ^ { j } . { } ^ { t } M ( u ) .\tag{45 DG}
\]

Let us verify, for example, (45 DG), the other verifications being analogous and slightly easier. Set \(x = \sum_{i} e_i x_{io}, u(x) = \sum_{k} y_{ko} f_k\) ; we have \(u(x) = u(\sum_{i} e_i x_{io}) = \sum_{i} x_{io}^{\mathbf{J}} u(e_i) = \sum_{i,k} x_{io}^{\mathbf{J}} u_{ki} f_k\) ; whence \(y_{ko} = \sum_{i} x_{io}^{\mathbf{J}} u_{ki}\) ; in order to place the two indices \(i\) next to each other, consider the transposed matrices \(^t M(x) = (x_{oi}')\) where \(x_{oi}' = x_{io}\) , and \(^t M(u) = (u_{ik}')\) where \(u_{ik}' = u_{ki}\) ; we then have \(y_{ko} = \sum_{i} x_{oi}' u_{ik}'\) ; since the right-hand side is the element with index \(k\) of the one-row matrix \(^t M(x)^{\mathbf{J}}\) . \(^t M(u)\) , formula (45 DG) is verified.

Remarks. -- 1) When A is commutative, (45 G) reduces to (45 D), and (45 DG) to (45 GD), by means of the formula \(t(M'M'') = 'M'''.'M'\) (cf. (42)), where both sides are here computed in A. 2) Let E, F, G be three left modules having finite bases, and \(u: \mathrm{E} \to \mathrm{F}\) , \(\nu: \mathrm{F} \to \mathrm{G}\) A-linear maps. It follows from (45 G) that one has

\[
{ } ^ { t } M ( \nu \circ u ) = { } ^ { t } M ( u ) . { } ^ { t } M ( \nu ) .\tag{46}
\]

Indeed, we have, for any \(x \in E\) ,

\[
\begin{array}{r l} ^ {t} M (x). ^ {t} M (\nu \circ u) & = ^ {t} M (\nu (u (x))) = ^ {t} M (u (x)). ^ {t} M (\nu) \\ & = ^ {t} M (x). ^ {t} M (u). ^ {t} M (\nu), \end{array}
\]

whence (46).

Recall that, in the case of right modules, we have

\[
M (\nu \circ u) = M (\nu) M (u).
\]

\begin{enumerate}
\def\labelenumi{\Roman{enumi}.}
\setcounter{enumi}{2}
\tightlist
\item
  From now on, we denote by A a ring, by B a ring (resp. a ring equipped with an antiautomorphism J, for which we set \(J' = J^{-1}\) ), by E a left A-module having a finite basis \((e_i)_{i \in I}\) and by F a right (resp. left) B-module having a finite basis \((f_k)_{k \in K}\) . We denote by \((e_i^*)\) and \((f_k^*)\) the dual bases of \(E^*\) and \(F^*\) . Unless expressly stated otherwise, the matrices considered are taken with respect to these bases.
\end{enumerate}

Let G be an (A, B)-bimodule (no. 1), Φ a bilinear (resp. sesquilinear on the right for J) map of E × F into G, and \(R = (\Phi(e_i, f_k))\) the matrix of Φ. Then, for \(x \in E\) and \(y \in F\) , formula (6) of no. 1 (resp. (8) of no. 2) is written, by virtue of the above conventions,

\[
\Phi (x, y) = ^ {t} M (x). R. M (y) \quad (\text { resp. } \Phi (x, y) = ^ {t} M (x). R. M (y) ^ {\mathrm{J}}),\tag{47}
\]

where the products are computed according to the maps which define the structure of the (A, B)-bimodule G; in particular, if A = B = G (in which case Φ is a form), the products are computed in A.

Let E' be a left A-module having a finite basis \((e_{s}^{\prime})_{s\in S}\) , F' a right (resp. left) A-module having a finite basis \((f_{t}^{\prime})_{t\in T}\) , \(u:E\to E^{\prime}\) and \(\nu:F\to F^{\prime}\) A-linear maps, and \(\Phi^{\prime}\) a bilinear map (resp. sesquilinear on the right for J) from \(E^{\prime}\times F^{\prime}\) into G. Denote by \(\Phi\) the inverse image of \(\Phi^{\prime}\) (relative to u and \(\nu\) ), U, V, R, \(R^{\prime}\) the matrices of u, \(\nu\) , \(\Phi\) , \(\Phi^{\prime}\) with respect to the bases considered. One then has

\[
R = ^ {t} U. R ^ {\prime}. V \quad (\text { resp. } R = ^ {t} U. R ^ {\prime}. V ^ {J}),\tag{48}
\]

the products being computed as in (47). Indeed, whatever \(x \in \mathbf{E}\) and \(y \in \mathbf{F}\) , we have by definition \(\Phi(x, y) = \Phi'(u(x), v(y))\) , whence, by (47),

\[
\begin{array}{c} ^ {t} M (x). R. M (y) = ^ {t} M (u (x)). R ^ {\prime}. M (\nu (y)) \\ (\text {resp.} ^ {t} M (x). R. M (y) ^ {\mathrm{J}} = ^ {t} M (u (x)). R ^ {\prime}. M (\nu (y)) ^ {\mathrm{J}}); \end{array}
\]

by (45 G) and (45 D) (resp. (45 G)) and (43) one deduces

\[
\begin{array}{r l} ^ {t} M (x). R. M (y) & = ^ {t} M (x). ^ {t} U. R ^ {\prime}. V. M (y) \\ (\text {resp.} ^ {t} M (x). R. M (y) ^ {J} & = ^ {t} M (x). ^ {t} U. R ^ {\prime}. ^ {t} (^ {t} M (y). ^ {t} V) ^ {J} \\ & = ^ {t} M (x). ^ {t} U. R ^ {\prime}. V ^ {J}. M (y) ^ {J}); \end{array}
\]

this proves our assertion.

\begin{enumerate}
\def\labelenumi{\Roman{enumi}.}
\setcounter{enumi}{3}
\tightlist
\item
  --- We assume here that the rings A and B are equal, and we denote by \(\Phi\) a bilinear (resp. sesquilinear on the right for J) form on E × F, and by \(R\) its matrix. Let us compute the matrices of the maps \(s_{\Phi}\) and \(d_{\Phi}\) associated with \(\Phi\) , which we shall denote by \(s\) and \(d\) to simplify. Since one has \(\Phi(x, y) = \langle y, s(x) \rangle = \langle x, d(y) \rangle\) by (23), no. 6 (resp. \(\Phi(x, y) = \langle x, d(y) \rangle = \langle y, s(x) \rangle^{J}\) by (24), no. 6), one has \(\Phi(e_i, f_k) = \langle f_k, s(e_i) \rangle = \langle e_i, d(f_k) \rangle\) (resp. \(\Phi(e_i, f_k) = \langle e_i, d(f_k) \rangle = \langle f_k, s(e_i) \rangle^{J}\) ), whence, by (44) and since ( \(e_i\) ) is the dual basis of ( \(e_i^*\) ) and ( \(f_k\) ) the dual basis of ( \(f_k^*\) ):
\end{enumerate}

\[
M (d) = R,   M (s) = ^ {t} R \quad (\text { resp. }   M (d) = R,   M (s) = ^ {t} R ^ {J ^ {\prime}}).\tag{49}
\]

Remarks. --- 1) When A is a field, the linear maps s and d have the same rank. We see here that their matrices \(M(s)\) and \(M(d)\) have the same rank; indeed, a matrix over A and its transpose have the same rank (chap. II, § 6, no. 7, prop. 3) and, when \(\Phi\) is sesquilinear, the equality of the ranks of \(R\) over A and of \(^tR\) over A \(^0\) (ibid.) and the fact that J' is an isomorphism of A \(^0\) onto A, imply the equality of the ranks of \(R\) and of \(^tR^{J'}\) over A.

\begin{enumerate}
\def\labelenumi{\arabic{enumi})}
\setcounter{enumi}{1}
\tightlist
\item
  If M and N are right A-modules having finite bases \((m_{i})\) and \((n_{k})\) , \(\Phi\) a sesquilinear form on the left for J on \(M \times N\) (no. 6, Remark), s and d its associated maps, and \(R = (\Phi(m_{i}, n_{k}))\) its matrix, formulas (26) of no. 6 show that one has
\end{enumerate}

\[
M (d) = R ^ {J ^ {\prime}}, \quad M (s) = ^ {t} R.
\]

Suppose now that the maps s and d associated with \(\Phi\) are bijective and let us compute the matrix \(\hat{R}\) of the inverse form of \(\Phi\) (no. 7). When \(\Phi\) is bilinear, \(\Phi\) is the inverse image of \(\hat{\Phi}\) with respect to the linear maps \(s: \mathrm{E} \to \mathrm{F}^*\) and \(d: \mathrm{F} \to \mathrm{E}^*\) ; one therefore has, by virtue of (48) and (49), \(R = R\) . \(\hat{R}\) . \(R\) , whence, since \(R\) is invertible ( \(d\) being bijective), \(\hat{R} = R^{-1}\) . This formula extends to the case where \(\Phi\) is sesquilinear, for, if one considers \(\Phi\) as a bilinear form on \(\mathrm{E} \times \mathrm{F}^{\mathrm{j}}\) , and if one identifies \((\mathrm{F}^{\mathrm{j}})^{*}\) with \((\mathrm{F}^{*})^{\mathrm{j}}\) (cf. no. 9), the inverse form of this bilinear form coincides with \(\hat{\Phi}\) considered as a bilinear form on \((\mathrm{F}^{*})^{\mathrm{j}} \times \mathrm{E}^{*}\) . In both cases the matrix of the inverse form of \(\Phi\) is the inverse of the matrix of \(\Phi\) .

Finally, let E' be a left A-module, F' a right (resp. left) A-module, both admitting finite bases \((e_{s}^{\prime})\) and \((f_{t}^{\prime})\) ; let \(\Phi^{\prime}\) be a bilinear form (resp. sesquilinear form with respect to J) on \(E^{\prime} \times F^{\prime}\) , and let \(R^{\prime}\) be its matrix. Suppose \(s_{\Phi}\) and \(d_{\Phi}\) bijective. Let \(u : E \to E^{\prime}\) and \(\nu : F \to F^{\prime}\) be linear maps, \(u^{*} : F^{\prime} \to F\) and \(\upsilon^{*} : E^{\prime} \to E\) their adjoints (no. 8, prop. 7); denote by U, V, \(U^{*}\) , \(V^{*}\) the matrices of u, \(\upsilon\) , \(u^{*}\) , \(\upsilon^{*}\) with respect to the given bases. One then has

\[
U ^ {*} = R ^ {- 1}. ^ {t} U. R ^ {\prime}, \quad {} ^ {t} V ^ {*} = R ^ {\prime}. V. R ^ {- 1}\tag{50}
\]

\[
(\text { resp. } U ^ {* \mathrm{J}} = R ^ {- 1}. ^ {t} U. R ^ {\prime}, ^ {t} V ^ {*} = R ^ {\prime}. V ^ {\mathrm{J}}. R ^ {- 1}).
\]

Indeed, whatever \(x \in E\) and \(y \in F'\) , one has \(\Phi'(u(x), y) = \Phi(x, u^*(y))\) (no. 8, def. 10). Hence, when \(\Phi\) is bilinear, by virtue of (47), \(^t M(u(x)). R'. M(y) = ^t M(x). R. M(u^*(y))\) ; this gives, by virtue of (45 G) and (45 D), \(^t M(x). ^t U. R'. M(y) = ^t M(x). R. U^*. M(y)\) , whence \(^t U. R' = R. U^*\) and the first announced formula, since, \(d\) being bijective, \(R\) is invertible. When \(\Phi\) is sesquilinear (47) and (45 G) give \(^t M(x). ^t U. R'. M(y)^J = ^t M(x). R. (^{t}M(y). ^{t}U^*)^{J}\) ; now, according to (43), we have \((^{t}M(y). ^{t}U^*)^{J} = ^{t}(^{u}U^{*J}. ^{u}M(y)^{J})\) , whence \((^{t}M(y). ^{t}U^*)^{J} = U^{*J}. M(y)^{J}\) ; it follows therefore \(^t M(x). ^{t}U. R'. M(y)^{J} = ^{t}M(x). R. U^{*J}. M(y)^{J}\) , whence \(^t U. R' = R. U^{*J}\) , and \(U^{*J} = R^{-1}. ^{t}U. R'\) . The verification of the formulas for \(V^*\) is analogous.

Exercises. --- 1) Let A be a commutative field, E a vector space over A admitting an infinite countable basis \((e_n)_{n \geqslant 1}\) . We define a bilinear form \(\Phi\) on E by setting \(\Phi(e_{i+1}, e_i) = 1\) for \(i \geqslant 1\) , \(\Phi(e_k, e_j) = 0\) for \(k \neq j + 1\) and \(j \geqslant 1\) . Show that the linear map \(d_\Phi\) associated on the right with \(\Phi\) is injective, but the linear map \(s_\Phi\) associated on the left with \(\Phi\) is not injective.

\begin{enumerate}
\def\labelenumi{\arabic{enumi})}
\setcounter{enumi}{1}
\item
  Let E be the Z-module direct sum of Z and Z/(2), and let E* be its dual (isomorphic to Z). Show that the bilinear form \((x, x') \to \langle x, x' \rangle\) on \(E \times E^*\) is such that the associated linear map on the right is injective, but the associated linear map on the left is not.
\item
  Give an example of a bilinear form \(\Phi\) defined on a product \(\mathbf{E} \times \mathbf{F}\) of two vector spaces, such that \(d_{\Phi}\) is bijective, \(s_{\Phi}\) injective but not bijective (take \(\mathbf{E}\) of infinite dimension and F equal to the dual \(\mathbf{E}^*\) of \(\mathbf{E}\) ; cf. chap. II, § 5, exerc. 3).
\item
  Let A be a ring equipped with an antiautomorphism J, let E be a left A-module, and let G be an (A, A)-bimodule, and \(\Phi\) a map from \(\mathrm{E} \times \mathrm{E}\) into G, sesquilinear on the right for J. Prove the identity (where \(\mathrm{Q}(x) = \Phi(x, x)\) ):
\end{enumerate}

\[
\begin{array}{r l} & {2 \Phi (x, y) (\mu^ {\mathrm{J}} \lambda^ {\mathrm{J}} - \lambda^ {\mathrm{J}} \mu^ {\mathrm{J}}) = \mathrm{Q} (x - \mu \lambda y) - \mathrm{Q} (x + \mu \lambda y) + \mu \mathrm{Q} (x + \lambda y)} \\ & {- \mu \mathrm{Q} (x - \lambda y) + \mathrm{Q} (x + \mu y) \lambda^ {\mathrm{J}} - \mathrm{Q} (x - \mu y) \lambda^ {\mathrm{J}} + \mu \mathrm{Q} (x - y) \lambda^ {\mathrm{J}} - \mu \mathrm{Q} (x + y) \lambda^ {\mathrm{J}}.} \end{array}
\]

\begin{enumerate}
\def\labelenumi{\arabic{enumi})}
\setcounter{enumi}{4}
\tightlist
\item
  Let K be a commutative field of characteristic 2, A a separable quadratic extension of K; we have \(\mathrm{A} = \mathrm{K}(\theta)\) , where \(\theta\) is a root of an irreducible polynomial \(\mathrm{X}^2 +\mathrm{X} + \beta\) of K{[}X{]} and the K-automorphism J of A, distinct from the identity, is such that \(\theta^{j} = \theta +1\) (chap. V, § 11, exerc. 8). Show that if E and G are vector spaces over A, \(\Phi\) a sesquilinear map (for J) from \(\mathrm{E}\times \mathrm{E}\) into G, we have, setting \(\mathrm{Q}(x) = \Phi (x,x)\) ,
\end{enumerate}

\[
\Phi (x, y) = \mathrm{Q} (\theta x + y) - \beta \mathrm{Q} (x) - \mathrm{Q} (y) - (\theta + 1) (\mathrm{Q} (x + y) - \mathrm{Q} (x) - \mathrm{Q} (y)).
\]

\begin{enumerate}
\def\labelenumi{\arabic{enumi})}
\setcounter{enumi}{5}
\tightlist
\item
  Let A be a field, E a vector space over A, \(\Phi\) a sesquilinear form on E, \(u\) an endomorphism of E.
\end{enumerate}

\begin{enumerate}
\def\labelenumi{\alph{enumi})}
\item
  For there to exist one and only one endomorphism \(u^*\) of E such that \(\Phi(u(x), y) = \Phi(x, u^*(y))\) for \(x, y\) in E, it is necessary and sufficient that \(d_\Phi\) be injective and that \(^t u(d_\Phi(E)) \subset d_\Phi(E)\) .
\item
  Give an example where E is infinite-dimensional and \(d_{\Phi}\) injective, but where \(u(d_{\Phi}(E))\) is not contained in \(d_{\Phi}(E)\) .
\end{enumerate}

\begin{enumerate}
\def\labelenumi{\arabic{enumi})}
\setcounter{enumi}{6}
\tightlist
\item
  Let E, \(E_{1}\) be two A-modules, \(\Phi\) (resp. \(\Phi_{1}\) ) a sesquilinear form on E (resp. \(E_{1}\) ). Assume that \(\Phi_{1}\) is nondegenerate and that there exists an element \(\alpha\in A\) and a bijection u from E onto \(E_{1}\) such that \(\Phi_{1}(u(x), u(y))=\Phi(x,y)\alpha\) for all x, y in E. Show that: \(1^{\circ}\Phi\) is nondegenerate; \(2^{\circ}u\) is linear; \(3^{\circ}\) if \(E_{1}\) is a faithful A-module, so is E, and \(\alpha\) is not a right zero divisor in A; \(4^{\circ}\) if \(\Phi_{1}\) takes values in A that are not left zero divisors, so is \(\Phi\) .
\end{enumerate}

¶ 8) Let A be a field, E₁, E₂ two nonzero vector spaces over A, Φ₁ (resp. Φ₂) a nondegenerate sesquilinear form on E₁ (resp. E₂) for an antiautomorphism J₁ (resp. J₂) of A. Let u be a linear map from E₁ onto E₂ such that the relation Φ₁(x, y) = 0 implies Φ₂(u(x), u(y)) = 0.

\begin{enumerate}
\def\labelenumi{\alph{enumi})}
\item
  Show that \(u\) is a bijection from \(\mathbf{E}_{1}\) onto \(\mathbf{E}_{2}\) . (If \(u(0)\) were not reduced to 0, show that there would exist in \(\mathbf{E}_{1}\) two vectors \(a\) , \(b\) such that \(u(a) \neq 0\) , \(u(b) = 0\) and \(\Phi_{1}(a, b) \neq 0\) ; if H is the hyperplane of \(x \in \mathbf{E}_{1}\) such that \(\Phi_{1}(a, x) = 0\) , note that one would have \(u(\mathrm{H}) = \mathrm{E}_{2}\) ).
\item
  Show that if dim \(E_{1} \geqslant 2\) , there exists \(\alpha \in A\) such that one has \(\Phi_{2}(u(x), u(y)) = \Phi_{1}(x, y) \alpha\) for all x, y in \(E_{1}\) . (For every \(y \in E_{1}\) , show that there exists an element \(m(y) \in A\) such that \(\Phi_{2}(u(x), u(y)) = \Phi_{1}(x, y)m(y)\) for all \(x \in E_{1}\) , and that if y and \(y'\) are linearly independent in \(E_{1}\) , one has \(m(y + y') = m(y) = m(y')\) ).
\end{enumerate}

\begin{enumerate}
\def\labelenumi{\arabic{enumi})}
\setcounter{enumi}{8}
\tightlist
\item
  Let A be a field, E, F be two left vector spaces over A, \(\Phi\) a nondegenerate sesquilinear form on \(E \times F\) for an antiautomorphism J of A.
\end{enumerate}

\begin{enumerate}
\def\labelenumi{\alph{enumi})}
\item
  Let M be a subspace of E, and N a subspace of F such that \(N \supset M^{0}\) and \(M \supset N^{0}\) . Show that if one of the spaces \(N/M^{0}\) , \(M/N^{0}\) is finite-dimensional, then so is the other, and the dimensions of these two spaces are equal.
\item
  Let M, \(M'\) be two subspaces of E such that \(M^{00} = M\) and \(M'\) be of finite dimension; show that we have \((M \cap M')^{0} = M^{0} + M'^{0}\) and \((M + M')^{00} = M + M'\) . (Applying a) to the subspaces \(M'\) and \(M^{0} + M'^{0}\) , show that \(\dim(M \cap M') = \text{codim}(M^{0} + M'^{0})\) ; applying a) to the subspaces \(M + M'\) and \(M^{0}\) , show that
\end{enumerate}

\[
\dim ((M + M ^ {\prime}) ^ {0 0} / M) = \dim ((M + M ^ {\prime}) / M).
\]

\begin{enumerate}
\def\labelenumi{\alph{enumi})}
\setcounter{enumi}{2}
\item
  If E = F and if M is a subspace of E such that \(E = M^{0} + M^{00}\) , show that E is the direct sum of \(M^{0}\) and \(M^{00}\) .
\item
  Let E be a vector space over a commutative field A admitting an infinite countable basis \((e_{n})_{n\geqslant0}\) , and let \(\Phi\) be the symmetric bilinear form on E such that \(\Phi(e_{n}, e_{n}) = 1\) for all n, \(\Phi(e_{i}, e_{j}) = 0\) for \(i \geqslant 1\) , \(j \geqslant 1\) and \(i \neq j\) , and \(\Phi(e_{0}, e_{n}) = 1\) for all \(n \geqslant 1\) . Show that \(\Phi\) is nondegenerate. Let M (resp. N) be the subspace of E spanned by the \(e_{2k}\) (resp. \(e_{2k-1}\) ) for \(k \geqslant 1\) , and let H = M + N, which is a hyperplane in E. Show that we have \(M^{0} = N\) , \(N^{0} = M\) , \(H^{00} = E \neq H\) , \((M \cap N)^{0} \neq M^{0} + N^{0}\) and \((M + N)^{00} \neq M + N\) , although \(M^{00} = M\) , \(N^{00} = N\) ; if L is the 2-dimensional subspace spanned by \(e_{0}\) and \(e_{1}\) , one has \((L \cap H)^{0} \neq L^{0} + H^{0}\) .
\end{enumerate}

¶ 10) Let E, E' be two left vector spaces over fields A, A' respectively, of dimension ≥3; let ℱ(E) (resp. ℱ(E')) be the lattice-ordered set (for the inclusion relation) formed by the finite-dimensional subspaces of E (resp. E').

\begin{enumerate}
\def\labelenumi{\alph{enumi})}
\item
  Let \(p\) be a map from \(\mathfrak{F}(\mathrm{E})\) into \(\mathfrak{F}(\mathrm{E}')\) such that for every \(\mathbf{M} \in \mathfrak{F}(\mathrm{E})\) , \(\dim p(\mathbf{M}) = \dim \mathbf{M}\) , and such that for every pair \((\mathbf{M}, \mathbf{N})\) of elements of \(\mathfrak{F}(\mathrm{E})\) , \(p(\mathbf{M} + \mathbf{N}) = p(\mathbf{M}) + p(\mathbf{N})\) . Show that \(p\) is injective; if \(p\) is bijective, there exists a bijective semilinear map \(u\) of \(\mathrm{E}\) into \(\mathrm{E}'\) such that one has \(u(\mathbf{M}) = p(\mathbf{M})\) for all \(\mathbf{M} \in \mathfrak{F}(\mathrm{E})\) (use exercise 10 of chap. II, 2nd ed., App. III).
\item
  Give an example where \(A' = A\) is commutative, \(E' = E\) is of finite dimension, and where there exists a map p from \(\mathfrak{F}(E)\) into itself, such that \(\dim p(M) = \dim M\) , \(p(M + N) = p(M) + p(N)\) , \(p(M \cap N) = p(M) \cap p(N)\) , but there exists no injective semilinear map u from E into itself such that \(u(M) = p(M)\) for \(M \in \mathfrak{F}(E)\) . (Consider the case where there exists an overfield \(A''\) of A of finite degree and isomorphic to A, for example \(\mathbf{A} = \mathbf{F}_{p}(\mathbf{X})\) , where \(p\) is prime; \(\mathrm{E}'' = \mathrm{A}'' \otimes_{\mathbf{A}} \mathrm{E}\) is then a vector space of the same dimension over \(\mathrm{A}''\) as \(\mathrm{E}\) over \(\mathrm{A}\) ; consider the map \(\mathrm{M} \to \mathrm{A}'' \otimes_{\mathbf{A}} \mathrm{M}\) .
\item
  We assume \(A' = A\) . Let \(\omega\) be a bijective map from \(\mathfrak{F}(E)\) onto the set \(\mathfrak{F}'(E')\) of subspaces of finite codimension of \(E'\) , such that \(\text{codim } \omega(M) = \dim M\) , and \(\omega(M + N) = \omega(M) \cap \omega(N)\) . Show that there exists a sesquilinear form \(\Phi\) on \(E \times E'\) , nondegenerate on the right, and such that \(\omega(M) = M^0\) for all \(M \in \mathfrak{F}(E)\) . (Use th. 1 of chap. II, § 4, no. 6).
\end{enumerate}

¶ 11) a) Let A be a ring artinian on the left and on the right (chap. VIII, § 2, no. 3). Show that the following conditions are equivalent: 1° the right annihilator (resp. left annihilator) of every left ideal (resp. right ideal) ≠ A is not reduced to 0; 2° the dual of every simple left A-module (resp. right A-module) is not reduced to 0; 3° the dual of every finitely generated left A-module (resp. right A-module) is not reduced to 0. A ring A is said to satisfy condition (N \(_{s}\) ) (resp. (N \(_{a}\) )) if it satisfies these conditions for left A-modules (resp. right A-modules).

\begin{enumerate}
\def\labelenumi{\alph{enumi})}
\setcounter{enumi}{1}
\item
  Let A be a ring artinian on the left and on the right satisfying condition (N \(_{a}\) ), E (resp. F) a free left A-module (resp. right A-module) having a countable basis over A, Φ a bilinear form on E × F such that s \(_{\Phi}\) is injective. Let M be a free submodule of E; show that there exists a free submodule N of F and a basis (e \(_{n}\) ) (resp. (f \(_{n}\) )) of M (resp. N) such that Φ(e \(_{i}\) , f \(_{j}\) ) = δ \(_{ij}\) ; moreover one can take for e \(_{1}\) any free element of M. (Note that if x is a free element of M, the image of F under the map y → Φ(x, y) is the whole ring A; then proceed by induction to construct the bases (e \(_{n}\) ) and (f \(_{n}\) )). Deduce that if the basis (e \(_{n}\) ) of M is finite, E is the direct sum of M and N \(^{0}\) , and that one has M \(^{00}\) = M.
\item
  We keep the hypotheses of \(b\) , and suppose in addition that A satisfies condition \((\mathrm{N}_{s})\) and that \(d_{\Phi}\) is injective. Show then that there exists a basis \((e_{n})\) in E and a basis \((f_{n})\) in F such that \(\Phi(e_{i}, f_{j}) = \delta_{ij}\) . (Use \(b\) ) by determining alternately by induction \(e_{n}\) and \(f_{n}\) ).
\end{enumerate}

¶ 12) Let E be a real Hilbert space of countable type, \(\Phi(x, y)\) the scalar product in E. Show that there does not exist a system of two algebraic bases \((e_{\lambda})\) , \((f_{\mu})\) of the vector space E over R such that we have \(\Phi(e_{\lambda}, f_{\mu}) = \delta_{\lambda y}\) for every pair of indices. (First note that the index set of these bases would have the cardinality of the continuum (Topological vector spaces, chap. II, § 3, exerc. 15); then consider an orthonormal (countable) basis \((a_n)\) of E and note that the subspace spanned by the \(a_n\) is contained in the subspace generated by a countable subfamily of \((e_{\lambda})\) .
% semantic-fragments: legacy-input-seam
\relax{}
\section{Discriminant of a sesquilinear form}

Throughout this paragraph, A denotes a commutative ring, J an automorphism of A, and E a free A-module of finite dimension n.

DEFINITION 1. --- Given a form \(\Phi\) sesquilinear for J on E and a system \(S = (x_1, \ldots, x_n)\) of n elements of E, one calls the discriminant of \(\Phi\) with respect to this system, and denotes it by \(D_{\Phi}(x_1, \ldots, x_n)\) or \(D_{\Phi}(S)\) the element det \((\Phi(x_i, x_j))\) of A.

If \((e_{1},\ldots,e_{n})\) is a basis of E, the discriminant of \(\Phi\) with respect to this basis is none other than the determinant of the matrix of \(\Phi\) with respect to this basis.

It follows from the definition of the extension of \(\Phi\) to \(\wedge\) E ( \(\S\) 1, no. 9) that we have

\[
\mathrm{D} _ {\Phi} (x _ {1}, \dots , x _ {n}) = \Phi_ {(n)} (x _ {1} \wedge \dots \wedge x _ {n}, x _ {1} \wedge \dots \wedge x _ {n}),\tag{1}
\]

where \(\Phi_{(n)}\) denotes the extension of \(\Phi\) to \(\wedge\) E. For every permutation \(\sigma \in \mathfrak{S}_n\) , we therefore have

\[
\mathrm{D} _ {\Phi} (x _ {\sigma (\mathbf {1})}, \dots , x _ {\sigma (n)}) = \mathrm{D} _ {\Phi} (x _ {\mathbf {1}}, x _ {\mathbf {2}}, \dots , x _ {n}).
\]

Example. --- Let B be an algebra over the ring A, such that B is a free A-module of finite dimension n. Then the map \((x, y) \to \mathrm{Tr}_{\mathrm{B}/\mathrm{A}}(xy)\) (chap.~VIII, § 12, no. 2) is a bilinear form on B. Given a system \((x_{1}, \ldots, x_{n})\) of n elements of B, the discriminant of this form with respect to this system is called the discriminant of the system \((x_{1}, \ldots, x_{n})\) over A, and is denoted \(\mathrm{D}_{\mathrm{B}/\mathrm{A}}(x_{1}, \ldots, x_{n})\) . We thus have

\[
\mathrm{D} _ {\mathrm{B/A}} (x _ {1}, \dots , x _ {n}) = \det (\operatorname{Tr} _ {\mathrm{B/A}} (x _ {i} x _ {j})).\tag{2}
\]

Remark. --- Let \((e_{1},\ldots,e_{n})\) be a basis of B over A, and \(e_{i}e_{j}=\sum_{k=1}^{n}c_{ijk}e_{k}\;(c_{ijk}\in\mathrm{A})\) the multiplication table of this basis (chap.~II, § 7, no. 2). Since the matrix of the A-linear endomorphism \(x\to e_{k}x\) of B with respect to \((e_{r})\) is \((c_{ksr})\) , one has \(\mathrm{Tr}_{\mathrm{B}/\mathrm{A}}(e_{k})=\sum_{r=1}^{n}c_{krr}\) , whence \(\mathrm{Tr}_{\mathrm{B}/\mathrm{A}}(e_{i}e_{j})=\sum_{k,r}c_{ijk}c_{krr}\) . It follows that one has

\[
\mathrm{D} _ {\mathbf {B} / \mathbf {A}} (e _ {1}, \dots , e _ {n}) = \det _ {i, j} (\sum_ {k, r} c _ {i j k} c _ {k r r}).\tag{3}
\]

PROPOSITION 1. — Let \(\Phi\) be a sesquilinear form for J on E, \((x_1,\ldots,x_n)\) be a system of \(n\) elements of E, and \((a_{ij})_{(i,j=1,\ldots,n)}\) a family of \(n^2\) elements of A; set \(y_j = \sum_{i=1}^{n} a_{ji} x_i\) . We then have

\[
\mathrm{D} _ {\Phi} (y _ {1}, \dots , y _ {n}) = \det (a _ {i j}). \det (a _ {i j}) ^ {\mathrm{J}}. \mathrm{D} _ {\Phi} (x _ {1}, \dots , x _ {n}).
\]

Indeed, as \(\Phi\) is a sesquilinear form, we have \(\Phi(y_i, y_j) = \sum_{k,m} a_{ik} \Phi(x_k, x_m) a_{jm}^{\mathfrak{J}}\) . Therefore, if we denote by \(A\) the matrix \((a_{ij})\) , the matrix \((\Phi(y_i, y_j))\) is equal to \(A \cdot (\Phi(x_i, x_j)) \cdot {}^t A^{\mathfrak{J}}\) . Since \(\det(^t A) = \det(A)\) and \(\det(A^{\mathfrak{J}}) = \det(A)^{\mathfrak{J}}\) , our assertion is proved.

In particular, if \((e_{i})\) and \((e_{i}^{\prime})\) are two bases of E, D and D' the discriminants of \(\Phi\) with respect to these bases, and a the determinant of the change-of-basis matrix from the basis \((e_{i})\) to the basis \((e_{i}^{\prime})\) , one has

\[
\mathrm{D} ^ {\prime} = a a ^ {\mathrm{J}} \mathrm{D}.\tag{4}
\]

It follows from prop. 1 that, if \((e_i)\) is a basis of E and \((x_i)\) any arbitrary system of \(n\) elements of E, \(\mathrm{D}_{\Phi}(e_1,\ldots ,e_n)\) divides \(\mathrm{D}_{\Phi}(x_1,\dots,x_n)\) . In particular, the discriminants of \(\Phi\) with respect to any two bases of E generate the same principal ideal of A.

Let \((\mathbf{E}_i)_{i\in I}\) be a finite family of free A-modules of finite dimensions, \(\Phi_i\) a sesquilinear form for J on \(\mathbf{E}_i\) , and \(\mathbf{B}_i\) a basis of \(\mathbf{E}_i\) . If \(\Phi\) denotes the direct sum of \(\Phi_i\) ( \(\S 1\) , no. 3) and B the basis of \(\prod E_i\) obtained by taking the union of \(\mathbf{B}_i\) , we obviously have

\[
\mathrm{D} _ {\Phi} (\mathrm{B}) = \prod_ {i \in \mathrm{I}} \mathrm{D} _ {\Phi_ {i}} (\mathrm{B} _ {i}).\tag{5}
\]

Let \(\Phi\) be a sesquilinear form for \(J\) on \(E\) , \(h\) a homomorphism of \(A\) into a commutative ring \(A'\) , \(\Phi'\) the sesquilinear form on \(A' \otimes_A E\) obtained by extension of \(\Phi\) ( \(\S 1\) , no. 4) and \((x_1, \ldots, x_n)\) an arbitrary system of elements of \(E\) . Since \(A' \otimes_A E\) is an \(A'\) -free module, \(D_{\Phi'}(1 \otimes x_1, \ldots, 1 \otimes x_n)\) is defined, and we obviously have

\[
\mathrm{D} _ {\Phi^ {\prime}} (1 \otimes x _ {1}, \dots , 1 \otimes x _ {n}) = h (\mathrm{D} _ {\Phi} (x _ {1}, \dots , x _ {n})).\tag{6}
\]

Example. --- Let B be an algebra over A which is a free A-module of finite dimension \(n\) , \((x_{1},\ldots,x_{n})\) a basis of B over A, and m an ideal of A. If we denote by \(h\) the canonical homomorphism from B onto B/mB, \((h(x_{1}),\ldots,h(x_{n}))\) is a basis of B/mB over A/m (chap.~I, § 6, no. 5, prop. 5), and B/mB is isomorphic to \((A/m)\otimes_{A}B\) . We therefore have

\[
\mathrm{D} _ {(\mathrm{B} / \mathfrak {m B}) / (\mathrm{A} / \mathfrak {m})} (h (x _ {1}), \dots , h (x _ {n})) = h (\mathrm{D} _ {\mathrm{B} / \mathrm{A}} (x _ {1}, \dots , x _ {n})).
\]

PROPOSITION 2. --- Assume that A is an integral domain. Let \(\Phi\) be a sesquilinear form for J on E and \((e_1, \ldots, e_n)\) a basis of E, such that \(\mathrm{D}_{\Phi}(e_1, e_2, \ldots, e_n) \neq 0\) .

\begin{enumerate}
\def\labelenumi{\alph{enumi})}
\item
  For a system \((x_{1},\ldots ,x_{n})\) of \(n\) elements of E to be free, it is necessary and sufficient that \(\mathrm{D}_{\Phi}(x_1,\dots,x_n)\) be \(\neq 0\) .
\item
  For a system \((x_{1},\ldots ,x_{n})\) of \(n\) elements of E to be a basis of E, it is necessary and sufficient that \(\mathrm{D}_{\Phi}(x_1,\dots,x_n)\) and \(\mathrm{D}_{\Phi}(e_1,\dots,e_n)\) be associated elements in A (cf.~Chap.~VI, § 1, no.~5).
\end{enumerate}

Let us set \(x_j = \sum_{i=1}^{n} a_{ji} e_i\) ( \(a_{ji} \in A\) ). Let us first prove \(a\) ). If \(D_{\Phi}(x_1, \ldots, x_n) = 0\) , one has \(\det(a_{ji})\) . \(\det(a_{ji})^{\mathfrak{J}} = 0\) (prop. 1) since \(D_{\Phi}(e_1, \ldots, e_n) \neq 0\) and A is an integral domain; we therefore have \(\det(a_{ji}) = 0\) , and the vectors \(x_j\) are linearly dependent (chap.~III, § 7, no. 1, th. 1, applied to the vector space \(K \otimes_A E\) , where K denotes the field of fractions of A). Conversely, if these vectors are linearly dependent, we have \(\det(a_{ji}) = 0\) ( \(ibid.\) ), whence \(D_{\Phi}(x_1, \ldots, x_n) = 0\) (prop. 1).

We now prove \(b\) ). If \(\mathrm{D}_{\Phi}(x_1, \ldots, x_n)\) and \(\mathrm{D}_{\Phi}(e_1, \ldots, e_n)\) are associated in A, prop. 1 shows that \(\det(a_{ij}) \cdot \det(a_{ij})^{\mathrm{J}}\) is invertible in A. Thus \(\det(a_{ij})\) is also invertible in A; therefore the matrix \((a_{ij})\) on A is invertible (chap.~III, § 6, no. 5, th. 2), and the endomorphism \(g\) of E defined by \(g(e_i) = x_i\) ( \(i = 1, \ldots, n\) ) is an automorphism; consequently \((x_1, \ldots, x_n)\) is a basis of E. The converse follows immediately from Prop. 1.

PROPOSITION 3. --- Let \(\Phi\) be a sesquilinear form for J on E, and S a basis of E. The following conditions are equivalent:

\begin{enumerate}
\def\labelenumi{\alph{enumi})}
\item
  The map \(s_{\Phi}\) from E into \(E^{*}\) associated with \(\Phi\) is bijective.
\item
  The map \(d_{\Phi}\) from E into \(E^{*}\) associated with \(\Phi\) is bijective.
\item
  The element \(D_{\Phi}(S)\) is invertible in A.
\end{enumerate}

Indeed the condition \(c\) expresses that the matrix of \(\Phi\) with respect to S is invertible (chap. III, § 6, no. 5, th. 2). Therefore \(c\) is equivalent to \(a\) ( \(§ 1\) , no. 10); likewise \(c\) is equivalent to \(b\) ).

PROPOSITION 4. --- Assume A is an integral domain. Let S be a basis of E. A necessary and sufficient condition for a sesquilinear form \(\Phi\) on E to be nondegenerate is that one has D \(_{\Phi}\) (S) ≠ 0.

Indeed, let K be the field of fractions of A, and let \(\Phi'\) be the extension of \(\Phi\) to the K-vector space \(K\otimes_{\mathbf{A}}E\) ; we identify E with a subset of this vector space. The relation \(D_{\Phi}(S)\neq0\) is then equivalent to \(D_{\Phi'}(S)\neq0\) by formula (6), which itself expresses that \(s_{\Phi'}\) is bijective (prop. 3), that is, \(\Phi'\) is nondegenerate ( \(\S1\) , no. 6, prop. 6). Now, for every \(x\in K\otimes_{\mathbf{A}}E\) , there exists \(a\in A\) such that \(ax\in E\) . Therefore, for \(\Phi\) to be degenerate, it is necessary and sufficient that \(\Phi'\) be so. This proves our assertion.

PROPOSITION 5. --- Let A be a field, B a commutative algebra of finite dimension n over A, and S a basis of B. For B to be separable (chap. VIII, § 7, no. 5, def. 1) it is necessary and sufficient that one have D \(_{B/A}\) (S) ≠ 0.

Let A' be the algebraic closure of A, and let B' be the algebra A' ⊗\_A B over A'. If B is separable, B' is semisimple (Chap. VIII, § 7, no. 5, cor. of Prop. 7) and is therefore a direct product of n fields isomorphic to A' (Chap. VIII, § 6, no. 4, cor. of Prop. 9). If S' denotes the canonical basis of B' (identified with A'	extsuperscript{n}), we have D\_\{B'/A'\}(S') = 1, hence D\_\{B'/A'\}(S) ≠ 0 (prop. 1) and D\_\{B/A\}(S) ≠ 0 (formula (6)).

Conversely, suppose we have \(D_{B/A}(S) \neq 0\) . To show that B is separable, it suffices to show that \(B'\) is semisimple, that is, that it admits no nilpotent element \(\neq 0\) . Now, if \(x'\) were a nonzero nilpotent element of \(B'\) , one could take it as the first element of a basis \(S'\) of \(B'\) , and one would then have \(\mathrm{Tr}_{\mathbf{B}'/\mathbf{A}'}(x'y') = 0\) for all \(y' \in S'\) since a nilpotent endomorphism has all its eigenvalues zero (chap.~VII, § 5, no. 3, cor. 3 of prop. 8), hence zero trace. It would then follow that \(\mathrm{D}_{\mathbf{B}'/\mathbf{A}'}(\mathrm{S}') = 0\) , whence \(\mathrm{D}_{\mathbf{B}'/\mathbf{A}'}(\mathrm{S}) = 0\) (prop. 1) and \(\mathrm{D}_{\mathbf{B}/\mathbf{A}}(\mathrm{S}) = 0\) (formula (6)), contrary to the hypothesis.

Remark. --- Suppose that B is an overfield of A. Let \(S = (x_{1}, \ldots, x_{n})\) be a basis of B, and \((s_{1}, \ldots, s_{n})\) the A-isomorphisms of B into the algebraic closure A' of A (each \(s_{j}\) being repeated \([B : A]_{i}\) times). Recall that, for every \(z \in B\) , one has \(\operatorname{Tr}_{\mathbf{B}/\mathbf{A}}(z) = \sum_{j=1}^{n} s_{j}(z)\) (chap.~VIII, § 12, no. 2, prop. 4).

It then follows from the product formula for determinants that one has

\[
(\det (s _ {j} (x _ {i}))) ^ {2} = \mathrm{D} _ {\mathrm{B/A}} (x _ {1}, \dots , x _ {n}).\tag{7}
\]

This formula shows that proposition 5 generalizes the separability condition given in chap.~V, § 7, no. 2, Remark.

PROPOSITION 6. --- Let \(\Phi\) be an A-bilinear form on E, and K a subring of A such that A is a free K-module of finite dimension q. If \((e_i)_{i=1,\ldots,n}\) is a basis of E over A and \((a_j)_{j=1,\ldots,q}\) a basis of A over K, then \((a_j e_i)\) is a basis of E over K. The map \(\Phi'\) from E \(\times\) E to K defined by \(\Phi'(x,y)=\mathrm{Tr}_{\mathrm{A/K}}(\Phi(x,y))\) is a K-bilinear form on E, and one has

\[
\mathrm{D} _ {\Phi^ {\prime}} (a _ {j} e _ {i}) = \mathrm{N} _ {\mathrm{A} / \mathrm{K}} (\mathrm{D} _ {\Phi} (e _ {1}, \dots , e _ {n})) \cdot (\mathrm{D} _ {\mathrm{A} / \mathrm{K}} (a _ {1}, \dots , a _ {q})) ^ {n}.\tag{8}
\]

The first two assertions being evident, it suffices to prove (7). By definition, the left-hand side is the determinant of the K-linear endomorphism u of E defined by

\[
u (a _ {j} e _ {i}) = \sum_ {r, s} \mathrm{Tr} _ {\mathbf {A} / \mathbf {K}} (a _ {j} a _ {r} \Phi (e _ {i}, e _ {s})) a _ {r} e _ {s}.
\]

Consider the A-linear endomorphism \(\nu\) of E defined by \(\nu(e_i) = \sum_s \Phi(e_i, e_s)e_s\) , and the K-linear endomorphism \(w\) of E defined by \(w(a_j e_i) = (\sum_r \mathrm{Tr}_{\mathrm{A/K}}(a_j a_r)a_r)e_i\) . We have

\[
w \left(v \left(a _ {j} e _ {i}\right)\right) = w \left(\sum_ {s} a _ {j} \Phi \left(e _ {i}, e _ {s}\right) e _ {s}\right) = \sum_ {\tau , s} \operatorname{Tr} _ {\mathrm{A} / \mathrm{K}} \left(a _ {j} \Phi \left(e _ {i}, e _ {s}\right) a _ {\tau}\right) a _ {r} e _ {s}
\]

since \(w(ae_s) = \sum_r \mathrm{Tr}_{\Lambda/\mathbb{K}}(aa_r)a_r e_s\) for all \(a \in A\) ; thus \(u\) is the composite map \(w \circ v\) . Thus, denoting by \(v_{\kappa}\) the map \(v\) considered as a K-linear map, we have \(\det(u) = \det(v_{\kappa})\det(w)\) . Now we have \(\det(v_{\kappa}) = N_{\Lambda/\mathbb{K}}(\det(v))\) (chap.~VIII, § 12, no. 2, prop. 7), and it is clear that \(\det(v) = D_{\Phi}(e_1, \ldots, e_n)\) . On the other hand, since each of the A-modules \(Ae_i\) ( \(i = 1, \ldots, n\) ) is stable under \(w\) and the determinant of the restriction of \(w\) to \(Ae_i\) is \(\det(\mathrm{Tr}_{\Lambda/\mathbb{K}}(a_j a_r)) = D_{\Lambda/\mathbb{K}}(a_1, \ldots, a_q)\) , one has \(\det(w) = (D_{\Lambda/\mathbb{K}}(a_1, \ldots, a_q))^n\) . Formula (8) therefore reduces to \(\det(u) = \det(v_{\kappa})\det(w)\) , a formula proved above.

COROLLARY ("Transitivity formula for discriminants"). — Let K be a commutative ring, A a commutative algebra admitting a finite basis \((a_{j})_{j=1,\ldots,q}\) over K, and E an algebra over A admitting a finite basis \((e_{i})_{i=1,\ldots,n}\) . Then \((a_{j}e_{i})\) is a basis of E over K, and one has

\[
\mathrm{D} _ {\mathrm{E} / \mathrm{K}} (a _ {j} e _ {i}) = \mathrm{N} _ {\mathrm{A} / \mathrm{K}} (\mathrm{D} _ {\mathrm{E} / \mathrm{A}} (e _ {1}, \dots , e _ {n})) \cdot (\mathrm{D} _ {\mathrm{A} / \mathrm{K}} (a _ {1}, \dots , a _ {q})) ^ {n}.
\]

Indeed, if one sets \(\Phi(x,y)=\mathrm{Tr}_{\mathrm{E/A}}(xy)\) , the K-bilinear form \(\Phi'\) of Prop. 6 is \(\Phi'(x,y)=\mathrm{Tr}_{\mathrm{E/K}}(xy)\) by the transitivity formula for traces (Chap. VIII, § 12, no. 2, cor. of Prop. 7).

Exercises. — 1) Let A be an algebra of finite rank over a commutative field K, having an identity element.

\begin{enumerate}
\def\labelenumi{\alph{enumi})}
\item
  Show that if the radical of A is not zero, the bilinear form \((x, y) \to \mathrm{Tr}_{\mathbf{A}/\mathbf{K}}(xy)\) on A is degenerate.
\item
  Assume that K has characteristic 0. Show that if A is a matrix algebra \(\mathbf{M}_{n}(\mathrm{K})\) , S the canonical basis of A over K, we have \(\mathrm{D}_{\mathrm{A}/\mathrm{K}}(\mathrm{S}) \neq 0\) .
\item
  From a) and b), deduce that, for an algebra A of finite rank over a field K of characteristic 0 to be absolutely semisimple, it is necessary and sufficient that the bilinear form \((x,y)\to\mathrm{Tr}_{\mathbf{A}/\mathbf{K}}(xy)\) be nondegenerate (or, equivalently, that \(\mathrm{D}_{\mathbf{A}/\mathbf{K}}(\mathrm{S})\neq0\) for every basis S of A over K).
\end{enumerate}

Let B be a ring, A a subring of B containing the identity element of B; B is therefore an (A, A)-bimodule; we denote by \(^s\text{B}\) (resp. \(^d\text{B}\) ) the set B considered as a left (resp. right) A-module, by \(^s\text{B}^*\) (resp. \(^d\text{B}^*\) ) the right (resp. left) dual A-module of \(^s\text{B}\) (resp. \(^d\text{B}\) ). For every \(x' \in ^s\text{B}^*\) and every \(b \in \text{B}\) , \(x \to \langle xb, x' \rangle\) is an A-linear form on \(^s\text{B}\) , hence an element of \(^s\text{B}^*\) denoted by \(bx'\) ; the map \((b, x') \to bx'\) defines on \(^s\text{B}^*\) a left B-module structure (cf.~Chap.~III, 2nd ed., App. II, no. 7).

\[
(\mathrm{A}, \mathrm{A})
\]

\[
(\mathrm{A}, \mathrm{A}) -
\]

\[
\Phi : (x, y) \rightarrow \varphi (x y)
\]

\(^{s}\) B × \(^{d}\) B into A be nondegenerate, it is necessary and sufficient that \(\varphi(0)\) contain no ideal (left or right) of B other than \(\{0\}\) . We then say that \(\varphi\) is a Frobenius homomorphism from B to A.

\begin{enumerate}
\def\labelenumi{\alph{enumi})}
\setcounter{enumi}{1}
\item
  Let \(\varphi\) a Frobenius homomorphism from B to A; show that the map \(d_{\Phi}\) right-associated to \(\Phi\) is an isomorphism of the left B-module \(B_s\) onto a submodule of the left B-module \(^sB^*\) . Show that \(d_{\Phi}\) is bijective in each of the following two cases: \(1^\circ\) A is a left and right artinian ring satisfying the conditions ( \(N_s\) ) and ( \(N_d\) ) ( \(\S 1\) , exerc. 11), and \(^sB\) and \(^dB\) are free A-modules of finite length (use exerc. 11 b) of \(\S 1\) ); \(2^\circ\) A is a commutative and involutive Artinian ring (chap.~VIII, \(\S 3\) , exerc. 11), contained in the center of B, and \(^{s}\) B is an A-module of finite length (use exerc. 11 of chap.~VIII, § 3).
\item
  Conversely, show that if B \(_{s}\) and \(^{s}\) B \(^{*}\) are isomorphic, there exists a Frobenius homomorphism from B to A when one is in one of the two cases considered in b) and that \(^{s}\) B and \(^{d}\) B have the same length (use exerc. 11 b) of § 1).
\item
  With the hypotheses and notations of a), show that the right annihilator (resp. left annihilator) of a left ideal I (resp. of a right ideal r) of B is the orthogonal l⁰ (resp. r⁰) for the form Φ of the submodule I (resp. r) of sB (resp. dB).
\item
  Let \(\varphi\) a Frobenius homomorphism from B to A. Show that if A is an involutive Artinian ring (chap.~VIII, § 3, exerc. 11) then so is B, when one moreover supposes that one or the other of the conditions of \(b\) ) is satisfied (use \(b\) ) and \(d\) ).
\end{enumerate}

¶ 3) Let A be a commutative field, B an algebra of finite rank over A, having an identity element.

\begin{enumerate}
\def\labelenumi{\alph{enumi})}
\item
  For B to be a Frobenius algebra, it is necessary and sufficient that there exist a Frobenius homomorphism from B to A (exerc. 2). (Use exerc. 2 c) and e) above, and exerc. 6 b) of chap.~VIII, § 13).
\item
  Let \(\varphi\) a Frobenius homomorphism from B to A; every A-linear form on B can then be written in a unique way \(x \to \varphi(b'x)\) (resp. \(x \to \varphi(xb'')\) ) where \(b'\) and \(b''\) belong to B; for this form to be a Frobenius homomorphism, it is necessary and sufficient that \(b'\) (resp. \(b''\) ) be invertible in B.
\item
  For every \(x \in B\) , let \(x^{\sigma}\) be the unique element (cf.~b)) such that \(\varphi(xy) = \varphi(yx^{\sigma})\) for all \(y \in B\) . Show that \(x \to x^{\sigma}\) is an A-automorphism of B. One says that the Frobenius algebra B is symmetric if \(\sigma\) is an inner automorphism of B; there is then a Frobenius homomorphism from B to A for which \(\sigma\) is the identity (cf.~b)). This is equivalent to saying that the (B, B)-bimodules B and \(^{s}B^{*} = ^{d}B^{*}\) (written B*) are isomorphic (exerc. 2 c)).
\item
  Let E be a left B-module of finite length, E' its dual, E'* the dual of E' considered as a vector space over A; E'* is endowed with a structure of left B-module by setting, for \(x' \in \mathrm{E}', x'' \in \mathrm{E}*\) , \(b \in \mathrm{B}\) , \(\langle x', bx'' \rangle = \langle x'b, x'' \rangle\) (chap.~III, 2nd ed., App. II, no. 7). For every \(x \in \mathrm{E}\) , let \(f_{\mathrm{E}}(x)\) (or simply \(f(x)\) ) be the element of E'* such that \(\langle x', f(x) \rangle = \varphi(\langle x, x' \rangle)\) for all \(x' \in \mathrm{E}'\) ; show that \(f\) is a semilinear bijection for the automorphism \(\sigma\) , from the left B-module E onto the left B-module E'* (use exerc. 10 of chap.~VIII, § 4). For \(\mathrm{E} = \mathrm{B}_s\) , one has (with the notations of exerc. 2 b)) \(d_{\Phi}(x^{\sigma}) = f_{\mathrm{B}_s}(x)\) for all \(x \in \mathrm{B}\) .
\end{enumerate}

\begin{enumerate}
\def\labelenumi{\arabic{enumi})}
\setcounter{enumi}{3}
\tightlist
\item
  \begin{enumerate}
  \def\labelenumii{\alph{enumii})}
  \tightlist
  \item
    Let G be a finite group. Show that the algebra B of the group G over any commutative field A is a symmetric Frobenius algebra (exerc. 3). (Consider the map \(\varphi\) from B to A which, to every element \(x = \sum_{s \in G} \xi_s.s\) , associates \(\varphi(x) = \xi_e\) , \(e\) denoting the identity element of G).
  \end{enumerate}
\end{enumerate}

\begin{enumerate}
\def\labelenumi{\alph{enumi})}
\setcounter{enumi}{1}
\tightlist
\item
  Let E be a vector space of finite dimension n over a commutative field A, and B the exterior algebra \(\wedge\) E. Show that B is a Frobenius algebra. (If \((e_i)_{1 \leqslant i \leqslant n}\) is a basis of E, consider the map which to every element \(x\) of B associates the coefficient of \(e_1 \wedge e_2 \wedge \ldots \wedge e_n\) in the expression of \(x\) using the basis of B corresponding to \((e_i)\) ). For the Frobenius algebra B to be symmetric, it is necessary and sufficient that \(n\) be even or A have characteristic 2.
\end{enumerate}

¶ 5) a) Show that the tensor product of two Frobenius algebras (resp. Frobenius and symmetric algebras) of finite rank over a commutative field K is a Frobenius algebra (resp. Frobenius and symmetric algebra) (cf.~§ 1, no. 9, prop. 9).

\begin{enumerate}
\def\labelenumi{\alph{enumi})}
\setcounter{enumi}{1}
\tightlist
\item
  Let B be an algebra of finite rank over K, L an extension of K of finite degree over K. Show that if the algebra \(\mathrm{B}_{(\mathrm{L})} = \mathrm{B} \otimes_{\mathrm{K}} \mathrm{L}\) over L is Frobenius (resp. Frobenius and symmetric), then so is B. (Use exerc. 2 c) and 3 c) above and exerc. 2 of chap.~VIII, § 2).
\end{enumerate}

\begin{enumerate}
\def\labelenumi{\arabic{enumi})}
\setcounter{enumi}{5}
\tightlist
\item
  Show that every absolutely semisimple algebra B of finite rank over a commutative field A is a symmetric Frobenius algebra. (Reduce to the case where B is simple; use exerc. 5 b) above, as well as prop. 9 of chap.~VIII, § 12, no. 3).
\end{enumerate}

\section{Hermitian forms and quadratic forms}

In all that follows in this Chapter, unless expressly stated otherwise, A denotes a ring and E a left A-module. One supposes A endowed with an involutive antiautomorphism J, denoted \(\alpha \to \overline{\alpha}\) ; one therefore has \((\overline{\alpha + \beta}) = \overline{\alpha} + \overline{\beta}\) , \((\overline{\alpha\beta}) = \overline{\beta}.\overline{\alpha}\) and \(\overline{\overline{\alpha}} = \alpha\) for all \(\alpha\) , \(\beta\) in A. Unless expressly stated otherwise, the sesquilinear forms considered are right sesquilinear ( \(\S 1\) , no. 2, def. 4) for this antiautomorphism.

\subsection{\texorpdfstring{Hermitian and \(\varepsilon\)-hermitian forms.}{Hermitian and epsilon-hermitian forms.}}

DEFINITION 1. --- Let \(\varepsilon\) be an element of the center of A. A sesquilinear form \(\Phi\) on E such that one has \(\Phi(x, y) = \varepsilon \overline{\Phi(y, x)}\) for all \(x\) and \(y\) in E is called a form \(\varepsilon\) -hermitian. A 1-hermitian form (resp. \((-1)\) -hermitian) is said to be hermitian (resp. antihermitian).

When J is the identity (which implies that A is commutative), a Hermitian (resp. antihermitian) form (for J) is none other than a symmetric (resp. antisymmetric) bilinear form (chap.~III, § 5, no. 1, def. 2). Recall that an alternating bilinear form (chap.~III, § 5, no. 2, def. 4) is antisymmetric; the converse is true if, in A, the relation \(2a = 0\) implies \(a = 0\) .

The orthogonality relation (§ 1, no. 3) with respect to an ε-hermitian form is evidently symmetric (cf.~exerc. 1).

If \(\alpha\) is an invertible element of A, the map T: \(\lambda \to \alpha^{-1}\overline{\lambda}\alpha\) is an antiautomorphism of A, and one readily verifies that the form \(\Phi\alpha\) is sesquilinear with respect to T. If, moreover, we have \(\alpha = \overline{\alpha}\) then T is involutive, and, if \(\Phi\) is \(\varepsilon\) -hermitian, so is \(\Phi\alpha\) ; indeed we have

\[
(\lambda^ {\mathrm{T}}) ^ {\mathrm{T}} = \alpha^ {- 1} (\overline {{\alpha^ {- 1} \lambda \alpha}}) \alpha = \alpha^ {- 1} \overline {{\alpha}} \lambda \overline {{\alpha}} ^ {- 1} \alpha = \lambda
\]

\[
\Phi (y, x) \alpha = \varepsilon \overline {{\Phi (x , y)}} \alpha = \varepsilon (\Phi (x, y) \alpha) ^ {\mathrm{T}}.
\]

In particular, when A is a field, the elements \(\alpha\) of the center of A such that \(\bar{\alpha} = \alpha\) form a subfield K of A, and the \(\varepsilon\) -hermitian forms on E (for J) form a vector space over K.

Remarks. --- 1) If \(\Phi\) is a form \(\varepsilon\) -hermitian on E, we have \(\Phi(x, y) = \varepsilon \Phi(x, y) \bar{\varepsilon}\) for all \(x, y\) in E. Therefore, if \(\Phi\) takes invertible values, we have \(\varepsilon \bar{\varepsilon} = 1\) .

\begin{enumerate}
\def\labelenumi{\arabic{enumi})}
\setcounter{enumi}{1}
\tightlist
\item
  If there exists an invertible element i of the center of A such that \(\bar{i}=i\varepsilon\) , then, for \(\Phi\) to be \(\varepsilon\) -hermitian, it is necessary and sufficient that \(i\Phi\) be Hermitian.
\end{enumerate}

The map \((y, x) \to \overline{\Phi(x, y)}\) being sesquilinear for J, for \(\Phi\) to be \(\varepsilon\) -hermitian, it is necessary and sufficient that one has \(\Phi(y, x) = \varepsilon \overline{\Phi(x, y)}\) when \(x\) and \(y\) range over a system of generators of E. In particular, if E admits a finite basis \((e_i)_{1 \leqslant i \leqslant n}\) for a sesquilinear form \(\Phi\) on E to be \(\varepsilon\) -hermitian, it is necessary and sufficient that its matrix \(R = (\rho_{ij}) = (\Phi(e_i, e_j))\) satisfies the relations \(\rho_{ji} = \varepsilon \overline{\rho_{ij}}\) for all \(i, j\) , that is to say \(^t R = \varepsilon \overline{R}\) ; a matrix \(R\) possessing this property is said to be \(\varepsilon\) -Hermitian. When \(\varepsilon = 1\) (resp. -1) one says that \(R\) is Hermitian (resp. antihermitian) with respect to the anti-automorphism J. When J is the identity (hence A is commutative), a Hermitian (resp. antihermitian) matrix \(R\) is such that \(^t R = R\) (resp. \(^t R = -R\) ); we then say that \(R\) is a symmetric (resp. antisymmetric) matrix. For \(\Phi\) to be an alternating form, it is necessary and sufficient that its matrix be antisymmetric and, in addition, that the diagonal terms of R all be zero; a matrix possessing these properties is said to be alternating.

Let \(\Phi\) a sesquilinear form on E, and let \(s_{\Phi}\) and \(d_{\Phi}\) be the maps from E into E* associated with \(\Phi\) to the left and to the right (§ 1, no. 6). For \(\Phi\) to be \(\varepsilon\) -hermitian, it is necessary and sufficient that \(\langle x, s_{\Phi}(y) \rangle = \bar{\varepsilon} \langle x, d_{\Phi}(y) \rangle\) for all elements \(x, y\) of E, hence that \(s_{\Phi} = \bar{\varepsilon} d_{\Phi}\) , or equivalently that \(\langle x, d_{\Phi}(y) \rangle = \varepsilon \langle x, s_{\Phi}(y) \rangle\) , hence that \(d_{\Phi} = \varepsilon s_{\Phi}\) .

Let \(\Phi\) be a form \(\varepsilon\) -hermitian such that the map \(d_{\Phi}\) from E into E* associated on the right to \(\Phi\) is bijective. For every endomorphism \(u\) of E we then have

\[
u ^ {* *} = \varepsilon \bar {\varepsilon} u.\tag{1}
\]

Indeed, whatever the elements \(x\) and \(y\) of \(\mathbf{E}\) , one has

\[
\begin{array}{r l} \Phi (x, u ^ {* *} (y)) & = \Phi (u ^ {*} (x), y) = \varepsilon \overline {{\Phi (y , u ^ {*} (x))}} = \varepsilon \overline {{\Phi (u (y) , x)}} \\ & = \varepsilon \Phi (x, u (y)) \bar {\varepsilon} = \Phi (x, \varepsilon \bar {\varepsilon} u (y)) \end{array}
\]

therefore \(u^{**}(x) = \varepsilon \bar{\varepsilon} u(x)\) since \(\Phi\) is nondegenerate.

If \(\Phi\) is a form \(\varepsilon\) -hermitian such that the maps \(s_{\Phi}\) and \(d_{\Phi}\) are bijective, then the inverse form \(\hat{\Phi}\) of \(\Phi\) ( \(\S 1\) , no. 7) is a \(\bar{\varepsilon}\) -hermitian form. Indeed, setting \(s = s_{\Phi}\) , \(d = d_{\Phi}\) for brevity, one deduces from \(d = \varepsilon s\) that \(s^{-1} = \bar{\varepsilon} d^{-1}\) , \(s\) being semilinear. Consequently, whatever \(u\) , \(\nu\) in E, one has

\[
\hat {\Phi} (u, \nu) = \Phi (s ^ {- 1} (u), d ^ {- 1} (\nu)) = \bar {\varepsilon} \Phi (d ^ {- 1} (u), d ^ {- 1} (\nu)),
\]

whence

\[
\widehat {\Phi} (\nu , u) = \overline {{\varepsilon \varepsilon}} \overline {{\Phi (d ^ {- 1} (u) , d ^ {- 1} (\nu))}} = \overline {{\varepsilon}} \widehat {\Phi} (u, \nu),
\]

since ε is in the center of A.

Finally, when the ring A is commutative, the canonical extensions of a \(\varepsilon\) -hermitian form \(\Phi\) to the tensor and exterior powers \(\bigotimes^{p}\mathbf{E}\) and \(\bigwedge^{p}\mathbf{E}\) of E are \(\varepsilon^{p}\) -hermitian forms, as follows immediately from formulas (35) and (37) of § 1, no. 9.

\subsection{Modules over a quadratic extension.}

Let K be a commutative ring. One takes for A the quadratic extension \(\mathrm{A} = \mathrm{K}(i)\) with \(i^{2} = -1\) , and for J the automorphism \(\lambda + i\mu \to \lambda - i\mu\) ( \(\lambda \in \mathbf{K}\) , \(\mu \in \mathbf{K}\) ) (chap.~II, § 7, no. 7). If E is an A-module, we denote by \(\mathbf{E}_0\) the K-module deduced from E by restriction of the ring of scalars, and by \(j\) the automorphism \(x \to ix\) of \(\mathbf{E}_0\) ; one evidently has \(j^2 = -I\) , where I is the identity map of \(\mathbf{E}_0\) . Conversely, let \(\mathbf{E}_0\) be a K-module and let \(j\) be an automorphism of \(\mathbf{E}_0\) such that \(j^2 = -I\) ; the map \(\lambda + i\mu \to \lambda I + \mu j\) is evidently a homomorphism from A into the ring \(\mathfrak{L}(\mathbf{E}_0)\) of endomorphisms of \(\mathbf{E}_0\) ; we have therefore defined on \(\mathbf{E}_0\) a structure of A-module, for which we have

\[
(\lambda + i \mu) x = \lambda x + \mu j (x) \quad (x \in \mathrm{E} _ {0}, \lambda \in \mathrm{K}, \mu \in \mathrm{K}).\tag{2}
\]

If E' is another A-module, \(E_{0}'\) the underlying K-module of \(E'\) , \(j'\) the automorphism \(x' \rightarrow ix'\) of \(E_{0}'\) then the A-linear maps f from E into \(E'\) are none other than the K-linear maps from \(E_{0}\) into \(E_{0}'\) such that \(f \circ j = j' \circ f\) . In particular, if one denotes by \(E^{*}\) and \((\mathrm{E}_{0})^{*}\) the respective duals of E and \(E_{0}\) , and if \(f_{1}\) and \(f_{2}\) are two maps from E into K, for the map \(x \to f_{1}(x) + if_{2}(x)\) from E into A to be A-linear, it is necessary and sufficient that \(f_{1}\) and \(f_{2}\) be in \((\mathrm{E}_{0})^{*}\) and that one has \(f_{1} \circ j + i(f_{2} \circ j) = if_{1} - f_{2}\) , that is to say \(f_{1} = f_{2} \circ j\) and \(f_{1} \circ j = -f_{2}\) . Since j is an automorphism of \(E_{0}\) and \(j^{2} = -I\) , these two conditions are equivalent. Eliminating \(f_{1}\) or \(f_{2}\) , we see that the formulas

\begin{enumerate}
\def\labelenumi{(\arabic{enumi})}
\setcounter{enumi}{2}
\tightlist
\item
\end{enumerate}

\[
f (x) = f _ {1} (x) - i f _ {1} (j (x))\tag{4}
\]

\[
f (x) = f _ {2} (j (x)) + i f _ {2} (x)
\]

\((x \in \mathbf{E}, f \in \mathbf{E}^{*}, f_{1} \in (\mathbf{E}_{0})^{*}, f_{2} \in (\mathbf{E}_{0})^{*})\) establish two bijective correspondences between \(\mathbf{E}^{*}\) and \((\mathbf{E}_{0})^{*}\) .

\subsection{Bilinear forms associated with a Hermitian form.}

We assume again here that the ring A is the quadratic extension \(\mathrm{A} = \mathrm{K}(i)\) (where \(i^{2} = -1\) ) of a commutative ring K, and that J is the automorphism \(\lambda + i\mu \to \lambda - i\mu\) of A ( \(\lambda \in K, \mu \in K\) ). Let E and \(E'\) be two A-modules, \(E_{0}\) and \(E_{0}'\) be the underlying K-modules of E and \(E'\) , j and \(j'\) the automorphisms \(x \to ix\) and \(x' \to ix'\) of E and \(E'\) (cf.~no. 2). A K-bilinear form f on \(E_{0} \times E_{0}'\) will be said to be invariant under j and \(j'\) if one has

\[
f (j (x), j ^ {\prime} (x ^ {\prime})) = f (x, x ^ {\prime})\tag{5}
\]

for \(x \in \mathrm{E}_0\) and \(x' \in \mathrm{E}_0'\) . Replacing \(x\) by \(j(x)\) , one sees that this condition is equivalent to

\[
f (x, j ^ {\prime} (x ^ {\prime})) = - f (j (x), x ^ {\prime})\tag{6}
\]

for all \(x \in \mathbf{E}_0\) and \(x' \in \mathbf{E}_0'\) .

PROPOSITION 1. --- Let \(\Phi_{1}\) (resp. \(\Phi_{2}\) ) be a K-bilinear form on \(\mathbf{E}_{0} \times \mathbf{E}_{0}^{\prime}\) , invariant under \(j\) and \(j^{\prime}\) . The map which associates to \(\Phi_{1}\) (resp. \(\Phi_{2}\) ) the map \(\Phi\) of \(\mathbf{E} \times \mathbf{E}^{\prime}\) into A defined by

\begin{enumerate}
\def\labelenumi{(\arabic{enumi})}
\setcounter{enumi}{6}
\tightlist
\item
\end{enumerate}

\[
\Phi (x, x ^ {\prime}) = \Phi_ {1} (x, x ^ {\prime}) + i \Phi_ {1} (x, j ^ {\prime} (x ^ {\prime}))\tag{8}
\]

\[
(\text { resp. } \Phi (x, x ^ {\prime}) = - \Phi_ {2} (x, j ^ {\prime} (x ^ {\prime})) + i \Phi_ {2} (x, x ^ {\prime}))
\]

\((x \in \mathrm{E}, x' \in \mathrm{E})\) , is an isomorphism of the K-vector space of K-bilinear forms on \(\mathrm{E}_0 \times \mathrm{E}_0'\) invariant under \(j\) and \(j'\) onto the K-vector space of sesquilinear forms on \(\mathrm{E} \times \mathrm{E}'\) . Suppose moreover \(\mathrm{E} = \mathrm{E}'\) ; for \(\Phi\) to be Hermitian, it is necessary and sufficient that \(\Phi_1\) be symmetric (resp. that \(\Phi_2\) be antisymmetric) (cf.~exerc. 4).

Indeed every map \(\Phi\) of E × E' into A can be written, in one and only one way, in the form \(\Phi = \Phi_1 + i\Phi_2\) , where \(\Phi_1\) and \(\Phi_2\) are maps of E × E' into K. For the partial map \(x \to \Phi(x, x')\) to be A-linear, it is necessary and sufficient, according to formula (3) (resp. (4)) of no. 2, that \(\Phi_1\) (resp. \(\Phi_2\) ) be K-linear in \(x\) and that one has

\begin{enumerate}
\def\labelenumi{(\arabic{enumi})}
\setcounter{enumi}{8}
\tightlist
\item
\end{enumerate}

\[
\Phi (x, x ^ {\prime}) = \Phi_ {1} (x, x ^ {\prime}) - i \Phi_ {1} (j (x), x ^ {\prime})\tag{10}
\]

\[
(\text { resp. } \Phi (x, x ^ {\prime}) = \Phi_ {2} (j (x), x ^ {\prime}) + i \Phi_ {2} (x, x ^ {\prime})).
\]

Similarly, for \(\overline{\Phi(x,x^{\prime})}\) to be A-linear in \(x^{\prime}\) , it is necessary and sufficient that \(\Phi_{1}\) (resp. \(\Phi_{2}\) ) be K-linear in \(x^{\prime}\) and that one has

\begin{enumerate}
\def\labelenumi{(\arabic{enumi})}
\setcounter{enumi}{10}
\tightlist
\item
\end{enumerate}

\[
\Phi (x, x ^ {\prime}) = \Phi_ {1} (x, x ^ {\prime}) + i \Phi_ {1} (x, j ^ {\prime} (x ^ {\prime}))\tag{12}
\]

\[
(\text { resp. } \Phi (x, x ^ {\prime}) = - \Phi_ {2} (x, j ^ {\prime} (x ^ {\prime})) + i \Phi_ {2} (x, x ^ {\prime})).
\]

It follows immediately that, for \(\Phi\) to be sesquilinear, it is necessary and sufficient that it be written in one or the other of the forms (9) and (11) (resp. (10) and (12)) with \(\Phi_{1}\) (resp. \(\Phi_{2}\) ) K-bilinear invariant under \(j\) and \(j'\) .

It follows from this that, for a sesquilinear form \(\Phi = \Phi_{1} + i\Phi_{2}\) to be zero, it is necessary and sufficient that \(\Phi_{1}\) (resp. \(\Phi_{2}\) ) be zero. Now, if \(E = E'\) , one has \(\Phi(y, x) = \Phi_{1}(y, x) + i\Phi_{2}(y, x)\) and \(\overline{\Phi(x, y)} = \Phi_{1}(x, y) - i\Phi_{2}(x, y)\) for these two expressions to be equal, in other words for \(\Phi\) to be Hermitian, it is therefore necessary and sufficient that \(\Phi_{1}\) be symmetric (resp. that \(\Phi_{2}\) be antisymmetric).

Remarks. --- 1) Formulas (7) and (8) show that, if \(x \in E\) , for \(\Phi(x, x') = 0\) to hold for all \(x' \in E'\) , it is necessary and sufficient that \(\Phi_1(x, x') = 0\) (resp. \(\Phi_2(x, x') = 0\) ) for all \(x' \in E'\) .

\begin{enumerate}
\def\labelenumi{\arabic{enumi})}
\setcounter{enumi}{1}
\tightlist
\item
  The adjoint of an endomorphism \(u\) of E with respect to \(\Phi\) ( \(\S 1\) , no. 8) is the same as the adjoint of \(u\) (considered as an endomorphism of \(\mathrm{E}_{0}\) ) with respect to \(\Phi_{1}\) (resp. \(\Phi_{2}\) ).
\end{enumerate}

\subsection{Quadratic forms.}

DEFINITION 2. --- Suppose the ring A is commutative. A map Q from E to A is said to be a quadratic form on E if

\begin{enumerate}
\def\labelenumi{\arabic{enumi})}
\item
  one has Q(αx) = α²Q(x) for α ∈ A and x ∈ E ;
\item
  the map \(\Phi : (x, y) \to \mathrm{Q}(x + y) - \mathrm{Q}(x) - \mathrm{Q}(y)\) of \(\mathbf{E} \times \mathbf{E}\) in A is a bilinear form.
\end{enumerate}

The bilinear form \(\Phi\) (which is necessarily symmetric) is called the bilinear form associated with Q. If \(\Phi\) is nondegenerate, Q is said to be nondegenerate.

Since Q(2x) = 4Q(x), it follows immediately from 2) that we have

\[
\Phi (x, x) = 2 \mathrm{Q} (x).\tag{13}
\]

In particular, if A is a ring of characteristic 2, the form \(\Phi\) is alternating.

We will say that two elements (resp. two subsets) of E are orthogonal with respect to Q if they are orthogonal with respect to the associated bilinear form Φ.

Let \((x_{i})_{i\in I}\) be a family of elements of E and \((a_{i})_{i\in I}\) a family of elements of A, all but finitely many of which are zero. By induction on the number of indices \(i\) for which \(a_{i}\neq0\) one shows easily that we have

\[
\mathrm{Q} (\sum_ {i} a _ {i} x _ {i}) = \sum_ {i} a _ {i} ^ {2} \mathrm{Q} (x _ {i}) + \sum_ {\{i, j \}} a _ {i} a _ {j} \Phi (x _ {i}, x _ {j}),\tag{14}
\]

the last summation being extended to the two-element subsets of I.

For every bilinear form \(f\) on \(\mathbf{E} \times \mathbf{E}\) one defines a quadratic form \(\mathbf{Q}\) by setting \(\mathrm{Q}(x) = f(x, x)\) ; the bilinear form \(\Phi\) associated with \(\mathbf{Q}\) is then defined by \(\Phi(x, y) = f(x, y) + f(y, x)\) for \(x, y\) in \(\mathbf{E}\) . Moreover, if one assumes that the scalar 2 has an inverse \(\frac{1}{2}\) in \(\mathbf{A}\) , there exists one and only one symmetric bilinear form \(f\) such that \(\mathrm{Q}(x) = f(x, x)\) , namely \(f = \frac{1}{2}\Phi\) ; the discriminant of \(f\) with respect to a system \(S = (x_1, \ldots, x_n)\) is also called the discriminant of \(\mathbf{Q}\) with respect to \(S\) . In this case, there is therefore a one-to-one correspondence between quadratic forms and symmetric bilinear forms on \(\mathbf{E}\) (cf. exercise 6).

In the case of a free module, we also have the following result:

PROPOSITION 2. --- Suppose that A is commutative and that E admits a basis \((e_i)_{i \in I}\) . Then, for every quadratic form Q on E, there exists a bilinear form \(f\) on E \(\times\) E such that Q(x) = f(x, x) for all \(x \in E\) . For every family \((b_{ij})_{(i,j) \in I \times I}\) of scalars such that \(b_{ij} = b_{ji}\) for \((i, j) \in I \times I\) , there exists one and only one quadratic form Q such that

\[
\mathrm{Q} (e _ {i}) = b _ {i i}, \quad \Phi (e _ {i}, e _ {j}) = b _ {i j} \text {   for   } i \neq j,\tag{15}
\]

where Φ denotes the bilinear form associated with Q; then Q is given by the formula

\[
\mathrm{Q} (\sum_ {i} a _ {i} e _ {i}) = \sum_ {\{i, j \}} b _ {i j} a _ {i} a _ {j},\tag{16}
\]

the last summation being extended over the subsets \(\{i, j\}\) of I having one or two elements.

Indeed, since formula (16) is only a transcription of formula (14), the uniqueness of a quadratic form Q satisfying (15) is proved. To prove its existence, we first note that there exists a family ( \(b_{ij}'\) ) of elements of A such that \(b_{ii}' = b_{ii}\) and \(b_{ij}' + b_{ji}' = b_{ij}\) for \(i \neq j\) ; for example, one obtains such a family by endowing I with a structure of totally ordered set (Ens., chap.~III, § 2, no. 3, th. 1) and setting \(b_{ij}' = b_{ij}\) for \(i<j\) and \(b_{ij}^{\prime}=0\) for \(i>j\) . Since the \(e_{i}\) form a basis of E, there exists a bilinear form \(f\) on \(\mathbf{E}\times\mathbf{E}\) such that \(f(e_{i},e_{j})=b_{ij}^{\prime}\) ; by setting \(\mathrm{Q}^{\prime}(x)=f(x,x)\) and denoting by \(\Phi^{\prime}\) the bilinear form associated with the quadratic form \(\mathrm{Q}^{\prime}\) , we obtain \(\mathrm{Q}^{\prime}(e_{i})=b_{ii}\) and \(\Phi^{\prime}(e_{i},e_{j})=f(e_{i},e_{j})+f(e_{j},e_{i})=b_{ij}\) . This proves our second assertion. As for the first, it follows immediately, because, by virtue of uniqueness, if a quadratic form Q satisfies (15), we have \(\mathrm{Q}(x)=\mathrm{Q}^{\prime}(x)=f(x,x)\) .

The module E equipped with the structure defined by a quadratic form Q is called a quadratic module. A homomorphism of the quadratic module (E, Q) into a quadratic module (E', Q') is a linear map u from E into E' such that Q = Q' ∘ u; if Φ and Φ' are the bilinear forms associated with Q and Q', we then have Φ(x, y) = Φ'(u(x), u(y)) for x ∈ E, y ∈ E; in other words, Φ' is the inverse image of Φ under u (§ 1, no. 1). Two quadratic forms Q and Q' on two A-modules E and E' are said to be equivalent if the corresponding quadratic modules are isomorphic.

Let \((\mathbf{E}_{i}, \mathbf{Q}_{i})_{i\in\mathbb{I}}\) be a family of quadratic modules, and let E be the direct sum of the modules \(E_{i}\) . One calls the external direct sum of the quadratic modules \((\mathbf{E}_{i}, \mathbf{Q}_{i})\) the quadratic module obtained by equipping E with the quadratic form Q defined by \(\mathrm{Q}(\sum_{i}x_{i}) = \sum_{i}\mathrm{Q}_{i}(x_{i})\) for \(x_{i} \in E_{i}\) . One also says that the quadratic form Q is the external direct sum of the quadratic forms \(Q_{i}\) .

If the forms \(Q_{i}\) are nondegenerate, so is Q.

Let Q be a quadratic form on the A-module E; if F is a submodule of E and if Q is constant on each class modulo F, the map \(\overline{Q}\) from E/F into A deduced from Q by passing to the quotient is evidently a quadratic form, and the canonical map from E onto E/F is a homomorphism for the structures of quadratic modules. For Q to be constant on each class modulo F, it is necessary and sufficient that one have \(Q(x + y) = Q(x)\) for \(x \in E\) and \(y \in F\) , that is, denoting by \(\Phi\) the bilinear form associated with Q, that one have \(Q(y) + \Phi(x, y) = 0\) for \(y \in F\) and \(x \in E\) . Setting x = 0, one sees that one has \(Q(y) = 0\) for \(y \in F\) , and hence \(\Phi(x,y)=0\) for \(x\in E\) and \(y\in F\) . In other words, if one calls the kernel of the quadratic module (E, Q) the set N of elements x of E such that \(Q(x)=0\) and \(\Phi(x,z)=0\) for all \(z\in E\) , for Q to be constant on each class modulo F, it is necessary and sufficient that F be contained in the kernel N of (E, Q). One verifies without difficulty that N is a submodule of E. For Q to be constant on each class modulo F, it is therefore necessary and sufficient that F be generated by elements of N.

One sees immediately that the kernel of the quadratic module E/N is \{0\}.

PROPOSITION 3. --- Let h be a homomorphism of A into a commutative ring A'. For every quadratic form Q on the A-module E, there exists one and only one quadratic form Q' on the A'-module A'⊗A E (chap.~III, 2nd ed., App. II, no. 10) such that one has

\[
\mathrm{Q} ^ {\prime} (1 \otimes x) = h (\mathrm{Q} (x))\tag{17}
\]

for all \(x \in \mathbf{E}\) . Moreover, the bilinear form \(\Phi'\) associated with Q' is obtained by extension of the ring of scalars from the bilinear form \(\Phi\) associated with Q.

We first show that, if there exists a quadratic form Q' satisfying (17), it is unique and the bilinear form \(\Phi'\) associated with Q' is obtained by extension of the ring of scalars from the form \(\Phi\) associated with Q. Indeed, this latter assertion follows from the fact that one has

\[
\begin{array}{r l} \Phi^ {\prime} (1 \otimes x, 1 \otimes y) & = Q ^ {\prime} (1 \otimes x + 1 \otimes y) - Q ^ {\prime} (1 \otimes x) - Q ^ {\prime} (1 \otimes y) \\ & = h (\Phi (x, y)) \end{array}\tag{18}
\]

for \(x \in \mathrm{E}, y \in \mathrm{E}\) . Formula (14) then shows that one has

\[
\mathrm{Q} ^ {\prime} \left(\sum_ {i} a _ {i} ^ {\prime} \otimes x _ {i}\right) = \sum_ {i} a _ {i} ^ {\prime 2} h (\mathrm{Q} (x _ {i})) + \sum_ {\{i, j \}} a _ {i} ^ {\prime} a _ {j} ^ {\prime} h (\Phi (x _ {i}, x _ {j}))\tag{19}
\]

for \(a_{i}^{\prime}\in\mathbf{A}^{\prime}\) and \(x_{i}\in\mathbf{E}\) , which proves the uniqueness of \(\mathbf{Q}^{\prime}\) .

To show the existence of Q', we shall first suppose that the module E admits a basis \((e_{i})_{i\in I}\) . There then exist elements \(b_{ij}\) of A such that \(b_{ij}=b_{ji}\) and \(\mathrm{Q}(\sum_{i}a_{i}e_{i})=\sum_{\{i,j\}}b_{ij}a_{i}a_{j}\) for \(a_{i}\in A\) (prop. 2). Since the elements \(1\otimes_{A}e_{i}\) form a basis of the

A'-module A' ⊗\_A E, one defines a quadratic form Q' on this latter module by setting

\[
\mathrm{Q} ^ {\prime} (\sum_ {i} a _ {i} ^ {\prime} \otimes e _ {i}) = \sum_ {\{i, j \}} a _ {i} ^ {\prime} a _ {j} ^ {\prime} h (b _ {i j})\tag{20}
\]

for \(a_i' \in A'\) ; whence, for \(x = \sum_{i} a_i e_i \in E\)

\[
Q ^ {\prime} (1 \otimes x) = Q ^ {\prime} (\sum_ {i} h (a _ {i}) \otimes e _ {i}) = \sum_ {\{i, j \}} h (a _ {i}) h (a _ {j}) h (b _ {i j}) = h (Q (x)),\tag{21}
\]

which proves the existence of Q' in this case.

Now let us turn to the general case. Let \((x_{i})_{i\in I}\) a system of generators of E, \(\mathbf{A}^{(1)}\) the module of formal linear combinations of elements of I (chap.~II, § 1, no. 8), and \((e_{i})_{i\in I}\) the canonical basis of \(\mathbf{A}^{(1)}\) . The linear map \(u\) of \(\mathbf{A}^{(1)}\) in E defined by \(u(e_{i}) = x_{i}\) is surjective since the elements \(x_{i}\) generate E. It follows (chap.~III, 2nd ed., App. II, no. 5, prop. 4) that the map \(1\otimes u\) of \(\mathbf{A}'\otimes\mathbf{A}^{(1)}\) in \(\mathbf{A}'\otimes\mathbf{E}\) is surjective, and its kernel P' is generated by elements of the form \(1\otimes p\) with \(u(p)=0\) Let then \(\mathrm{Q}_{1}^{\prime}\) be the extension to \(\mathbf{A}'\otimes_{\mathbf{A}}\mathbf{A}^{(1)}\) of the quadratic form \(\mathrm{Q}_{1}=\mathrm{Q}\circ u\) on \(\mathbf{A}^{(1)}\) . If \(p\) is an element of \(\mathbf{A}^{(1)}\) such that \(u(p)=0\) , one has \(\mathrm{Q}_{1}^{\prime}(1\otimes p)=h(\mathrm{Q}_{1}(p))=0\) and (denoting by \(\Phi_{1}^{\prime}\) the bilinear form associated with \(\mathrm{Q}_{1}^{\prime}\) ) \(\Phi_{1}^{\prime}(1\otimes p,1\otimes x)=h(\Phi(u(p),u(x)))=0\) for all \(x\in\mathbf{A}^{(1)}\) . Therefore, if \(u(p)=0\) ( \(p\in\mathbf{A}^{(1)}\) ), then \(1\otimes p\) belongs to the kernel of the quadratic module \(\mathbf{A}'\otimes_{\mathbf{A}}\mathbf{A}^{(1)}\) and consequently, as we saw above, there exists a quadratic form Q' on \(\mathbf{A}'\otimes_{\mathbf{A}}\mathbf{E}\) such that \(\mathrm{Q}_{1}^{\prime}=\mathrm{Q}^{\prime}\circ(1\otimes u)\) . Since \(u\) is surjective, we see that Q' satisfies condition (17). QED.

The quadratic form Q′, whose existence and uniqueness are guaranteed by Prop. 3, is called the extension of Q to A′ ⊗\_A E (with respect to h). One also says that Q' is obtained from Q by extension of the ring of scalars.

Exercises. --- 1) Let A be a field, E a left vector space over A, \(\Phi\) a sesquilinear form on E (for an antiautomorphism J of A). We assume that the rank (finite or infinite) of \(\Phi\) is \(\geqslant 2\) and that the relations \(\Phi(x, y) = 0\) and \(\Phi(y, x) = 0\) are equivalent.

\begin{enumerate}
\def\labelenumi{\alph{enumi})}
\item
  Show that there exists \(\lambda\neq0\) in A such that one has \(\Phi(y,x)=\lambda(\Phi(x,y))^{\mathfrak{J}}\) (Use exerc. 8 of § 1).
\item
  Show that there exists \(\alpha\in A\) such that the sesquilinear form \(\Phi\alpha\) (for the antiautomorphism \(\xi\to\alpha^{-1}\xi^{J}\alpha\) ) is Hermitian or alternating.
\end{enumerate}

(First note that one has \(\xi^{J^2} = \lambda^{-1}\xi\lambda^{-J}\) and \(\lambda\lambda^J = \lambda^J\lambda = 1\) . Then distinguish two cases according as \(\xi + \xi^J\lambda^{-1} = 0\) for all \(\xi \in A\) or not; in the second case, show that every element \(\neq 0\) of the form \(\alpha = \xi + \xi^J\lambda^{-1}\) answers the question).

\begin{enumerate}
\def\labelenumi{\arabic{enumi})}
\setcounter{enumi}{1}
\tightlist
\item
  Let \(\Phi\) be a sesquilinear form \(\varepsilon\) -hermitian on a finite-dimensional vector space E over a field A.
\end{enumerate}

\begin{enumerate}
\def\labelenumi{\alph{enumi})}
\item
  Show that for every vector subspace M of E, we have \((\mathbf{M}^{\mathbf{0}})^{\mathbf{0}} = \mathbf{M} + \mathbf{E}^{\mathbf{0}}\) and \(\dim M^{\mathbf{0}} + \dim M = \dim E + \dim (M \cap E^{\mathbf{0}})\) .
\item
  If \(M_{1}\) , \(M_{2}\) are two vector subspaces of E, show that one has dim \((\mathbf{M}_{1} \cap \mathbf{M}_{2}^{0}) + \dim (\mathbf{M}_{2} + \mathbf{M}_{1}^{0}) = \dim \mathbf{E} + \dim (\mathbf{M}_{1} \cap \mathbf{E}^{0})\) (consider the canonical map from E to \(E/E^{0}\) ).
\end{enumerate}

\begin{enumerate}
\def\labelenumi{\arabic{enumi})}
\setcounter{enumi}{2}
\tightlist
\item
  Let K be a commutative ring, \(f\) a monic polynomial in K{[}X{]}, of degree \(n \geqslant 1\) ; let A be the quotient algebra K{[}X{]}/(f) admitting as a basis over K the elements 1, \(\xi, \xi^2, \ldots, \xi^{n-1}\) (chap.~IV, § 1, no. 5, prop. 4). Show that the data of an A-module E is equivalent to the data of a K-module \(E_0\) and of a K-endomorphism \(j\) of \(E_0\) such that \(f(j) = 0\) . For every
\end{enumerate}

\(u \in \mathrm{E}^*\) , set \(u(x) = \sum_{k=0} u_k(x)\xi^k\) ; show that if \(\alpha_0 = f(0)\) is invertible in K, the map \(u \to u_0\) is a K-isomorphism of \(\mathbf{E}^*\) on \((\mathbf{E}_0)^*\) whose reciprocal isomorphism will be made explicit.

¶ 4) a) Let A be a ring (commutative or not), σ an automorphism of A such that there exists an invertible element γ ∈ A satisfying γ\^{}σ = γ, and such that one has ξ\^{}σ² = γξγ\^{}-1 for all ξ ∈ A. Let B be a left A-module having a basis of two elements (e₁, e₂); show that one defines on B a ring structure by taking as multiplication in B the composition law

\[
(\xi e _ {1} + \eta e _ {2}) (\xi^ {\prime} e _ {1} + \eta^ {\prime} e _ {2}) = (\xi \xi^ {\prime} + \eta \eta^ {\prime \sigma} \gamma) e _ {1} + (\eta \xi^ {\prime \sigma} + \xi \eta^ {\prime}) e _ {2}.
\]

For this ring structure, \(e_1\) is the identity element (identified with the identity element 1 of A); if we set \(e_2 = \rho\) , one has \(\rho^2 = \gamma\) and \(\rho\xi = \xi^\sigma\rho\) for all \(\xi \in A\) . Moreover, B is a right A-module, of which 1 and \(\rho\) form a basis. If A is a field, a necessary and sufficient condition for B to be a field is that \(\gamma\) is not of the form \(\lambda^\sigma\lambda\) (where \(\lambda \in A\) ). (Cf. Chap. VIII, § 12, Exerc. 8).

\begin{enumerate}
\def\labelenumi{\alph{enumi})}
\setcounter{enumi}{1}
\tightlist
\item
  Let J be an involutive antiautomorphism of A. Suppose there exists an invertible element \(\delta\in\mathsf{A}\) satisfying the following conditions:
\end{enumerate}

\[
\delta^ {J} = \delta , \quad \delta^ {\sigma} \delta = \gamma \gamma^ {J}, \quad (\xi^ {J}) ^ {\sigma} = \delta (\xi^ {\sigma}) ^ {J} \delta^ {- 1} \quad \text { for   all } \quad \xi \in A.\tag{1}
\]

Show that one can then extend J to an involutive antiautomorphism of B (still denoted J) by setting

\[
(\xi + \eta \rho) ^ {J} = \xi^ {J} + \gamma^ {- 1} \delta^ {\sigma} (\eta^ {J}) ^ {\sigma} \rho .\tag{2}
\]

If in addition A and B are fields, and if \(\sigma\) is not an inner automorphism, show that conditions (1) are necessary for the existence of an extension of J to an involutive antiautomorphism of B (by setting \(\rho^{J} = \alpha + \beta\rho\) , with \(\alpha, \beta\) in A, show that one necessarily has \(\alpha = 0\) by writing the condition \((\rho\xi)^{J} = \xi^{J}\rho^{J}\) for all \(\xi \in A\) ).

\begin{enumerate}
\def\labelenumi{\alph{enumi})}
\setcounter{enumi}{2}
\tightlist
\item
  We assume that conditions (1) are satisfied and that the involutive antiautomorphism J is extended to B by (2). Let F be a unitary B-module, E the unitary A-module obtained from F by restricting the ring of scalars to A; the map \(j: x \to \rho x\) is then a bijective semilinear map from E onto itself, relative to the automorphism \(\sigma\) of A and such that \(j^{2}(x) = \gamma x\) . Let \(\Phi\) be a Hermitian form on F (for J); for \(x \in E\) , \(y \in E\) , set \(\Phi(x, y) = \Phi_{1}(x, y) + \Phi_{2}(x, y)\rho\) , where \(\Phi_{1}(x, y) \in A\) , \(\Phi_{2}(x, y) \in A\) . Show that \(\Phi_{1}\) is a Hermitian form on E (with respect to J), such that
\end{enumerate}

\[
\Phi_ {1} (j (x), j (y)) = (\Phi_ {1} (x, y)) ^ {\sigma} \delta ;
\]

we have \(\Phi_{2}(x,y)=\Phi_{1}(x,j(y))\delta^{-1}\) , \(\Phi_{2}\) is a sesquilinear form on E for the antiautomorphism (in general non-involutive) \(\xi\to(\xi^{\mathrm{J}})^{\sigma}\) , such that

\[
\Phi_ {2} (j (x), j (y)) = (\Phi_ {2} (x, y)) ^ {\sigma} \delta^ {\sigma}
\]

and

\[
\Phi_ {2} (y, x) = \gamma^ {\mathrm{J}} \delta^ {- 1} ((\Phi_ {2} (x, y)) ^ {\mathrm{J}}) ^ {\sigma}.
\]

Converses. For \(\Phi\) to be nondegenerate, it is necessary and sufficient that \(\Phi_{1}\) or \(\Phi_{2}\) be nondegenerate. Special case where B is a quaternion algebra over a commutative ring K, corresponding to a pair \((\alpha, \beta)\) of elements of K, \(\beta\) being invertible, A the subalgebra K + Ku of B, and J and \(\sigma\) the map \(\xi \to \bar{\xi}\) (chap.~II, § 7, no. 8).

\begin{enumerate}
\def\labelenumi{\alph{enumi})}
\setcounter{enumi}{3}
\tightlist
\item
  Let u be an automorphism of the A-module E. For u to be an automorphism of the B-module F, leaving the form \(\Phi\) invariant, it is necessary and sufficient that u satisfy any two of the following three conditions:
\end{enumerate}

\(1^{\circ} u\) leaves \(\Phi_{1}\) invariant;

\(2^{o} u\) leaves \(\Phi_{2}\) invariant;

\(3^{o} u\) commutes with \(j\) .

\begin{enumerate}
\def\labelenumi{\arabic{enumi})}
\setcounter{enumi}{4}
\tightlist
\item
  Let A be a commutative ring, E an A-module, \((x_{i})_{i\in I}\) a system of generators of E; let R be the submodule of the module \(\mathrm{A}^{(\mathrm{I})}\) formed by the elements \((y_{i})_{i\in\mathrm{I}}\) such that \(\sum_{i}y_{i}x_{i}=0\) , and let \((a_{\lambda})_{\lambda\in\mathrm{L}}\) be a system of generators of R (with \(a_{\lambda}=(a_{\lambda i})_{i\in\mathrm{I}}\) ). Let \((b_{ij})\) be a family of elements of A \((i\in\mathrm{I},j\in\mathrm{I})\) . For there to exist a quadratic form Q such that \(\mathrm{Q}(x_{i})=b_{ii}\) and \(\Phi(x_{i},x_{j})=b_{ij}\) for \(i\neq j\) ( \(\Phi\) denoting the bilinear form associated with Q), it is necessary and sufficient that \(b_{ij}=b_{ji}\) for all i, j in I, and that, for all \(\lambda\in L\) and \(i\in I\) one has
\end{enumerate}

\[
\sum_ {\{i, j \}} b _ {i j} a _ {\lambda i} a _ {\lambda j} = \sum_ {j \neq i} b _ {i j} a _ {\lambda j} + 2 b _ {i i} a _ {\lambda i} = 0;\tag{3}
\]

we then have \(\mathrm{Q}(\sum_{i}a_{i}x_{i})=\sum_{\{i,j\}}b_{ij}a_{i}a_{j}\) . Deduce from this a new proof of prop. 3 of no. 4. (Note that the \(x_{i}^{\prime}=1\otimes x_{i}\) form a system of generators of \(A^{\prime}\otimes_{A}E\) , and that the \(A^{\prime}\) -module \(A^{\prime}\otimes_{A}E\) is isomorphic to \(A^{\prime(I)}/R^{\prime}\) , where \(A^{\prime(I)}\) is identified with \(A^{\prime}\otimes_{A}A^{(I)}\) and \(R^{\prime}\) is generated by the image of R under the canonical map from \(A^{(I)}\) into \(A^{\prime(I)}\) ).

\begin{enumerate}
\def\labelenumi{\arabic{enumi})}
\setcounter{enumi}{5}
\tightlist
\item
  Let A be a commutative ring of characteristic 2, E a free A-module, \(\mathcal{A}\) (resp. \(\mathcal{S},\mathcal{Q}\) ) the A-module of alternating (resp. symmetric bilinear, quadratic) forms on E. We have \(\mathcal{A}\subset\mathcal{S}\) ; one also defines a linear map \(\omega\) of \(\mathcal{S}\) in \(\mathcal{Q}\) , and a linear map \(\theta\) of \(\mathcal{Q}\) in \(\mathcal{A}\) as follows: for every bilinear form \(\Phi\in\mathcal{S}\) , \(\omega(\Phi)\) is the quadratic form \(x\to\Phi(x,x)\) , and for every quadratic form \(Q\in\mathcal{Q}\) , \(\theta(Q)\) is the bilinear form associated with Q, which is alternating. Show that \(\omega(0)=\mathcal{A}\) , \(\theta(\mathcal{Q})=\mathcal{A}\) and \(\bar{\theta}(0)=\omega(\mathcal{S})\) .
\end{enumerate}

¶ 7) Let A be a commutative ring, and E, F two A-modules. A map Q from E to F is said to be quadratic if it satisfies the following conditions: \(1^{\circ} \mathrm{Q}(\alpha x) = \alpha^{2} \mathrm{Q}(x)\) for \(\alpha \in \mathrm{A}\) , \(x \in \mathrm{E}\) ; \(2^{\circ}\) the map \((x, y) \to \mathrm{Q}(x + y) - \mathrm{Q}(x) - \mathrm{Q}(y)\) of \(\mathrm{E} \times \mathrm{E}\) into F is bilinear. If \(f\) is a linear map from an A-module \(\mathrm{E}_{1}\) into E, \(\mathrm{Q} \circ f\) is a quadratic map from \(\mathrm{E}_{1}\) into F.

\begin{enumerate}
\def\labelenumi{\alph{enumi})}
\tightlist
\item
  Let E be an A-module, \(\mathrm{A}^{(\mathrm{E})}\) the module of formal linear combinations of elements of E with coefficients in A (chap. II, § 1, no. 8), and for every \(x \in \mathrm{E}\) , let \(\varepsilon_x\) be the corresponding element of the standard basis of \(\mathrm{A}^{(\mathrm{E})}\) . Let \(\Gamma^2(\mathrm{E})\) be the quotient of \(\mathrm{A}^{(\mathrm{E})} \times (\mathrm{E} \otimes_{\mathbf{A}} \mathrm{E})\) by the submodule R generated by the elements ( \(\varepsilon_{x+y} - \varepsilon_x - \varepsilon_y, -x \otimes y\) ) and ( \(\varepsilon_{\lambda x} - \lambda^2\varepsilon_x, 0\) ), for \(x \in \mathrm{E}\) , \(y \in \mathrm{E}\) , \(\lambda \in \mathrm{A}\) . For every \(x \in \mathrm{E}\) , set \(\gamma(x) = \varphi(\varepsilon_x, 0)\) , denoting by \(\varphi\) the canonical mapping from \(\mathrm{A}^{(\mathrm{E})} \times (\mathrm{E} \otimes \mathrm{E})\) on \(\Gamma^2(\mathrm{E})\) ; we say that \(\gamma\) is the canonical mapping from E into \(\Gamma^2(\mathrm{E})\) . Show that \(\gamma\) is a quadratic map from E to \(\Gamma^2(\mathrm{E})\) and that, for every quadratic map Q from E into an A-module F, there exists one and only one linear map \(q\) of \(\Gamma^2(\mathrm{E})\) in F such that \(Q = q \circ \gamma\) (in other words, \((\Gamma^2(\mathrm{E}), \gamma)\) is a solution of a universal mapping problem; cf.~Sets, chap.~IV, § 3).
\end{enumerate}

For every pair of A-modules E, E′ and every linear map f from E into E′, show that, if γ′ denotes the canonical map from E′ into Γ²(E′), there exists one and only one linear map \(\bar{f}\) from Γ²(E) to Γ²(E') such that γ'∘f = \(\bar{f} \circ \gamma\) .

\begin{enumerate}
\def\labelenumi{\alph{enumi})}
\setcounter{enumi}{1}
\item
  Assume that E is the direct sum of two submodules M, N. Define a canonical isomorphism of \(\Gamma^{2}(\mathrm{E})\) onto the direct sum of modules \(\Gamma^{2}(\mathrm{M})\) , \(\Gamma^{2}(\mathrm{N})\) and \(\mathbf{M} \otimes \mathbf{N}\) (show that this direct sum is a solution of the same universal mapping problem as \(\Gamma^{2}(\mathrm{E})\) ).
\item
  Let F be a submodule of E, j the canonical injection of F into E. Define a canonical isomorphism of \(\Gamma^{2}(\mathrm{E}/\mathrm{F})\) onto
\end{enumerate}

\[
\Gamma^ {2} (\mathrm{E}) / (\bar {j} (\Gamma^ {2} (\mathrm{F})) + \psi (\mathrm{E} \times \mathrm{F})),
\]

where \(\psi(x, y) = \varphi(0, x \otimes j(y))\) for \(x \in \mathrm{E}, y \in \mathrm{F}\) . (Same method).

\begin{enumerate}
\def\labelenumi{\alph{enumi})}
\setcounter{enumi}{3}
\item
  Let A' be a commutative ring, and h a homomorphism from A into A'. Define a canonical isomorphism of \(\Gamma^{2}(\mathbf{A}'\otimes_{\mathbf{A}}\mathbf{E})\) onto \(\mathbf{A}'\otimes_{\mathbf{A}}\symbf{\Gamma}^{2}(\mathbf{E})\) (same method).
\item
  There exists one and only one linear map \(s\) of \(\Gamma^{2}(\mathrm{E})\) into \(\mathrm{E}\otimes\mathrm{E}\) such that \(s(\gamma(x))=x\otimes x\) for all \(x\in\mathrm{E}\) ; show that if \(\mathrm{E}\) is a free module, s is an isomorphism onto the submodule of symmetric tensors of order 2 on E.
\item
  Assume that A = Z and that E is a finite cyclic group of order n. Show that \(\Gamma^{2}(E)\) is a cyclic group of order n if n is odd, and of order 2n if n is even. (First note that if a is a generator of E, \(\gamma(a)\) is a generator of \(\Gamma^{2}(E)\) , and that \(\gamma(-ha) = \gamma(ha)\) for every integer h; deduce from this that if n is odd, \(n\gamma(a) = 0\) by taking \(h = (n - 1)/2\) ; similarly show that \(2n\gamma(a) = 0\) if n is even. Finally, prove that if n is odd (resp. even), there exists a quadratic map Q from E into a cyclic group of order n (resp. 2n) sending a onto a generator of this group.
\end{enumerate}

\begin{enumerate}
\def\labelenumi{\arabic{enumi})}
\setcounter{enumi}{7}
\tightlist
\item
  Let A be a commutative field, and let E, F be two vector spaces over A. Let g be a map from E to F, such that there exist three maps a, b, c from E × E to F, satisfying the identity
\end{enumerate}

\[
g (\lambda x + \mu y) = \lambda^ {2} a (x, y) + \lambda \mu b (x, y) + \mu^ {2} c (x, y)\tag{4}
\]

for all x, y in E, λ, μ in A.

\begin{enumerate}
\def\labelenumi{\alph{enumi})}
\tightlist
\item
  Show that we have \(a(x, y) = g(x)\) , \(c(x, y) = g(y)\) , \(b(y, x) = b(x, y)\) and \(b(\lambda x, y) = \lambda b(x, y)\) ; moreover, if we set
\end{enumerate}

\[
d (x, y, z) = b (x + y, z) - b (x, z) - b (y, z)
\]

show that we have \(d(x,y,z)=d(y,z,x)=d(z,x,y)\) and conclude from this that \(d(\lambda x,\mu y,\nu z)=\lambda\mu\nu d(x,y,z)\) .

\begin{enumerate}
\def\labelenumi{\alph{enumi})}
\setcounter{enumi}{1}
\tightlist
\item
  From a), deduce that if \(\mathbf{A} \neq \mathbf{F}_2\) one necessarily has \(d(x, y, z) = 0\) and consequently that \(g\) is a quadratic map (exerc. 7). On the contrary, if \(\mathbf{A} = \mathbf{F}_2\) and if dim \(\mathrm{E} \geqslant 3\) show that there exist maps \(g\) from E to A, satisfying (4) and for which \(d(x, y, z)\) is not identically zero.
\end{enumerate}

\begin{enumerate}
\def\labelenumi{\arabic{enumi})}
\setcounter{enumi}{8}
\tightlist
\item
  Let A be a commutative ring, E an A-module having a basis of n elements, Φ a symmetric bilinear form on E. Let e be an element of \(\bigwedge^{n}\) E forming a basis of this module, \(\Delta\) the discriminant of \(\Phi\) with respect to \(e\) , \(\varphi_{p}\) the canonical isomorphism of \(\bigwedge^{p}\) E onto \(\bigwedge^{n-p}\) E* relative to \(e\) (chap.~III, §8, no. 5). For \(x \in \bigwedge^{p}\) E, let \(d_{(p)}(x) = \bar{\varphi}_{n-p}^{-1}(d_{\Phi(p)}(x))\) ; show that the map \(d_{(p)}\) of \(\bigwedge^{p}\) E into \(\bigwedge^{n-p}\) E possesses the following properties:
\end{enumerate}

\begin{enumerate}
\def\labelenumi{\alph{enumi})}
\item
  For every \(x \in \wedge \mathrm{E}\) , \(d_{(n - p)}(d_{(p)}(x)) = (-1)^{p(n - p)}\Delta x\) .
\item
  For every pair of elements \(x, y\) of \(\wedge\) E, we have
\end{enumerate}

\[
x \wedge d _ {(p)} (y) = \Phi_ {(p)} (x, y) e \quad \text { and } \quad \Phi_ {(n - p)} (d _ {(p)} (x), d _ {(p)} (y)) = \Delta \Phi_ {(p)} (x, y).
\]

\begin{enumerate}
\def\labelenumi{\alph{enumi})}
\setcounter{enumi}{2}
\item
  We further assume that A is a field and that \(\Phi\) is nondegenerate. Then, if \(x\) is a decomposable \(p\) -vector \(\neq 0\) corresponding to a subspace F of E (chap. III, § 7, no. 3), \(d_{(p)}(x)\) is a decomposable \((n-p)\) -vector \(\neq 0\) corresponding to the subspace F \(^{0}\) orthogonal to F.
\item
  Extend the previous results to the case where (A being a commutative ring), \(\Phi\) is a sesquilinear form \(\varepsilon\) -hermitian for an involutive automorphism \(J \neq 1\) of A.
\end{enumerate}

\begin{enumerate}
\def\labelenumi{\arabic{enumi})}
\setcounter{enumi}{9}
\tightlist
\item
  \begin{enumerate}
  \def\labelenumii{\alph{enumii})}
  \tightlist
  \item
    Let A be a commutative ring, E an A-module having a basis of 3 elements, \(\Phi\) a symmetric bilinear form on E. With the notations of exerc. 9, for any two elements \(x\) , \(y\) of E, one sets \(x \overline{\wedge} y = d_{(2)}(x \wedge y)\) , and one says that this element is the vector product
  \end{enumerate}
\end{enumerate}

\[
\stackrel {{3}} {\wedge} \mathrm{E})
\]

of \(x\) and of \(y\) (relative to \(\Phi\) and to the basis \(e\) of \(\wedge E\) ). Show that \((x, y) \to x \overline{\wedge} y\) is an alternating bilinear map from \(E \times E\) into \(E\) , and that \(x \overline{\wedge} y\) is orthogonal to \(x\) and to \(y\) .

\begin{enumerate}
\def\labelenumi{\alph{enumi})}
\setcounter{enumi}{1}
\tightlist
\item
  Let \(\alpha, \beta\) be two invertible elements of A, B the quaternion algebra over A corresponding to the pair \((\alpha, \beta)\) (chap.~II, § 7, no. 8), 1, u, v, w the canonical basis of B over A; let E be the submodule of B having u, v, w as basis. Show that if x, y are two quaternions belonging to E, one has
\end{enumerate}

\[
x y = \Phi (x, y) + x \overline {{{{\wedge}}}} y
\]

where \(\Phi\) is a symmetric bilinear form on E, such that the linear maps associated with \(\Phi\) are bijective, and \(x\overline{\wedge}y\) is the vector product of x and y relative to the form \(\Phi\) and to the basis \(\alpha^{-1}\beta^{-1}u\wedge v\wedge w\) of \(\bigwedge^{3}\mathbf{E}\) .

\begin{enumerate}
\def\labelenumi{\arabic{enumi})}
\setcounter{enumi}{10}
\tightlist
\item
  Let \(\Phi\) be a nondegenerate \(\varepsilon\) -hermitian sesquilinear form on a finite-dimensional vector space E. One says that a vector subspace M of E is weakly orthogonal to a vector subspace N (relative to \(\Phi\) ) if one of the two subspaces M, N° contains the other.
\end{enumerate}

\begin{enumerate}
\def\labelenumi{\alph{enumi})}
\item
  Show that the relation “M is weakly orthogonal to N” is symmetric.
\item
  If M and N are weakly orthogonal, \(M^{0}\) and \(N^{0}\) are weakly orthogonal.
\item
  If M and N are weakly orthogonal and if M ∩ N = \{0\}, M and N are orthogonal.
\end{enumerate}

\section{Totally isotropic subspaces. Witt's theorem}

In this section one assumes, unless expressly stated otherwise, that A is a field. One denotes by \(\Phi\) , either a form \(\varepsilon\) -hermitian on E (with respect to the involutive antiautomorphism \(\lambda \to \overline{\lambda}\) of A), or the symmetric bilinear form associated with a quadratic form Q on E (A being assumed commutative in the latter case).

\subsection{Isotropic subspaces.}

DEFINITION 1. --- Given a module E over the ring A, an element x of E is said to be isotropic if \(\Phi(x, x) = 0\) . A submodule F of E is said to be

\begin{enumerate}
\def\labelenumi{\arabic{enumi})}
\item
  isotropic if there exists an element \(x \neq 0\) of F orthogonal to F;
\item
  totally isotropic if the restriction of \(\Phi\) to F is zero.
\end{enumerate}

When the module E is equipped with a quadratic form Q, an element of E will be said to be isotropic (resp. a submodule of E will be said to be isotropic, or totally isotropic) if this element is isotropic (resp. if this submodule is isotropic, or totally isotropic) with respect to the bilinear form associated with Q.

An isotropic vector is none other than a vector orthogonal to itself. To say that a submodule F is isotropic means that \(F \cap F^{0} \neq \{0\}\) , or equivalently that the restriction of \(\Phi\) to F is degenerate; a non-isotropic submodule G of E is therefore a submodule such that the restriction of \(\Phi\) to G is nondegenerate. For a submodule F of E to be totally isotropic, it is necessary and sufficient that one has \(F \subset F^{0}\) . If F is a totally isotropic submodule of E, the same is true of every submodule \(F'\) contained in F. The sum of a family of totally isotropic submodules that are pairwise orthogonal is a totally isotropic submodule. The set of totally isotropic submodules of E, ordered by inclusion, is evidently inductive; it follows that every totally isotropic submodule is contained in a maximal totally isotropic submodule.

PROPOSITION 1. --- Suppose that A is a field. Let F be a finite-dimensional non-isotropic subspace of E; then E is the direct sum of F and F⁰.

Indeed, since the restriction \(\Phi'\) of \(\Phi\) to F is nondegenerate by hypothesis, the map \(d_{\Phi'}\) from F to its dual F* associated on the right to \(\Phi'\) is injective, hence bijective since F and F* are two spaces of the same finite dimension. Consequently, for every \(y \in E\) , there exists one and only one element \(y_0\) of F such that one has \(\Phi(x, y) = \Phi(x, y_0)\) for all \(x \in F\) , that is to say \(y - y_0 \in F^0\) ; this proves that E is the direct sum of F and F \(^0\) .

COROLLARY. --- If F is a finite-dimensional subspace of E, and if Φ is nondegenerate, the following conditions are equivalent:

\begin{enumerate}
\def\labelenumi{\alph{enumi})}
\item
  F is non-isotropic.
\item
  \(\mathbf{F}^0\) is non-isotropic.
\item
  E is the direct sum of F and F \(^{0}\) .
\end{enumerate}

Proposition 1 indeed shows that \(a)\) implies \(c)\) , and \(c)\) implies \(a)\) and \(b)\) . Finally, if \(\mathbf{F}^0\) is non-isotropic, we have \(\mathbf{F} \cap \mathbf{F}^0 \subset \mathbf{F}^0 \cap \mathbf{F}^{00} = \{0\}\) ; hence \(\mathbf{F}\) is non-isotropic, which shows that \(b)\) implies \(a)\) .

DEFINITION 2. --- Let Q be a quadratic form on E. An element x of E is said to be singular (relative to Q) if Q(x) = 0. A submodule F of E is said to be:

\begin{enumerate}
\def\labelenumi{\arabic{enumi})}
\item
  singular if there exists an element \(x \neq 0\) of F which is singular and orthogonal to F;
\item
  totally singular if the restriction of Q to F is zero.
\end{enumerate}

The kernel of the quadratic module (E, Q) (§ 3, no. 4) consists of the singular elements of E°; for a submodule F to be singular, it is necessary and sufficient that its kernel be ≠ \{0\} . As Φ(x, y) = Q(x + y) − Q(x) − Q(y), every totally singular submodule ≠ \{0\} is singular. Since Φ(x, x) = 2Q(x), every singular vector is isotropic and every singular (resp. totally singular) submodule is isotropic (resp. totally isotropic); the converse holds if A is a field of characteristic ≠ 2. Every submodule contained in a totally singular submodule is itself totally singular. The sum of a family of totally singular submodules that are pairwise orthogonal is a totally singular submodule. The set of totally singular submodules of E, ordered by inclusion, is inductive; hence every totally singular submodule of E is contained in a maximal totally singular submodule.

\subsection{Witt decomposition.}

In addition to the conventions already in force since the beginning of the present paragraph, we shall add the following:

Condition (T). --- For all \(x \in E\) , there exists \(\alpha \in A\) such that \(\Phi(x, x) = \alpha + \varepsilon\bar{\alpha}\) .

This condition is always satisfied when \(\Phi\) is alternating, or when \(\varepsilon = 1\) and A is a field of characteristic \(\neq 2\) , taking then \(\alpha = \frac{1}{2}\Phi(x, x)\) (cf. Exercises 1 and 14).

Lemma 1. --- Let \(\Phi\) be a form \(\varepsilon\) -hermitian satisfying (T) (resp. the bilinear form associated with a quadratic form Q) on E, and let F be a totally isotropic (resp. totally singular) subspace of E, not reduced to 0. For every \(x \in E\) not orthogonal to F and all \(\alpha \in A\) , there exists \(y \in F\) such that

\[
\Phi (x + y, x + y) = \alpha + \varepsilon \bar {\alpha} \quad (\text { resp. } Q (x + y) = \alpha).
\]

Indeed, let us set \(\Phi(x, x) = \beta + \varepsilon\bar{\beta}\) (resp. \(Q(x) = \beta\) ). For \(y \in F\) we then have \(\Phi(x + y, x + y) = (\beta + \Phi(x, y)) + \varepsilon(\overline{\beta + \Phi(x, y)})\) since \(\Phi(y, y) = 0\) (resp. \(Q(x + y) = \beta + \Phi(x, y)\) since \(Q(y) = 0\) ). Since \(x\) is not orthogonal to \(F\) , the affine linear function \(y \to \Phi(x, y) + \beta\) on \(F\) is not constant; it therefore takes the value \(\alpha\) for some element \(y\) of \(F\) , which thus answers the question.

A Witt decomposition of E is any decomposition of E as a direct sum of three subspaces F, F', G such that F and F' are totally isotropic (resp. totally singular) and G is non-isotropic and orthogonal to F + F'; if E is finite-dimensional, the matrix of Φ with respect to a basis of E adapted to a Witt decomposition of E takes the form

\[
\left( \begin{array}{c c c} 0 & U & 0 \\ \varepsilon \overline {{U}} & 0 & 0 \\ 0 & 0 & V \end{array} \right)\tag{1}
\]

We say that \(\Phi\) is a neutral form if it is nondegenerate and if E is finite-dimensional and is the direct sum of two totally isotropic (resp. totally singular) subspaces. The direct sum of two neutral forms is a neutral form.

PROPOSITION 2. --- Let Φ be a nondegenerate ε-hermitian form satisfying (T) (resp. the bilinear form associated with a nondegenerate quadratic form Q), and let F be a totally isotropic (resp. totally singular) subspace of finite dimension r.

\begin{enumerate}
\def\labelenumi{\alph{enumi})}
\item
  If \(F'\) is a totally isotropic subspace of dimension r such that \(F' \cap F^{0} = \{0\}\) , then \(F + F'\) is non-isotropic and, for any basis \((f_{i})\) of F, there exists a basis \((f_{i}')\) of \(F'\) such that \(\Phi(f_{i}, f_{j}') = \delta_{ij}\) (Kronecker index) for \(i, j = 1, \ldots, r\) .
\item
  If G is a totally isotropic (resp. totally singular) subspace of dimension \(\leqslant r\) such that \(G \cap F^{0} = \{0\}\) , there exists a totally isotropic (resp. totally singular) subspace \(F' \supset G\) of dimension r such that \(F' \cap F^{0} = \{0\}\) .
\end{enumerate}

Let \(\Psi\) be the restriction of \(\Phi\) to \(F \times F'\) ; for \(x' \in F'\) , the relation \(\langle \Phi(x, x') = 0\) for all \(x \in F\rangle\) implies \(x = 0\) since \(F' \cap F^0 = \{0\}\) . The assertion \(a)\) then follows from the corollary of Proposition 6 of §1, no. 6, except for the fact that \(F + F'\) is non-isotropic. Now the subspace \(H = (F + F') \cap (F + F')^0\) is equal to \((F + F') \cap F^0 \cap F'^0\) . Since \(F \subset F^0\) , one has \((F + F') \cap F^0 = F + (F' \cap F^0) = F\) , whence \(H = F \cap F'^0\) ; hence \(H = \{0\}\) since we have seen that \(\Psi\) is nondegenerate. This indeed proves that \(F + F'\) is non-isotropic.

To demonstrate \(b\) ), we proceed by descending induction on \(s = \dim G\) . It thus suffices for us to prove that, if \(s < r\) , there exists a totally isotropic (resp. totally singular) subspace \(G'\) containing \(G\) , of dimension \(s + 1\) , and such that \(G' \cap F^0 = \{0\}\) . Since \(\dim G < \dim F\) , the restriction of \(\Phi\) to \(F \times G\) is degenerate, and since \(G \cap F^0\) is zero, \(F \cap G^0\) is nonzero. If one then had \(G + F^0 \supset G^0\) one would deduce, by taking the orthogonal subspaces and noting that \(F = F^{00}\) and \(G = G^{00}\) ( \(\S 1\) , no. 6, cor. 1 of prop. 4), that \(G^0 \cap F \subset G\) , whence

\[
\mathrm{G} ^ {0} \cap \mathrm{F} \subset \mathrm{G} \cap \mathrm{F} \subset \mathrm{G} \cap \mathrm{F} ^ {0} = \{0 \},
\]

which is impossible. There then exists an element \(x\) of \(G^0\) such that \(x \notin G + F^0\) ; since \(F \subset F^0\) , one can add to \(x\) a vector of \(G^0 \cap F\) without modifying these properties; since \(G^0 \cap F\) is totally isotropic (resp. totally singular) and \(\neq \{0\}\) , lemma 1 shows that one can choose \(x\) isotropic (resp. singular).

Then the subspace \(G' = G + Ax\) is of dimension \(s + 1\) , and is totally isotropic (resp. totally singular); moreover we have \(G' \cap F^{0} = \{0\}\) because, if \(y = z + ax\) ( \(z \in G, a \in A\) ) is in \(F^{0}\) , we have a = 0, since otherwise \(x \in F^{0} + G\) contrary to the choice of x, whence \(y = z \in G \cap F^{0} = \{0\}\) and y = 0. Consequently the subspace \(G'\) answers the question.

COROLLARY 1. --- If F is a totally isotropic (resp. totally singular) subspace of dimension r, there exists a totally isotropic (resp. totally singular) subspace F' of dimension r such that F ∩ F' = \{0\} and such that F + F' is non-isotropic.

It suffices to take G = \{0\} in prop. 2, b).

COROLLARY 2. --- Two neutral ε-hermitian forms on spaces of the same dimension over A are equivalent.

Remark. --- Under the conditions of corollary 1, E is the direct sum of F + F' and of the orthogonal of F + F'. One thus has a Witt decomposition of E. By prop. 2, a), there exist bases of F and F' such that, in the matrix (1) of Φ, the block U is the identity matrix 1.

PROPOSITION 3. --- Let \(\Phi\) be a nondegenerate \(\varepsilon\) -hermitian form satisfying (T) (resp. the bilinear form associated with a nondegenerate quadratic form Q). Let \(F_{1}\) and \(F_{2}\) be two maximal totally isotropic (resp. totally singular) subspaces of E, one of the two being finite-dimensional. Set \(F = F_{1} \cap F_{2}\) . Let \(S_{i}\) ( \(i = 1, 2\) ) be a supplement of F in \(F_{i}\) ; set \(S = S_{1} + S_{2}\) . There then exist two subspaces G and H of E such that

\begin{enumerate}
\def\labelenumi{\alph{enumi})}
\item
  The subspaces G + F, S, and H are non-isotropic and pairwise orthogonal;
\item
  E is the direct sum of F, S, G, and H;
\item
  There is no nonzero isotropic (resp. singular) vector in H;
\item
  G is totally isotropic (resp. totally singular).
\end{enumerate}

Additionally, \(F_{1}\) and \(F_{2}\) are both finite-dimensional and we have \(\dim \mathrm{F}_1 = \dim \mathrm{F}_2, \dim \mathrm{G} = \dim \mathrm{F}, \dim \mathrm{S}_1 = \dim \mathrm{S}_2, \operatorname{codim} \mathrm{H} = 2 \dim \mathrm{F}_1.\)

First, note that if N is a maximal totally isotropic (resp. totally singular) subspace, then every isotropic (resp. singular) vector x orthogonal to N belongs to N, for otherwise \(N + Ax\) would contradict the maximality of N. Thus, if, for i = 1, 2, \(x_{i}\) is an isotropic (resp. singular) vector of \(F_{i}^{0}\) , one has \(x_{i} \in F_{i}\) On the other hand, if y is an element of \(S_{1}\) orthogonal to \(S_{2}\) it is orthogonal to \(F_{1}\) since \(F_{1}\) is totally isotropic, hence to F, and consequently to \(F_{2} = S_{2} + F\) . Since y is isotropic (resp. singular) and is orthogonal to \(F_{2}\) , one has

\[
y \in S _ {1} \cap F _ {2} = S _ {1} \cap F _ {1} \cap F _ {2} = S _ {1} \cap F = \{0 \}.
\]

We therefore have \(S_1 \cap S_2^0 = \{0\}\) , and likewise \(S_2 \cap S_1^0 = \{0\}\) . Since one of the two subspaces \(F_1, F_2\) , for example \(F_1\) , is finite-dimensional, \(S_1\) is finite-dimensional, hence \(S_1^0\) has finite codimension (§ 1, no. 6, cor. 1 of prop. 4), and consequently \(S_2\) is finite-dimensional since \(S_2 \cap S_1^0 = \{0\}\) ; moreover this shows that one has dim \(S_2 \leqslant \text{codim } S_1^0 = \text{dim } S_1\) ; likewise dim \(S_1 \leqslant \text{dim } S_2\) , whence dim \(S_1 = \text{dim } S_2\) . Prop. 2 a) then shows that \(S = S_1 + S_2\) is non-isotropic.

This being so, the orthogonal N of S is non-isotropic (no. 1, cor. of prop. 1) and contains F; cor. 1 of prop. 2 therefore shows that there exists a subspace G totally isotropic (resp. totally singular) of N such that dim G = dim F, that G ∩ F = \{0\} and that G + F is non-isotropic. Thus d) is satisfied by G. One then satisfies a) and b) by taking for H the orthogonal of G + F in N. As for c), one notes that, since H is orthogonal to F₁ = S₁ + F, there is no nonzero isotropic (resp. singular) vector in H by virtue of what was remarked at the beginning of the proof and of the fact that H ∩ F₁ = \{0\}. Finally, certain of the assertions concerning dimensions were proved along the way; the others follow from them trivially.

COROLLARY 1. --- Under the hypotheses of prop. 3, two maximal totally isotropic (resp. totally singular) subspaces of finite dimension have the same dimension. For every maximal totally isotropic (resp. totally singular) subspace F of finite dimension, there exists another one \(\mathbf{F}'\) such that \(\mathbf{F} \cap \mathbf{F}' = \{0\}\) , and under these conditions \(\mathbf{F} + \mathbf{F}'\) is non-isotropic.

If \(F \cap F' = \{0\}\) , one will have \(G = \{0\}\) with the notations of prop. 3, and \(F + F'\) will be non-isotropic. The other assertions follow trivially from prop. 3 and cor. 1 of prop. 2.

COROLLARY 2. --- Let Q be a nondegenerate quadratic form on a vector space E of finite dimension n over an algebraically closed field A; there then exists a basis \((e_i)_{1 \leqslant i \leqslant n}\) of E such that

\begin{enumerate}
\def\labelenumi{(\arabic{enumi})}
\setcounter{enumi}{1}
\tightlist
\item
\end{enumerate}

\[
\mathrm{Q} (\sum_ {i = 1} ^ {n} x _ {i} e _ {i}) = \sum_ {i = 1} ^ {\nu} x _ {i} x _ {i + \nu} \quad \text {   if   } n = 2 \nu ,\tag{3}
\]

\[
\mathrm{Q} \left(\sum_ {i = 1} ^ {n} x _ {i} e _ {i}\right) = \sum_ {i = 1} ^ {\nu} x _ {i} x _ {i + \nu} + x _ {2 \nu + 1} ^ {2} \quad \text {   if   } n = 2 \nu + 1.
\]

Indeed, let \(F_{1}\) and \(F_{2}\) be two maximal totally singular subspaces such that \(F_{1} \cap F_{2} = \{0\}\) (cor. 1) and let q be their dimension. One then has \(G = \{0\}\) with the notations of prop. 3. Taking a basis \((e_{i})_{1 \leqslant i \leqslant q}\) of \(F_{1}\) and a basis \((e_{i})_{q+1 \leqslant i \leqslant 2q}\) of \(F_{2}\) such that \(\Phi(e_{i}, e_{j+q}) = \delta_{ij}\) for \(i, j = 1, \ldots, q\) (prop. 2 a)), we see that it suffices to show that dim \(H \leqslant 1\) . Now, if \(x \in H\) , \(y \in H\) and if \(x \neq 0\) , the equation \(Q(y - ax) = Q(y) - a\Phi(x, y) + a^{2}Q(x) = 0\) has at least one solution \(a_{0}\) since \(Q(x) \neq 0\) and one has \(y = a_{0}x\) since every singular vector of H is zero.

DEFINITION 3. --- Suppose that E is finite-dimensional and that \(\Phi\) is a nondegenerate \(\varepsilon\) -hermitian form satisfying (T) (resp. the bilinear form associated with a nondegenerate quadratic form Q). The index of \(\Phi\) (resp. Q) is the common dimension of the maximal totally isotropic (resp. totally singular) subspaces of E.

If n is the dimension of E and v the index of \(\Phi\) (resp. Q), Prop. 3 shows that one has

\[
n \geqslant 2 v.\tag{4}
\]

Moreover, since every totally isotropic (resp. totally singular) subspace is contained in a maximal totally isotropic (resp. totally singular) subspace, the maximal totally isotropic (resp. totally singular) subspaces are precisely those of dimension v. The assertion that \(\Phi\) (resp. Q) has index 0 means that every isotropic (resp. singular) vector of E is zero. In a space of even dimension \(n\) the neutral forms are those of index \(\frac{1}{2} n\) ; there is no neutral form in a space of odd dimension. Prop. 3 shows that every form is a direct sum of a neutral form and a form of index 0.

PROPOSITION 4. --- Let Q be a nondegenerate quadratic form on E such that there exists a vector \(x \neq 0\) of E with Q(x) = 0. For every element a of A, there then exists \(y \in E\) such that Q(y) = a.

Indeed, by cor. 1 of prop. 2, there exists a subspace \(\mathbf{G} = \mathbf{F} + \mathbf{F}'\) ( \(\mathbf{F}, \mathbf{F}'\) : totally singular subspaces of dimension 1) of \(\mathbf{E}\) of dimension 2, such that the restriction of \(\mathbf{Q}\) to \(\mathbf{G}\) is neutral. If \(\{e, e'\}\) ( \(e \in \mathbf{F}, e' \in \mathbf{F}'\) ) is a basis of \(\mathbf{G}\) , one has

\[
\mathrm{Q} (x e + x ^ {\prime} e ^ {\prime}) = b x x ^ {\prime} \quad (x \in \mathrm{A}, x ^ {\prime} \in \mathrm{A}, b \in \mathrm{A}, b \neq 0).
\]

It therefore suffices to take for y the vector \(ae + b^{-1}e'\) .
% semantic-fragments: legacy-input-seam
\relax{}
\subsection{Witt's Theorem.}

Given two vector spaces E, E' over A equipped respectively with two sesquilinear forms \(\Phi\) , \(\Phi'\) (resp. with two quadratic forms Q, Q'), one calls a metric homomorphism from E to E' any linear map \(u\) from E to E' such that \(\Phi'(u(x), u(y)) = \Phi(x, y)\) (resp. Q'(u(x)) = Q(x)) for \(x \in E\) , \(y \in E\) . If E and E' have the same finite dimension and if \(\Phi\) (resp. Q) is nondegenerate, every metric homomorphism \(u\) from E to E' is an isomorphism, since \(u(x) = 0\) implies \(\Phi(x, y) = 0\) for all \(y \in E\) , hence \(x = 0\) ; thus \(u\) is injective, hence bijective since E and E' have the same finite dimension.

THEOREM 1 (Witt). — Let E and E' be two finite-dimensional vector spaces, equipped respectively with two nondegenerate ε-hermitian forms Φ and Φ' satisfying condition (T) of no. 2 (resp. with two nondegenerate quadratic forms Q and Q'), and isomorphic for these structures. Given any subspace F of E, every injective metric homomorphism of F into E' extends to a metric isomorphism of E onto E'.

Using the given isomorphism from E onto \(\mathbf{E}'\) , one sees that it suffices to show that every injective metric homomorphism \(u\) from F into E extends to a metric automorphism of E. Note that if, for \(i = 1, 2\) , \(\mathbf{F}_i\) is a subspace of E and \(u_i\) a metric homomorphism from \(\mathbf{F}_i\) into E, such that \(\mathbf{F}_1 \cap \mathbf{F}_2 = \{0\}\) and \(\Phi(u_1(x_1), u_2(x_2)) = \Phi(x_1, x_2)\) for \(x_i \in \mathbf{F}_i\) ( \(i = 1, 2\) ), then the homomorphism \(\nu : x_1 + x_2 \to u_1(x_1) + u_2(x_2)\) of \(\mathbf{F}_1 + \mathbf{F}_2\) into E which extends \(u_1\) and \(u_2\) is metric: indeed, for any \(x_i, y_i\) in \(\mathbf{F}_i\) ( \(i = 1, 2\) ), the expansion of each of the expressions \(\Phi(x_1 + x_2, y_1 + y_2)\) and \(\Phi(u_1(x_1) + u_2(x_2), u_1(y_1) + u_2(y_2))\) (resp. \(\mathrm{Q}(x_1 + x_2)\) and \(\mathrm{Q}(u_1(x_1) + u_2(x_2))\) ) contains four (resp. three) terms equal each to each according to the hypotheses made. Moreover, if \(u_1\) and \(u_2\) are injective and if \(u_1(\mathbf{F}_1) \cap u_2(\mathbf{F}_2) = \{0\}\) , then \(\nu\) is injective.

\begin{enumerate}
\def\labelenumi{\arabic{enumi})}
\tightlist
\item
  Let us first prove Witt's theorem in the case where the set of fixed points of u is a hyperplane U of F. The set of vectors of the form \(u(x) - x\) with \(x \in F\) is then a line D. If \(F'\) is a subspace orthogonal to D such that \(F' \cap F = F' \cap u(F) = \{0\}\) , one will have \(\Phi(u(x), y) = \Phi(x, y)\) for \(x \in F\) and \(y \in F'\) ; our initial remark therefore applies to u and to the identity map of \(F'\) into E, showing that u extends to \(F + F'\) while leaving the points of \(F'\) fixed; the set of vectors of the form \(u(x) - x (x \in F + F')\) is still the line D. Now we have, for \(x \in F, y \in F\)
\end{enumerate}

\[
\Phi (u (x), u (y) - y) = \Phi (u (x), u (y)) - \Phi (u (x), y) = \Phi (x - u (x), y),
\]

which, when \(x \in \mathrm{U}\) (that is, when \(u(x) = x\) ), shows that \(x \in \mathrm{D}^0\) ; in other words, we have \(\mathrm{U} \subset \mathrm{D}^0\) . We distinguish two cases:

\begin{enumerate}
\def\labelenumi{\alph{enumi})}
\item
  \(F \subset D^{0}\) . Formula (5) shows that \(u(F)\) is not contained in \(D^{0}\) , hence \(F \cap D^{0} = u(F) \cap D^{0} = U\) . One can then take for \(F'\) a supplement of U in \(D^{0}\) ; since \(F + F'\) contains the hyperplane \(D^{0}\) and is distinct from it, we have \(F + F' = E\) , and we have found in this case the desired extension of u to E.
\item
  \(\mathbf{F} \subset \mathbf{D}^0\) . Formula (5) shows that \(u(\mathbf{F}) \subset \mathbf{D}^0\) , and hence that
\end{enumerate}

D ⊂ D⁰ ; the line D is therefore isotropic (resp. singular, since we have Q(u(x) - x) = Q(u(x)) - Φ(x, u(x)) + Q(x) = 2Q(x) - Φ(x, x) = 0 for x ∈ F). We shall show that, under these conditions, there exists a subspace F' of D⁰ which is a supplement of F and of u(F) in D⁰. This is immediate if F = u(F). Otherwise, let x and y be vectors such that x ∈ F, x ∉ U, y ∈ u(F), y ∉ U; we then have F = U + Ax, u(F) = U + Ay, and F does not contain x + y, otherwise y = (x + y) - x would belong to F ∩ u(F) = U; we see similarly that x + y does not belong to u(F); thus the line A(x + y) is a supplement of F and of u(F) in the subspace F + u(F); it then suffices to set F' = A(x + y) + G, where G is a supplement of F + u(F) in D⁰. This being so, we have F + F' = u(F) + F' = D⁰, and, in this case, what was said at the beginning of 1) shows that there exists an extension of u to the hyperplane D⁰ of E, and that D⁰ is stable for this extension.

We are therefore reduced to the case where F is the hyperplane D⁰ and where u is an automorphism of F. Let us prove that, for every z ∈ E, there exists z' ∈ E such that

\[
\Phi (u (x), z ^ {\prime}) = \Phi (x, z)\tag{6}
\]

for all \(x \in \mathbf{F}\) ; indeed the linear form \(x \to \Phi(u^{-1}(x), z)\) on \(\mathbf{F}\) is the restriction of a linear form on \(\mathbf{E}\) , a form which is of the type \(x \to \Phi(x, z')\) since \(\Phi\) is nondegenerate; hence (6) is valid. Moreover, if \(z \notin \mathbf{F}\) there exists a vector \(z' \in \mathbf{E}\) satisfying (6) and such that \(\Phi(z', z') = \Phi(z, z)\) (resp. \(Q(z') = Q(z)\) ). Indeed, formula (6) remains valid if one adds to \(z'\) an element \(u(y) - y\) ( \(y \in \mathbf{F}\) ) of \(\mathbf{D}\) since \(\mathbf{F} = \mathbf{D}^0\) , and Lemma 1 of No. 2 allows us to conclude since \(z\) is not orthogonal to \(\mathbf{D}\) Our initial remark then shows that there exists a metric homomorphism \(v\) of \(\mathbf{F} + A z = E\) into \(E\) which extends \(u\) and which transforms \(z\) to \(z'\) . Since \(\Phi\) is nondegenerate, \(v\) is the metric automorphism of \(E\) sought.

\begin{enumerate}
\def\labelenumi{\arabic{enumi})}
\setcounter{enumi}{1}
\tightlist
\item
  In the general case, we will argue by induction on \(r = \dim F\) The case \(r = 0\) is trivial. Let then \(r > 0\) , that is to say \(F \neq \{0\}\) and let U be a hyperplane of F. The restriction \(u_0\) of \(u\) to U extends, by the induction hypothesis, to a metric automorphism \(v_0\) of E. If \(v_0\) extends \(u\) , the theorem is proved. Otherwise, U is the set of invariant elements by \(v_{0}^{-1}u\) , and there exists, by 1), a metric automorphism \(v_{1}\) of E extending \(v_{0}^{-1}u\) . The automorphism \(v_{0}v_{1}\) is then the desired extension of u. QED.
\end{enumerate}

COROLLARY 1. --- Let, for \(i=1,2\) , \(\mathbf{E}_{i}\) be a finite-dimensional vector space, \(\Phi_{i}\) a nondegenerate \(\varepsilon\) -hermitian form on \(\mathbf{E}_{i}\) satisfying (T) (resp. \(\mathbf{Q}_{i}\) a nondegenerate quadratic form on \(\mathbf{E}_{i}\) ), \(\mathbf{E}_{i}^{\prime}\) and \(\mathbf{E}_{i}^{\prime\prime}\) two orthogonal subspaces of \(\mathbf{E}_{i}\) of which \(\mathbf{E}_{i}\) is the direct sum. If the forms \(\Phi_{1}\) and \(\Phi_{2}\) (resp. \(\mathbf{Q}_{1}\) and \(\mathbf{Q}_{2}\) ) are equivalent, and if their restrictions to \(\mathbf{E}_{1}^{\prime}\) and \(\mathbf{E}_{2}^{\prime}\) are equivalent, then so are their restrictions to \(\mathbf{E}_{1}^{\prime\prime}\) and \(\mathbf{E}_{2}^{\prime\prime}\) .

Indeed, let \(u\) be a metric isomorphism of \(\mathbf{E}_{1}^{\prime}\) onto \(\mathbf{E}_{2}^{\prime}\) . By th. 1, \(u\) extends to a metric isomorphism \(\nu\) of \(\mathbf{E}_{1}\) onto \(\mathbf{E}_{2}\) . Since \(\Phi_{i}\) is nondegenerate, \(\mathbf{E}_{i}^{\prime\prime}\) is the orthogonal of \(\mathbf{E}_{i}^{\prime}\) in \(\mathbf{E}_{i}\) , hence \(\nu\) maps \(\mathbf{E}_{1}^{\prime\prime}\) on \(\mathbf{E}_{2}^{\prime\prime}\) . QED.

COROLLARY 2. --- Under the hypotheses of Theorem 1, the group of metric automorphisms of E acts transitively on the totally isotropic (resp. totally singular) subspaces of E of a given dimension. Moreover, if F is a totally isotropic (resp. totally singular) subspace of E, then every bijective linear map from F onto F is induced by a metric automorphism of E.

COROLLARY 3. --- Let Q be a nondegenerate quadratic form on a finite-dimensional vector space E over an algebraically closed field A. The group of metric automorphisms of E acts transitively on the non-isotropic subspaces of E of a given dimension.

This follows immediately from th. 1 and cor. 2 of prop. 3.

Exercises. -- 1) a) Let K be a field of characteristic 2, J : \(\xi \to \bar{\xi}\) an involutive antiautomorphism of K, Z the center of K. Show that if the restriction of J to Z is not the identity, every element \(\mu\) of K such that \(\bar{\mu} = \mu\) is of the form \(\lambda + \bar{\lambda}\) (note that there is an element \(\rho \neq 0\) in Z which can be written in the form \(\zeta + \bar{\zeta}\) , with \(\zeta \in \mathbb{Z}\) ); every Hermitian form on a vector space over K then satisfies condition (T).

\begin{enumerate}
\def\labelenumi{\alph{enumi})}
\setcounter{enumi}{1}
\tightlist
\item
  Give examples of fields of characteristic 2 admitting an involutive antiautomorphism \(\xi \to \overline{\xi}\) distinct from the identity map, and for which there exist elements \(\mu = \bar{\mu}\) which are not of the form \(\lambda + \bar{\lambda}\) (cf. Chap. VIII, § 11, Exerc. 4).
\end{enumerate}

\begin{enumerate}
\def\labelenumi{\arabic{enumi})}
\setcounter{enumi}{1}
\tightlist
\item
  Let A be a field, E a vector space over A, \(\Phi\) (resp. Q) a nondegenerate \(\varepsilon\) -hermitian form on E, satisfying condition (T) (resp. a nondegenerate quadratic form on E, \(\Phi\) then denoting the symmetric bilinear form associated with Q).
\end{enumerate}

\begin{enumerate}
\def\labelenumi{\alph{enumi})}
\item
  Show that for a plane P ⊂ E to be isotropic (resp. singular) and not totally isotropic (resp. not totally singular), it is necessary and sufficient that it contain only one isotropic (resp. singular) line (cf. exerc. 14 e)).
\item
  We assume that dim \(E \geqslant 3\) and that there exist isotropic vectors in E \(\neq 0\) Show that if P is a non-totally isotropic plane in E, there exists a non-isotropic subspace \(V \subset E\) of dimension 3, containing isotropic vectors \(\neq 0\) , and such that \(P \subset V\) .
\end{enumerate}

\begin{enumerate}
\def\labelenumi{\arabic{enumi})}
\setcounter{enumi}{2}
\item
  With the same hypotheses as in exerc. 2, show that if dim E ≥ 3, every isotropic line in E is the intersection of two non-isotropic planes.
\item
  The hypotheses are those of exerc. 2, and it is further assumed that E is of finite dimension.
\end{enumerate}

\begin{enumerate}
\def\labelenumi{\alph{enumi})}
\item
  If the index v of \(\Phi\) (resp. Q) is \(\geqslant 1\) , show that for every isotropic (resp. singular) vector \(a \neq 0\) in E, there exists a basis \((e_i)\) of E consisting of isotropic (resp. singular) vectors, such that \(e_1 = a\) (cf. exerc. 14 e)).
\item
  Let V, W be two totally isotropic (resp. totally singular) subspaces of the same dimension \(r \leqslant v\) ; show that there exist two maximal totally isotropic (resp. totally singular) subspaces \(V_{1}\) , \(W_{1}\) , such that \(V \subset V_{1}\) , \(W \subset W_{1}\) , and \(V_{1} \cap W_{1} = V \cap W\) . (If U = V ∩ W, reason in \(U^{0}/U\) ).
\item
  Let V, W, \(V_{1}\) , \(W_{1}\) be four totally isotropic (resp. totally singular) subspaces of the same dimension, such that \(V + W\) and \(V_{1} + W_{1}\) are non-isotropic. Show that there exists a metric automorphism u of E such that \(u(V) = V_{1}\) and \(u(W) = W_{1}\) .
\item
  Let \(f\) be a linear form on E, \(\alpha\) an element of A of the form \(\lambda + \varepsilon\bar{\lambda}\) (resp. an element of A). We consider the sesquilinear form on E
\end{enumerate}

\[
(x, y) \rightarrow \Phi_ {1} (x, y) = \Phi (x, y) + f (x) \alpha \overline {{f (y)}}
\]

(resp. the quadratic form

\[
x \rightarrow \mathrm{Q} _ {1} (x) = \mathrm{Q} (x) + \alpha (f (x)) ^ {2}).
\]

Show that if \(\Phi_{1}\) (resp. \(Q_{1}\) ) is nondegenerate and if \(\nu_{1}\) denotes its index, we have \(|\nu_{1} - \nu| \leqslant 1\) .

\begin{enumerate}
\def\labelenumi{\arabic{enumi})}
\setcounter{enumi}{4}
\tightlist
\item
  \begin{enumerate}
  \def\labelenumii{\alph{enumii})}
  \tightlist
  \item
    Let B be a ring, \(\xi \to \overline{\xi}\) an involutive antiautomorphism of B, \(\varepsilon\) an element of the center of B such that \(\varepsilon \varepsilon = 1\) . Show that if \(\beta\) is an invertible element of B such that \(\beta + \varepsilon \overline{\beta} \neq 0\) , there exists an invertible element \(\mu \neq 1\) in B, such that \(\mu (\beta + \varepsilon \overline{\beta}) \mu = \beta + \varepsilon \overline{\beta}\) (Show that one can take \(\mu\) such that \(\mu \beta \overline{\mu} = \beta\) ).
  \end{enumerate}
\end{enumerate}

\begin{enumerate}
\def\labelenumi{\alph{enumi})}
\setcounter{enumi}{1}
\tightlist
\item
  Let A be a field, and let E be a vector space over A, \(\Phi\) a nondegenerate ε-hermitian sesquilinear form on E, satisfying (T). Show that if Φ is not alternating, then for every non-isotropic hyperplane H of E, there exists a metric automorphism of E, distinct from the identity, leaving every element of H invariant (use a)).
\end{enumerate}

\begin{enumerate}
\def\labelenumi{\arabic{enumi})}
\setcounter{enumi}{5}
\item
  Given the hypotheses of exercise 2, let \(a\) be an isotropic vector \(\neq 0\) in E (resp. a nonsingular isotropic vector (note that such vectors exist only if A has characteristic 2)). Let \(\lambda \in \mathbf{A}\) such that \(\lambda + \varepsilon \bar{\lambda} = 0\) (resp. \(\lambda = (\mathrm{Q}(a))^{-1}\) ); show that the transvection \(x \to x + \Phi(x, a) \lambda a\) (chap. II, § 6, exerc. 7) is a metric automorphism of E; conversely.
\item
  Under the hypotheses of Exercise 2, let G be the group of metric automorphisms of E. Show that the only semilinear bijections of E onto itself that commute with every element of G are the homotheties of E, except in the following three cases: dim E = 2, G is the group of metric automorphisms corresponding to a quadratic form of index 1 on E, and A is one of the three fields \(\mathbf{F}_{2}, \mathbf{F}_{3}\) or \(\mathbf{F}_{4}\) (Use Exercises 5, 6, and 3; examine separately the case of a quadratic form on a vector space of dimension 2).
\end{enumerate}

*8) Let A be a field, E a vector space of finite dimension \textgreater0 over A, Φ a nondegenerate ε-hermitian sesquilinear form on E, satisfying (T). Let M(Φ) be the group of multipliers of the similarities of E for Φ (§ 6, no. 5).

\begin{enumerate}
\def\labelenumi{\alph{enumi})}
\item
  Let \(V_{1}\) , \(V_{2}\) two vector subspaces of E, of the same dimension, and let \(\Phi_{1}\) , \(\Phi_{2}\) be the restrictions of \(\Phi\) to \(V_{1}\) , \(V_{2}\) respectively. For there to exist a similarity u such that \(u(V_{1}) = V_{2}\) , it is necessary and sufficient that there exist \(\alpha \in M(\Phi)\) such that \(\Phi_{2}\) is equivalent to \(\alpha\Phi_{1}\) (use Witt's theorem).
\item
  Let (F, F', G) be a Witt decomposition of E (no. 2), and let \(\Phi_0\) be the restriction of \(\Phi\) to the non-isotropic subspace G. Show that one has M( \(\Phi\) ) = M( \(\Phi_0\) ) if G ≠ \{0\}. (Use Witt's theorem and prop. 2 of no. 2).
\item
  Show that if the index v of Φ is such that dim E = 2v, M(Φ) is the group of elements ζ ≠ 0 of the center of A such that ζ̄ = ζ. If dim E = 2v + 1, M(Φ) is the group of elements of the form ρρ̄, where ρ runs through the multiplicative group of elements ≠ 0 of the center of A (use Witt's theorem). (Cf. § 10, exerc. 18).*
\end{enumerate}

*9) Let A be a commutative field, E a vector space over A, Q a nondegenerate quadratic form on E. One calls a similarity for Q any automorphism \(u\) of E such that there exists an element \(\alpha \neq 0\) of A for which \(Q(u(x)) = \alpha Q(x)\) for every \(x \in E\) ; \(u\) is then also a similarity for the bilinear form associated with Q. Assuming the dimension of E finite and \(>0\) , state and prove for the similarities relative to Q the analogues of the results of exerc. 8.*

*10) Let A be a field, E a vector space over A of dimension \(\geqslant 2\) , \(\Phi_1\) (resp. \(\Phi_2\) ) a nondegenerate sesquilinear form \(\varepsilon_1\) -hermitian (resp. \(\varepsilon_2\) -hermitian) on E for an involutive antiautomorphism \(J_1\) (resp. \(J_2\) ) of A, satisfying condition (T). Show that if the group of metric automorphisms of E for \(\Phi_1\) is a subgroup of the group of similarities for \(\Phi_2\) , there exists \(\alpha \in A\) such that \(\Phi_2 = \Phi_1\alpha\) (use Exercises 5 b) and 6).

Prove the analogous property when A is assumed commutative, and that \(\Phi_{1}\) and \(\Phi_{2}\) are replaced in the statement by two nondegenerate quadratic forms \(Q_{1}, Q_{2}\) on E.*

\begin{enumerate}
\def\labelenumi{\arabic{enumi})}
\setcounter{enumi}{10}
\tightlist
\item
  Let A be a ring, J an involutive antiautomorphism of A, and E an A-module admitting a finite basis \((e_{i})\) , \(\Phi\) a nondegenerate \(\varepsilon\) -hermitian form on E, R the matrix of \(\Phi\) with respect to \((e_{i})\) ; the group of metric automorphisms of \(\Phi\) is identified with the group G of invertible matrices U such that \(^{t}U.R.U^{j}=R\) .
\end{enumerate}

\begin{enumerate}
\def\labelenumi{\alph{enumi})}
\tightlist
\item
  Assume that there exists a matrix \(P\) such that \(R = {}^t P + \varepsilon P^J\) . Show that for every matrix \(S\) such that \({}^t S + \varepsilon S^J = 0\) , and that \(P + S\) is invertible, \(U = (^{t}P^{J} - \varepsilon^{J}S)^{-1}(P + S)\) belongs to G, and \(\varepsilon I + U\) is invertible. Conversely (show that for every matrix \(U \in G\) such that \(\varepsilon I + U\) is invertible, we have
\end{enumerate}

\[
\varepsilon (\varepsilon I + ^ {t} U) ^ {- 1} R + \varepsilon^ {\mathrm{J}} R (\varepsilon^ {\mathrm{J}} I + U ^ {\mathrm{J}}) ^ {- 1} = R).
\]

\begin{enumerate}
\def\labelenumi{\alph{enumi})}
\setcounter{enumi}{1}
\tightlist
\item
  Show that the condition of a) is satisfied when \(\Phi\) satisfies condition (T). Case where in A the equation \(2\xi = \alpha\) admits one and only one solution for every \(\alpha \in \mathbf{A}\) .
\end{enumerate}

¶ 12) Let A be a commutative field, E a vector space of finite dimension n over A, Φ a nondegenerate ε-hermitian sesquilinear form on E.

\begin{enumerate}
\def\labelenumi{\alph{enumi})}
\item
  Let \(u\) be an endomorphism of E; show that if the \(r_i\) ( \(1 \leqslant i \leqslant m\) ) are the similarity invariants of \(u\) (chap. VII, § 5, no. 1, def. 1), the similarity invariants of the adjoint \(u^*\) of \(u\) with respect to \(\Phi\) are the polynomials \(\bar{r}_i\) ( \(1 \leqslant i \leqslant m\) ), where \(\bar{r}_i\) is deduced from \(r_i\) by applying to each coefficient the automorphism J (cf. chap. VII, § 5, exerc. 2). For every monic irreducible polynomial \(p \in \mathrm{A}[\mathrm{X}]\) dividing the minimal polynomial of \(u\) , let \(\mathrm{F}_k(u, p)\) be the kernel of \((p(u))^k\) in E, and let \(\mathrm{F}(u, p)\) be the union of the \(\mathrm{F}_k(u, p)\) for all integers \(k > 0\) . Show that if \(p\) and \(q\) are two distinct monic irreducible polynomials dividing the minimal polynomial of \(u\) , the subspaces \(\mathrm{F}(u, p)\) and \(\mathrm{F}(u^*, \bar{q})\) are orthogonal (use Bezout's identity). Finally, if G is a vector subspace of E such that \(u(\mathrm{G}) \subset \mathrm{G}\) , one has \(u^*(\mathrm{G}^0) \subset \mathrm{G}^0\) .
\item
  We assume that \(uu^{*} = u^{*}u\) (the case where one says that \(u\) is a normal endomorphism for \(\Phi\) ; cf.§ 7, no. 3); show that one then has \(u^{*}(F_{k}(u,p)) \subset F_{k}(u,p)\) for all \(k\) and consequently \(u^{*}(F(u,p)) \subset F(u,p)\) . If one sets \(G(p,\bar{q}) = F(u,p) \cap F(u^{*},\bar{q})\) , show that E is the direct sum of the subspaces \(G(p,\bar{q})\) , and that \(G(p,\bar{q})\) and \(G(p_{1},\bar{q}_{1})\) are orthogonal if \(p \neq q_{1}\) or if \(p_{1} \neq q\) ; in particular \(G(p,\bar{q})\) is totally isotropic if \(p \neq q\) . Show that \(G(p,\bar{p})\) is reduced to 0 or non-isotropic, and that if \(p \neq q\) , no nonzero vector of \(G(p,\bar{q})\) is orthogonal to \(G(q,\bar{p})\) (using the fact that \(\Phi\) is nondegenerate); deduce that if \(p \neq q\) , \(G(p,\bar{q})\) and \(G(q,\bar{p})\) are totally isotropic subspaces of the same dimension, and that \(G(p,\bar{q}) + G(q,\bar{p})\) is non-isotropic.
\item
  Assume that J is not the identity or that A does not have characteristic 2, and that \(u^{*} = u\) . Let M be the set of non-isotropic subspaces \(\mathbf{M} \subset \mathrm{G}(p, \overline{p})\) stable under \(u\) (hence submodules of the A{[}X{]}- module \(E_{u}\) (chap. VII, § 5, no. 1)). Show that if M is a minimal element of M, then M is an indecomposable submodule of \(E_{u}\) (chap. VII, § 4, no. 7). (Assume that M is a direct sum of an indecomposable submodule \(M_{1}\) and of a submodule \(M_{2} \neq \{0\}\) , the minimal polynomials \(p^{h}\) and \(p^{k}\) of the restrictions of u to \(M_{1}\) and \(M_{2}\) respectively being such that \(h \geqslant k\) . Note then that \(M_{1}\) is necessarily isotropic and that every \(z \neq 0\) in \(M_{1}\) such that \(p(u).z = 0\) is orthogonal to \(M_{1}\) (using the fact that every submodule of \(M_{1}\) is monogenic); writing that \(z = (p(u))^{h-1}.x\) and that z is not orthogonal to \(M_{2}\) , and deduce that one necessarily has k = h. Then show that there exists an indecomposable submodule \(N_{2}\) of \(M_{2}\) such that \(p^{h}\) be the minimal polynomial of the restriction of u to \(N_{2}\) , and that \(M_{1} + N_{2}\) be non-isotropic; conclude from this that \(M_{2} = N_{2}\) . Finally, if \(y \in M_{2}\) is not orthogonal to z, consider the submodule P of M generated by \(w = x + \lambda y\) , where \(\lambda \in A\) , and show that one can take \(\lambda\) such that P is non-isotropic, by proving that one has \(\Phi((p(u))^{h-1}.w,w) \neq 0\) ; which leads to a contradiction).
\end{enumerate}

\(d)\) Deduce from \(c)\) that \(\mathrm{G}(p,\bar{p})\) is a direct sum of indecomposable submodules \(\mathrm{H}_{i}\) pairwise orthogonal. If \(p^{h}\) is the minimal polynomial of the restriction of \(u\) to \(\mathrm{H}_{i}\) , and if \(d\) is the degree of \(p\) , show that there exists in \(\mathrm{H}_{i}\) a totally isotropic subspace of dimension \(d\) . {[}h/2{]}. Case where E contains no isotropic vector \(\neq 0\) (cf.§ 7, no. 3).

\begin{enumerate}
\def\labelenumi{\alph{enumi})}
\setcounter{enumi}{4}
\item
  State and prove the properties analogous to those of c) and d) when one has \(u^{*} = -u\) or \(u^{*}u = 1\) .
\item
  Give an example where n = 4, Φ is symmetric and of index 2, p = \(\bar{p} = X - 1\) , u is normal, E = G(p, \(\bar{p}\) ), but E is not a direct sum of minimal submodules of M, and where there exists an eigenvector of u that is not an eigenvector of \(u^{*}\) (cf.§ 7, no. 3).
\end{enumerate}

\begin{enumerate}
\def\labelenumi{\arabic{enumi})}
\setcounter{enumi}{12}
\item
  The hypotheses are those of Exercise 2, and one further assumes that E admits a countable basis \((e_n)\) . Let F be a totally isotropic (resp. totally singular) subspace of E such that \(F^{00} = F\) . Show that there exists a totally isotropic (resp. totally singular) subspace \(F'\) such that: \(1^\circ F \cap F' = \{0\}\) ; \(2^\circ\) there exists a basis \((a_m)_{m \in I}\) of F and a basis \((a'_m)_{m \in I}\) of \(F'\) (I interval of N starting at 0) such that \(\Phi(a_i, a_j') = \delta_{ij}\) for every pair of indices; \(3^\circ (F + F')^{00} = F + F'\) and E is the direct sum of \(F + F'\) and of \(G = (F + F')^0\) . (Construct by induction an increasing sequence \((L_n)\) of non-isotropic subspaces, with union E, such that dim \(L_{n+1} - \dim L_n = 2\) , and apply prop. 2 of no. 2 to each of the \(L_n\) ; to construct this sequence, one will consider, for every \(n\) , the smallest integer \(k\) such that \(e_k \notin L_n\) , and one will use exerc. 9 b) of § 1).
\item
  Let A be a field of characteristic 2, E a vector space of finite dimension n over A, Φ a nondegenerate Hermitian form on E, not necessarily satisfying condition (T).
\end{enumerate}

\begin{enumerate}
\def\labelenumi{\alph{enumi})}
\item
  Show that the set V of \(x \in E\) such that \(\Phi(x, x)\) is of the form \(\alpha + \bar{\alpha}\) is a vector subspace of E.
\item
  Let \(V_{1}=V\cap V^{0}\) , \(q=\dim V_{1}\) , \(V_{2}\) a supplement of \(V_{1}\) with respect to V, \(V_{3}\) a supplement of \(V_{1}\) with respect to \(V^{0}\) . Show that there exists a basis \((e_{i})_{1\leqslant i\leqslant2q}\) of \((\mathrm{V}_{2}+\mathrm{V}_{3})^{0}=\mathrm{V}_{2}^{0}\cap\mathrm{V}_{3}^{0}\) such that the vectors \(e_1, \ldots, e_q\) form a basis of \(V_1\) and that one has \(\Phi(e_i, e_{q+j}) = \delta_{ij}\) for \(1 \leqslant i \leqslant q\) , \(1 \leqslant j \leqslant q\) .
\item
  Let \(\mathbf{G}(\Phi)\) be the group of metric automorphisms of E (for \(\Phi\) ). Show that for every \(u \in \mathbf{G}(\Phi)\) , one has \(u(x) = x\) for all \(x \in V^0\) .
\item
  For every \(u \in \mathbf{G}(\Phi)\) , one has \(u(V) = V\) ; let \(u_{v}\) be the restriction of \(u\) to \(V\) , and let \(G_{v}\) be the group formed by the \(u_{v}\) . Show that: 1° the kernel of the homomorphism \(u \to u_{v}\) of \(\mathbf{G}(\Phi)\) onto \(G_{v}\) is commutative; 2° if \(\Phi_{2}\) is the restriction of \(\Phi\) to \(V_{2}\) and \(\mathbf{G}(\Phi_{2})\) the group of metric automorphisms of \(V_{2}\) for \(\Phi_{2}\) , there exists a homomorphism from \(G_{v}\) onto \(\mathbf{G}(\Phi_{2})\) whose kernel is commutative (use \(b\) ) and \(c\) ).
\item
  We assume that A is commutative and J is the identity; let E be a vector space of dimension 3 over A, \((e_{i})_{1\leqslant i\leqslant3}\) a basis of E, \(\Phi\) the nondegenerate symmetric form on E whose matrix with respect to \((e_{i})\) is
\end{enumerate}

\[
\left( \begin{array}{c c c} 1 & 0 & 0 \\ 0 & 0 & 1 \\ 0 & 1 & 0 \end{array} \right).
\]

Show that all isotropic vectors in E are contained in the hyperplane spanned by \(e_{2}\) and \(e_{3}\) (cf. exercise 4(a)). Give an example of a non-isotropic plane containing only one isotropic line (cf. exercise 2(a)). Show that there exists no automorphism \(u \in \mathbf{G}(\Phi)\) such that \(u(e_{1}) = e_{1} + e_{2}\) , even though one has \(\Phi(e_{1}, e_{1}) = \Phi(e_{1} + e_{2}, e_{1} + e_{2})\) .

\section{Special properties of alternating bilinear forms}

\subsection{Reduction of alternating bilinear forms.}

THEOREM 1. --- Let A be a principal (commutative) ring, E a free A-module of finite dimension n, and \(\Phi\) an alternating bilinear form on E. Then there exists a basis \((e_i)_{1 \leqslant i \leqslant n}\) of E and an even integer \(2r \leqslant n\) , such that

\[
1 ^ {0} \quad \Phi (e _ {1}, e _ {2}) = \alpha_ {1}, \Phi (e _ {3}, e _ {4}) = \alpha_ {2}, \dots , \Phi (e _ {2 r - 1}, e _ {2 r}) = \alpha_ {r}
\]

where the \(\alpha_{i}\) are elements \(\neq 0\) of A, and where \(\alpha_{i}\) divides \(\alpha_{i+1}\) for \(i=1,\ldots,r-1\) .

2° All the other elements \(\Phi(e_{i}, e_{j})\) where \(i \leqslant j\) are zero.

The ideals \(\mathrm{A}\alpha_{i}\) ( \(i=1,\ldots,r\) ) are uniquely determined by the preceding conditions. The submodule \(\mathrm{E}^{0}\) of E orthogonal to E is generated by \(e_{2r+1},\ldots,e_{n}\) .

We shall proceed by induction on the dimension n of E. The theorem is obvious for n = 0. If \(\Phi = 0\) , the theorem is also obvious; we may therefore assume \(\Phi \neq 0\) . Let f denote the linear map \(d_{\Phi}\) from E into \(\mathbf{E}^*\) associated on the right to \(\Phi\) ( \(\S 1, \text{no. 1}\) ); then \(f(\mathbf{E})\) is a submodule not reduced to 0 of the module \(\mathbf{E}^*\) , which is a free module of dimension \(n\) . Let \(\mathrm{A}\alpha_{1}\) be the largest invariant factor of \(f(\mathbf{E})\) with respect to \(\mathbf{E}^*\) (chap. VII, \(\S 4, \text{no. 2}\) , th. 1); we know (loc. cit.) that there exists a basis \((e_{1}', a_{2}', \ldots, a_{n}')\) of \(\mathbf{E}^*\) and an element \(f(e_{2}) \in f(\mathbf{E})\) such that \(f(e_{2}) = \alpha_{1}e_{1}'\) . Let \((e_{1}, a_{2}, \ldots, a_{n})\) the basis of E (identified with the bidual \(\mathbf{E}^{**}\) ) dual of \((e_{1}', a_{2}', \ldots, a_{n}')\) ; we have

\[
\Phi (e _ {1}, e _ {2}) = - \Phi (e _ {2}, e _ {1}) = \left\langle e _ {1}, f (e _ {2}) \right\rangle = \alpha_ {1}.\tag{1}
\]

Let P be the submodule \(Ae_{1} + Ae_{2}\) of E. We shall see that E is the direct sum of \(Ae_{1}\) , \(Ae_{2}\) and of the submodule \(P^{0}\) orthogonal to P. For this it suffices to prove that, for every \(x \in E\) , there exist elements \(\xi_{1}\) , \(\xi_{2}\) of A, uniquely determined and such that \(x - \xi_{1}e_{1} - \xi_{2}e_{2} \in P^{0}\) that is, such that

\[
\Phi (e _ {1}, x - \xi_ {1} e _ {1} - \xi_ {2} e _ {2}) = 0, \quad \Phi (e _ {2}, x - \xi_ {1} e _ {1} - \xi_ {2} e _ {2}) = 0.
\]

According to (1) these conditions are written

\[
\left\langle e _ {1}, f (x) \right\rangle = \xi_ {2} \alpha_ {1}, \left\langle e _ {2}, f (x) \right\rangle = - \xi_ {1} \alpha_ {1}.
\]

But we know (loc. cit.) that the image of \(f(E)\) under any linear form on \(E^{*}\) is contained in the ideal \(A\alpha_{1}\) , in other words all the values \(\Phi(x,y)=\langle x,f(y)\rangle\) belong to \(A\alpha_{1}\) ; whence the existence and uniqueness of \(\xi_{1}\) and \(\xi_{2}\) . Thus \(P^{0}\) is a free module of rank \(n-2\) ; there therefore exists, in \(P^{0}\) , by the induction hypothesis, a basis \((e_{3},e_{4},\ldots,e_{n})\) satisfying the conditions of the statement. To show that the basis \((e_{1},\ldots,e_{n})\) of E thus obtained also satisfies these conditions, it suffices to prove that \(\alpha_{1}\) divides \(\alpha_{2}\) ; now this follows from the fact that all the values \(\Phi(x,y)\) are multiples of \(\alpha_{1}\) . It is then clear that \(e_{2r+1},\ldots,e_{n}\) generate \(E^{0}\) . Finally, if \((e_{i}^{\prime})\) is the dual basis of \((e_{i})\) , one has \(f(e_{2j-1})=-\alpha_{j}e_{2j}^{\prime}\) and \(f(e_{2j})=\alpha_{j}e_{2j}^{\prime}\) for \(j=1,\ldots,r\) and \(f(e_{k})=0\) for \(k=2r+1,\ldots,n\) ; the ideals \(A\alpha_{1},A\alpha_{1},A\alpha_{2},A\alpha_{2},\ldots,A\alpha_{r},A\alpha_{r}\) are therefore the invariant factors of \(f(E)\) with respect to \(E^{*}\) , which proves their uniqueness (chap. VII, § 4, no. 2, th. 1).

COROLLARY 1. --- Let A be a commutative field, E a vector space of finite dimension n over A, and Φ an alternating bilinear form on E. Then there exists a basis \((e_i)_{1\leqslant i\leqslant n}\) of E and an even integer \(2r\leqslant n\) such that

\[
\Phi (\sum_ {i = 1} ^ {n} \xi_ {i} e _ {i}, \sum_ {i = 1} ^ {n} \eta_ {i} e _ {i}) = \sum_ {j = 1} ^ {r} (\xi_ {2 j - 1} \eta_ {2 j} - \xi_ {2 j} \eta_ {2 j - 1}).\tag{2}
\]

In particular \(\Phi\) is of even rank \(2r\) .

A basis satisfying (2) is called a symplectic basis for \(\Phi\) .

Remark. --- Note that this corollary is also an immediate consequence of prop. 3 of § 4, no. 2 and of its cor. 1, for with the notations of that proposition, one necessarily has H = \{0\} since Φ is alternating.

COROLLARY 2. --- Let A be a commutative field, E a vector space of finite dimension n over A. For every bivector \(z \in \bigwedge^{2} \mathrm{E}\) there exists a basis \((e_i)_{1 \leqslant i \leqslant n}\) of E such that

\[
z = e _ {1} \wedge e _ {2} + e _ {3} \wedge e _ {4} + \dots + e _ {2 r - 1} \wedge e _ {2 r} \quad (2 r \leqslant n).
\]

It suffices indeed to note that z is canonically identified with an alternating bilinear form on E* (chap. III, § 8, no. 2), and to apply cor. 1 to this form.

Translating cor. 1 into matrix language, we obtain:

COROLLARY 3. --- Let A be a commutative field, R an alternating square matrix over A. The rank of R is an even number 2r, and there exists an invertible matrix P over A such that

\[
{ } ^ { t } P . R . P = \left( \begin{array} { c c c } 0 & I _ { r } & 0 \\ - I _ { r } & 0 & 0 \\ 0 & 0 & 0 \end{array} \right) .
\]

Remark. --- If A is any commutative ring and R is an alternating square matrix of odd order n over A, then det R = 0. This follows from Cor. 3 when A is a field. A direct proof in the case where A is a field of characteristic ≠ 2 is the following: since \(^tR = -R\) , we have det \(R = \det ^t R = (-1)^n\) det R, hence 2 det \(^tR = 0\) . This being so, since the determinant of an alternating matrix ( \(\alpha_{ij}\) ) is a polynomial with integer coefficients in the \(\alpha_{ij}\) such that i \textless{} j, the principle of extension of algebraic identities (chap. IV, § 2, no. 5, Scholium) then shows that our assertion is true for an arbitrary commutative ring A.

\subsection{Pfaffian of a skew-symmetric matrix.}

Let A be a commutative field of characteristic 0, and \(R = (\alpha_{ij})\) an alternating matrix of even order \(2m\) on A. Let E denote the vector space \(\mathbf{A}^{2m}\) , \((e_i)\) ( \(i = 1, \ldots, 2m\) ) its canonical basis, and \(e\) the element \(e_1 \wedge e_2 \wedge \ldots \wedge e_{2m}\) of \(\bigwedge^{2m}\) E. The \(m\) -th exterior power of the bivector \(u = \sum_{i < j} \alpha_{ij} e_i \wedge e_j \in \bigwedge^2 E\) is of the form \(\alpha.e\) , where \(\alpha\) is an element of A that we will compute. The element \(\bigwedge^m u\) is a sum of terms of the form

\[
\alpha_ {h _ {1} k _ {1}} \alpha_ {h _ {2} k _ {2}} \dots \alpha_ {h _ {m} k _ {m}} e _ {h _ {1}} \wedge e _ {k _ {1}} \wedge e _ {h _ {2}} \wedge e _ {k _ {2}} \wedge \dots \wedge e _ {h _ {m}} \wedge e _ {k _ {m}}\tag{3}
\]

with \(h_j < k_j\) for \(j = 1, \ldots, m\) . Such a term is zero if it contains two equal \(e_j\) , that is, if the set \(\{h_1, k_1, \ldots, h_m, k_m\}\) is not exactly \(\{1, 2, \ldots, 2m\}\) . Moreover, if, in (3), one simultaneously interchanges \(e_{h_r}\) and \(e_{h_{r+1}}\) on the one hand, \(e_{k_r}\) and \(e_{k_{r+1}}\) on the other hand, the product does not change; hence it does not change under any permutation performed on the pairs \((h_1, k_1), \ldots, (h_m, k_m)\) . Consider then the sets (and not the sequences) \(S = \{(h_1, k_1), \ldots, (h_m, k_m)\}\) of pairs \((h_j, k_j)\) such that \(1 \leqslant h_j < k_j \leqslant 2m\) for \(j = 1, 2, \ldots, m\) ; let \(\mathscr{F}\) be the set of these pairs. For \(S \in \mathscr{F}\) , set

\[
1 ^ {0}) \varepsilon (S) = 0 \text {   si   } \left\{h _ {1}, k _ {1}, \dots , h _ {m}, k _ {m} \right\} \neq \left\{1, 2, \dots , 2 m \right\};
\]

2°) in the contrary case, \(\varepsilon(S)=1\) or \(\varepsilon(S)=-1\) depending on whether the permutation that maps \(h_{j}\) to 2j-1 and \(k_{j}\) to 2j ( \(j=1,\ldots,m\) ) is even or odd.

The preceding remarks then prove that \(\overset{m}{\wedge} u\) is equal to

\[
m! \sum_ {S \in \mathfrak {S}} \varepsilon (S) (\prod_ {(h, k) \in S} \alpha_ {h k}) e.\tag{4}
\]

We now introduce \(m(2m-1)\) indeterminates \(X_{hk}\) indexed by the pairs \((h,k)\) such that \(1\leqslant h<k\leqslant2m\) and let us call P the polynomial over \(\mathbf{Z}\) with respect to the \(X_{hk}\) defined by

\[
\mathrm{P} ((X _ {h k})) = \sum_ {\mathbf {S} \in \mathfrak {S}} \varepsilon (\mathrm{S}) (\prod_ {(h, k) \in \mathrm{S}} X _ {h k}).\tag{5}
\]

We therefore have

\[
\stackrel {m} {\wedge} u = m! \mathrm{P} ((\alpha_ {h k})). e.\tag{6}
\]

DEFINITION 1. --- Given an alternating matrix \(R = (\alpha_{ij}) (i, j = 1, \ldots, 2m)\) of even order \(2m\) on an arbitrary commutative ring A, one calls the Pfaffian of \(R\) , and denotes it \(\mathrm{Pf}(R)\) , the element \(\mathrm{P}((\alpha_{hk}))\) of A, where \(1 \leqslant h < k \leqslant 2m\) .

Example. --- Suppose:

\[
\alpha_ {1 2} = - \alpha_ {2 1} = \beta_ {1}, \alpha_ {3 4} = - \alpha_ {4 3} = \beta_ {2}, \dots , \alpha_ {2 m - 1, 2 m} = - \alpha_ {2 m, 2 m - 1} = \beta_ {m},
\]

all the others \(\alpha_{ij}\) being zero (cf. th. 1). Then the Pfaffian of \(R = (\alpha_{ij})\) is \(\beta_1\beta_2\ldots\beta_m\) .

PROPOSITION 1. --- Let R be an alternating matrix of even order 2m over a commutative ring A, and let P be a square matrix of order 2m over A. Then

\[
\operatorname{Pf} (^ {t} P. R. P) = (\det P) \operatorname{Pf} (R).\tag{7}
\]

Indeed, suppose first that A is a field of characteristic 0 and set \(R = (\alpha_{ij})\) , \(P = (\beta_{st})\) . Associate with \(R\) the bivector

\[
u = \sum_ {i <   j} \alpha_ {i j} e _ {i} \wedge e _ {j} = \frac {1}{2} \sum_ {1 \leqslant i, j \leqslant 2 m} \alpha_ {i j} e _ {i} \wedge e _ {j}
\]

of \(\wedge\) A \(^{2m}\) , where (e \(_i\) ) denotes the canonical basis of A \(^{2m}\) ; consider \(^t P\) as the matrix, with respect to the basis (e \(_i\) ), of an endomorphism \(f\) of A \(^{2m}\) . Then the bivector ( \(\overset{2}{\wedge} f\) )(u) is associated with the matrix \(^t P.R.P\) since it is equal to \(\frac{1}{2}\sum_{i,j,s,t}\beta_{is}\alpha_{ij}\beta_{jt}e_s\wedge e_t\) . Since the extension \(\wedge f\) of \(f\) to the exterior algebra \(\wedge\) A \(^{2m}\) is an endomorphism of this algebra (chap. III, § 5, no. 9), we have \(\overset{m}{\wedge}((\overset{2}{\wedge}f)(u)) = (\overset{2m}{\wedge}f)(\overset{m}{\wedge}u)\) ; since \(\overset{2m}{\wedge} f\) is the homothety with ratio det \(f\) , it therefore follows from (6) and from def. 1 that we have \(m!\) Pf( \(^t P.R.P\) ) = \(m!\) (det \(P\) )Pf( \(R\) ), whence (7) in the case under consideration. The general case follows by noting that both sides of (7) are polynomials with integer coefficients in the entries of the matrices \(R\) and \(P\) (chap. IV, § 2, no. 5, Scholium).

PROPOSITION 2. --- For every alternating matrix R of even order 2m over a commutative ring A, we have

\[
\det R = (\operatorname{Pf} (R)) ^ {2}.\tag{8}
\]

Indeed, since both sides of (8) are polynomials with integer coefficients in the elements of \(R\) , the principle of extension of algebraic identities (chap. IV, § 2, no. 5, Scholium) shows that it suffices to carry out the proof in the case where A is a field of characteristic 0 and where det \(R \neq 0\) . If \(P\) is an invertible square matrix of order \(2m\) on A, we have det( \(^{t}P.R.P\) ) = ( \(\det P\) ) \(^{2}\) ( \(\det R\) ) and \(\mathrm{Pf}(^{t}P.R.P)\) = ( \(\det P\) ) \(\mathrm{Pf}(R)\) (prop. 1), so that it suffices to prove (8) for \(^{t}P.R.P\) instead of \(R\) . By cor. 1 of th. 1, one can, by a suitable choice of \(P\) , suppose that the matrix \(^{t}P.R.P\) is of the form ( \(\alpha_{ij}\) ) where

\[
\alpha_ {1 2} = - \alpha_ {2 1} = 1, \dots , \alpha_ {2 m - 1, 2 m} = - \alpha_ {2 m, 2 m - 1} = 1,
\]

all the others \(\alpha_{ij}\) being zero (cf. Example). Now the determinant of this matrix equals 1, and its Pfaffian likewise; this completes the proof.

\subsection{Symplectic group.}

Suppose the ring A is commutative. If \(\Phi\) is an alternating bilinear form on E, the automorphisms of the module E leaving \(\Phi\) invariant are called symplectic automorphisms (or symplectic transformations) relative to \(\Phi\) , and they form a group called the symplectic group associated with \(\Phi\) ; it is sometimes denoted \(\mathbf{Sp}(\Phi)\) .

Consider in particular, on the module \(E = A^{2m}\) the alternating bilinear form \(\Phi_{0}\) whose matrix with respect to the canonical basis \((e_{i})\) of E is

\[
R _ {m} = \left( \begin{array}{c c} 0 & I _ {m} \\ - I _ {m} & 0 \end{array} \right).
\]

The symplectic automorphisms and the symplectic group relative to \(\Phi_{0}\) are simply called symplectic automorphisms and symplectic group in \(2m\) variables (over A), respectively; this group is denoted \(\mathbf{Sp}(2m, A)\) or \(\mathbf{Sp}_{2m}(A)\) . Every matrix \(A\) of a symplectic automorphism with respect to the canonical basis \((e)_i\) is called a symplectic matrix. Such a matrix is invertible, and, by formula (48) of § 1, no. 10, satisfies the relation

\[
{ } ^ { t } A . R _ { m } . A = R _ { m } .\tag{9}
\]

Conversely, if a square matrix A of order 2m over A satisfies (9), it is symplectic: it suffices indeed to prove that it is invertible; but (9) implies Pf(R \(_{m}\) ) = Pf( \(^{t}\) A . R \(_{m}\) . A) = (det A)Pf(R \(_{m}\) ) by prop. 1 of no. 2, hence det A = 1. We have at the same time proved the following proposition:

PROPOSITION 3. --- The determinant of a symplectic matrix is equal to 1.

If A is a commutative field and \(\Phi\) a nondegenerate alternating bilinear form on a vector space E of even dimension \(2m\) over A, the symplectic group associated with \(\Phi\) is isomorphic to \(\mathbf{Sp}(2m, \mathrm{A})\) , by cor. 1 of th. 1.

Exercises. --- ¶ 1) Let A be a principal commutative ring, E a free A-module of finite dimension n, Φ an alternating bilinear form on E; the ideals Aαi (1 ≤ i ≤ r) defined in th. 1 of no. 1 are called the invariant factors of Φ.

\begin{enumerate}
\def\labelenumi{\alph{enumi})}
\item
  Let F be a submodule of E, and let \(\Phi_{\mathrm{F}}\) the restriction of \(\Phi\) to F × F. Show that if A \(\beta_{i}\) ( \(1 \leqslant i \leqslant s\) ) are the invariant factors of \(\Phi_{\mathrm{F}}\) (where \(\beta_{i}\) divides \(\beta_{i+1}\) ), we have \(s \leqslant r\) and \(\beta_{i}\) is a multiple of \(\alpha_{i}\) for \(1 \leqslant i \leqslant s\) . (Reduce to the case where \(r = s = n/2\) , and use exerc. 9 b) and 9 c) of chap. VII, § 4).
\item
  Let \(E_{1}\) be a second free A-module of finite dimension, \(\Phi_{1}\) an alternating bilinear form on \(E_{1}, A\gamma_{1}, \ldots, A\gamma_{s}\) its invariant factors ( \(\gamma_{i}\) dividing \(\gamma_{i+1}\) ). In order that \(\Phi_{1}\) be the inverse image of \(\Phi\) by a linear mapping of \(E_{1}\) into E, it is necessary and sufficient that \(s \leqslant r\) and \(\gamma_{i}\) be a multiple of \(\alpha_{i}\) for \(1 \leqslant i \leqslant s\) . (Use a) and prop. 4 of chap. VII, § 4, no. 5).
\item
  Let F, G be two submodules of E, such that \(F^{0}\) (resp. \(G^{0}\) ) be a supplement of F (resp. G) in E. If the restrictions of \(\Phi\) to F and G are equivalent, show that the same holds for the restrictions of \(\Phi\) to \(F^{0}\) and to \(G^{0}\) , and that there exists an automorphism of E leaving \(\Phi\) invariant and mapping F to G.
\item
  Give an example of two submodules F, G of E, of dimension 2, such that F and G admit supplements in E and such that the restrictions \(\Phi_{\mathrm{F}}\) and \(\Phi_{\mathrm{G}}\) are equivalent, but that there exists no automorphism of E leaving \(\Phi\) invariant and mapping F to G (take \(n = 4\) ).
\end{enumerate}

\begin{enumerate}
\def\labelenumi{\arabic{enumi})}
\setcounter{enumi}{1}
\tightlist
\item
  Let \(\Phi\) be an alternating bilinear form on a finite-dimensional vector space E. Show that for every subspace M of E, the difference dim M - dim (M ∩ M⁰) is even. (First consider the case where Φ is nondegenerate).
\end{enumerate}

¶ 3) Let E be a vector space of even dimension n = 2m over a commutative field A; let Φ and Ψ be two alternating bilinear forms on E; suppose Ψ is nondegenerate. Let u and v be the linear maps from E to E* associated on the right with Φ and Ψ respectively; v is an isomorphism from E onto E*; set w = v⁻¹ ∘ u, so that w is an endomorphism of E.

\begin{enumerate}
\def\labelenumi{\alph{enumi})}
\item
  Set \(\mathbf{M}_{\mathbf{0}} = \mathbf{E}\) and by induction \(\mathbf{M}_{k+1} = w(\mathbf{M}_{k})\) for \(k > 0\) . Show that if \(\mathbf{M}_{k}^{\prime}\) is the orthogonal subspace to \(\mathbf{M}_{k}\) for \(\Phi\) , \(\mathbf{M}_{k+1}\) is the orthogonal subspace to \(\mathbf{M}_{k}^{\prime}\) for \(\Psi\) .
\item
  Let \(n_{0}\) be the dimension of \(\mathbf{M}_{0}^{\prime}\) ; we set \(m_{0}=0\) , and for \(k\geqslant1\) , we denote by \(m_{k}\) the dimension of \(\mathbf{M}_{k}\cap\mathbf{M}_{0}^{\prime}\) . Show that for \(k\geqslant1\) the dimension of \(\mathbf{M}_{k}\) is \(n-n_{0}-(m_{1}+\ldots+m_{k-1})\) and that of \(\mathbf{M}_{k}^{\prime}\) is \(n_{0}+(m_{1}+\ldots+m_{k})\) .
\item
  Show that, for all \(k \geqslant 0\) , the dimension of \(\mathbf{M}_k \cap \mathbf{M}_k'\) is \(m_k + m_{k+1} + \ldots + m_{2k}\) , and that of \(\mathbf{M}_k' \cap \mathbf{M}_{k+1}\) is \(m_{k+1} + m_{k+2} + \ldots + m_{2k+1}\) (apply exerc. 2 b) of § 3 to compute the dimensions of \(\mathbf{M}_h \cap \mathbf{M}_{2k-h}'\) and of \(\mathbf{M}_h' \cap \mathbf{M}_{2k+1-h}\) by induction on \(h\) ).
\item
  Deduce from c) that the numbers \(m_{k}\) are even (use exerc. 2).
\item
  Conclude from d) that the number of elementary divisors of the endomorphism w, corresponding to the characteristic root \(\lambda = 0\) , and having a given degree, is even (cf. chap. VII, § 5, exerc. 20).
\end{enumerate}

\begin{enumerate}
\def\labelenumi{\arabic{enumi})}
\setcounter{enumi}{3}
\item
  Let E be a vector space over a commutative field A, admitting a countable basis \((e_n)_{n\geqslant 1}\) , \(\Phi\) a nondegenerate alternating form on E. Show that there exists in E a basis \((a_n)\) such that \(\Phi(a_{2n-1}, a_{2n}) = 1\) for all \(n\geqslant 1\) , and \(\Phi(a_i, a_j) = 0\) for any other pair of indices such that \(i < j\) (reason as in exercise 13 of § 4).
\item
  For any alternating matrix \(X = (x_{ij})\) of even order \(n = 2m\) over a commutative ring, and for every index \(i\) , show that one has \(\operatorname{Pf}(X) = \sum_{j=1}^{n} (-1)^{i+j-1} \operatorname{Pf}(X_{ij}) x_{ij}\) , where \(X_{ij}\) is the matrix of order \(n - 2\) obtained by deleting from \(X\) the rows and columns with indices \(i\) and \(j\) .
\item
  Let M be a square matrix of order m over a commutative ring. Show that if we set
\end{enumerate}

\[
R = \left( \begin{array}{c c} 0 & M \\ - ^ {t} M & 0 \end{array} \right)
\]

we have \(\operatorname{Pf}(R) = \det M\) (first prove it when \(M\) is invertible, using formula (7)).

\begin{enumerate}
\def\labelenumi{\arabic{enumi})}
\setcounter{enumi}{6}
\item
  Let A be a commutative field, \(\mathfrak{A}(A)\) the set of alternating matrices of order \(2m\) on A; let I be a mapping from \(\mathfrak{A}(A)\) into A such that for every matrix \(R\in\mathfrak{A}(A)\) and for every matrix \(P\) of order \(2m\) over A, one has \(\mathrm{I}(^{t}PRP)=(\det P)^{h}\mathrm{I}(R)\) , where \(h\) is a rational integer. Show that \(\mathrm{I}(R)=c(\mathrm{Pf}R)^{h}\) , where \(c\in\mathrm{A}\) (using formula (7) and th. 1 of no. 1).
\item
  Let \(P\) , \(Q\) be two alternating square matrices of even order \(2m\) on a commutative ring A. Let \(\varphi(X) = \mathrm{Pf}(P - XQ)\) ; show that, if \(Q\) is invertible, one has \(\varphi(Q^{-1}P)=0\) . (First consider the case where A is a field; then deduce from exerc. 3 that the minimal polynomial of the matrix \(Q^{-1}P\) divides \(\varphi(X)\) , by passing to an algebraically closed extension of A, and noting that \(\varphi^{2}\) is the characteristic polynomial of \(Q^{-1}P\) up to a scalar factor).
\end{enumerate}

¶ 9) Let E be a vector space of dimension n over a commutative field A, Φ an alternating bilinear form on E, Ψ a sesquilinear Hermitian form (for an involutive automorphism of A) on E, and P and Q the respective matrices of Φ and Ψ with respect to the same basis of E.

\begin{enumerate}
\def\labelenumi{\alph{enumi})}
\item
  We assume \(\Psi\) nondegenerate. Show that the number of elementary divisors of \(Q^{-1}P\) corresponding to the characteristic root 0 and having the same even degree, is an even number (method of exerc. 3).
\item
  We assume that \(\Phi\) be nondegenerate (which implies that \(n\) is even). Show that the number of elementary divisors of \(P^{-1}Q\) corresponding to the characteristic root 0, and having the same odd degree, is an even number (same method).
\end{enumerate}

\begin{enumerate}
\def\labelenumi{\arabic{enumi})}
\setcounter{enumi}{9}
\tightlist
\item
  Let \(\omega_{2m}(\mathbf{F}_q)\) be the order of the symplectic group \(\mathbf{Sp}(2m, \mathbf{F}_q)\) over the finite field \(\mathbf{F}_q\) . Show that if \(h_{2m}\) is the number of pairs of vectors \((x, y)\) of \(\mathbf{F}_q^{2m}\) such that \(\Phi_0(x, y) = 1\) (notations of no. 3), one has \(\omega_{2m}(\mathbf{F}_q) = h_{2m}\omega_{2m-2}(\mathbf{F}_q)\) ; deduce from this
\end{enumerate}

\[
\omega_ {2 m} (\mathbf {F} _ {q}) = (q ^ {2 m} - 1) q ^ {2 m - 1} (q ^ {2 m - 2} - 1) q ^ {2 m - 3} \dots (q ^ {2} - 1) q.
\]

\begin{enumerate}
\def\labelenumi{\arabic{enumi})}
\setcounter{enumi}{10}
\tightlist
\item
  Let A be a commutative field. Show that every transformation \(u\) belonging to the symplectic group \(\mathbf{Sp}(2m, \mathrm{A})\) is a product of transvections belonging to this group (called symplectic transvections; cf.§ 4, exerc. 6). (Reason by induction on \(m\) , by showing that if \(x, y\) are two non-orthogonal vectors of \(\mathrm{E} = \mathrm{A}^{2m}\) , there exists a product \(\nu\) of symplectic transvections such that \(\nu u\) leaves \(x\) and \(y\) invariant). Deduce a new proof of prop. 3 of no. 3.
\end{enumerate}

*12) Let A be a commutative field, E a vector space of even dimension \(n=2m\) over A, \(\Phi\) a nondegenerate alternating bilinear form on E. Show that, for every similarity \(u\) for the form \(\Phi\) ( \(\S 6\) , no. 5), with multiplier \(\alpha\) , we have det \(u=\alpha^{m}\) (using formula (7)).*

¶ 13) We assume that A is a commutative field of characteristic 0, E a vector space of dimension 2m over A, Φ a nondegenerate alternating bilinear form on E. We identify the inverse form \(\widehat{\Phi}\) of Φ with a bivector \(\Gamma\in\bigwedge E\) , so that for any symplectic basis \((e_{i})_{1\leqslant i\leqslant2m}\) of E (for \(\Phi\) ), indexed such that \(\Phi(e_{i},e_{j})=\Phi(e_{m+i},e_{m+j})=0\) , \(\Phi(e_{i},e_{m+j})=\delta_{ij}(1\leqslant i\leqslant m,1\leqslant j\leqslant m)\) , one has

\[
\Gamma = e _ {1} \wedge e _ {m + 1} + e _ {2} \wedge e _ {m + 2} + \dots + e _ {m} \wedge e _ {2 m}.
\]

We say that a nonzero decomposable \(p\) -vector \(z \in \bigwedge E\) is isotropic (resp. totally isotropic) for \(\Phi\) if the vector subspace \(V_z\) corresponding to \(z\) (chap. III, § 7, no. 3) is isotropic (resp. totally isotropic).

\begin{enumerate}
\def\labelenumi{\alph{enumi})}
\item
  If z is a decomposable p-vector \(\neq0\) , 2r the dimension of a supplement of \(V_{z}\cap V_{2}^{0}\) with respect to \(V_{z}\) , show that \(m-p+r\) is the largest of the integers h such that one has \(z\wedge\Gamma^{h}\neq0\) , denoting by \(\Gamma^{h}\) the h-th power of \(\Gamma\) in the exterior algebra \(\wedge E\) (using prop. 2 of § 4, no. 2).
\item
  Show that every bivector \(x \in \bigwedge E\) can be written in the form \(\lambda \Gamma + x_{1}\) , where \(\lambda\) is a scalar and \(x_{1}\) a linear combination of totally isotropic decomposable bivectors. (Reduce to the case where \(x\) is decomposable, and note that if \((e_{i})\) is a symplectic basis, \((e_{1} + e_{2}) \wedge (e_{m+1} - e_{m+2})\) is totally isotropic).
\item
  If \(p \leqslant m\) , show that every \(p\) -vector \(z \in \hat{\Lambda}\) E can be written \(z = x \wedge \Gamma + z_1\) , where \(x\) is a \((p - 2)\) -vector and \(z_1\) a linear combination of \(p\) -vectors that are decomposable and totally isotropic. (Reduce to the case where \(z\) is decomposable, and reason by induction on \(p\) . One can thus reduce to the case where, \((e_i)\) being a symplectic basis, one has \(z = e_1 \wedge e_2 \wedge \ldots \wedge e_{p-1} \wedge e_{m+p-1}\) , and consider the bivectors \((e_{p-1} + e_{p+i}) \wedge (e_{m+p-1} - e_{m+p+i})\) for \(0 \leqslant i \leqslant m - p\) ).
\item
  For \(1 \leqslant p \leqslant m\) , show that every \((m+p)\) -vector \(z \in \wedge\) E can be written \(z = y \wedge \Gamma^p\) , where \(y\) is a \((m-p)\) -vector. (Reduce to the case where \(z\) is decomposable; if \(2r\) is the dimension of a supplement of \(\mathrm{V}_z \cap \mathrm{V}_z^0\) with respect to \(\mathrm{V}_z\) , distinguish two cases, according as \(r = p\) or \(r > p\) ; in the second case, reason by induction on \(r\) , in the same manner as in \(c)\) . Deduce from this that the map \(y \to y \wedge \Gamma^p\) of \(\wedge^{m-p}\) E into \(\wedge^{m+p}\) E is bijective (note that the two spaces have the same dimension). If y is a decomposable \((m-p)\) -vector, show that \(z = y \wedge \Gamma^{p}\) is decomposable and that \(V_{z} = V_{y}^{0}\) .
\end{enumerate}

\begin{enumerate}
\def\labelenumi{\alph{enumi})}
\setcounter{enumi}{4}
\tightlist
\item
  Deduce from d) that, for \(p \leqslant m\) , the subspace \(\hat{E}\) is the direct sum of the homogeneous component of degree p of the two-sided ideal c generated by \(\Gamma\) in \(\hat{E}\) , and of the subspace \(R_{p}\) of p-vectors \(z_{1}\) such that
\end{enumerate}

\(z_{1} \wedge \Gamma^{m-p+1} = 0\) (for \(z \in \bigwedge E\) , apply \(d\) ) to the \((2m - p + 2)\) -vector \(z \wedge \Gamma^{m-p+1}\) ). Using \(c\) ), prove that \(R_{p}\) is generated by the decomposable totally isotropic \(p\) -vectors, and by using \(d\) ), show

that the map \(x \to x \wedge \Gamma\) of \(\wedge E\) into \(\wedge E\) is injective, and that \(R_p\) is of dimension \(\binom{2m}{p} - \binom{2m}{p-2}\) .

\begin{enumerate}
\def\labelenumi{\alph{enumi})}
\setcounter{enumi}{5}
\tightlist
\item
  Show that if \(p \leqslant m\) , a necessary and sufficient condition for a p-vector z to be of the form \(x \wedge \Gamma\) is that, for every decomposable totally isotropic m-vector u, one has \(z \wedge u = 0\) . (Show that, if this condition is satisfied, and if \(z \in R_{p}\) , one has \(z \wedge y = 0\) for all \((2m - p)\) -vector y; for this, put y in the form \(x \wedge \Gamma^{m-p}\) where x is a p-vector, then apply c) to x, and note that if \(x_{1}\) is a decomposable totally isotropic p-vector, \(x_{1} \wedge \Gamma^{m-p}\) is a decomposable \((2m - p)\) -vector which can be written \(u_{1} \wedge \varphi_{1}\) , where \(u_{1}\) is a decomposable totally isotropic \(m\) -vector).
\end{enumerate}

Let a be the annihilator of the ideal c in \(\wedge\) E; a is the direct sum of the subspaces \(R_{m}\) , \(R_{m-1} \wedge \Gamma, \ldots, R_{1} \wedge \Gamma^{m-1}\) and \(K\Gamma^{m}\) . Show that the annihilator of a in \(\wedge\) E is equal to c (cf.§ 2, exerc. 4 b)).

\begin{enumerate}
\def\labelenumi{\alph{enumi})}
\setcounter{enumi}{6}
\tightlist
\item
  For every symplectic automorphism u of E (for \(\Phi\) ), let \(\bar{u}\) be the canonical extension of u to an automorphism of the algebra \(\wedge\) E (chap. III, § 5, no. 9). Show that the only elements of \(\wedge\) E invariant under all the automorphisms \(\bar{u}\) are the linear combinations of 1, \(\Gamma\) , \(\Gamma^{2}\) , \(\ldots\) , \(\Gamma^{m}\) with coefficients in K. (If ( \(e_{i}\) ) is a symplectic basis, write that an element of \(\wedge\) E is invariant under the symplectic transvections (exerc. 11) corresponding to the hyperplanes orthogonal to the \(e_{i}\) ; then consider the symplectic transformations \(u_{ij}\) defined by \(u_{ij}(e_{i}) = e_{j}\) , \(u_{ij}(e_{j}) = e_{i}\) , \(u_{ij}(e_{m+i}) = e_{m+j}\) , \(u_{ij}(e_{m+j}) = e_{m+i}\) , \(u_{ij}(e_{k}) = e_{k}\) for every other index k).
\end{enumerate}

\begin{enumerate}
\def\labelenumi{\arabic{enumi})}
\setcounter{enumi}{13}
\item
  Let A be a commutative field of characteristic 2, E the vector space \(\mathrm{A}^{2m}\) ; we identify the symplectic group \(\mathbf{Sp}(2m,\mathrm{A})\) with the group of symplectic matrices \(U\) , thus satisfying the relation \(^t U.R.U = R\) , where \(R = \left( \begin{array}{cc}0 & I_m\\ I_m & 0 \end{array} \right)\) . We set \(D = \left( \begin{array}{cc}0 & 0\\ I_m & 0 \end{array} \right)\) ; show that every symplectic matrix \(U\) such that \(I + U\) is invertible can be written in a unique way in the form \((^t D + S)^{-1}(D + S)\) , where \(S\) is a symmetric matrix such that \(D + S\) is invertible, and conversely (cf.§ 4, exerc. 11).
\item
  Let A be a commutative field, E a vector space of dimension 2m over A, Φ a nondegenerate alternating bilinear form on E. a) Extend to endomorphisms u of E such that \(u^{*}u = uu^{*}\) ( \(u^{*}\) being the adjoint of u with respect to Φ) the results of § 4, exerc. 12 a) and b).
\end{enumerate}

\begin{enumerate}
\def\labelenumi{\alph{enumi})}
\setcounter{enumi}{1}
\item
  We assume that \(u^{*} = u\) . With the notations of exerc. 12 of § 4, show that if M is a minimal element of the set \(\mathfrak{M}\) of non-isotropic subspaces contained in \(G(p, p)\) and stable under \(u\) , then either M is an indecomposable submodule of \(E_{u}\) , or it is a direct sum of two such isomorphic submodules (reason as in exerc. 12 c) of § 4). Show by examples (with \(p(X) = X - 1\) ) that the two cases can occur.
\item
  We assume \(u^* = -u\) , with A not of characteristic 2. With the same notations, let \(p^h\) be the minimal polynomial of the restriction of \(u\) to M, and let \(d\) be the degree of \(p\) . Show that if \(d(h - 1)\) is odd, M is an indecomposable submodule of \(E_u\) ; if on the contrary \(d(h - 1)\) is even, M is either indecomposable or a direct sum of two isomorphic indecomposable submodules.
\item
  We assume that \(u^{*}u = 1\) (in other words that \(u \in \mathbf{Sp}(\Phi)\) ). With the notations of c), if \(p(X)\) does not divide \(X^{d(h-1)} - 1\) , or if \(p(X) = X - 1\) and \((-1)^{h} \neq -1\) , M is an indecomposable submodule of \(E_{u}\) ; otherwise, M is either indecomposable or a direct sum of two isomorphic indecomposable submodules.
\end{enumerate}

\section{Special properties of Hermitian forms}

\subsection{Orthogonal bases.}

DEFINITION 1. --- Let Φ be a Hermitian form on E. A basis ( \(e_{i}\) ) of E is said to be orthogonal for Φ if any two elements of this basis are orthogonal for Φ.

If moreover \(\Phi(e_i, e_i) = 1\) for all \(i\) , the basis \((e_i)\) is said to be orthonormal.

Let \((e_i)\) be an orthogonal basis; if we set \(\Phi(e_i, e_i) = \alpha_i\) , one has

\[
\Phi (\sum_ {i} \xi_ {i} e _ {i}, \sum_ {i} \eta_ {i} e _ {i}) = \sum_ {i} \xi_ {i} \alpha_ {i} \overline {{\eta}} _ {i}.\tag{1}
\]

Lemma 1. --- Suppose that A is a field and \(\Phi\) a Hermitian form \(\neq 0\) on E. If all vectors of E are isotropic, then A is a commutative field of characteristic 2, the antiautomorphism J is the identity, and \(\Phi\) is alternating.

Indeed, if one expands \(\Phi(x+y, x+y)=0\) , one obtains, taking the hypotheses into account, \(\Phi(x,x)=\Phi(y,y)=0\) , the relation \(\Phi(x,y)=-\overline{\Phi(x,y)}\) for all \(x,y\) in E. Since \(\Phi\) is \(\neq 0\) , there exist \(x,y\) in E such that \(\Phi(x,y)=1\) . Writing that \(\Phi(\lambda x,y)=-\overline{\Phi(\lambda x,y)}\) , one obtains \(\bar{\lambda}=-\lambda\) for all \(\lambda\in A\) . Taking first \(\lambda=1\) , one sees that A is of characteristic 2; the relation \(\bar{\lambda}=-\lambda\) then shows that J is the identity, hence A is commutative and \(\Phi\) is alternating.

THEOREM 1. --- Suppose that A is a field and E a vector space of finite dimension n over A. Then, for every Hermitian form Φ on E, E admits an orthogonal basis, except if the following conditions are simultaneously satisfied:

\begin{enumerate}
\def\labelenumi{(\Alph{enumi})}
\setcounter{enumi}{2}
\tightlist
\item
  A is commutative of characteristic 2, the antiautomorphism J is the identity, Φ is alternating and nonzero.
\end{enumerate}

Reason by induction on \(n\) , the result being evident for \(n=0\) . We may suppose \(\Phi\neq0\) . If (C) is not satisfied, Lemma 1 shows that there exists an element \(x\in\mathbf{E}\) such that \(\Phi(x,x)\neq0\) . Let H be the subspace of E orthogonal to \(x\) ; it has dimension \(\geqslant n-1\) , and, since \(x \notin H\) , \(H\) is exactly of dimension \(n-1\) . If the restriction \(\Psi\) of \(\Phi\) to \(H\) does not satisfy (C), then by the induction hypothesis there exists a basis \((e_2, \ldots, e_n)\) of \(H\) which is orthogonal for \(\Psi\) ; by setting \(e_1 = x\) one obtains an orthogonal basis \((e_1, e_2, \ldots, e_n)\) of \(E\) . It remains to examine the case where \(A\) is a commutative field of characteristic 2, where \(J\) is the identity, and where \(\Psi\) is alternating and \(\neq 0\) . There then exists \(y\) , \(z\) in \(H\) such that \(\Psi(y, z) \neq 0\) ; set \(e_1 = x + y\) ; for \(x + \lambda z (\lambda \in A)\) to be orthogonal to \(e_1\) , it is necessary and sufficient that we have \(0 = \Phi(x + y, x + \lambda z) = \Phi(x, x) + \lambda \Psi(y, z)\) , a condition which determines \(\lambda\) in a unique way; the scalar \(\lambda\) having been chosen thus, we have \(\Phi(x + \lambda z, x + \lambda z) = \Phi(x, x) \neq 0\) , hence the restriction \(\Psi'\) of \(\Phi\) to the subspace \(H'\) of \(E\) orthogonal to \(e_1\) is not alternating; one may consequently apply the induction hypothesis to \(H'\) , which proves the theorem.

When (C) is satisfied, there evidently exists no orthogonal basis for \(\Phi\) .

COROLLARY 1. --- With the notations of th. 1, one supposes moreover that (C) is not satisfied, that \(\Phi\) is nondegenerate and that, for every \(x \in E\) , there exists \(\rho \in A\) such that \(\Phi(x, x) = \rho \bar{\rho}\) . Then E admits an orthonormal basis for \(\Phi\) .

Indeed, let \((e_i)\) ( \(i = 1, \ldots, n\) ) be an orthogonal basis of E. Put \(\Phi(e_i, e_i) = \alpha_i\) . We have \(\alpha_i \neq 0\) for \(i = 1, \ldots, n\) since \(\Phi\) is nondegenerate. By hypothesis, there exist elements \(\beta_i\) of A such that \(\alpha_i = \beta_i \overline{\beta_i}\) for \(i = 1, \ldots, n\) ; we have \(\beta_i \neq 0\) . By setting \(f_i = \beta_i^{-1} e_i\) , one has \(\Phi(f_i, f_i) = \beta_i^{-1} \alpha_i \overline{\beta_i^{-1}} = 1\) for all \(i\) , and \(\Phi(f_i, f_j) = 0\) for \(i \neq j\) . Hence \((f_i)\) is an orthonormal basis.

Remark. --- The last hypothesis of the corollary is satisfied when J is the identity and every element of A is the square of an element of A (for example, when A is algebraically closed).

COROLLARY 2. --- Let A be a field and R a Hermitian matrix of order n and rank r over A. Then, unless the following condition holds:

(C') A is commutative of characteristic 2, J is the identity, R is alternating and nonzero,

there exists an invertible matrix P of order n over A such that

\[
{ } ^ { t } P . R . \overline { { P } } = \left( \begin{array} { c c c c c } \alpha _ { 1 } & 0 \dots 0 \dots 0 \\ 0 & \alpha _ { 2 } \dots 0 \dots 0 \\ \cdot & \cdot & \cdot & \cdot & \cdot \\ 0 & 0 \dots \alpha _ { r } \dots 0 \\ 0 & 0 \dots 0 \dots 0 \\ \cdot & \cdot & \cdot & \cdot & \cdot \\ 0 & 0 \dots 0 \dots 0 \end{array} \right)
\]

\(where \overline{\alpha}_{i} = \alpha_{i} \neq 0 \text{ for } i = 1, \ldots, r.\)

PROPOSITION 1. --- Suppose that A is a commutative field. Let \(\Phi\) be a Hermitian form on E, and \((x_{n})\) ( \(n=1,2,\ldots\) ) a sequence (finite or infinite) of linearly independent vectors of E such that, for every \(n\) , the subspace \(\mathrm{E}_{n}=\mathrm{A}x_{1}+\cdots+\mathrm{A}x_{n}\) is non-isotropic. Let \(\mathrm{D}_{jn}\) ( \(j\leqslant n\) ) be the cofactor of \(\Phi(x_{j},x_{n})\) in the matrix \((\Phi(x_{s},x_{t}))_{(s,t=1,\ldots,n)}\) . We then have \(\mathrm{D}_{nn}\neq0\) for all \(n\) . Set

\[
e _ {n} = \sum_ {j = 1} ^ {n} \mathrm{D} _ {n n} ^ {- 1} \mathrm{D} _ {j n} x _ {j}.\tag{2}
\]

Then, for every \(n\) , \((e_{1},\ldots,e_{n})\) is an orthogonal basis of \(\mathrm{E}_{n}\) and we have

\[
\Phi (e _ {n}, e _ {n}) = \mathrm{D} _ {n n} ^ {- 1} \mathrm{D} _ {n + 1, n + 1}.\tag{3}
\]

Indeed, since the restriction of \(\Phi\) to \(\mathbf{E}_{n-1}\) is nondegenerate, we have \(\mathrm{D}_{nn} \neq 0\) ( \(\S 2\) , prop. 3); note that we have \(\mathrm{D}_{11} = 1\) since the determinant of the empty matrix is equal to 1. Formulas (2) imply first that we have \(e_n \equiv x_n\) (mod. \(\mathbf{E}_{n-1}\) ) for all \(n\) , hence that the \(e_n\) are linearly independent, and that \((e_1, \ldots, e_n)\) is a basis of \(\mathbf{E}_n\) . For every \(j < n\) , one has

\[
\Phi (e _ {n}, x _ {j}) = \mathrm{D} _ {n n} ^ {- 1} \sum_ {k = 1} ^ {n} \mathrm{D} _ {k n} \Phi (x _ {k}, x _ {j}) = 0
\]

(chap. III, § 6, no. 1, formula (12)); hence \(e_n\) is orthogonal to \(\mathrm{E}_{n-1}\) , and in particular to \(e_j\) for \(j < n\) . On the other hand, we have

\[
\begin{array}{r l} \Phi (e _ {n}, e _ {n}) & = \Phi (e _ {n}, \sum_ {j = 1} ^ {n} \mathrm{D} _ {n n} ^ {- 1} \mathrm{D} _ {j n} x _ {j}) = \Phi (e _ {n}, x _ {n}) = \Phi (\sum_ {j = 1} ^ {n} \mathrm{D} _ {n n} ^ {- 1} \mathrm{D} _ {j n} x _ {j}, x _ {n}) \\ & = \mathrm{D} _ {n n} ^ {- 1} \sum_ {j = 1} ^ {n} \mathrm{D} _ {j n} \Phi (x _ {j}, x _ {n}) = \mathrm{D} _ {n n} ^ {- 1} \mathrm{D} _ {n + 1, n + 1} \end{array}
\]

(chap. III, § 6, no. 1, formula (10)). This proves our assertions.

With the notations of prop. 1, one says that the sequence \((e_n)\) is obtained from the sequence \((x_n)\) by the Gram-Schmidt orthogonalization process.

PROPOSITION 2. --- Let \(\Phi\) be a Hermitian form on E, and \((e_i)\) ( \(i = 1, \ldots, n\) ) an orthogonal basis (resp. orthonormal basis) of E for \(\Phi\) . Then, for every \(p \geqslant 0\) , the basis of \(\bigotimes^{p} \mathbf{E}\) formed by the \(e_{i_1} \otimes \cdots \otimes e_{i_p}\) and the basis \((e_h)\) of \(\bigwedge^{p} \mathbf{E}\) (where H ranges over the set of subsets with \(p\) elements of {[}1, n{]}; cf. chap. III, § 5, no. 6) are orthogonal (resp. orthonormal) for the extensions of \(\Phi\) to \(\bigotimes^{p} \mathbf{E}\) and \(\bigwedge^{p} \mathbf{E}\) respectively (§ 1, no. 9). If, moreover, the maps associated with \(\Phi\) are bijective, the basis \((e_i')\) of \(\mathbf{E}^*\) dual of \((e_i)\) is orthogonal (resp. orthonormal) for the inverse form \(\hat{\Phi}\) of \(\Phi\) (§ 1, no. 7).

The assertions concerning \(\widehat{\otimes}\) E and \(\wedge\) E follow immediately from formulas (35) and (37) of § 1, no. 9. The one relating to the inverse form follows from the fact that the matrix of \(\widehat{\Phi}\) with respect to \((e_i')\) is the inverse of the matrix of \(\Phi\) with respect to \((e_i)\) ( \(§ 1\) , no. 10).

\subsection{Unitary group and orthogonal group.}

Let \(\Phi\) be a Hermitian form on E; the automorphisms of the A-module E that leave \(\Phi\) invariant are called unitary automorphisms (or unitary transformations) relative to \(\Phi\) , and their group is called the unitary group associated with \(\Phi\) ; it is denoted \(\mathbf{U}(\Phi)\) . Given a quadratic form \(Q \neq 0\) on E, the automorphisms of the A-module E which leave Q invariant are called orthogonal automorphisms (or orthogonal transformations) relative to Q; their group is called the orthogonal group associated with Q; it is denoted by \(\mathbf{O}(Q)\) .

Every orthogonal transformation for a quadratic form Q is unitary for the bilinear form associated with Q. The converse is true when the scalar 2 is not equal to 0 or a zero divisor in A (§ 3, no. 4, (13)), for example if A is a field of characteristic ≠ 2.

Consider, in particular, on the module \(E = A^{n}\) the Hermitian form \(\Phi_{0}\) whose matrix with respect to the canonical basis \((e_{i})\) of E is the identity matrix \(I_{n}\) . The unitary automorphisms associated with \(\Phi_{0}\) are simply called unitary automorphisms (or transformations) in n variables; their group is called the unitary group in n variables and is sometimes denoted \(\mathbf{U}(n, \mathbf{A})\) or \(\mathbf{U}_{n}(\mathbf{A})\) . The matrix U of a unitary automorphism with respect to \((e_{i})\) is called a unitary matrix. Such a matrix is invertible, and satisfies, according to formula (48) of § 1, no. 10, the relation

\[
{ } ^ { t } U . \overline { { U } } = I _ { n } ;\tag{4}
\]

Conversely, if A is a commutative ring or a field, a matrix U satisfying (4) is invertible, and is then unitary.

When J is the identity and 2 is neither equal to 0 nor a divisor of 0 in A, one uses the terms orthogonal group in n variables, orthogonal automorphism (or orthogonal transformation) in n variables, and orthogonal matrix instead of the preceding terms, and one writes \(\mathbf{O}(n,\mathrm{A})\) (resp. \(\mathbf{O}_n(\mathrm{A})\) ) instead of \(\mathbf{U}(n,\mathrm{A})\) (resp. \(\mathbf{U}_n(\mathrm{A})\) ). Relation (4) is then written

\[
{ } ^ { t } U . U = I _ { n }\tag{5}
\]

and, since A is commutative, it is a necessary and sufficient condition for U to be an orthogonal matrix.

PROPOSITION 3. --- Suppose that A is a commutative field and that E is of finite dimension \textgreater{} 0. Let Φ be a nondegenerate Hermitian form on E. The map u → det u is a homomorphism from the unitary group U(Φ) associated with Φ onto the multiplicative subgroup H of A consisting of the elements ρ such that ρρ̄ = 1 (a subgroup reduced to \{1, -1\} when J is the identity).

Indeed, let \(u\) be an element of \(\mathbf{U}(\Phi)\) , \(U\) its matrix with respect to a basis of E, and \(R\) the matrix of \(\Phi\) with respect to this basis. The relation \(R = {}^t U.R.\overline{U}\) ( \(\S 1\) , no. 10, formula (48)) shows that one has \((\det U)(\det \overline{U}) = 1\) since \(R\) is invertible; whence \((\det u)(\overline{\det u}) = 1\) . The homomorphism \(u \to \det u\) maps \(\mathbf{U}(\Phi)\) onto H. Indeed, when A is of characteristic 2 and J is the identity, H is reduced to the element 1. Otherwise there exists an orthogonal basis \((e_i)(i = 1,\ldots,n)\) of E (th. 1); for every \(\rho \in A\) such that \(\rho \bar{\rho} = 1\) , let \(u\) be the automorphism of E defined by \(u(e_1) = \rho e_1\) and \(u(e_i) = e_i\) for \(i = 2,\ldots,n\) ; then \(u\) is unitary and det \(u = \rho\) , whence the proposition.

Under the conditions of prop. 3, the kernel of the homomorphism \(u \to \det u\) is a normal subgroup of \(\mathbf{U}(\Phi)\) , which is called the special unitary group associated with \(\Phi\) ; it is sometimes denoted \(\mathbf{SU}(\Phi)\) .

When J is the identity and A is not of characteristic 2, this group is also called the special orthogonal group associated with \(\Phi\) (or with the quadratic form \(Q(x) = \Phi(x, x)\) ) and is sometimes denoted SO(Q).

If E = A \(^{n}\) and if Φ is the form whose matrix with respect to the canonical basis of E is the identity matrix, one uses the notations SU(n, A) or SU \(_{n}\) (A) and SO(n, A) or SO \(_{n}\) (A).

\subsection{Orthogonal projectors and involutions.}

Throughout this no., one assumes that the scalar 2 is \(invertible\) in A (for example that A is a field of characteristic \(\neq 2\) ), and that \(\Phi\) is a Hermitian form \(nondegenerate\) on E. One denotes by \(\frac{1}{2}\) the inverse of 2.

Lemma 2. --- For an endomorphism u of E to be such that \(u^{2}=1\) , it is necessary and sufficient that \(\frac{1}{2}(1-u)\) be a projector in E; then u is the difference of the two projectors \(\frac{1}{2}(1+u)\) and \(\frac{1}{2}(1-u)\) .

Indeed, in the ring \(\mathfrak{L}(\mathrm{E})\) , the relation \(\left(\frac{1}{2}(1 - u)\right)^2 = \frac{1}{2}(1 - u)\) is equivalent to \(u^2 = 1\) . The rest is trivial.

An endomorphism \(u\) of E such that \(u^2 = 1\) (which is then necessarily an automorphism of E equal to its inverse) is called an involution. Set \(v = \frac{1}{2}(1 - u)\) , \(U^- = v(E)\) , \(U^+ = v^{-1}(0)\) ( \(= w(E)\) by setting \(w = \frac{1}{2}(1 + u)\) ); one knows that E is the direct sum of \(U^+\) and of \(U^-\) (chap. VIII, § 1, no. 1), and one has \(u(x) = x\) in \(U^+\) , \(u(x) = -x\) in \(U^-\) . When A is a field and E is of finite dimension, it follows, since A is of characteristic \(\neq 2\) , that the only eigenvectors \(\neq 0\) of \(u\) are the elements ≠ 0 in U⁺ or in U⁻; they correspond respectively to the eigenvalues +1 and −1.

PROPOSITION 4. --- Let \(u \in \mathbf{GL}(\mathrm{E})\) be an involution. The following properties are equivalent:

\begin{enumerate}
\def\labelenumi{\alph{enumi})}
\item
  u belongs to the unitary group associated with \(\Phi\) ;
\item
  the submodules \(U^{+} = \frac{1}{2}(1 + u)(E)\) and \(U^{-} = \frac{1}{2}(1 - u)(E)\) are orthogonal (and consequently non-isotropic).
\end{enumerate}

Moreover, if A is a field and E is finite-dimensional, properties a) and b) are equivalent to:

\begin{enumerate}
\def\labelenumi{\alph{enumi})}
\setcounter{enumi}{2}
\tightlist
\item
  \(u = u^{*}\) .
\end{enumerate}

Indeed, for \(x \in \mathrm{U}^{+}\) and \(y \in \mathrm{U}^{-}\) , the relation \(\Phi(u(x), u(y)) = \Phi(x, y)\) gives \(2\Phi(x, y) = 0\) , hence \(a)\) implies \(b)\) . Conversely, we obviously have \(\Phi(u(x), u(y)) = \Phi(x, y)\) when \(x\) and \(y\) are both in \(\mathrm{U}^{+}\) or both in \(\mathrm{U}^{-}\) , and, given \(b)\) , this relation is still true when one of them is in \(\mathrm{U}^{+}\) and the other in \(\mathrm{U}^{-}\) ; since E is the direct sum of \(\mathrm{U}^{+}\) and \(\mathrm{U}^{-}\) , we see that \(b)\) implies \(a)\) . Finally, when E is a finite-dimensional vector space, the adjoint \(u^{*}\) is defined since \(\Phi\) is nondegenerate; the relation \(a)\) is equivalent to \(uu^{*} = 1\) ( \(§ 1, n^{\circ} 8\) , cor. of prop. 8); since \(u^{2} = 1\) by hypothesis, \(a)\) and \(c)\) are equivalent.

COROLLARY 1. --- Suppose that A is a field and that E is finite-dimensional. The map \(u \to \frac{1}{2}(1 + u)(E)\) is a bijection from the set of involutions u belonging to the unitary group associated with \(\Phi\) onto the set of non-isotropic subspaces of E; the subspace U \(^+\) corresponding to u is the set of elements of E invariant under u.

By prop. 4, it suffices to show that every non-isotropic subspace M of E is the set of vectors invariant under an involution \(u \in \mathbf{U}(\Phi)\) , and that this involution is unique. Now E is the direct sum of M and of \(\mathbf{M}^0\) ( \(\S 4\) , no. 1, cor. of prop. 1), and one necessarily has \(u(x) = x\) for \(x \in \mathbf{M}\) and \(u(x) = -x\) for \(x \in \mathbf{M}^0\) by virtue of prop. 4; these relations determine \(u\) uniquely, and the endomorphism \(u\) thus determined evidently answers the question (prop. 4).

One says that the involution \(u\) thus determined is the symmetry with respect to the non-isotropic subspace M.

COROLLARY 2. --- Suppose that A is a field and that E is of finite dimension. For a projector v in E to be such that v(E) and \(v^{-1}(0)\) are orthogonal (and consequently non-isotropic), it is necessary and sufficient that v = v*.

It suffices to apply prop. 4 to the involution \(u = 1 - 2v\) .

A projector satisfying the condition of corollary 2 is called an orthogonal projector for \(\Phi\) .

\subsection{Symmetries in the orthogonal group.}

Unless expressly stated otherwise, one assumes, in this no., that A is a commutative field of characteristic ≠ 2, and that Φ is the symmetric bilinear form associated with a nondegenerate quadratic form Q on E. Recall that one has Φ(x, x) = 2Q(x) for x ∈ E (§ 3, no. 4).

Let H be a non-isotropic hyperplane in E, and u the symmetry with respect to H (no. 3). Let \(a \neq 0\) be a vector orthogonal to H; one has by hypothesis \(u(a) = -a\) . Every vector \(x \in E\) can be written in one and only one way in the form \(x = \lambda a + y\) with \(\lambda \in A\) and \(y \in H\) ; since \(a\) and \(y\) are orthogonal, one has \(\Phi(x, a) = \lambda \Phi(a, a)\) , whence, since \(a\) is non-isotropic ( \(\S 4\) , no. 1, cor. of prop. 1), \(\lambda = \Phi(x, a) \Phi(a, a)^{-1}\) . This being so, one has

\[
u (x) = \lambda u (a) + u (y) = - \lambda a + y = x - 2 \lambda a,
\]

whence

\[
u (x) = x - 2 \Phi (x, a) \Phi (a, a) ^ {- 1}. a = x - \Phi (x, a) Q (a) ^ {- 1}. a.\tag{6}
\]

One notes that the last member of (6) retains a meaning when A is a field of characteristic 2, and a is a non-singular vector of E; one verifies immediately that one still has then Q(u(x)) = Q(x) for all x ∈ E, in other words u ∈ O(Q). One also says that the involution u thus defined is the symmetry with respect to the hyperplane orthogonal to a (cf. exerc. 28).

PROPOSITION 5. --- Suppose the vector space E is of finite dimension n. The orthogonal group \(\mathbf{O}(\mathrm{Q})\) associated with Q is then generated by the symmetries with respect to the non-isotropic hyperplanes of E.

The proposition being evident for \(n=0\) , we shall reason by induction on \(n\) . Let \(u\) be an orthogonal transformation of E, and let \(x\) be a non-isotropic vector of E (Lemma 1); we distinguish three cases:

\begin{enumerate}
\def\labelenumi{\alph{enumi})}
\item
  Suppose first that \(u(x) = x\) . Then the hyperplane H orthogonal to \(x\) is non-isotropic, and we have \(u(\mathrm{H}) = \mathrm{H}\) . The restriction \(u'\) of \(u\) to H therefore belongs to the orthogonal group \(\mathbf{O}(Q')\) associated with the restriction \(Q'\) of Q to H. The induction hypothesis implies, since \(Q'\) is nondegenerate, that one has \(u' = v_1' \ldots v_m'\) , where \(v_i'\) is a symmetry with respect to a hyperplane \(L_i\) of H. The endomorphism \(v_i\) of E which extends \(v_i'\) and is such that \(v_i(x) = x\) is then the symmetry with respect to the hyperplane \(Ax + L_i\) of E. We obviously have \(u = v_1v_2 \ldots v_m\) .
\item
  Suppose secondly that \(u(x) = -x\) . If we denote by \(s\) the symmetry with respect to the hyperplane H orthogonal to \(x\) , and if we set \(\nu = su\) , one has \(\nu(x) = x\) , and we are reduced to the case \(a\) ).
\item
  Let us finally pass to the general case, and set \(y = u(x)\) so that \(Q(y) = Q(x)\) . Under these conditions, the vectors \(x - y\) and \(x + y\) cannot both be isotropic, for from the relations \(Q(x - y) = 0\) and \(Q(x + y) = 0\) , one would deduce, by adding memberwise, \(2(Q(x) + Q(y)) = 0\) (§ 3, no. 4, def. 2), whence \(4Q(x) = 0\) , contrary to the hypothesis. Suppose, for example, that \(a = x - y\) is not isotropic; we then have
\end{enumerate}

\[
\Phi (y, a) = \mathrm{Q} (y + a) - \mathrm{Q} (y) - \mathrm{Q} (a) = \mathrm{Q} (x) - \mathrm{Q} (y) - \mathrm{Q} (a) = - \mathrm{Q} (a);
\]

consequently, if one denotes by \(s\) the symmetry with respect to the hyperplane orthogonal to \(a\) , formula (6) proves that \(s(y) = y + a = x\) ; by setting \(\nu = su\) , one has \(\nu(x) = x\) , and we are reduced to the case \(a\) . If \(a = x - y\) is isotropic and \(b = x + y\) not isotropic; one sees in the same way that we are reduced to the case \(b\) .

\subsection{The group of similarities.}

Let \(\Phi\) be a Hermitian form on E. An automorphism \(u\) of the A-module E is called a similarity (relative to \(\Phi\) ) if there exists an invertible element \(\alpha\) of A such that one has

\[
\Phi (u (x), u (y)) = \alpha \Phi (x, y)\tag{7}
\]

for all x, y in E. The similarities form a group Γ.

When \(\Phi\) takes values that are regular elements of A (for example when A is a field and \(\Phi \neq 0\) ), the element \(\alpha\) of A satisfying (7) is uniquely determined by \(u\) ; it is called the multiplier of the similarity \(u\) . Changing \(x\) to \(\lambda x\) in (7), we then see that \(\alpha\) belongs to the center of A; exchanging \(x\) and \(y\) in (7), we see moreover that \(\overline{\alpha} = \alpha\) . If, for \(u \in \Gamma\) , we denote by \(\alpha(u)\) the multiplier of \(u\) , the map \(u \to \alpha(u)\) is a homomorphism of \(\Gamma\) into the multiplicative group of invertible elements of the center of A. The kernel of this homomorphism is the unitary group associated with \(\Phi\) , which is therefore a normal subgroup of \(\Gamma\) . Let \(\beta\) be an invertible element of the center of A, \(v\) the homothety of ratio \(\beta\) , and \(w\) a unitary automorphism of E; then \(vw = wv\) is a similarity of E, and its multiplier is \(\beta\overline{\beta}\) . Conversely, let \(u\) be a similarity whose multiplier is of the form \(\beta\overline{\beta}\) ( \(\beta\) denoting an invertible element of the center of A); then \(uv^{-1}\) is a unitary automorphism \(w\) , hence \(u\) is of the form \(vw\) .

Suppose now that A is a field, E is a finite-dimensional vector space, and that \(\Phi\) is nondegenerate. For every similarity \(u\) with multiplier \(\alpha\) we have

\[
\Phi (x, \alpha y) = \alpha \Phi (x, y) = \Phi (u (x), u (y)) = \Phi (x, u ^ {*} (u (y))),
\]

therefore \(u^*u\) is the homothety with ratio \(\alpha\) . If A is commutative, and if \(n\) denotes the dimension of E, it follows from this and from formula (50) of § 1, no. 10 that one has

\[
(\det u) \overline {{(\det u)}} = \alpha^ {n}.\tag{8}
\]

Let us then distinguish two cases:

1°) The integer n is odd, i.e., \(n = 2q + 1\) . Then, setting \(\rho = \alpha^{-q}(\det u)\) , one has \(\alpha = (\det u) \overline{(\det u)} \alpha^{-2q} = \rho \bar{\rho}\) . Thus u is the product of the homothety with ratio \(\rho\) and a unitary automorphism.

2°) The integer n is even, say n = 2q. Then, setting \(\rho = \alpha^{-q}(\det u)\) , one has \(\rho\bar{\rho} = 1\) . In particular, when J is the identity, one has \((\det u)^{2} = \alpha(u)^{2q}\) ; the similarities u such that det \(u = \alpha(u)^{q}\) (resp. det \(u = -\alpha(u)^{q}\) ) are called direct (resp. inverse); the direct similarities form a normal subgroup of index 2 of \(\Gamma\) ;

the homotheties with ratio ≠ 0 are direct similarities; the same is true of orthogonal transformations with determinant 1 (no. 2); orthogonal transformations with determinant -1 are inverse similarities.

The preceding definitions and results remain valid for ε-hermitian forms (§ 3, no. 1), and in particular for alternating forms.

Let A be a commutative field and Q a quadratic form \(\neq 0\) on E. One calls a similarity (relative to Q) any automorphism \(u\) of E such that there exists a nonzero element \(\alpha\) of A (called the multiplier of \(u\) ) for which \(Q(u(x)) = \alpha Q(x)\) for every \(x \in E\) . It is clear that \(u\) is then a similarity with multiplier \(\alpha\) relative to the bilinear form associated with Q; the converse is true when the characteristic of A is \(\neq 2\) .
% semantic-fragments: legacy-input-seam
\relax{}
\subsection{Hermitian geometry.}

DEFINITION 2. --- Let A be a field, L an affine space over A, and T the space of translations of L (chap.~II, 2 \(^{e}\) ed., App. II). If T is equipped with a nondegenerate Hermitian form Φ, then L is said to be a Hermitian space over A, and Φ is called the metric form of L.

If J is the identity (which implies that A is commutative), one says instead that L is a Euclidean space.

If \(a\) and \(b\) are two points of L, set \(e(a, b) = \Phi(b - a, b - a)\) . Let \(c\) be a third point of L. For \(b - a\) and \(c - a\) to be orthogonal, it is necessary, according to formula (17) of § 1, no. 5, that one have \(e(b, c) = e(a, b) + e(a, c)\) and this condition is sufficient when J = 1 and A is not of characteristic 2 (“Pythagorean theorem”).

Two linear varieties of L are said to be orthogonal if their directions (chap.~II, 2 \(^{e}\) ed., App. II, no. 3) are orthogonal. A linear manifold of L is said to be isotropic (resp. totally isotropic) if its direction is isotropic (resp. totally isotropic). A vector of T is said to be orthogonal to a linear manifold of L if it is orthogonal to the direction of this manifold.

Let V be a linear variety in L, and let x be a point of L. The set of points y of L such that y − x is orthogonal to V is a linear variety W passing through x; we say that W is the totally orthogonal variety (or, more simply, orthogonal) to V passing through \(x\) . If L is finite-dimensional, the dimension of W is equal to the codimension of V. Moreover, if V is non-isotropic, the directions of V and W are supplementary (§ 4, no. 1, cor. of prop. 1); then W meets V in a single point \(x_1\) ; taking an origin in V, one sees immediately that, for fixed V, the map \(x \to x_1\) is an idempotent affine linear map; it is called the orthogonal projection of L onto V; the linear map associated with it (chap.~II, 2 \(^{e}\) ed., App. II, no. 4) is the orthogonal projector of T onto the direction of V (no. 3).

DEFINITION 3. --- Let L be a Hermitian space over a field A, and T the space of translations of L. One calls a displacement (resp. similarity) of L any affine bijection u of L onto L such that the linear map v associated with u in T (chap.~II, 2 \(^{e}\) ed., App. II, no. 4) is unitary (resp. is a similarity).

The group of translations is a distinguished subgroup of the affine group; it is therefore a distinguished subgroup of the group of similarities and of the group of displacements. For every \(a \in L\) , let \(G_a\) be the group of similarities (resp. displacements) leaving \(a\) fixed; if one identifies L with T by taking \(a\) as origin, \(G_a\) is the group of similarities (resp. the unitary group) of T. Every similarity (resp. displacement) \(u\) can be written, in one and only one way, in the form \(u = u_1 t_1\) where \(u_1 \in G_a\) and \(t_1 \in T\) , and also in the form \(u = t_2 u_2\) where \(u_2 \in G_a\) and \(t_2 \in T\) ; moreover, one has \(u_2 = u_1\) and \(t_2 = u_1 t_1 u_1^{-1}\) (chap.~II, 2 \(^\text{e}\) ed., App. II, no. 4).

Let \(u\) be a similarity in L, \(\varrho\) the associated similarity in T. The multiplier of \(\varrho\) is also called the multiplier of \(u\) (no. 5). If we denote by \(\alpha(u)\) this multiplier, the map \(u \to \alpha(u)\) is a homomorphism from the group of similarities of L into the multiplicative group of invertible elements of the center of A; its kernel is the group of displacements, which is therefore a normal subgroup of the group of similarities. When A is commutative and L is finite-dimensional, there is, between the determinant det \(u\) (equal by definition to det \(\varrho\) ) and \(\alpha(u)\) , the same relations as in no. 5. The displacements \(u\) such that det \(u = 1\) form a distinguished subgroup of the group of displacements; this subgroup has index 2 if A is a commutative field of characteristic ≠ 2 and J is the identity.

PROPOSITION 6. --- Let L be a finite-dimensional Hermitian space over A, whose metric form has index 0. Every similarity u of L, with multiplier \(\mu \neq 1\) then admits one and only one fixed point.

Indeed, let \(a\) be a point of L. There exists a similarity \(\nu\) of L leaving \(a\) fixed and a translation \(t\) of L such that \(u = tv\) . To say that \(b\) is a fixed point of \(u\) amounts to saying that \(\nu(b) - b = t\) . To show that this equation admits one and only one solution \(b\) , let us identify L with its space of translations T by taking \(a\) as origin. It then suffices to prove that the endomorphism \(\nu - 1\) of T is invertible, in other words that the relation \(\nu(x) - x = 0\) ( \(x \in \mathrm{T}\) ) implies \(x = 0\) . Now, if \(\nu(x) - x = 0\) , one has \(\Phi(x, x) = \Phi(\nu(x), \nu(x)) = \mu \Phi(x, x)\) , hence \(\Phi(x, x) = 0\) since \(\mu \neq 1\) ; this entails \(x = 0\) since \(\Phi\) has index 0. QED.

Suppose that A is a field of characteristic \(\neq 2\) Every displacement \(u\) of L such that \(u^2 = 1\) admits at least one fixed point, for example the midpoint \(\frac{1}{2}(x + u(x))\) of two homologous points; taking this point as origin, one sees that the unitary automorphism of T associated with \(u\) is a symmetry (no. 3). Let V be a non-isotropic linear variety in L; one says that a displacement \(u\) is the symmetry with respect to V if, taking an origin in V, \(u\) is identified with the symmetry of T with respect to V. It is equivalent to say that \(u(x)\) is obtained as follows: denoting by \(x_1\) the orthogonal projection of \(x\) on V, we have \(u(x) - x = 2(x_1 - x)\) .

Exercises. --- 1) Suppose that A is a commutative field. Given a Hermitian matrix R of order n over A, the principal minors of order r of R are called the minors obtained by deleting from R n - r rows and the n - r columns with the same indices.

\begin{enumerate}
\def\labelenumi{\alph{enumi})}
\setcounter{enumi}{1}
\tightlist
\item
  From a), deduce that if R has rank r, there exists a permutation \(\sigma \in S_{n}\) such that, if one applies the same permutation \(\sigma\) to the rows and columns of R and if we denote by S the matrix thus obtained, by \(\Delta_{k}\) the principal minor of order k of S obtained by deleting from S the rows and columns with index \textgreater{} k , one has the following two properties: \(1^{\circ} \Delta_{r} \neq 0\) ; \(2^{\circ}\) there is no index k \textless{} r such that \(\Delta_{k} = \Delta_{k+1} = 0\) .
\end{enumerate}

\begin{enumerate}
\def\labelenumi{\arabic{enumi})}
\setcounter{enumi}{1}
\tightlist
\item
  We assume that A is a commutative field, and that E is of finite dimension n. Let \(\Phi\) be a Hermitian sesquilinear form on E, satisfying condition (T) of § 4, no. 2, \(R = (\alpha_{ij})\) the matrix of \(\Phi\) with respect to a basis ( \(e_i\) ) of E.
\end{enumerate}

\begin{enumerate}
\def\labelenumi{\alph{enumi})}
\item
  If \(\Phi\) has rank \(r\) and if the principal minor (exerc. 1) obtained by deleting in \(R\) the rows and columns with indices \(>r\) is not zero, show that there exists a new basis \((f_i)\) of E such that \(e_i = f_i\) for \(1 \leqslant i \leqslant r\) and that the matrix of \(\Phi\) with respect to \((f_i)\) is obtained by replacing by 0 in \(R\) all the \(\alpha_{ij}\) such that \(i > r\) or \(j > r\) (consider the subspace E \(^\circ\) orthogonal to E).
\item
  From a), deduce that if \(\Phi\) has rank \(n\) , and if the cofactor \(\Delta_{n-1}\) of \(\alpha_{nn}\) in the determinant \(\Delta = \det R\) is nonzero, there exists a new basis \((f_i)\) of E such that \(f_i = e_i\) for \(1 \leqslant i \leqslant n-1\) , and such that one has
\end{enumerate}

\[
\Phi (x, y) = \Phi \left(\sum_ {i = 1} ^ {n} \xi_ {i} f _ {i}, \sum_ {i = 1} ^ {n} \eta_ {i} f _ {i}\right) = \sum_ {i = 1} ^ {n - 1} \sum_ {j = 1} ^ {n - 1} \alpha_ {i j} \xi_ {i} \bar {\eta} _ {j} + \frac {\Delta}{\Delta_ {n - 1}} \xi_ {n} \bar {\eta} _ {n}
\]

(consider the Hermitian form whose matrix with respect to \((e_i)\) is obtained by replacing \(\alpha_{nn}\) by \(\alpha_{nn} - \frac{\Delta}{\Delta_{n-1}}\) in \(R\) ).

\begin{enumerate}
\def\labelenumi{\alph{enumi})}
\setcounter{enumi}{2}
\tightlist
\item
  We assume that \(\Phi\) has rank \(n\) , that \(\Delta_{n-1} = 0\) , but the principal minor \(\Delta_{n-2}\) of \(R\) obtained by deleting the rows and columns with indices \(n-1\) and \(n\) in \(R\) is nonzero. Show that there exists a new basis \((f_i)\) of E such that \(f_i = e_i\) for \(1 \leqslant i \leqslant n-2\) , and such that one has
\end{enumerate}

\[
\Phi (x, y) = (\sum_ {i = 1} ^ {n} \xi_ {i} f _ {i}, \sum_ {i = 1} ^ {n} \eta_ {i} f _ {i}) = \sum_ {i = 1} ^ {n - 2} \sum_ {j = 1} ^ {n - 2} \alpha_ {i j} \xi_ {i} \bar {\eta} _ {j} + \xi_ {n - 1} \bar {\eta} _ {n} + \xi_ {n} \bar {\eta} _ {n - 1}.
\]

(If H is the hyperplane spanned by \(e_{1},\ldots,e_{n-1}\) which is isotropic, note that the line orthogonal to H is not contained in the subspace generated by \(e_{1},\ldots,e_{n-2}\) , and use prop. 2 of § 4, no. 2).

\begin{enumerate}
\def\labelenumi{\arabic{enumi})}
\setcounter{enumi}{2}
\item
  Let A be a finite field, E a finite-dimensional vector space over A, Φ a nondegenerate Hermitian sesquilinear form on E, relative to an automorphism J ≠ 1 of A. Show that E admits an orthonormal basis for Φ (cf.~chap.~V, § 11, no. 5, cor. of th. 3).
\item
  Let A be a finite field of characteristic \(\neq 2\) , E a finite-dimensional vector space of dimension \(n\) over A.
\end{enumerate}

\begin{enumerate}
\def\labelenumi{\alph{enumi})}
\item
  Show that for every nondegenerate symmetric bilinear form \(\Phi\) on E, there exists an orthogonal basis \((e_i)\) of E such that \(\Phi(e_i, e_i) = 1\) for \(1 \leqslant i \leqslant n - 1\) , \(\Phi(e_n, e_n) = \Delta\) (discriminant of \(\Phi\) with respect to \((e_i)\) ). (Note that if \(\alpha\beta \neq 0\) , the equation \(\alpha\xi^2 + \beta\eta^2 = \gamma\) always admits solutions \((\xi, \eta)\) in A if \(\gamma \neq 0\) (chap.~V, § 11, exerc. 4)).
\item
  For two nondegenerate symmetric bilinear forms on E to be equivalent, it is necessary and sufficient that the ratio of their discriminants (with respect to the same basis of E) be a square in A. Deduce that, if n is odd, for every symmetric bilinear form \(\Phi\) nondegenerate on E, there exists an orthogonal basis with respect to which the matrix of \(\Phi\) is of the form \(\lambda I_{n} (\lambda \in \mathbf{A})\) ; the index of \(\Phi\) is then \((n - 1)/2\) .
\item
  If \(n=2m\) is even, show that the index of a nondegenerate symmetric bilinear form \(\Phi\) on E is \(m\) if \((-1)^{m}\Delta\) is a square in A, \(m-1\) otherwise.
\end{enumerate}

\begin{enumerate}
\def\labelenumi{\arabic{enumi})}
\setcounter{enumi}{4}
\tightlist
\item
  Let A be a commutative field of characteristic \(\neq 2\) Let I be a polynomial with coefficients in A, in the \(n(n+1)/2\) indeterminates \(X_{ij}\) ( \(1 \leqslant i \leqslant j \leqslant n\) ); for every symmetric matrix \(R = (\alpha_{ij})\) over a commutative extension field A' of A, denote by I(R) the element of A' obtained by substituting \(\alpha_{ij}\) to the indeterminate \(X_{ij}\) ( \(i \leqslant j\) ) in I.
\end{enumerate}

We assume that I is such that, for the matrix \(U = (u_{ij})\) with \(u_{ij} = \mathrm{X}_{ij}\) for \(i \leqslant j\) , \(u_{ij} = \mathrm{X}_{ji}\) for \(i > j\) and the square matrix \(P = (\mathrm{Y}_{ij})\) of order \(n\) (where the \(\mathrm{Y}_{ij}\) are \(n^2\) other indeterminates), one has

\[
\mathrm{I} (\mathbf {t} P U P) = (\det P) ^ {\hbar} \mathrm{I} (U)
\]

where h is an integer \textgreater{} 0. Show that h is even and that I(U) = γ(det U)k, where h = 2k and γ ∈ A. (Using th. 1, show that for every symmetric matrix R over the algebraic closure Ω of A, one has (I(R))² = λ(det R)h, where λ ∈ Ω, and use the fact that the polynomial det U with respect to the Xij is not a square, by considering the terms of this polynomial containing an Xii).

*6) Let A be a complete non-discrete valued field, commutative and of characteristic \(\neq 2\) (Top. gen., chap.~IX, § 3, no. 2), \(\Phi\) a non-degenerate Hermitian form on a finite-dimensional vector space E of dimension \(n\) over A, \(R = (\alpha_{ij})\) the matrix of \(\Phi\) with respect to a basis \((e_i)\) of E. Show that there exists \(\varepsilon > 0\) such that, for every Hermitian matrix \(R' = (\alpha_{ij}')\) satisfying the conditions \(|\alpha_{ij}' - \alpha_{ij}| \leqslant \varepsilon\) for every pair \((i, j)\) , the form \(\Phi'\) having \(R'\) as matrix with respect to the basis \((e_i)\) is equivalent to \(\Phi\) . (Reduce to the case where \(R\) is diagonal; use exerc. 2 b) relying on the following lemma: there exists a number \(a > 0\) such that for \(|\eta| \leqslant a\) , there exists in A an element \(\xi\) such that \(\xi^2 = 1 - \eta\) . To prove this lemma, use the binomial series for \((1 - x)^{1/2})\cdot *\)

¶ 7) Let A be a non-orderable commutative field (chap.~VI, § 2, exerc. 8) of characteristic ≠ 2, E a vector space of finite dimension n \textgreater{} 0 over A, Q a non-degenerate quadratic form on E, \((e_i)\) an orthogonal basis for Q, so that \(\mathrm{Q}\left(\sum_{i=1}^{n}\xi_{i}e_{i}\right)=\sum_{i=1}^{n}\alpha_{i}\xi_{i}^{2}\) . For \(1\leqslant r\leqslant n\) , set \(\mathrm{Q}_{r}(\xi_{1},\ldots,\xi_{r})=\sum_{i=1}^{r}\alpha_{i}\xi_{i}^{2}\) , and denote by \(M_{r}\) the set of values of \(Q_{r}\) as the \(\xi_{i}\left(1\leqslant i\leqslant r\right)\) range over A.

\begin{enumerate}
\def\labelenumi{\alph{enumi})}
\item
  Show that if, for some index \(r\) , one has \(\mathbf{M}_r = \mathbf{M}_{r+1}\) , then it follows that \(\mathbf{M}_r = \mathbf{A}\) (note that every element of A is a sum of squares (chap.~VI, § 2, exerc. 7)).
\item
  Suppose that the subgroup S of the multiplicative group A*, formed by the squares of elements of A, has finite index s in A*. Deduce from a) that if n \textgreater{} s, every non-degenerate quadratic form on E has index \textgreater{} 0 (note that every set M \(_{r}\) is the union of 0 and of classes mod. S). *Application to the case where A is a p-adic field Q \(_{p}\) (Top. gen., chap.~III, § 5, exerc. 35).*
\end{enumerate}

\begin{enumerate}
\def\labelenumi{\arabic{enumi})}
\setcounter{enumi}{7}
\item
  Let A be a commutative field of characteristic \(\neq 2\) , E a finite-dimensional vector space \(n\) over A, Q a non-degenerate quadratic form of index 0 on E. Let A' be an algebraic extension of A, of finite odd degree, E' the vector space over A' obtained by extending to A' the field of scalars of E. Show that the extension Q' of Q to E' (§ 3, no. 4, prop. 3) still has index 0. (Reduce to the case where A' = A{[}X{]}/(f), f being an irreducible polynomial of odd degree m over A. Let ( \(e_i\) ) be an orthogonal basis of E for Q, and \(\rho_i = Q(e_i)\) ; show that, in A{[}X{]}, a relation of the form \(\sum_{i} \rho_i(g_i(X))^2 = f(X)h(X)\) , where the \(g_i\) are polynomials, not all zero, of degree \(\leqslant m - 1\) , is impossible; for this observe that h would necessarily have odd degree, and consider an irreducible factor of h, of odd degree).
\item
  Let A be a field, E a vector space over A admitting a countable basis \((e_n)_{n\geqslant 1}\) , \(\Phi\) a non-degenerate Hermitian sesquilinear form on E, satisfying condition (T) ( \(\S 4\) , no. 2).
\end{enumerate}

\begin{enumerate}
\def\labelenumi{\alph{enumi})}
\item
  Show that if the conditions (C) of th. 1 are not simultaneously satisfied, there exists in E an orthogonal basis for \(\Phi\) (reason as in exerc. 4 of § 5).
\item
  Suppose moreover that A is commutative, and that there exists an integer s such that on every finite-dimensional vector space of dimension \textgreater s over A, every non-degenerate Hermitian sesquilinear form has index \textgreater0 (cf.~exerc. 7). Show that there then exists in E an orthonormal basis for Φ. (Reason as in a), observing that for every element of A of the form α = λ + λ̄, and every non-degenerate Hermitian form Ψ on a finite-dimensional space F of dimension \textgreater s, there exists z ∈ F such that Ψ(z, z) = α (cf.~§ 4, no. 2, prop. 4).)
\end{enumerate}

¶ 10) a) Let A be a principal ideal ring in which there is only one maximal ideal Aπ, such that 2 is not divisible by π (chap.~VII, § 1, exerc. 4). Let E be a free module over A, of dimension n.~Show that every symmetric bilinear form Φ on E admits an orthogonal basis. (Let r be the largest exponent such that πr divides all elements Φ(x, y); show that there exists a ∈ E such that Φ(a, a) = απr, where α is invertible in A; deduce that E is the direct sum of F = Aa and of the submodule F⁰ orthogonal to F.)

\begin{enumerate}
\def\labelenumi{\alph{enumi})}
\setcounter{enumi}{1}
\item
  Give an example (for n = 2) where \(\Phi\) is non-degenerate and where there exists a non-isotropic submodule F of E, of rank 1, admitting a complement in E but such that \(F^{0}\) is not a complement of F.
\item
  Let \((e_i)\) be an orthogonal basis for \(\Phi\) , and \(\alpha_i = \Phi(e_i, e_i)\) . Show that the ideals \(\mathbf{A}\alpha_i\) are, up to order, independent of the orthogonal basis considered (cf.~§ 5, th. 1).
\end{enumerate}

These ideals are said to be the invariant factors of the form \(\Phi\) . Give an example of two forms having the same invariant factors and not equivalent (take two forms whose quotient of discriminants is not a square).

\begin{enumerate}
\def\labelenumi{\alph{enumi})}
\setcounter{enumi}{3}
\item
  Let F be a submodule of E, \(\Phi_{\mathrm{F}}\) the restriction of \(\Phi\) to \(\mathbf{F} \times \mathbf{F}\) , \(\mathbf{A}\alpha_{i}\) ( \(1 \leqslant i \leqslant r\) ) the nonzero invariant factors of \(\Phi\) , arranged so that \(\alpha_{i}\) divides \(\alpha_{i+1}\) , \(\mathrm{A}\beta_{i}\) ( \(1 \leqslant i \leqslant s\) ) the nonzero invariant factors of \(\Phi_{\mathrm{F}}\) , arranged so that \(\beta_{i}\) divides \(\beta_{i+1}\) . Show that we have \(s \leqslant r\) and \(\beta_{i}\) is a multiple of \(\alpha_{i}\) for \(1 \leqslant i \leqslant s\) (same method as in exerc. 1 a) of § 5).
\item
  We assume \(\Phi\) nondegenerate; let F, G be two non-isotropic submodules of E such that \(F^{0}\) (resp. \(G^{0}\) ) is a supplement of F (resp. G). One assumes that the restrictions of \(\Phi\) to F and to G are equivalent; show that there then exists an automorphism \(u\) of E, leaving invariant \(\Phi\) , and such that \(u(F) = G\) . (Using \(a\) ), one can reduce to the case where \(F = Aa\) , \(G = Ab\) , \(\Phi(a, a) = \Phi(b, b)\) . Let \((c_{j})\) be a basis of \(G^{0}\) , and let \(b'\) , \(c_{j}'\) ( \(1 \leqslant j \leqslant n - 1\) ) be the components of \(b\) and \(c_{j}\) respectively in \(F^{0}\) ; show that there exist scalars \(\mu_{j}\) ( \(1 \leqslant j \leqslant n - 1\) ) such that the elements \(d_{j} = c_{j}' + \mu_{j}b'\) satisfy the relations \(\Phi(d_{j}, d_{k}) = \Phi(c_{j}, c_{k})\) for every pair of indices; to do this, note that for every \(\lambda \in A\) one of the elements \(1 \pm \lambda\) is invertible in A.)
\end{enumerate}

\begin{enumerate}
\def\labelenumi{\arabic{enumi})}
\setcounter{enumi}{10}
\item
  Let A be a principal ideal ring of characteristic 0, in which there is only one principal ideal \(\pi\) , such that 2 is divisible by \(\pi\) . If \((e_1, e_2)\) is the standard basis of \(E = A^2\) , \(\Phi\) the symmetric bilinear form on E defined by \(\Phi(\xi_1 e_1 + \xi_2 e_2, \eta_1 e_1 + \eta_2 e_2) = \xi_1 \eta_2 + \xi_2 \eta_1\) , show that there does not exist an orthogonal basis of E for \(\Phi\) .
\item
  Let A be the finite field \(\mathbf{F}_{q^{2}}\) , J the involutive automorphism \(\xi \to \xi^{q}\) of A, whose \(\mathbf{F}_{q}\) is the field of invariants. If E is a vector space of dimension \(n\) over A, \(\Phi\) a non-degenerate Hermitian sesquilinear form (with respect to J) on E, show that the order of the unitary group \(\mathbf{U}(\Phi)\) is equal to
\end{enumerate}

\[
(q ^ {n} - (- 1) ^ {n}) q ^ {n - 1} (q _ {-} ^ {n - 1} - (- 1) ^ {n - 1}) q ^ {n - 2} \dots (q ^ {2} - 1) q (q + 1)
\]

(using a method analogous to that of exerc. 10 of § 5, by means of exerc. 3).

\begin{enumerate}
\def\labelenumi{\arabic{enumi})}
\setcounter{enumi}{12}
\tightlist
\item
  Let A be the finite field \(F_{q}\) (q not a multiple of 2), E a vector space of dimension n over A, and Q a non-degenerate quadratic form on E. Show that:
\end{enumerate}

\begin{enumerate}
\def\labelenumi{\alph{enumi})}
\tightlist
\item
  If n is odd, the order of the group SO(Q) is
\end{enumerate}

\[
(q ^ {n - 1} - 1) q ^ {n - 2} (q ^ {n - 3} - 1) q ^ {n - 4} \dots (q ^ {2} - 1) q.
\]

\begin{enumerate}
\def\labelenumi{\alph{enumi})}
\setcounter{enumi}{1}
\tightlist
\item
  If n = 2m is even, the order of the group SO(Q) is equal to
\end{enumerate}

\[
(q ^ {2 m - 1} - \varepsilon q ^ {m - 1}) (q ^ {2 m - 2} - 1) q ^ {2 m - 3} \dots (q ^ {2} - 1) q
\]

where \(\varepsilon = 1\) if \((-1)^{m}\Delta\) is a square in A, \(\varepsilon = -1\) otherwise, \(\Delta\) denoting the discriminant of Q with respect to an arbitrary basis of E. (Method analogous to that of exerc. 12, using exerc. 3 of §6 and exerc. 5 of chap.~V, §11.)

\begin{enumerate}
\def\labelenumi{\arabic{enumi})}
\setcounter{enumi}{13}
\tightlist
\item
  Assume that A is a commutative field, E a finite-dimensional vector space \(n \geqslant 2\) over A, \(\Phi\) a non-degenerate Hermitian sesquilinear form on E, satisfying condition (T) ( \(\S 4\) , no. 2). Show that the only endomorphisms \(w\) of E commuting with all automorphisms \(u\) belonging to the special unitary group \(\mathbf{SU}(\Phi)\) are the homotheties, except when one simultaneously has \(n = 2\) , J = 1, A being of characteristic \(\neq 2\) . (If \(n \geqslant 3\) , write that \(w\) commutes with the involutions \(u \in \mathbf{SU}(\Phi)\) , and use exerc. 3 of \(\S 4\) ; if \(n = 2\) and J \(\neq 1\) , write that \(w\) commutes with the elements of \(\mathbf{SU}(\Phi)\) whose matrix is of the form \(\begin{pmatrix} \lambda & 0 \\ 0 & \lambda^{-1} \end{pmatrix}\) with respect to an orthogonal basis of E.
\end{enumerate}

¶ 15) Let A be a commutative field of characteristic ≠ 2, E a vector space of dimension n ≥ 1 over A, Q a non-degenerate quadratic form on E. For every automorphism u ∈ O(Q), let w = u − 1, and let r be the rank of w, and W = ̄w(0).

\begin{enumerate}
\def\labelenumi{\alph{enumi})}
\item
  Show that \(w(E)\) is the subspace \(W^{0}\) orthogonal to \(W\) .
\item
  Show that if \(n=2\) , \(r=2\) , \(u\) is the product of two reflections with respect to lines of E. (Establish that if \(w(x)\) is isotropic for every non-isotropic vector \(x\in\mathrm{E}\) , \(w(x)\) is isotropic for every \(x\in\mathrm{E}\) ; we will consider separately the case where A has at least 5 elements and the case \(\mathrm{A}=\mathbf{F}_{3}\) .)
\item
  We assume n and r are arbitrary. Show that if \(w(\mathrm{E})\) is not totally isotropic, u is a product of r symmetries with respect to hyperplanes of E, and cannot be expressed as a product of a smaller number of symmetries. (Reduce to the case where W is totally isotropic, and proceed by induction on n and r. If \(W \neq \{0\}\) show that there exists a vector \(a \in W^{0}\) such that \(w(a)\) is not isotropic, arguing by contradiction and using the fact that a plane all of whose lines except at most one are isotropic is necessarily totally isotropic; then take the symmetry s with respect to the hyperplane orthogonal to \(w(a)\) , and consider the automorphism su. If \(W = \{0\}\) take \(a \in E\) such that \(w(a)\) is not isotropic, and, with the same meaning for s, consider again the automorphism su, and use b).)
\item
  We assume that \(w(E)\) is totally isotropic. If \(s\) is a symmetry with respect to a non-isotropic hyperplane H; show that the subspace of vectors invariant under \(su\) is H ∩ W, hence is of dimension \(n-r-1\) , and deduce that \(su\) cannot be the product of fewer than \(r+1\) symmetries with respect to hyperplanes. Then deduce from \(c\) ) that \(u\) is a product of \(r+2\) symmetries with respect to hyperplanes, but cannot be the product of a smaller number of symmetries.
\item
  Deduce from c) and d) that every orthogonal automorphism is a product of at most n symmetries with respect to hyperplanes.
\item
  Show that if n is odd (resp. even), for every automorphism \(u \in \mathbf{0}(\mathrm{Q})\) of determinant 1 (resp. -1), there exists a vector \(x \neq 0\) invariant under u (use e)).
\end{enumerate}

\begin{enumerate}
\def\labelenumi{\arabic{enumi})}
\setcounter{enumi}{15}
\tightlist
\item
  Given the same hypotheses as in exerc. 15, show that if \(n \geqslant 3\) , the group SO(Q) is generated by the symmetries with respect to the non-isotropic subspaces of E of dimension \(n - 2\) (reason as in prop. 5 of no. 4).
\end{enumerate}

¶ 17) The hypotheses are the same as in exerc. 15.

\begin{enumerate}
\def\labelenumi{\alph{enumi})}
\item
  Show that, for \(n \geqslant 2\) the commutator subgroup \(\Omega(Q)\) of the orthogonal group \(\mathbf{O}(Q)\) is generated by the elements \((st)^{2}\) , where \(s\) and \(t\) range over the set of symmetries with respect to hyperplanes (use prop. 5 of no. 4, and note that for any group \(\Gamma\) the subgroup generated by the squares of the elements of \(\Gamma\) contains the commutator subgroup of \(\Gamma\) ).
\item
  Show that if \(n \geqslant 3\) the commutator subgroup of \(\mathbf{SO}(\mathbf{Q})\) is generated by the squares of the elements of \(\mathbf{SO}(\mathbf{Q})\) (use exerc. 16); deduce that this group is identical to \(\Omega(\mathbf{Q})\) , and that the quotient group \(\mathbf{SO}(\mathbf{Q})/\Omega(\mathbf{Q})\) is a commutative group all of whose elements have order 2.
\item
  A plane P ⊂ E is said to be hyperbolic if it is non-isotropic and contains isotropic lines (necessarily two in number). An automorphism \(u \in \mathbf{0}(Q)\) is said to be hyperbolic if there exists a hyperbolic plane P such that \(u(x) = x\) for all \(x \in \mathrm{P}^0\) ; one then says that \(u\) is a hyperbolic transformation associated with P. Show that if Q has index \(\geqslant 1\) , every \(u \in \mathbf{0}(Q)\) is a product of hyperbolic transformations (use prop. 5 of no. 4 and exerc. 4 a) of § 4). Deduce that if P is a hyperbolic plane, every \(u \in \mathbf{0}(Q)\) can be written \(u = tv\) , where \(t\) is a hyperbolic transformation associated with P and \(v \in \Omega(Q)\) .
\end{enumerate}

¶ 18) Let A be a commutative field, E a vector space of dimension n over A, Φ a nondegenerate Hermitian sesquilinear form on E, satisfying condition (T) (§ 4, no. 2). Let V be a vector subspace of E, H \(_{v}\) the subgroup of the unitary group U(Φ) consisting of unitary automorphisms u such that u(V) = V.

\begin{enumerate}
\def\labelenumi{\alph{enumi})}
\item
  Show that, when V is not a totally isotropic subspace of dimension \(n/2\) , the image of \(\mathrm{H}_{\mathrm{v}}\) by the map \(u \to \det u\) is the subgroup of \(\mathbf{A}^*\) consisting of \(\rho \in \mathbf{A}\) such that \(\rho \bar{\rho} = 1\) .
\item
  If \(n\) is even and \(V\) is a totally isotropic subspace of dimension \(n/2\), show that the image of \(H_{v}\) by the map \(u \rightarrow \det u\) is the subgroup of A* consisting of elements of the form \(\bar{\lambda}/\lambda\) (using prop. 2 of §4, no. 2).
\item
  Let V, W be two vector subspaces of E such that the restrictions of \(\Phi\) to V and to W are equivalent. Show that there exists \(u \in \mathbf{S}\mathbf{U}(\Phi)\) such that \(u(\mathrm{V}) = \mathrm{W}\) in the following cases:
\end{enumerate}

\(^{1}\) J is distinct from the identity (use th. 3 of chap.~V, § 11, no. 5).

\(2^{\circ} \text{ J} = 1, \text{ A is of characteristic } \neq 2, \text{ V and W are not totally isotropic subspaces of dimension } n/2.\)

\begin{enumerate}
\def\labelenumi{\alph{enumi})}
\setcounter{enumi}{3}
\tightlist
\item
  Assume J = 1, A has characteristic ≠ 2, n = 2m is even, and Φ is a nondegenerate symmetric bilinear form of index m. Let V, W be two totally isotropic subspaces of dimension m in E; show that if dim(V ∩ W) = q, then for every orthogonal automorphism u such that u(V) = W, one has det u = \((-1)^{m-q}\) (using b) and prop. 2 of § 4, no. 2). Deduce that the set of totally isotropic subspaces of dimension m is the union of two intransitivity classes N \(_{1}\) , N \(_{2}\) for the group SU(Φ); if V and W are in the same class (resp. in different classes), the dimension of V ∩ W has the same parity as m (resp. does not have the same parity as m). For a similarity u (for Φ) to be direct, it is necessary and sufficient that \(u(N_{1}) = N_{1}\) (using exerc. 4 c) of § 4.)
\end{enumerate}

\begin{enumerate}
\def\labelenumi{\arabic{enumi})}
\setcounter{enumi}{18}
\tightlist
\item
  Let A be a field, E a vector space of finite dimension \textgreater{} zero over A, \(\Phi\) a non-degenerate and non-alternating \(\varepsilon\)-hermitian sesquilinear form on E. Let u be an automorphism of E such that we have
\end{enumerate}

\[
\Phi (u (x), u (x)) = \alpha \Phi (x, x)
\]

for all \(x \in \mathbf{E}\) , with \(\alpha \in \mathbf{A}\) . Show that \(u\) is a similarity with multiplier \(\alpha\) unless the following conditions hold simultaneously: A is commutative and of characteristic 2, and J is the identity (use exerc. 8 of § 1).

\begin{enumerate}
\def\labelenumi{\arabic{enumi})}
\setcounter{enumi}{19}
\item
  Let A be a field, and let L be a finite-dimensional Hermitian space over A; we assume that the metric form \(\Phi\) of L satisfies condition (T) ( \(\S\) 4, no. 2). Show that if the index of \(\Phi\) is \(>0\) there may be similarities of L, with multiplier \(\neq 1\) , and which admit no fixed point (use the reasoning of prop. 6 of no. 6, and prop. 2 of \(\S\) 4, no. 2).
\item
  Let A be a commutative field of characteristic ≠ 2, L a finite-dimensional Euclidean space over A, and Φ the metric form of L.
\end{enumerate}

\begin{enumerate}
\def\labelenumi{\alph{enumi})}
\tightlist
\item
  Show that every bijection u from L onto itself, such that
\end{enumerate}

\[
\Phi (u (x) - u (y), u (x) - u (y)) = \Phi (x - y, x - y)
\]

for all \(x, y\) in \(L\) is a displacement (use exerc. 7 of § 1).

\begin{enumerate}
\def\labelenumi{\alph{enumi})}
\setcounter{enumi}{1}
\tightlist
\item
  Show that the group of displacements is generated by the symmetries with respect to the non-isotropic hyperplanes of the affine space L (using prop. 5 of no. 4, reduce to proving that every non-isotropic translation is the product of two such symmetries).
\end{enumerate}

\begin{enumerate}
\def\labelenumi{\arabic{enumi})}
\setcounter{enumi}{21}
\item
  In a Hermitian space L, two linear varieties are said to be perpendicular if their directions are weakly orthogonal subspaces (§ 3, exerc. 11). Suppose L is of finite dimension; let \(V_{1}\) , \(V_{2}\) be two linear varieties, \(W_{1}\) , \(W_{2}\) their respective directions. Suppose that \(p = \dim(W_{1} + W_{2}) < n\) ; show that if \(W_{1} + W_{2}\) is not isotropic, there exists at least one linear variety U of dimension n - p, perpendicular to \(V_{1}\) and to \(V_{2}\) , and meeting each of the varieties \(V_{1}\) , \(V_{2}\) in exactly one point; moreover, if \(q = \dim(W_{1} \cap W_{2})\) , the union of all linear varieties U having the preceding properties is a linear variety of dimension \(n - p + q\) .
\item
  Let A be a commutative field of characteristic \(\neq 2\) , E a finite-dimensional vector space \(n + 1 \geqslant 2\) on A, Q a quadratic form on E, \(\Phi\) the symmetric bilinear form associated with Q. The set C of \(x \in E\) such that \(Q(x) = 0\) is called the isotropic cone with vertex 0 and equation \(Q(x) = 0\) . If it is not reduced to 0, the image S of \(C - \{0\}\) in the projective space \(\mathbf{P}(\mathrm{E})\) , under the canonical map \(\pi\) of \(\mathrm{E}-\{0\}\) on \(\mathbf{P}(\mathrm{E})\) (chap.~II, \(2^{\mathrm{e}}\) ed., App. III), is called a projective quadric (resp. projective conic if \(n=2\) ) with homogeneous equation \(\mathrm{Q}(x)=0\) . S is said to be degenerate if Q is degenerate. Two projective linear varieties \(\mathrm{V}_{1}, \mathrm{V}_{2}\) of \(\mathbf{P}(\mathrm{E})\) are said to be conjugate with respect to S if \(\overline{\pi}(\mathrm{V}_{1})\) and \(\overline{\pi}(\mathrm{V}_{2})\) are orthogonal (for \(\Phi\) ). The polar \(\mathrm{V}^{0}\) of a projective linear variety \(\mathrm{V}\subset\mathbf{P}(\mathrm{E})\) with respect to S is the variety such that \(\overline{\pi}(\mathrm{V}^{0})\cup\{0\}\) is the totally orthogonal subspace (for \(\Phi\) ) to \(\overline{\pi}(\mathrm{V})\cup\{0\}\) ; if V is a hyperplane and if S is non-degenerate, \(\mathrm{V}^{0}\) is reduced to a point, called the pole of V. A projective linear variety V is said to be tangent to S if \(\overline{\pi}(\mathrm{V})\cup\{0\}\) is an isotropic subspace (for \(\Phi\) ).
\end{enumerate}

We assume in what follows that S is nonempty and nondegenerate.

\begin{enumerate}
\def\labelenumi{\alph{enumi})}
\item
  Show that the intersection of S and a projective linear variety V is empty or is a quadric in V; for this quadric to be degenerate, it is necessary and sufficient that V be tangent to S.
\item
  Show that the tangent hyperplane to S at a point \(z \in S\) is the union of the lines passing through z and tangent to S.
\end{enumerate}

\[
z ^ {\prime}
\]

z with respect to a and b, that is, the point of D such that \(\begin{bmatrix} a & b \\ z' & z \end{bmatrix} = -1\) (chap.~II, 2nd ed., App. III, exerc. 4); show that \(z'\) belongs to the polar hyperplane of z with respect to S, and that there exist n of these points forming a projectively free family in P(E) and belonging to S (cf.~§ 4, exerc. 4 a)).

\begin{enumerate}
\def\labelenumi{\alph{enumi})}
\setcounter{enumi}{3}
\item
  We assume that n = 3 and that \(\Phi\) has maximum index v = 2. The set of lines contained in S is then the union of two sets \(N_{1}\) , \(N_{2}\) such that every line of \(N_{1}\) meets every line of \(N_{2}\) , but two distinct lines of \(N_{1}\) (resp. \(N_{2}\) ) do not meet (exerc. 18 d)). Let D, \(D'\) be two distinct lines belonging to \(N_{1}\) ; for every \(z \in D\) there exists one and only one line \(\Delta \in N_{2}\) passing through z; if \(u(z)\) is the point where \(\Delta\) meets \(D'\) , show that u is a projective linear map from D onto \(D'\) .
\item
  Assuming still n = 3, let D, D', D'' be three lines of P(E), no two of which intersect. Show that the union of the lines meeting D, D', and D'' is a nondegenerate quadric.
\end{enumerate}

\begin{enumerate}
\def\labelenumi{\arabic{enumi})}
\setcounter{enumi}{23}
\tightlist
\item
  The hypotheses and notation are those of exerc. 23, the quadric S being assumed nonempty and nondegenerate.
\end{enumerate}

\begin{enumerate}
\def\labelenumi{\alph{enumi})}
\item
  Show that the subgroup \(\Gamma\) of the projective group PGL(E) consisting of the projective linear bijections mapping S to itself, is the canonical image of the group of similarities with respect to Q. (Use exerc. 23 c) above, exerc. 2 a) of § 4, and exerc. 8 of § 1.)
\item
  Let \(a\) be a point of \(\mathbf{P}(\mathrm{E})\) not belonging to S, and let \(\Phi_{1}\) be the restriction of \(\Phi\) to the hyperplane orthogonal to \(\overline{\pi}(a)\) in E. Show that the subgroup of \(\Gamma\) leaving \(a\) invariant is isomorphic to the quotient of the orthogonal group \(\mathbf{U}(\Phi_1)\) by its center.
\item
  Let \(b\) be a point of S, F the (isotropic) hyperplane orthogonal to \(\overline{\pi}(b)\) in E, M a (non-isotropic) supplement of \(\overline{\pi}(b)\) with respect to F, and \(\Phi_{2}\) the restriction of \(\Phi\) to M. Show that the subgroup of \(\Gamma\) leaving \(b\) invariant is isomorphic to the group of similarities of a Euclidean space L of dimension \(n-1\) , having as metric form the inverse form ( \(\S\) 1, no. 7) of \(\Phi_{2}\) . (Note that if a similarity transformation for \(\Phi\) maps the line \(\overline{\pi}(b)\) in itself, it maps F into itself, and is entirely determined by its restriction to F).
\end{enumerate}

\begin{enumerate}
\def\labelenumi{\arabic{enumi})}
\setcounter{enumi}{24}
\tightlist
\item
  Let A be a commutative field of characteristic \(\neq 2\) , L a finite-dimensional affine space of dimension \(n \geqslant 2\) over A. We identify L with the complement of a projective hyperplane \(\mathrm{H}_{0}\) (“hyperplane at infinity”) of a projective space \(\mathbf{P}(\mathrm{E})\) of dimension \(n\) (chap.~II, \(2^{\mathrm{e}}\) ed., App. III, no. 4). A nonempty set S ⊂ L is said to be an affine quadric (resp. an affine conic if \(n = 2\) ) if S is the intersection of L and a projective quadric (resp. conic) in \(\mathbf{P}(\mathrm{E})\) (exerc. 23).
\end{enumerate}

\begin{enumerate}
\def\labelenumi{\alph{enumi})}
\item
  Show that if there exists a nondegenerate projective quadric \(\overline{\mathrm{S}}\subset\mathbf{P}(\mathrm{E})\) such that \(S=L\cap\overline{S}\) , this quadric is the only one having these properties, except when \(n=2\) , \(A=F_{3}\) and that S is reduced to 2 elements (note that apart from this exceptional case, for every point \(z\in H_{0}\) not belonging to \(\overline{S}\) there exists a line passing through \(z\) and meeting S in two distinct points). One then says that S is a nondegenerate affine quadric. Two affine linear varieties \(V_{1},V_{2}\) contained in L are said to be conjugate with respect to S if the projective linear varieties \(\overline{V}_{1},\overline{V}_{2}\) such that \(V_{i}=L\cap\overline{V}_{i}\) ( \(i=1,2\) ) are conjugate with respect to \(\overline{S}\) ; one defines similarly the polar (when it is not contained in \(H_{0}\) ) or the pole of an affine linear variety with respect to S, and the affine linear varieties tangent to S.
\item
  Assume that S is nondegenerate; show that one can choose an origin a in L such that, by identifying L in this way with a vector
\end{enumerate}

space, there exists a basis \((e_i)\) of L such that S is the set of \(x = \sum_{i=1}^{\infty} \xi_i e_i\) satisfying an equation of one of the two forms

\[
\begin{array}{c} \alpha_ {1} \xi_ {1} ^ {2} + \dots + \alpha_ {n} \xi_ {n} ^ {2} = 1 \\ \alpha_ {1} \xi_ {1} ^ {2} + \dots + \alpha_ {n - 1} \xi_ {n - 1} ^ {2} + \xi_ {n} = 0. \end{array}
\]

In the first case, the point \(a\) is well determined and is the pole with respect to \(\bar{S}\) of the hyperplane at infinity \(H_{0}\) (called the center of S). (Distinguish two cases according as \(H_{0}\) is or is not tangent to \(\bar{S}\) ; using th. 1 of § 6, no. 1 and prop. 2 of § 4, no. 2.)

\begin{enumerate}
\def\labelenumi{\arabic{enumi})}
\setcounter{enumi}{25}
\item
  Let A be an algebraically closed commutative field of characteristic \(\neq 2\) , E a finite-dimensional vector space over A, Q a non-degenerate quadratic form on E. Let \(u \in \mathbf{0}(Q)\) ; with the notations of exerc. 12 of § 4, we have G(p, p) = \{0\} except for \(p(X) = X - 1\) and \(p(X) = X + 1\) . Let M be a minimal element of the set of non-isotropic subspaces contained in \(G(p, p)\) and stable under \(u\) , and let \(p^h\) be the minimal polynomial of the restriction of \(u\) To M. Show that if \(h\) is odd, M is an indecomposable submodule of \(E_u\) and that if \(h\) is even, M is the direct sum of two isomorphic indecomposable submodules of \(E_u\) (To see that if \(h = 2k\) is even, M cannot be indecomposable, show that \(N = p^k(u)(M)\) would then be totally isotropic; if \((e_i)_{1 \leqslant i \leqslant 2k}\) is a basis of M such that \(u(e_i) = \varepsilon e_i + e_{i+1}\) for \(i \leqslant 2k - 1\) , \(u(e_{2k}) = \varepsilon e_{2k}\) (with \(\varepsilon = \pm 1\) ), show that \(e_k\) cannot be orthogonal to \(e_{k+1}\) , and deduce that the relation \(Q(u(e_k)) = Q(e_k)\) leads to a contradiction).
\item
  Let A be a commutative field of characteristic 2, E a vector space over A, of finite dimension n, Q a quadratic form on E, Φ the associated bilinear form, which is alternating, hence of even rank 2m (§ 5, no. 1, cor. 1 of th. 1).
\end{enumerate}

\begin{enumerate}
\def\labelenumi{\alph{enumi})}
\item
  Show that if \(E^{0}\) is the subspace of E (of dimension n - 2m) orthogonal to E for \(\Phi\) , one has \(Q(\lambda x + \mu y) = \lambda^{2}Q(x) + \mu^{2}Q(y)\) for all x, y in \(E^{0}\) ; in other words, the restriction \(Q_{0}\) of Q to \(E^{0}\) is a semilinear map from \(E^{0}\) (considered as a vector space over A) into A (considered as a vector space over the subfield \(A^{2}\) ), relative to the isomorphism \(\xi \to \xi^{2}\) from A to \(A^{2}\) . Let q be the dimension (over A) of the kernel \(E^{0} \cap \overline{Q}(0)\) of Q, and let \(E_{1}\) be a supplement of \(E^{0} \cap \overline{Q}(0)\) with respect to \(E^{0}\) ; we have \(n - 2m - q \leqslant [A : A^{2}]\) .
\item
  Deduce from a) that there exists a basis \((e_i)_{1 \leqslant i \leqslant n}\) of E, whose \(2m\) first vectors form a basis of a complement \(\mathbf{E}_2\) of \(\mathbf{E}^0\) in
\end{enumerate}

E, the \(n - 2m - q\) following ones a basis of \(\mathbf{E}_1\) , such that one has, for \(x = \sum_{i=1}^{n} \xi_i e_i\)

\[
\mathrm{Q} (x) = \sum_ {i = 1} ^ {m} \left(\alpha_ {i} \xi_ {i} ^ {2} + \xi_ {i} \xi_ {m + i} + \beta_ {i} \xi_ {m + i} ^ {2}\right) + \sum_ {i = 2 m + 1} ^ {n - q} \gamma_ {i} \xi_ {i} ^ {2}
\]

the \(\gamma_{i}\) ( \(2m + 1 \leqslant i \leqslant n - q\) ) being elements of A linearly independent with respect to \(A^{2}\) .

\begin{enumerate}
\def\labelenumi{\alph{enumi})}
\setcounter{enumi}{2}
\item
  We call the index of Q the maximum dimension of totally singular subspaces V of E such that \(V \cap E^{0} = \{0\}\) . Show that if v is the index of Q, one can take the basis \((e_{i})\) of E having the properties stated in b) so that \(\alpha_{i} = \beta_{i} = 0\) for \(1 \leqslant i \leqslant v\) and that the restriction of Q to the subspace of \(E_{2}\) generated by \(e_{v+1}, \ldots, e_{m}, e_{m+v+1}, \ldots, e_{2m}\) is a nondegenerate quadratic form of index 0.
\item
  We assume \(q=0\) ; let \(\mathbf{O}(\mathrm{Q})\) the group of automorphisms of E leaving Q invariant. If \(u\in\mathbf{O}(\mathrm{Q})\) , show that \(u(x)=x\) for all \(x\in\mathrm{E}^{0}\) . For every \(x\in\mathrm{E}_{2}\) , let \(u(x)=u_{0}(x)+u_{2}(x)\) , where \(u_{0}(x)\in\mathrm{E}^{0}\) and \(u_{2}(x)\in\mathrm{E}_{2}\) ; show that \(u_{2}\) belongs to the symplectic group \(\mathbf{Sp}(\Phi_{2})\) (where \(\Phi_{2}\) is the restriction of \(\Phi\) to \(\mathrm{E}_{2}\) ) and one has \(\mathrm{Q}(u_{2}(x))+\mathrm{Q}(x)\in\mathrm{Q}(\mathrm{E}^{0})\) . Conversely, for every automorphism \(u_{2}\in\mathbf{Sp}(\Phi_{2})\) such that \(\mathrm{Q}(u_{2}(x))+\mathrm{Q}(x)\in\mathrm{Q}(\mathrm{E}^{0})\) for all \(x\in\mathrm{E}_{2}\) , show that there exists one and only one linear map \(u_{0}\) of \(\mathrm{E}_{2}\) in E \(^{0}\) such that the linear map equal to \(u_{0} + u_{2}\) in E \(_{2}\) , and to the identity in E \(^{0}\) , belongs to O(Q).
\item
  Assume that A is a perfect field (A² = A) and that q = 0. Deduce from b) that every vector subspace of E of dimension ≥ 3 contains at least one vector x such that Q(x) = 0. If n is odd, one necessarily has m = v and n = 2m + 1, so that there exists a basis (ei) of E with respect to which one has
\end{enumerate}

\[
\mathrm{Q} \left(\sum_ {i = 1} ^ {n} \xi_ {i} e _ {i}\right) = \xi_ {1} \xi_ {m + 1} + \dots + \xi_ {m} \xi_ {2 m} + \xi_ {2 m + 1} ^ {2},
\]

and (with the notations of \(d\) ) \(\mathbf{O}(\mathrm{Q})\) is isomorphic to \(\mathbf{Sp}(\Phi_{2})\) ; all quadratic forms such that \(q=0\) are then equivalent. If \(n\) is even, one necessarily has \(n=2m\) , \(\nu=m\) or \(\nu=m-1\) , and there exists a basis \((e_{i})\) of E with respect to which one has

\[
\mathrm{Q} \left(\sum_ {i = 1} ^ {n} \xi_ {i} e _ {i}\right) = \xi_ {1} \xi_ {m + 1} + \dots + \xi_ {m - 1} \xi_ {2 m - 1} + \xi_ {m} \xi_ {2 m} + \lambda \left(\xi_ {m} ^ {2} + \xi_ {2 m} ^ {2}\right)\tag{1}
\]

where \(\lambda\in A\) . Let \(A_{1}\) be the field obtained by adjoining to A the roots of the polynomial \(\lambda X^{2}+X+\lambda\) show that this field is independent of the basis \((e_{i})\) with respect to which Q can be written in the form (1), and that for two quadratic forms (such that \(q=0\) ) to be equivalent, it is necessary and sufficient that the quadratic extensions of A which correspond to them in this way be identical (use Witt's theorem). Case where A is a finite field of characteristic 2.

¶ 28) Let A be a commutative field of characteristic 2, distinct from \(\mathbf{F}_{2}\) , E a vector space of dimension \(n = 2m\) on A, Q a nondegenerate quadratic form on E.

\begin{enumerate}
\def\labelenumi{\alph{enumi})}
\item
  Show that the orthogonal group \(\mathbf{O}(\mathrm{Q})\) is generated by the symmetries (which here are none other than the transvections belonging to \(\mathbf{O}(\mathrm{Q})\) ( \(\S 4\) , exerc. 6)) (reason as in exerc. 11 of \(\S 5\) ). Deduce from this that the commutator subgroup of \(\mathbf{O}(\mathrm{Q})\) is generated by the squares of the elements of \(\mathbf{O}(\mathrm{Q})\) (cf.~exerc. 17).
\item
  We assume that Q is of maximum index; let V, W be two totally singular subspaces of E (§ 4, no. 1) of dimension m. Let u be a symmetry \(x \to x + \frac{\Phi(x, a)}{\mathrm{Q}(a)} a\) (§ 4, exerc. 6); let k be the dimension of V ∩ W. Show that the dimension of V ∩ u(W) is \(k + 1\) if a is orthogonal to V ∩ W, \(k - 1\) in the contrary case (in the first case, note that \(a = x + y\) , where \(x \in V\) , \(y \in W\) and show that \(u(y) = x\) ; in the second, note that u cannot leave invariant any singular vector not orthogonal to a).
\item
  We again assume that the index of Q is arbitrary. Show that the subgroup \(\mathbf{SO}(\mathrm{Q})\) of \(\mathbf{O}(\mathrm{Q})\) consisting of the automorphisms of E that are products of an even number of symmetries, is a normal subgroup of index 2 of \(\mathbf{O}(\mathrm{Q})\) . (Show that the product of an odd number of symmetries cannot be the identity, by considering the extension of Q to the
\end{enumerate}

vector space E' obtained by extending the scalar field of E to its algebraic closure; then use b).) (Cf. § 9, exerc. 9.)

\begin{enumerate}
\def\labelenumi{\alph{enumi})}
\setcounter{enumi}{3}
\item
  If \(V_{1}\) , \(\bar{V}_{2}\) are two totally singular subspaces of E, of the same dimension \textless m, show that there exists an automorphism \(u \in \mathbf{SO}(Q)\) such that \(u(V_{1}) = V_{2}\) . Conversely, if \(V_{1}\) and \(V_{2}\) are two totally singular subspaces of dimension m, then for there to exist an automorphism \(u \in \mathbf{SO}(Q)\) such that \(u(V_{1}) = V_{2}\) , it is necessary and sufficient that the dimension of \(V_{1} \cap V_{2}\) have the same parity as m (argue as in exerc. 18, using b)).
\item
  A plane P ⊂ E is said to be hyperbolic if it is non-isotropic and contains singular lines (necessarily two in number). A transformation \(u \in \mathbf{0}(Q)\) is said to be hyperbolic if there exists a hyperbolic plane P such that \(u(x) = x\) for all \(x \in \mathrm{P}^0\) ; one then says that \(u\) is a hyperbolic transformation associated with P. Show that if Q has index \textgreater0, every \(u \in \mathbf{0}(Q)\) is a product of hyperbolic transformations (use \(a\) ). Deduce that if P is a hyperbolic plane, every transformation \(u \in \mathbf{0}(Q)\) can be written \(u = sv\) , where \(s\) is a hyperbolic transformation associated with P and \(v\) belongs to the commutator subgroup of \(\mathbf{0}(Q)\) .
\end{enumerate}

\begin{enumerate}
\def\labelenumi{\arabic{enumi})}
\setcounter{enumi}{28}
\tightlist
\item
  Given the hypotheses of exerc. 28, suppose further that Q has maximum index m; let \((e_{i})\) be a symplectic basis of E (for the alternating form \(\Phi\) associated with Q) consisting of singular vectors (§ 4, no. 2, prop. 2), so that the matrix of \(\Phi\) with respect to this basis is the matrix denoted by R in exerc. 14 of § 5. Using the notation of that exercise, show that, for a symplectic matrix \((^{t}D + S)^{-1}(D + S)\) to be the matrix of an automorphism \(u \in \mathbf{O}(\mathbb{Q})\) , it is necessary and sufficient that S be alternating (write that every vector \(u(e_{i})\) is singular, noting that one has \((^{t}D + S).u(e_{i}) = (D + S).e_{i}\) ).
\end{enumerate}

\section{Hermitian forms and ordered fields}

Throughout this paragraph, K denotes a maximal ordered field (hence commutative and of characteristic zero; cf.~Chap.~VI, § 2), and we assume that we are in one of the following three cases:

1°) A = K, J is the identity;

2°) A is the field \(\mathrm{K}(i)\) , obtained by adjoining to K a square root i of -1, and, for every \(\lambda \in A\) , \(\overline{\lambda}\) is the conjugate of \(\lambda\) (chap.~II, § 7, no. 7).

3°) A is the quaternion algebra over K corresponding to the pair \((-1,-1)\) (or, as we shall say for short, the field of quaternions over K), and for every \(\lambda\in A\) , \(\overline{\lambda}\) is the conjugate of \(\lambda\) (cf.~chap.~II, § 7, no. 8 and chap.~VIII, § 11, no. 2).

If \(\Phi\) is a Hermitian form on E, so in all cases we have, \(\Phi(x, x) \in K\) for all \(x \in E\) , since \(\Phi(x, x) = \overline{\Phi(x, x)}\) .

\subsection{Positive Hermitian forms.}

DEFINITION 1. --- A Hermitian form \(\Phi\) on E is said to be positive (resp. negative) if \(\Phi(x, x) \geqslant 0\) (resp. \(\Phi(x, x) \leqslant 0\) ) for all \(x \in \mathbf{E}\) .

When A = K, one also says that the quadratic form \(Q(x) = \frac{1}{2} \Phi(x, x)\) to which \(\Phi\) is associated is positive (resp. negative).

Suppose E is finite-dimensional over A, and let \((e_{i})\) ( \(i = 1, \ldots, n\) ) be an orthogonal basis of E ( \(\S 6, no. 1\) , th. 1). For a Hermitian form \(\Phi\) on E to be positive, it is necessary and sufficient that \(\Phi(e_{i}, e_{i}) \geqslant 0\) for \(i = 1, \ldots, n\) . Let \(\Phi\) be a nondegenerate positive Hermitian form on E; since every positive element of K is a square, hence of the form \(\rho\bar{\rho} (\rho \in A)\) , there exist in E orthonormal bases for \(\Phi\) ( \(\S 6, no. 1\) , cor. 1 of th. 1) .

PROPOSITION 1. --- Suppose E is finite-dimensional, and A = K or A = K(i). If \(\Phi\) is a nondegenerate positive Hermitian form on E, then the extensions of \(\Phi\) to \(\bigotimes^{p}\mathbf{E}\) and \(\bigwedge^{p}\mathbf{E}\) ( \(p > 0\) ), as well as the inverse form of \(\Phi\) , are nondegenerate positive Hermitian forms.

This follows immediately from the existence of an orthonormal basis of E, and from prop. 2, § 6, no. 1.

Prop. 1 remains true if one replaces everywhere “nondegenerate positive” by “positive”.

PROPOSITION 2. --- Let Φ be a positive Hermitian form on E. For x, y in E, one has

\[
\Phi (x, y) \overline {{{{\Phi (x , y)}}}} \leqslant \Phi (x, x) \Phi (y, y).\tag{1}
\]

The inequality is indeed immediate when the vectors \(x\) and \(y\) are proportional. Suppose therefore \(x\) and \(y\) linearly independent. Let A' be the (commutative) subfield K( \(\Phi(x, y)\) ) of A, F the vector plane over A' generated by \(x\) and \(y\) , and \(\Phi_{\mathbb{F}}\) the restriction of \(\Phi\) to F; this takes its values in A'. By prop. 1, the discriminant of \(\Phi_{\mathbb{F}}\) with respect to the basis \((x, y)\) is \(\geqslant 0\) . Now this discriminant is \(\Phi(x, x)\Phi(y, y) - \Phi(x, y)\overline{\Phi(x, y)}\) . QED.

COROLLARY. --- The set of isotropic vectors of E is the subspace \(E^{0}\) orthogonal to E for \(\Phi\) . For \(\Phi\) to be nondegenerate, it is necessary and sufficient that one have \(\Phi(x, x) > 0\) for all \(x \neq 0\) .

PROPOSITION 3. --- Suppose E is finite-dimensional and A is commutative (A = K or A = K(i)). Let \(\Phi\) be a Hermitian form on E, X its matrix with respect to a basis ( \(x_j\) ) ( \(j = 1, \ldots, n\) ) of E; for every subset H of {[}1, n{]}, denote by \(X_{\mathrm{H,H}}\) the minor of X obtained by deleting the rows and columns with indices \(j \notin \mathrm{H}\) (chap.~III, § 6, no. 3).

\begin{enumerate}
\def\labelenumi{\alph{enumi})}
\item
  If \(\Phi\) is nondegenerate positive, one has \(X_{\mathrm{H,H}} > 0\) for every subset H of \([1, n]\) .
\item
  Conversely, if, setting \(H_{j} = (1, j)\) , one has \(X_{H_{j}, H_{j}} > 0\) for \(j = 1, \ldots, n\) , \(\Phi\) is nondegenerate positive.
\end{enumerate}

Suppose first \(\Phi\) nondegenerate positive. The elements \((x_{j})\) ( \(j\in\mathbb{H}\) ) form a basis of a subspace F of E, and the minor \(X_{\mathrm{H,H}}\) is the discriminant of the restriction \(\Phi_{\mathrm{F}}\) of \(\Phi\) to F with respect to this basis; now, since \(\Phi_{\mathrm{F}}(x,x)>0\) for all \(x\neq0\) in F, \(\Phi_{\mathrm{F}}\) is nondegenerate positive (cor. of prop. 2); one therefore has \(X_{\mathrm{H,H}}>0\) (prop. 1). To prove \(b\) ), note that, with the notations of prop. 1, § 6, no. 1, the minor \(X_{\mathrm{Hj,Hj}}\) is equal to D \(_{j+1,j+1}\) ; there therefore exists (prop. 1, § 6, no. 1) an orthogonal basis \((e_{j})\) ( \(j=1,\ldots,n\) ) of E such that \(\Phi(e_{j},e_{j})>0\) for \(j=1,\ldots,n\) ; consequently \(\Phi\) is nondegenerate positive.

Remark. --- It follows from the existence of orthonormal bases that two nondegenerate positive Hermitian forms on two vector spaces of the same finite dimension are equivalent (§ 1, no. 6). Let then L be a finite-dimensional Hermitian space over A, whose metric form is nondegenerate positive, and V₁ and V₂ two linear varieties of the same dimension in L (§ 6, no. 6); since the restrictions of the metric form to the directions T₁ and T₂ of V₁ and V₂ on the one hand, and to the orthogonal subspaces T₁ and T₂ on the other hand, are equivalent, there exists a unitary automorphism \(u\) of the space T of translations of L such that \(u(\mathrm{T}_1) = \mathrm{T}_2\) and \(u(\mathrm{T}_1^0) = \mathrm{T}_2^0\) ; there therefore exists a displacement \(v\) of L such that \(v(\mathrm{V}_1) = \mathrm{V}_2\) . Let \((a, b), (a', b')\) two pairs of points of L; for there to exist a displacement \(v\) of L such that \(v(a) = a'\) and \(v(b) = b'\) it is necessary and sufficient (with the notation of § 6, no. 6) that one have \(e(a, b) = e(a', b')\) ; the element \(\sqrt{e(a, b)}\) of A (chap.~VI, § 2, no. 4) is called the distance of \(a\) and \(b\) in the Hermitian space L.

\subsection{The law of inertia.}

THEOREM 1 (“law of inertia”). --- Suppose that A satisfies the hypotheses at the beginning of this paragraph, and that E is of finite dimension n.~Let Φ be a Hermitian form on E. Then:

\begin{enumerate}
\def\labelenumi{\alph{enumi})}
\item
  There exists a decomposition of E as a direct sum of the subspace \(E^{0}\) orthogonal to E, and of two subspaces \(E^{+}\) and \(E^{-}\) such that the restriction of \(\Phi\) to \(E^{+}\) (resp. \(E^{-}\) ) is positive (resp. negative) and non-degenerate.
\item
  There exists an orthogonal basis \((e_i)_{1\leqslant i\leqslant n}\) of E such that
\end{enumerate}

\[
\Phi (\sum_ {i = 1} ^ {n} \xi_ {i} e _ {i}, \sum_ {i = 1} ^ {n} \eta_ {i} e _ {i}) = \sum_ {i = 1} ^ {s} \xi_ {i} \overline {{\eta}} _ {i} - \sum_ {i = s + 1} ^ {s + t} \xi_ {i} \overline {{\eta}} _ {i}.\tag{2}
\]

\begin{enumerate}
\def\labelenumi{\alph{enumi})}
\setcounter{enumi}{2}
\item
  The dimensions s of \(E^{+}\) and t of \(E^{-}\) are the same for all direct-sum decompositions satisfying the conditions stated in a); the integer s (resp. t) is the maximum of the dimensions of subspaces F of E such that the restriction of \(\Phi\) to F is positive (resp. negative) and non-degenerate.
\item
  The rank of \(\Phi\) is \(s + t\) .
\item
  If \(\Phi\) is non-degenerate, its index is equal to inf(s, t) (§ 4, no. 2, def. 2).
\end{enumerate}

Let, in fact, \((x_{i})\) ( \(i = 1, \ldots, n\) ) be an orthogonal basis of E (§ 6, no. 1, th. 1); recall that one has \(\Phi(x, x) \in K\) for all \(x \in E\) . One may assume that \(\Phi(x_{i}, x_{i}) > 0\) for \(i = 1, \ldots, s\) , \(\Phi(x_{i}, x_{i}) < 0\) for \(i = s + 1, \ldots, s + t\) and \(\Phi(x_{i}, x_{i}) = 0\) for \(i = s + t + 1, \ldots, n\) . This proves \(a\) ), for one takes as \(E^{+}\) (resp. \(E^{-}\) ) the subspace generated by \(x_{1}, \ldots, x_{s}\) (resp. by \(x_{s+1}, \ldots, x_{s+t}\) ), and \(E^{0}\) is then generated by \(x_{s+t+1}, \ldots, x_{n}\) . From this we deduce \(b\) ) by noting that \(\Phi(x_{i}, x_{i})\) is of the form \(\rho\bar{\rho}\) (resp.

\(-\rho\bar{\rho}) \quad (\rho\in\Lambda^{*}) \quad \text{for} \quad i=1,\ldots,s \quad (\text{resp. } i=s+1,\ldots,s+t).\) To prove c), consider a subspace P of E such that the restriction of \(\Phi\) to P is positive and nondegenerate; we then have \(\mathrm{P}\cap(\mathrm{E}^{-}+\mathrm{E}^{0})=\{0\}\) and the sum \(P+E^{-}+E^{0}\) is therefore direct; we conclude that dim \(P\leqslant\dim E^{+}=s\) and this proves c). Assertion d) follows immediately from a).

Finally suppose that \(\Phi\) is nondegenerate, set \(q = \inf (s, t)\) , and let us show that \(q\) is the index of \(\Phi\) . With the notations of \(b\) ), the vectors \(e_i + e_{s+i}\) (resp. \(e_i - e_{s+i}\) ) ( \(i = 1, \ldots, q\) ) generate a totally isotropic subspace F (resp. F'). Since K has characteristic 0, F + F' is generated by \(e_1, \ldots, e_q\) , \(e_{s+1}, \ldots, e_{s+q}\) and is therefore non-isotropic. The restriction of \(\Phi\) to the subspace H = (F + F')⁰ is then positive (or negative) and nondegenerate, and H therefore contains no isotropic vector \(\neq 0\) . Since F⁰ contains F and H, and dim(F + H) = codim F' = codim F, we have F⁰ = F + H. Thus an isotropic vector z orthogonal to F necessarily lies in F, because, in the direct sum F + H, the component of z in H is isotropic. Consequently, F is a maximal totally isotropic subspace, and this proves e).

DEFINITION 2. --- With the notation of Th. 1, the pair \((s,t)\) is called the signature of \(\Phi\) .

\subsection{Reduction of a form with respect to a positive Hermitian form.}

In this section we will assume that E is of finite dimension n over A, and we will denote by Φ a positive nondegenerate Hermitian form on E.

Since the linear maps from E into \(\mathbf{E}^*\) associated with \(\Phi\) are bijective, whatever the elements \(x, y, z\) of E with \(y \neq 0\) , there exists one and only one element \(t\) of E such that \(\Phi(x, z) = \overline{\Phi(t, y)}\) In particular, if \(\mathrm{A} = \mathrm{K}\) or \(\mathrm{A} = \mathrm{K}(i)\) and if \(u\) is a semilinear map for J from E into itself (chap.~II, App. I, no. 1), then for every \(x \in \mathrm{E}\) there exists one and only one element \(u^*(x)\) such that for all \(y \in \mathrm{E}\) , \(\Phi(x, u(y)) = \overline{\Phi(u^*(x), y)}\) We see immediately that \(u^*\) is a semilinear map for J from E into itself; it is called the adjoint of \(u\) .

Remarks. --- 1) When J is the identity, we recover the notion of the adjoint of a homomorphism defined in § 1, no. 8.

\begin{enumerate}
\def\labelenumi{\arabic{enumi})}
\setcounter{enumi}{1}
\tightlist
\item
  Assume always that A = K or A = K(i). Let E \(^{j}\) be the right A-module defined in § 1, no. 2, def. 5. The form Φ is a bilinear form on E × E \(^{j}\) , Φ \(^{j}\) is a bilinear form on E \(^{j}\) × E, and u is an A-linear map from E to E \(^{j}\) . The adjoint of this homomorphism, in the sense of § 1, no. 8, is an A-linear map from E into E \(^{j}\) ; we see at once that this coincides with the map u* defined above.
\end{enumerate}

We say that an endomorphism \(u\) of E is normal (for \(\Phi\) ) if one has \(uu^{*}=u^{*}u\) .
% semantic-fragments: legacy-input-seam
\relax{}
\paragraph{Examples of normal endomorphisms:}

\begin{enumerate}
\def\labelenumi{\arabic{enumi})}
\item
  unitary automorphisms for \(\Phi\) ( \(\S 6\) , no 2), which are characterized by the relation \(u^{-1} = u^{*}\) ( \(\S 1\) , no 8, cor. of prop. 8);
\item
  endomorphisms \(u\) such that \(u^{*}=u\) ; these are called Hermitian endomorphisms.
\end{enumerate}

For every Hermitian endomorphism \(u\) of E, set \(\Phi_{u}(x,y)=\Phi(u(x),y)\) ; we have

\[
\Phi_ {u} (y, x) = \Phi (u (y), x) = \Phi (y, u (x)) = \overline {{\Phi (u (x) , y)}} = \overline {{\Phi_ {u} (x , y)}}.
\]

which shows that \(\Phi_{u}\) is a Hermitian form on E. Conversely, let \(\Psi\) be a Hermitian form on E; since the map \(s_{\Phi}\) from E to E* associated with \(\Phi\) is bijective, there exists, for every \(x \in \mathbf{E}\) , a unique element \(u(x)\) of E such that \(\Psi(x, y) = \Phi(u(x), y)\) ; one verifies easily that \(u\) is a Hermitian endomorphism of E. Thus \(u \to \Phi_{u}\) is a bijection, called canonical, from the set of Hermitian endomorphisms of E onto the set of Hermitian forms on E.

Suppose that A = K or A = K(i); given a bilinear form Ψ on E, there exists, for every \(x \in E\) , a unique element \(u(x)\) , a unique element of E such that Ψ(x, y) = \(\overline{\Phi(u(x), y)}\) ; one sees immediately that u is a semilinear map (for J) from E into itself. Since \(\overline{\Phi(u(x), y)} = \Phi(y, u(x)) = \overline{\Phi(u^{*}(y), x)}\) , one sees that, for Ψ to be symmetric (resp. alternating), it is necessary and sufficient that one have \(u^{*} = u\) (resp. \(u^{*} = -u\) ).

THEOREM 2. --- Let S be a set of endomorphisms of E (resp. of semilinear maps from E into E when A = K or A = K(i)) stable under the map u → u*. Then, if V is a subspace of E stable under S, its orthogonal V⁰ is stable under S. Moreover, E is a direct sum of subspaces stable under S, minimal among the subspaces ≠ \{0\} and stable under S, and pairwise orthogonal.

Indeed, let V be a subspace of E stable under S; for any \(x \in \mathrm{V}^0\) , \(y \in \mathrm{V}\) and \(u \in \mathrm{S}\) , one has \(u^*(y) \in \mathrm{V}\) , whence \(\Phi(y, u(x)) = \Phi(u^*(y), x) = 0\) (resp. \(\Phi(y, u(x)) = \overline{\Phi(u^*(y), x)} = 0\) ), and consequently \(u(x) \in \mathrm{V}^0\) ; thus \(\mathrm{V}^0\) is stable under S, which proves our first assertion. For the second, we proceed by induction on the dimension. \(n\) of E, the case \(n = 0\) being trivial. For \(n \neq 0\) there exists a subspace \(\mathrm{V} \neq \{0\}\) of E stable under S and minimal, for example a stable subspace \(\neq \{0\}\) of minimal dimension. It then suffices to apply to \(\mathrm{V}^0\) the induction hypothesis, since the adjoint of the restriction of \(u\) to \(\mathrm{V}^0\) (with respect to the restriction of \(\Phi\) ) is identical to the restriction to \(\mathrm{V}^0\) of the adjoint of \(u\) .

COROLLARY 1. --- Suppose that A = K or that A = K(i). Let B be a subalgebra of \(\mathfrak{L}_{\mathbf{A}}(\mathbf{E})\) stable under the map \(u \to u^{*}\) . Then E is a semisimple B-module, and is a direct sum of simple submodules that are pairwise orthogonal. The algebra B is semisimple.

Indeed, since every sub-B-module V of E admits a complement, for example \(V^{0}\) the B-module V is semisimple (chap.~VIII, § 3, no. 3, prop. 7). Since every sub-B-module \(\neq\{0\}\) and minimal of E is simple, E is a direct sum of pairwise orthogonal simple submodules. Finally B is a semisimple algebra, since it admits a semisimple and faithful module E whose contramodule is finitely generated (chap.~VIII, § 5, no.~1, prop.~3).

COROLLARY 2. --- With the hypotheses and notations of cor.~1, suppose moreover that the algebra B is commutative. Then all its elements are (normal) semisimple endomorphisms of E. When A = K (resp. A = K(i)), E is a direct sum of pairwise orthogonal simple B-submodules, which are vector spaces of dimension 1 or 2 (resp. 1) over A.

The first assertion follows from chap.~VIII, § 9, no.~1, prop.~2, since B is semisimple. On the other hand every simple B-module is isomorphic to a B-module of the form B/m, where m is a maximal ideal of B, hence to a commutative field L of finite degree over A; when A = K (resp. A = K(i)), the field L is necessarily isomorphic to K or K(i) (resp. K(i)), since K(i) is algebraically closed (chap.~VI, § 2, no.~6, th.~3).

PROPOSITION 4. --- Let u be a normal endomorphism of E. When A is equal to K(i) or to the field of quaternions over K, there exists an orthonormal basis (for Φ) of E consisting of eigenvectors of u. When A = K, u is semisimple and E is a direct sum of subspaces stable under u, pairwise orthogonal, and of dimension 1 or 2.

Let us first examine the case where A is commutative (A = K or A = K(i)). Then the subalgebra B = A{[}u, u*{]} of ℒ\_A(E) is commutative since u is normal; it is stable under the map v → v* by virtue of formulas (32) and (33) of § 1, no.~8. The assertion relative to the case A = K then follows immediately from cor.~2 of th.~2. When A = K(i), this corollary also shows that E is a direct sum of vector subspaces Ax\_i (i = 1, . . ., n) of dimension 1, pairwise orthogonal and stable under u; if we set \(e_{i} = (\Phi(x_{i}, x_{i}))^{-1/2}x_{i}\) , \((e_{i})\) is the desired orthonormal basis.

When A is the field of quaternions over K, it will likewise suffice for us to prove, by virtue of th.~2, that every minimal element of the set of subspaces ≠ \{0\} of E stable under u and u* is of dimension 1. Now such a subspace V necessarily contains an eigenvector x ≠ 0 of u (*), as one sees by noting that the field of quaternions A contains K(i) as an algebraically closed subfield, and by restricting to K(i) the field of scalars of V. It therefore remains only to show that the eigenvector \(x\) of \(u\) is also an eigenvector of \(u^{*}\) . Set \(u(x) = ax (a \in A)\) ; we then have \(\Phi(u(x), x) = a\Phi(x, x) = \Phi(x, x)a = \Phi(x, \bar{a}x)\) since \(\Phi(x, x)\) belongs to the center of A, and on the other hand \(\Phi(u(x), x) = \Phi(x, u^*(x))\) ; it follows that we have \(\Phi(x, u^*(x) - \bar{a}x) = 0\) , and we can therefore write \(u^*(x) = \bar{a}x + z\) , where \(z\) is a vector orthogonal to \(x\) . We therefore have

\[
\Phi (u ^ {*} (x), u ^ {*} (x)) = \bar {a} a \Phi (x, x) + \Phi (z, z) = \Phi (u (x), u (x)) + \Phi (z, z).
\]

Now, since \(u\) is normal, we have

\[
\begin{array}{r l} \Phi (u (x), u (x)) & = \Phi (x, u ^ {*} u (x)) = \Phi (x, u u ^ {*} (x)) \\ & = \Phi (u u ^ {*} (x), x) = \Phi (u ^ {*} (x), u ^ {*} (x)). \end{array}
\]

Consequently we have \(\Phi(z, z) = 0\) , which, by hypothesis, implies \(z = 0\) and \(u^{*}(x) = \bar{a}x\) .Q.E.D.

Remark. --- It follows from th. 2 that the eigenspaces corresponding to two distinct eigenvalues of \(u\) are orthogonal.

PROPOSITION 5. --- Let u be a Hermitian endomorphism of E. The eigenvalues of u belong to K, and there exists an orthonormal basis of E consisting of eigenvectors of u.

When A is equal to K(i) or to the field of quaternions over K, it suffices, by virtue of prop. 4, to prove the first assertion; now, if x is a nonzero eigenvector of u and a is the corresponding eigenvalue, we have, by virtue of the hypothesis u = u*,

\[
a \Phi (x, x) = \Phi (u (x), x) = \Phi (x, u (x)) = \Phi (x, x) \bar {a};
\]

since \(\Phi(x, x)\) is a nonzero element of the center of A, it follows that \(a = \bar{a}\) , hence \(a \in K\) . Consequently a Hermitian matrix has all its eigenvalues in K. Suppose now that A is equal to K. The matrix \(M\) of \(u\) with respect to an orthonormal basis of E is then symmetric; it is therefore a Hermitian matrix if one considers it as a matrix over \(K(i)\) . The first part of the proof then shows that this matrix has all its eigenvalues in K, and hence admits, if \(E \neq \{0\}\) , eigenvectors \(\neq 0\) in E. It follows that every stable subspace for \(u\) which is minimal is necessarily of dimension 1, and the conclusion follows immediately, as in prop. 4.

PROPOSITION 6. --- a) Let Ψ be a Hermitian form (resp. symmetric bilinear form if A = K or A = K(i)) on E. There exists a basis of E which is orthonormal for Φ and orthogonal for Ψ.

\begin{enumerate}
\def\labelenumi{\alph{enumi})}
\setcounter{enumi}{1}
\tightlist
\item
  Suppose A = K or A = K(i), and let Ψ be an alternating bilinear form on E. There exists a basis of E which is orthonormal for Φ and with respect to which the matrix of Ψ is of the form
\end{enumerate}

\[
\left( \begin{array}{c c c c} 0 & a _ {1} & 0 & 0 \ldots 0 \\ - a _ {1} & 0 & 0 & 0 \ldots 0 \\ 0 & 0 & 0 & a _ {2} \ldots 0 \\ 0 & 0 & - a _ {2} & 0 \ldots 0 \\ \ldots \ldots \end{array} \right).
\]

where the \(a_{i}\) are \(\geqslant 0\) in K.

When \(\Psi\) is Hermitian, our assertion follows immediately from prop. 5 and the canonical correspondence between Hermitian forms on E and Hermitian endomorphisms for \(\Phi\) . In the other two cases, let \(u\) be the semilinear map from E to itself defined at the beginning of this n° by the formula \(\Psi(x, y) = \overline{\Phi(u(x), y)}\) . Then \(u^2\) is an A-linear map; when \(\Psi\) is symmetric (resp. alternating), one has \(u = u^*\) (resp. \(u = -u^*\) ), hence \(u^2\) is Hermitian; whence also

\[
\begin{array}{r l} \Phi (u ^ {2} (x), x) & = \overline {{\Phi (x , u ^ {2} (x))}} = \Phi (u ^ {*} (x), u (x)) = \Phi (u (x), u (x)) \\ & (\text { resp. } - \Phi (u (x), u (x))) \end{array}
\]

this shows that the Hermitian form \((x, y) \to \Phi(u^{2}(x), y)\) is positive (resp. negative). Then apply th. 2, and let V be a minimal element of the family of subspaces \(\neq\{0\}\) of E stable under u and under \(u^{*}\) . Since \(u^{2}\) is a Hermitian endomorphism, V contains an eigenvector \(x \neq 0\) of \(u^{2}\) ; set \(u^{2}(x) = ax\) , where \(a \in K\) (prop. 5).

When \(\Psi\) is symmetric, the inequality \(\Phi(u^{2}(x), x) \geqslant 0\) shows that one has \(a \geqslant 0\) . Set \(y = a^{1/2}x + u(x)\) ; we have \(u(y) = a^{1/2}u(x) + ax = a^{1/2}y\) , and Ay is stable under \(u\) and \(u^{*} = u\) . If \(y = 0\) , Ax is stable under \(u\) and \(u^{*}\) ; in all cases, V is a subspace of dimension 1, and our conclusion follows immediately from th. 2.

When \(\Psi\) is alternating, one has \(a \leqslant 0\) ; set

\[
y = (- a) ^ {1 / 2} x + u (x), \quad z = (- a) ^ {1 / 2} x - u (x).
\]

\[
\text { One has } u (y) = - (- a) ^ {1 / 2} z, u (z) = (- a) ^ {1 / 2} y \text { and }
\]

\[
\Phi (y, z) = - a \Phi (x, x) - \Phi (u (x), u (x)) = - \Phi (a x, x) + \Phi (u ^ {2} (x), x) = 0.
\]

If \(y = z = 0\) , one has \(a = 0\) and \(u(x) = 0\) , hence \(Ax\) is stable under \(u\) and \(u^{*}\) , \(V\) is of dimension 1, and the matrix of the restriction of \(\Psi\) to \(V\) is zero since \(\Psi\) is alternating. Otherwise, \(y\) and \(z\) are both \(\neq 0\) , they generate \(V\) , and since they are orthogonal, \(V\) is of dimension 2 and the matrix of the restriction of \(\Psi\) to \(V\) is of the form \(\begin{pmatrix} 0 & b \\ -b & 0 \end{pmatrix}\) . Finally, by virtue of the formula \(\Psi(x', y') = \overline{\Phi(u(x'), y')}\) , one has \(\Psi(x', y') = 0\) when \(x'\) and \(y'\) belong to two subspaces stable under \(u\) and orthogonal for \(\Phi\) . This proves the existence of an orthonormal basis of \(E\) (for \(\Phi\) ) with respect to which the matrix of \(\Psi\) has the indicated form. Q.E.D.

Exercises. --- 1) Suppose the hypotheses at the beginning of § 7 are satisfied; let \(\Phi\) be a positive Hermitian form on E.

\begin{enumerate}
\def\labelenumi{\alph{enumi})}
\item
  For us to have \(\Phi(x,x)\Phi(y,y)=\Phi(x,y)\overline{\Phi(x,y)}\) , it is necessary and sufficient that \(x\) and \(y\) be linearly dependent, or that the plane spanned by \(x\) and \(y\) be isotropic.
\item
  We assume that, for all \(x \in E\) , \(\Phi(x, x)\) is a square in K. Show that for any two vectors x, y in E, we have
\end{enumerate}

\[
\sqrt {\Phi (x + y , x + y)} \leqslant \sqrt {\Phi (x , x)} + \sqrt {\Phi (y , y)}.
\]

If \(\Phi\) is nondegenerate, the two sides of this inequality can be equal only if \(\alpha x + \beta y = 0\) , where \(\alpha\) and \(\beta\) are two elements of K, not both zero, and such that \(\alpha \beta \leqslant 0\) .

\begin{enumerate}
\def\labelenumi{\arabic{enumi})}
\setcounter{enumi}{1}
\tightlist
\item
  We assume \(\mathbf{A} = \mathbf{K}\) or \(\mathbf{A} = \mathbf{K}(i)\) . Let \(X\) be a square Hermitian matrix of order \(n\) on \(\mathbf{A}\) such that for every nonempty subset \(\mathbf{H}\) of \([1, n]\) one has \(X_{\mathbf{H},\mathbf{H}} \geqslant 0\) (notations of prop. 3 of no. 1).
\end{enumerate}

\begin{enumerate}
\def\labelenumi{\alph{enumi})}
\item
  Let \(\lambda\) be an element \(>0\) of K; show that the matrix \(X + \lambda I\) is positive nondegenerate (use prop. 3 of no. 1, arguing by induction on \(n\) ).
\item
  Deduce that the Hermitian matrix X is positive.
\end{enumerate}

\begin{enumerate}
\def\labelenumi{\arabic{enumi})}
\setcounter{enumi}{2}
\item
  We assume that \(\mathbf{A} = \mathbf{K}\) or \(\mathbf{A} = \mathbf{K}(i)\) and that E is finite-dimensional. Let \(\Phi_1, \Phi_2\) be two positive Hermitian forms on E, and let \(V = (\alpha_{ij})\) , \(W = (\beta_{ij})\) be the matrices of these forms with respect to the same basis \((e_i)\) of E. Show that the Hermitian form \(\Phi\) whose matrix \((\gamma_{ij})\) with respect to \((e_i)\) is such that \(\gamma_{ij} = \alpha_{ij}\beta_{ij}\) for every pair of indices, is positive; moreover, if \(\Phi_1\) and \(\Phi_2\) are nondegenerate, the same holds for \(\Phi\) . (In the computation of \(\Phi(x, x)\) express the \(\alpha_{ij}\) in terms of the values \(\Phi_1(c_i, c_i)\) for an orthogonal basis \((c_i)\) of E relative to \(\Phi_1\) .)
\item
  Suppose that A = K or A = K(i). Let R be a Hermitian matrix of order n over A; we say that R has signature \((s, t)\) if the Hermitian form on \(A^{n}\) having R as its matrix with respect to the canonical basis, has signature \((s, t)\) . We denote by \(\Delta_{k}\) the principal minor of R obtained by deleting from R the rows and columns with index \textgreater{} k ; we assume that \(\Delta_{s+t} \neq 0\) and that, for no index \(k < s + t\) , \(\Delta_{k}\) and \(\Delta_{k+1}\) are not simultaneously zero (cf.~§ 6, exerc. 1 b)). Show that if \(\Delta_{k} = 0\) for a \(k < s + t\) , \(\Delta_{k-1}\) and \(\Delta_{k+1}\) have opposite signs (method of exerc. 1 a) of § 6) and that the number s - t is equal to
\end{enumerate}

\[
\operatorname{sgn} \Delta_ {1} + \operatorname{sgn} \left(\Delta_ {1} \Delta_ {2}\right) + \dots + \operatorname{sgn} \left(\Delta_ {s + t - 1} \Delta_ {s + t}\right)
\]

(use exercise 2 of § 6).

\begin{enumerate}
\def\labelenumi{\arabic{enumi})}
\setcounter{enumi}{4}
\item
  We assume that A = K or A = K(i); let R, S be two Hermitian matrices over A, with signatures \((s, t)\) and \((s', t')\) respectively (exerc. 4); show that the matrix \(R \otimes S\) is a Hermitian matrix of signature \((ss' + tt', st' + s't)\) .
\item
  Let K be a maximal ordered field, L a finite-dimensional simple algebra over K. If \((s, t)\) is the signature of the symmetric bilinear form \((x, y) \to \mathrm{Tr}_{\mathrm{L/K}}(xy)\) on L, show that we have:
\end{enumerate}

\(s-t=m\) if L is isomorphic to a matrix algebra of order \(m\) over K;

s - t = 0 if L is isomorphic to a matrix algebra over the field K(i);

\(s - t = -2m\) if L is isomorphic to a matrix algebra of order \(m\) over the field of quaternions over K.

\begin{enumerate}
\def\labelenumi{\arabic{enumi})}
\setcounter{enumi}{6}
\tightlist
\item
  Assume that A satisfies the conditions at the beginning of §7 and that E has finite dimension n. Let Φ be a positive nondegenerate Hermitian form on E.
\end{enumerate}

\begin{enumerate}
\def\labelenumi{\alph{enumi})}
\item
  Show that every similarity for \(\Phi\) is written in a unique way as the product of a homothety with ratio \(\geqslant 0\) in K, and of a unitary transformation.
\item
  For every basis \((a_{i})_{1\leqslant i\leqslant n}\) of E, show that there exists one and only one orthonormal basis \((e_{i})_{1\leqslant i\leqslant n}\) of E satisfying the following conditions: \(1^{\circ}\) for all \(m\) such that \(1\leqslant m\leqslant n\) , the subspace spanned by \(a_1,\ldots ,a_m\) is identical to the subspace spanned by \(e_1,\ldots ,e_m\) ; \(2^{\circ}\) we have \(\Phi (a_i,e_i) > 0\) in K, for every index \(i\) (cf. § 6, no. 1, prop. 1).
\item
  Deduce from \(b\) that for every invertible square matrix \(M\) of order \(n\) on A, there exists one and only one pair of matrices \((L, U)\) of order \(n\) such that \(U\) is a unitary matrix, that \(L = (\lambda_{ij})\) has only zeros below its diagonal and diagonal entries \(\lambda_{ii}\) belonging to K and \(>0\) and finally that \(M = LU\) .
\end{enumerate}

Assume A = K; let L be a finite-dimensional Hermitian space over Λ, whose metric form is positive non-degenerate. Let M, N be two subsets of L, and u a bijection from M onto N such that one has \(e(u(a), u(b)) = e(a, b)\) for any points a, b in M, show that there exists a displacement whose restriction to M is equal to u (argue by induction on the dimension of the affine linear variety spanned by M, as in the Gram-Schmidt orthogonalization process). Does the proposition extend to the case where A equals K(i) or the field of quaternions over K?

\begin{enumerate}
\def\labelenumi{\arabic{enumi})}
\setcounter{enumi}{8}
\item
  Assume the conditions of No. 3 are satisfied, and in addition that A is the field of quaternions over K. Let u be a normal endomorphism of E. Show that E is the direct sum of subspaces \(F_{k}\) ( \(1 \leqslant k \leqslant r\) ), pairwise orthogonal, such that in each of the \(F_{k}\) there exists an orthonormal basis ( \(e_{ik}\) ) ( \(1 \leqslant i \leqslant n_{k}\) ) and that one has \(u(e_{ik}) = \lambda_{k}e_{ik}\) for \(1 \leqslant i \leqslant n_{k}\) , two \(\lambda_{k}\) with distinct indices not being transformed into each other by an inner automorphism of A. Moreover, if ( \(F_{k}'\) ) is a second decomposition of E having the same properties, with a system ( \(\lambda_{k}'\) ) of eigenvalues, one has \(F_{k}' = F_{k}\) (up to a permutation of the indices) and \(\lambda_{k}' = \alpha_{k}\lambda_{k}\alpha_{k}^{-1}\) ; the set of eigenvectors of u for the eigenvalue \(\lambda_{k}\) is the subspace over the commutant \(A_{k}\) of \(\lambda_{k}\) in A, generated by the \(e_{ik}\) ( \(1 \leqslant i \leqslant n_{k}\) ).
\item
  Assume the conditions of No. 3 are satisfied, and in addition that A is equal to K(i) or to the field of quaternions over K. Let u be a normal endomorphism of E.
\end{enumerate}

\begin{enumerate}
\def\labelenumi{\alph{enumi})}
\item
  Show that for a vector subspace F of E to be such that \(u(\mathrm{F}) \subset \mathrm{F}\) , it is necessary and sufficient that F be generated by eigenvectors of u; we then have \(u(\mathrm{F}^{0}) \subset \mathrm{F}^{0}\) and \(u^{*}(\mathrm{F}) \subset \mathrm{F}\) .
\item
  For \(u\) to be unitary (resp. for \(u^{*} = u\) , \(u^{*} = -u\) ), it is necessary and sufficient that for every eigenvalue \(\lambda\) of \(u\) one has \(\lambda \bar{\lambda} = 1\) (resp. \(\bar{\lambda} = \lambda\) , \(\bar{\lambda} = -\lambda\) ).
\item
  We say that a Hermitian endomorphism \(u\) of E is positive (resp. positive and nondegenerate) if the Hermitian form \((x, y) \to \Phi(u(x), y)\) canonically corresponding to it is positive (resp. positive and nondegenerate); for this it is necessary and sufficient that all eigenvalues of \(u\) be \(\geqslant 0\) (resp. \(>0\) ).
\item
  Show that for every integer m\textgreater{} 0, there exists a normal endomorphism v of E such that \(v^{m}=u\) . If u is positive Hermitian, there exists a unique positive Hermitian endomorphism v such that \(v^{m}=u\) , and there exists a polynomial \(f\in K[X]\) such that \(v=f(u)\) ; this last result is also valid when A = K.
\end{enumerate}

\begin{enumerate}
\def\labelenumi{\arabic{enumi})}
\setcounter{enumi}{10}
\tightlist
\item
  Assume the conditions of no. 3 are satisfied, and in addition that A = K. Let u be a normal endomorphism of E.
\end{enumerate}

\begin{enumerate}
\def\labelenumi{\alph{enumi})}
\item
  Let V be a minimal element of the set of subspaces of E not reduced to 0, stable under u and \(u^{*}\) ; show that if V has dimension 2, the restriction of u to V is a direct similarity with multiplier \textgreater{} 0.
\item
  Show that every subspace of E stable under u is also stable under \(u^{*}\) . (Let \(E_{0}\) be the vector space obtained from E by extending the field of scalars to \(\mathrm{K}(i)\) ; \(\Phi\) is the restriction to E of a positive nondegenerate Hermitian form on \(E_{0}\) and u the restriction to E of a normal endomorphism \(u_{0}\) of \(E_{0}\) (for this form); moreover, E is the set of \(x \in E_{0}\) invariants under an involutive semilinear bijection j of \(E_{0}\) , and we have \(u_{0}j = ju_{0}\) . Then apply exerc. 10 a).)
\end{enumerate}

\begin{enumerate}
\def\labelenumi{\arabic{enumi})}
\setcounter{enumi}{11}
\item
  Assume the conditions of No. 3 are satisfied, and in addition that A is equal to K(i) or to the field of quaternions over K. Show that for every endomorphism u of E, there exists an orthonormal basis of E with respect to which the matrix of u has only zeros below the diagonal. (Proceed by induction on the dimension of E, considering an eigenvector of u.)
\item
  Let A be a field satisfying the conditions at the beginning of § 7, and let E, F be two finite-dimensional vector spaces over A, \(\Phi\) (resp. \(\Psi\) ) a positive non-degenerate Hermitian form on E (resp. F). Show that, for every linear map \(u\) from E to F, the adjoint \(u^{*}\) of \(u\) (for \(\Phi\) and \(\Psi\) ; cf. § 1, no. 8) is such that \(u^{*}u\) and \(uu^{*}\) are positive Hermitian endomorphisms (exercise 10 c)) of E and F, respectively.
\item
  Assume the conditions of No. 3 are satisfied. Let u be an endomorphism of E, \(h_{1}\) , \(h_{2}\) the two positive Hermitian endomorphisms of E, such that we have \(h_{1}^{2}=u^{*}u\) , \(h_{2}^{2}=uu^{*}\) (Exercises 13 and 10 d)).
\end{enumerate}

\begin{enumerate}
\def\labelenumi{\alph{enumi})}
\item
  Show that there exists a unitary endomorphism \(\nu\) such that \(u = \nu h_{1} = h_{2}\nu\) and in particular that \(h_{1}\) and \(h_{2}\) are similar (note that \(\bar{u}(0) = \bar{h}_{1}^{-1}(0)\) and that if V is the subspace orthogonal to \(\bar{u}(0)\) , one has \(\Phi(u(x), u(x)) = \Phi(h_{1}(x), h_{1}(x))\) for all \(x \in \mathrm{V}\) . For \(\nu\) to be uniquely determined, it is necessary and sufficient that \(u\) be bijective. In order for one to be able to take \(\nu\) commuting with \(h_{1}\) , it is necessary and sufficient that \(u\) be normal.
\item
  From a), deduce that every square matrix M of order n over A can be written as UDV, where U and V are unitary matrices and D is a diagonal matrix whose entries are \(\geqslant 0\) and have as their squares the eigenvalues of MM*.
\end{enumerate}

\begin{enumerate}
\def\labelenumi{\arabic{enumi})}
\setcounter{enumi}{14}
\tightlist
\item
  Assume the conditions of No. 3 are satisfied. Show that every positive Hermitian matrix H over A can be written in the form \(LL^{*}\) , where \(L = (\lambda_{ij})\) has only zeros below its diagonal and diagonal terms belonging to K and \(\geqslant 0\) ; moreover, L is uniquely determined by these conditions when H is invertible (cf.~exerc. 10 d) and 7 c).
\end{enumerate}

¶ 16) We assume that the conditions of no. 3 are satisfied and, in addition, that A = K(i). Let u be an endomorphism of E.

\begin{enumerate}
\def\labelenumi{\alph{enumi})}
\item
  The set of eigenvalues of \(u\) is contained in the set U of values of \(\Phi(x, u(x))\) when \(x\) ranges over the set of elements of E such that \(\Phi(x, x) = 1\) .
\item
  A subset C of A = K(i) is said to be convex if, for every pair of elements \((\xi, \eta) \in \mathrm{C}^{2}\) and every \(\tau \in K\) such that \(0 \leqslant \tau \leqslant 1\) we have \(\tau\xi + (1 - \tau)\eta \in C\) . Show that the set U is convex. (Reduce to the case where n = 2; by writing u in the form \(\nu + iw\) , where \(\nu\) and w are Hermitian, and replacing u if necessary by \(\lambda u\) , where \(\lambda \in A\) and \(\lambda\overline{\lambda} = 1\) show that everything reduces to proving the following property. Let \(f(\xi_{1}, \xi_{2}) = \xi_{1}\overline{\xi}_{1} + \xi_{2}\overline{\xi}_{2}\) , \(g(\xi_{1}, \xi_{2}) = a\xi_{1}\overline{\xi}_{1} + b\xi_{2}\overline{\xi}_{2}\) ( \(a \in K, b \in K\) ), \(h(\xi_{1}, \xi_{2}) = \alpha\xi_{1}\overline{\xi}_{1} + \beta\xi_{1}\overline{\xi}_{2} + \beta\xi_{2}\overline{\xi}_{1} + \gamma\xi_{2}\overline{\xi}_{2}\) ( \(\alpha \in K, \gamma \in K, \beta \in A\) ); let \((\eta_{1}, \eta_{2}) \in A^{2}\) , \((\zeta_{1}, \zeta_{2}) \in A^{2}\) such that \(f(\eta_{1}, \eta_{2}) = f(\zeta_{1}, \zeta_{2}) = 1\) , \(g(\eta_{1}, \eta_{2}) = g(\zeta_{1}, \zeta_{2}) = 1\) , \(h(\eta_{1}, \eta_{2}) > 0\) , \(h(\zeta_{1}, \zeta_{2}) < 0\) ; there then exists \((\theta_{1}, \theta_{2}) \in A^{2}\) such that \(f(\theta_{1}, \theta_{2}) = 1\) , \(g(\theta_{1}, \theta_{2}) = 1\) and \(h(\theta_{1}, \theta_{2}) = 0\) . One will begin by noting that for every pair \((\xi_{1}, \xi_{2})\) , there exist \(\mu \in A\) such that \(\mu\overline{\mu} = 1\) and \(\beta\mu\xi_{1}\overline{\xi}_{2} + \beta\overline{\mu}\xi_{2}\overline{\xi}_{1} = 0\) , which will make it possible to reduce to the case where \(\beta = 0\) ; use prop. 5 of chap.~VI, § 2, no.~5.)
\item
  Show that if \(u\) is normal, U is the smallest convex set containing all the eigenvalues of \(u\) . Give an example where \(u\) is not normal but where U still possesses the preceding property. (Take as eigenvalues of \(u\) the elements \(\pm 1 \pm i\) as simple eigenvalues, and 0 as a double eigenvalue.)
\end{enumerate}

¶ 17) We assume the conditions of no.~3 are satisfied; for every \(\xi\in\mathbf{A}\) , we denote by \(|\xi|\) the element \(\rho\geqslant0\) of K such that \(\rho^{2}=\xi\bar{\xi}\) (absolute value of \(\xi\) ). For every square matrix \(M=(\alpha_{ij})\) over A, we set \(f(M)=\max_{i}|\alpha_{ii}|\) , \(g(M)=\max_{i,j}|\alpha_{ij}|\) , and denote by \(\varphi(M)\) the largest absolute value of the eigenvalues of \(M\) (it will be shown that this definition makes sense when A is the field of quaternions over K).

\begin{enumerate}
\def\labelenumi{\alph{enumi})}
\tightlist
\item
  Let \(A, B, D\) be three square matrices of order \(n\) over A. Show that if \(D\) is diagonal, then
\end{enumerate}

\[
g ^ {2} (A D B ^ {*}) \leqslant f ^ {2} (D) f (A A ^ {*}) f (B B ^ {*})\tag{1}
\]

(use inequality (1) of no. 1). Deduce that

\[
g (A B B ^ {*} A ^ {*}) \leqslant f (A A ^ {*}) \varphi (B B ^ {*})\tag{2}
\]

(apply prop. 5 to the Hermitian matrix \(BB^{*}\) ). Deduce that, for \(m\) arbitrary square matrices \(A_{i}\) ( \(1 \leqslant i \leqslant m\) ) over A, we have

\begin{enumerate}
\def\labelenumi{(\arabic{enumi})}
\setcounter{enumi}{2}
\tightlist
\item
\end{enumerate}

\[
g ^ {2} (A _ {1} A _ {2} \dots A _ {m}) \leqslant f (A _ {1} A _ {1} ^ {*}) \varphi (A _ {2} A _ {2} ^ {*}) \dots \varphi (A _ {m - 1} A _ {m - 1} ^ {*}) f (A _ {m} ^ {*} A _ {m})\tag{4}
\]

\[
\varphi^ {2} (A _ {1} A _ {2} \dots A _ {m}) \leqslant \varphi (A _ {1} A _ {1} ^ {*}) \varphi (A _ {2} A _ {2} ^ {*}) \dots \varphi (A _ {m} A _ {m} ^ {*})\tag{5}
\]

\[
g ^ {2} (A _ {1} A _ {2} \dots A _ {m}) \leqslant \varphi (A _ {1} A _ {1} ^ {*}) \dots \varphi (A _ {m} A _ {m} ^ {*}).
\]

(For (3), reason by induction starting from (2). For (4), use exerc. 12; finally deduce (5) from (3).)

Show that inequality (3) no longer necessarily holds when one replaces \(A_{m}^{*}A_{m}\) by \(A_{m}A_{m}^{*}\) (observe that one can have \(f(A^{*}A)\neq f(A A^{*})\) ), or when one replaces \(\varphi(A_{i}A_{i}^{*})\) by \(f(A_{i}A_{i}^{*})\) for \(2\leqslant i\leqslant m-1\) (take all the \(A_{i}\) equal to the square matrix all of whose elements are 1).

\begin{enumerate}
\def\labelenumi{\alph{enumi})}
\setcounter{enumi}{1}
\tightlist
\item
  Let u be a normal endomorphism of E; if λ is an eigenvalue of u (an element of K(i) when A = K), then \(\lambda\bar{\lambda}\) is an eigenvalue of u*u. For every normal matrix M (matrix of a normal endomorphism of E with respect to an orthonormal basis), we therefore have φ²(M) = φ(MM*). Deduce that one also has g(M) ≤ φ(M), and more generally, if M₁, . . ., Mₘ are normal,
\end{enumerate}

\begin{enumerate}
\def\labelenumi{(\arabic{enumi})}
\setcounter{enumi}{5}
\tightlist
\item
\end{enumerate}

\[
g (M _ {1} \dots M _ {m}) \leqslant \varphi (M _ {1}) \dots \varphi (M _ {m})\tag{7}
\]

\[
\varphi (M _ {1} \dots M _ {m}) \leqslant \varphi (M _ {1}) \dots \varphi (M _ {m})
\]

(use (4) and (5)).

\begin{enumerate}
\def\labelenumi{\alph{enumi})}
\setcounter{enumi}{2}
\tightlist
\item
  Show that if \(H\) is a positive Hermitian matrix over A, we have \(f(H) = g(H) \leqslant \varphi(H)\) (use (3) and b), by writing \(H = AA^{*}\) ).
\end{enumerate}

¶ 18) We assume the conditions of no. 3 are satisfied.

\begin{enumerate}
\def\labelenumi{\alph{enumi})}
\tightlist
\item
  For every endomorphism u of E, let \((e_{i})\) be an orthonormal basis of E consisting of eigenvectors of \(u^{*}u\) , and let \(\rho_{i}\) be the eigenvalue of \(u^{*}u\) corresponding to the vector \(e_{i}\) ; we set \(s(u) = (\sum_{i=1}^{n} \rho_{i})^{1/2}\) (square root of \(\mathrm{Tr}(u^{*}u)\) when A is commutative); we have \(s(u^{*}) = s(u)\) . If u, v are two endomorphisms of E, and U, V their matrices with respect to the same orthonormal basis of E, show that for the matrix \(UV^{*} + VU^{*}\) , with elements in K, one has \((\mathrm{Tr}(UV^{*} + VU^{*}))^{2} \leqslant 4s(u)s(v)\) (note that \((u^{*} + \lambda v^{*})(u + \lambda v)\) is positive Hermitian for every \(\lambda \in K\) ); deduce that \(s(u + v) \leqslant s(u) + s(v)\) . If in addition A is commutative, one has
\end{enumerate}

\[
\left| \operatorname{Tr} (u v) \right| ^ {2} \leqslant s (u) s (v) \quad \text { and } \quad \left| \operatorname{Tr} (u) \right| \leqslant \sqrt {n} s (u)
\]

(write \(u\) as a product using Exerc. 14 b)).

\begin{enumerate}
\def\labelenumi{\alph{enumi})}
\setcounter{enumi}{1}
\tightlist
\item
  We further assume A is commutative (hence equal to K or K(i)). If \(H_{1}, H_{2}\) are two positive Hermitian square matrices, show that one has \(\mathrm{Tr}(H_{1}H_{2}) \geqslant 0\) (reduce to the case where \(H_{2}\) is a diagonal matrix). Deduce that we have (using the notation of exercise 17)
\end{enumerate}

\[
\left| \operatorname{Tr} (H _ {1} H _ {2}) \right| \leqslant \varphi (H _ {1}) \operatorname{Tr} (H _ {2}) \leqslant \operatorname{Tr} (H _ {1}) \operatorname{Tr} (H _ {2}).
\]

From these inequalities, conclude that for any two endomorphisms \(u, \nu\) of E, we then have \(s(u\nu) \leqslant s(u)s(\nu)\) .

\begin{enumerate}
\def\labelenumi{\alph{enumi})}
\setcounter{enumi}{2}
\tightlist
\item
  Assuming A is still commutative, let \(\prod_{i=1}^{n}(x-\lambda_i)\) be the decomposition into linear factors of the characteristic polynomial of an arbitrary endomorphism \(u\) of E; show that one has \(\sum_{i=1}^{n}|\lambda_i|^2\leqslant(s(u))^2\) and that, for the two sides of this inequality to be equal, it is necessary and sufficient that \(u\) is normal (use Exerc. 12).
\end{enumerate}

\begin{enumerate}
\def\labelenumi{\arabic{enumi})}
\setcounter{enumi}{18}
\item
  Assume the conditions of no. 3 are satisfied. Let M be a square matrix of order n over A; show that for every square submatrix N of M (obtained by deleting from M a certain number of rows and the columns with the same indices as those rows), one has (with the notations of exercise 17) \(\varphi^{2}(N) \leqslant \varphi(MM^{*})\) (apply formula (4) of exerc. 17 appropriately). If in particular M is a normal matrix, \(\varphi(N) \leqslant \varphi(M)\) (cf. exercise 17 b)).
\item
  We assume that the conditions of No. 3 are satisfied and, in addition, that A is equal to K or to K(i). Apply the results of Exercises 17 to 19 to the extension of Φ to the p-th exterior powers (§ 1, No. 9) and to the p-th exterior powers of the endomorphisms or matrices under consideration. In particular, show that, if \(\prod_{i=1}^{n}(X-\lambda_i)\) and \(\prod_{i=1}^{n}(X-\rho_i^2)\) are the decompositions into linear factors of the characteristic polynomials of an
\end{enumerate}

endomorphism \(u\) of E and of the endomorphism \(u^*u\) and if we suppose \(|\lambda_i| \geqslant |\lambda_{i+1}|\) and \(\rho_i \geqslant \rho_{i+1} \geqslant 0\) for \(1 \leqslant i \leqslant n-1\) , one has

\[
\mid \lambda_ {1} \lambda_ {2} \dots \lambda_ {h} \mid \leqslant \rho_ {1} \rho_ {2} \dots \rho_ {h}
\]

for \(1 \leqslant h \leqslant n - 1\) and \(|\lambda_1\lambda_2 \ldots \lambda_n| = \rho_1\rho_2 \ldots \rho_n\) .

\begin{enumerate}
\def\labelenumi{\arabic{enumi})}
\setcounter{enumi}{20}
\tightlist
\item
  We suppose the conditions of No. 3 are satisfied. Let u be a Hermitian endomorphism of E and let \(\prod_{i=1}^{n}(X-\lambda_{i})\) be the decomposition into linear factors of its characteristic polynomial; we assume \(\lambda_{i}\geqslant\lambda_{i+1}\) for \(1\leqslant i\leqslant n-1\) .
\end{enumerate}

\begin{enumerate}
\def\labelenumi{\alph{enumi})}
\item
  Show that the largest (resp. the smallest) of the eigenvalues \(\lambda_{i}\) in K is equal to the largest (resp. the smallest) of the values of \(\Phi(u(x), x)\) when \(x\) ranges over the set of \(x \in E\) such that \(\Phi(x, x) = 1\) . (Reason directly, or apply Exerc. 16 c) by reducing to the case where A = K(i) ( \(\S\) 3, exercise 4).)
\item
  Let \(\Psi_{\mathrm{v}}\) be the restriction to a vector subspace V of E of the Hermitian form \(\Psi\) associated with \(u\) , and let \(u_{\mathrm{v}}\) be the Hermitian endomorphism of V associated with \(\Psi_{\mathrm{v}}\) . Show that \(\lambda_{k}\) is the smallest of the largest eigenvalues of the \(u_{\mathrm{v}}\) as V ranges over the set of vector subspaces of E of dimension \(n - k + 1\) (using Prop. 5).
\end{enumerate}

\begin{enumerate}
\def\labelenumi{\arabic{enumi})}
\setcounter{enumi}{21}
\tightlist
\item
  We assume that the conditions of no. 3 are satisfied, and in addition that A is equal to K(i) or to the field of quaternions over K. Let u, v be two normal endomorphisms of E; let (E \(_{i}\) ) \(_{1\leqslant i\leqslant r}\) (resp. (F \(_{j}\) ) \(_{1\leqslant j\leqslant s}\) ) be the decomposition of E into a direct sum of pairwise orthogonal subspaces, such that in E \(_{i}\) (resp. F \(_{j}\) ) there is an orthonormal basis formed by eigenvectors of u (resp. v) for the same eigenvalue \(\lambda_{i}\) (resp. \(\mu_{j}\) ), and that \(\lambda_{h}\) and \(\lambda_{i}\) (resp. \(\mu_{j}\) and \(\mu_{k}\) ) are not transformed into one another by an inner automorphism of A if \(h\neq i\) (resp. \(j\neq k\) ) (cf. prop. 4 and exerc. 9). For an endomorphism w of E to satisfy uw = wν, it is necessary and sufficient that for every j ( \(1\leqslant j\leqslant s\) ), w(F \(_{j}\) ) be contained in one of the E \(_{i}\) and that the image under w of every eigenvector of v relative to the eigenvalue \(\mu_{j}\) be an eigenvector of u relative to the eigenvalue \(\mu_{j}\) (which implies in particular that \(\mu_{j}\) and \(\lambda_{i}\) are transformed into one another by an inner automorphism of A). Deduce that if this is so, then one has u*w = wν*.
\end{enumerate}

¶ 23) We assume that the conditions of no. 3 are satisfied. Let u and v be two endomorphisms of E.

\begin{enumerate}
\def\labelenumi{\alph{enumi})}
\item
  We assume that u and uv are normal. For vu to be normal, it is necessary and sufficient that v and \(u^{*}u\) commute. (To see that the condition is necessary, use the relation \(u(vu) = (uv)u\) and exerc. 22; to see that it is sufficient, use exerc. 14 a) and 10 d).
\item
  We assume \(u, v\) and \(uv\) normal; let \(\beta\) be the largest eigenvalue of \(uu^*\) , and let \(F\) be the subspace of \(E\) formed by the eigenvectors of \(uu^*\) relative to the eigenvalue \(\beta\) . Show that \(v(F) \subset F\) . (Note that for every normal endomorphism \(u\) and every \(x \in E\) , one has
\end{enumerate}

\[
\Phi (u (x), u (x)) = \Phi (u ^ {*} (x), u ^ {*} (x)).
\]

Reduce to the case where A is equal to K(i) or to the field of quaternions over K; F then admits a basis formed by eigenvectors for u (and \(u^{*}\) ); note that for such a vector z, we have \(\Phi(u^{*}u\nu(z),\nu(z))=\beta\Phi(\nu(z),\nu(z))\) Deduce from this that \(\nu u\) is normal (argue by induction on the number of distinct eigenvalues of \(u^{*}u\) ). If h (resp. \(h'\) ) is the positive Hermitian endomorphism such that \(h^{2}=uu^{*}\) (resp. \(h'^{2}=\nu\nu^{*}\) ) and if we set \(u=hu_{1}\) , \(\nu=h'\nu_{1}\) , where \(u_{1}\) and \(\nu_{1}\) are unitary, h commuting with \(u_{1}\) and \(h'\) with \(\nu_{1}\) (exerc. 14 a)), show that the pairs \((h,h')\) , \((h,\nu_{1})\) and \((h',u_{1})\) are permutable; and conversely. Deduce that \(u^{m}\nu^{n}\) , \(\nu^{n}u^{m}\) , \(u\nu^{*}\) and \(\nu^{*}u\) are then normal (m and n integers \textgreater0 arbitrary).

¶ 24) Assume the conditions of no. 3 are satisfied. Let Γ be a group of automorphisms of E such that every \(u \in \Gamma\) is normal. Show that there exists a decomposition of E into a direct sum of subspaces \(E_k\) ( \(1 \leqslant k \leqslant r\) ) pairwise orthogonal, such that the restriction to \(E_k\) of every \(u \in \Gamma\) is of the form \(\lambda_k v_k\) , where \(\lambda_k\) is an element \textgreater0 of K and where \(v_k\) is a unitary endomorphism of \(E_k\) (decompose each \(u \in \Gamma\) in the form \(h\nu\) , where \(v\) is unitary, \(h\) positive Hermitian, and \(h^2 = uu^*\) (Exerc. 14a); use Exerc. 23b and apply Cor. 2 of Thm. 2 to the algebra generated by the Hermitian endomorphisms \(h\) corresponding to the \(u \in \Gamma\) ). Deduce that if Γ is finite, the elements \(u \in \Gamma\) are unitary.

\begin{enumerate}
\def\labelenumi{\arabic{enumi})}
\setcounter{enumi}{24}
\tightlist
\item
  Suppose that A = K (maximal ordered field). Let L be an n-dimensional Euclidean space over A, whose metric form is positive non-degenerate.
\end{enumerate}

Let S be a nondegenerate affine quadric in L (§ 6, exerc. 25). If S admits a center a and if one takes a as origin in L, show that there exists an orthonormal basis \((e_{i})\) for L such that, with respect to this basis, S has the equation \(\lambda_{1}\xi_{1}^{2} + \cdots + \lambda_{n}\xi_{n}^{2} = 1\) . Moreover, if two orthonormal bases have this property, the elements \(\lambda_{i} \in K\) which correspond to them are the same up to order.

If S admits no center, show that there exists a point b of S and (taking b as origin) an orthonormal basis for L such that, with respect to this basis, S has equation \(\lambda_{1}\xi_{1}^{2} + \cdots + \lambda_{n-1}\xi_{n-1}^{2} + \xi_{n} = 0\) . (If \(\overline{S}\) is the projective quadric such that \(S = L \cap \overline{S}\) , c the pole (at infinity) with respect to \(\overline{S}\) of the hyperplane at infinity \(H_{0}\) , determine b by the condition that the line passing through b and c be perpendicular to the tangent hyperplane to S at the point b).

¶ 26) a) Let K be a commutative field, S a subset of K such that K is a maximal S-orderable field (chap.~VI, § 2, exerc. 8). Let f be a polynomial in K{[}X{]}, L its splitting field. Show that if, for every (total) order structure on K compatible with its ring structure, L admits an ordered extension structure of K, then necessarily L = K. (Argue by contradiction: proceeding as in exerc. 8 e) of chap.~VI, § 2, reduce to the case where one would have {[}L : K{]} = 2. Conclude by noting that if b ∈ K is not a square, there exists a total order structure on K, compatible with its ring structure, for which b \textless{} 0 (cf.~chap.~VI, § 2, n° 3, lemma of th. 1)).

\begin{enumerate}
\def\labelenumi{\alph{enumi})}
\setcounter{enumi}{1}
\tightlist
\item
  Extend the results of n° 3 (prop. 6 excepted) to the case where K is a maximal S-orderable field (*). (Begin by proving the analogue of prop. 5, using a). Then establish the analogue of cor. 2 of th. 2 for the case where the commutative algebra B consists of Hermitian endomorphisms, using prop. 10 of chap.~VIII, § 9, n° 4. Pass to the case of a normal endomorphism u for A = K(i) by noting that one can then write u = v + iw, where v and w are Hermitian and commute. Finally, if A is the quaternion field over K, E₀ the space E considered as a vector space of dimension 2n over K(i), and if one sets Φ(x, y) = Φ₁(x, y) + Φ₂(x, y)j, where Φ₁ and Φ₂ take their values in K(i), note that if u is normal for Φ, it is also normal for Φ₁.)
\end{enumerate}

¶ 27) Let K be a maximal ordered field, E a vector space of dimension n over K, Φ a non-degenerate symmetric bilinear form on E, of signature (s, t) distinct from (n, 0) and (0, n). Let (ei) be an orthogonal basis for Φ, such that Φ(ei, ei) = 1 for 1 ≤ i ≤ s, Φ(ei, ei) = -1 for s + 1 ≤ i ≤ n.~For every orthogonal transformation u ∈ U(Φ), let

\[
U = \left( \begin{array}{c c} M & N \\ P & Q \end{array} \right)
\]

the matrix of \(u\) with respect to \((e_{i})\) , written in the form of a square array of matrices corresponding to the partition of \([1, n]\) into \([1, s]\) and \([s + 1, n]\) .

\begin{enumerate}
\def\labelenumi{\alph{enumi})}
\tightlist
\item
  Prove the relations
\end{enumerate}

\[
\begin{array}{r l} ^ {t} M. M - ^ {t} P. P & = I _ {s} \\ ^ {t} Q. Q - ^ {t} N. N & = I _ {t} \\ ^ {t} M. N - ^ {t} P. Q & = 0. \end{array}
\]

\begin{enumerate}
\def\labelenumi{\alph{enumi})}
\setcounter{enumi}{1}
\item
  Let R be a matrix over K with t rows and s columns, such that \(I_{s} - ^{t}R \cdot R\) is the matrix of a positive non-degenerate Hermitian form (over \(K^{s}\) ). Show that det \((M + \lambda NR)\) does not change sign for \(-1 \leqslant \lambda \leqslant 1\) in K (show, using a), that \(^{t}(M + \lambda NR)(M + \lambda NR)\) is the matrix of a positive non-degenerate symmetric form).
\item
  Set \(\sigma(u) = \sigma(U) = \operatorname{sgn}(\det M)\) ; show that, for any two elements \(u\) , \(\nu\) of the orthogonal group \(\mathbf{U}(\Phi)\) , one has \(\sigma(u\nu) = \sigma(u)\sigma(\nu)\) (using \(a\) ) and \(b\) ). Deduce that the commutator subgroup of the special orthogonal group \(\mathbf{SU}(\Phi)\) is distinct from \(\mathbf{SU}(\Phi)\) (cf.~§ 10, exerc. 9).
\end{enumerate}

\section{Types of quadratic forms}

In this paragraph we assume that A is a commutative field.

\subsection{Types of quadratic forms.}

Given a quadratic form Q (§ 3, n° 4) on a vector space E over A, we shall say that E is the space of definition of Q and that dim(E) is the dimension of Q. Given two quadratic forms Q, Q' on vector spaces E, E' over A, we shall denote by Q τ Q' their direct sum (§ 3, n° 4). Recall that the direct sum of two neutral forms is neutral (§ 4, n° 2).

Introduce the following relation:

“Q and Q' are non-degenerate quadratic forms of finite dimensions over A, and there exist neutral quadratic forms, N, N' such that Q τ N is equivalent to Q' τ N'”.

This relation, which we shall denote by \(Q \sim Q'\) , is manifestly reflexive and symmetric. It is also transitive: indeed, if \(Q\) , \(Q'\) , \(Q''\) are quadratic forms such that \(Q \sim Q'\) and \(Q' \sim Q''\) , there exist neutral quadratic forms \(M\) , \(M'\) , \(N\) , \(N'\) such that \(Q \tau M\) is equivalent to \(Q' \tau M'\) and \(Q' \tau N\) is equivalent to \(Q'' \tau N'\) ; then \(Q \tau (M \tau N)\) is equivalent to \((Q \tau M) \tau N\) , hence to \((Q' \tau M') \tau N\) , and also to \((Q' \tau N) \tau M'\) , hence again to \((Q'' \tau N') \tau M'\) and to \(Q'' \tau (N' \tau M')\) ; since \(M \tau N\) and \(N' \tau M'\) are neutral, one indeed has \(Q \sim Q''\) . The relation \(Q \sim Q'\) is therefore an equivalence relation between \(Q\) and \(Q'\) . It is clear that, if \(Q\) and \(Q'\) are two non-degenerate quadratic forms of finite dimensions and equivalent, one has \(Q \sim Q'\) .

For every quadratic form Q on A, non-degenerate and of finite dimension, we set

\[
\theta (Q) = \tau_ {x} (X \sim Q),\tag{1}
\]

and we shall say that \(\theta(Q)\) is the type of \(Q\) . If \(Q\) and \(Q'\) are two quadratic forms on \(A\) , non-degenerate and of finite dimensions, the relations \(Q \sim Q'\) and \(\theta(Q) = \theta(Q')\) are equivalent.

PROPOSITION 1. --- Let Q and Q' be two nondegenerate quadratic forms on A, of finite dimension. For Q and Q' to be equivalent, it is necessary and sufficient that they have the same dimension and the same type.

The condition is obviously necessary. Suppose it is satisfied. Then there exist neutral forms N, N′ such that Q ⊥ N and Q′ ⊥ N′ are equivalent. Since these two forms have the same dimension, so do N and N′, which are consequently equivalent (§ 4, no. 2, cor. 2 of prop. 2). Hence Q and Q′ are equivalent by virtue of Witt's theorem (§ 4, no. 3, cor. 1 of th. 1).

PROPOSITION 2. --- The relation “there exists a nondegenerate quadratic form Q on A, of finite dimension, such that X = θ(Q)” is collectivizing in X (Set Theory, Chap. II, § 1, no. 4).

Let indeed V be a vector space of infinite dimension over A, S the set of nondegenerate quadratic forms defined on the finite-dimensional subspaces of V, and W the set of \(\theta(Q)\) for \(Q \in \mathfrak{S}\) . It is clear that every nondegenerate quadratic form \(Q'\) of finite dimension over A is equivalent to at least one element of \(\mathfrak{S}\) ; whence \(\theta(Q') \in \mathfrak{W}\) , which proves our assertion.

\subsection{Group of types of quadratic forms.}

We shall equip the set \(\mathfrak{W}\) of types of nondegenerate quadratic forms of finite dimensions over A with a structure of a commutative group. We will define an addition in \(\mathfrak{W}\) by the formula

\[
\mathrm{T} + \mathrm{T} ^ {\prime} = \theta (\mathrm{T} _ {\tau} \mathrm{T} ^ {\prime})\tag{2}
\]

\[
(\mathrm{T}, \mathrm{T} ^ {\prime} \text {   in   } \mathfrak {W}),
\]

This addition is commutative since \(T^{\prime} \tau T\) is equivalent to \(T \tau T^{\prime}\) . It is associative because, if T, \(T^{\prime}\) and \(T^{\prime\prime}\) are elements of W, we have

\[
\begin{array}{r l} & (\mathrm{T} + \mathrm{T} ^ {\prime}) + \mathrm{T} ^ {\prime \prime} \sim (\mathrm{T} + \mathrm{T} ^ {\prime}) \tau \mathrm{T} ^ {\prime \prime} \sim (\mathrm{T} \tau \mathrm{T} ^ {\prime}) \tau \mathrm{T} ^ {\prime \prime} \\ & \qquad \sim \mathrm{T} \tau (\mathrm{T} ^ {\prime} \tau \mathrm{T} ^ {\prime \prime}) \sim \mathrm{T} \tau (\mathrm{T} ^ {\prime} + \mathrm{T} ^ {\prime \prime}) \sim \mathrm{T} + (\mathrm{T} ^ {\prime} + \mathrm{T} ^ {\prime \prime}), \end{array}
\]

whence \((\mathrm{T} + \mathrm{T}' ) + \mathrm{T}'' = \mathrm{T} + (\mathrm{T}' + \mathrm{T}'' )\) since two elements of W that have the same type are equal. Moreover, the addition just defined has an identity element: indeed, it is clear that all neutral forms have the same type \(T_{0}\) , namely that of the zero form of dimension zero; one sees immediately that \(T_{0}\) is the sought neutral element. Finally, the existence, for every \(T \in W\) , of an element opposite to T follows immediately from the following proposition:

PROPOSITION 3. — Let Q be a nondegenerate quadratic form of finite dimension on a vector space V over A. Denote by −Q the quadratic form on V defined by \((-Q)(x) = -Q(x) (x \in V)\) . Then the form \(Q \tau (-Q)\) is neutral.

Indeed, the restriction of \(Q \tau (-Q)\) to the diagonal \(D\) of \(V \times V\) is zero. The index of this form is therefore \(\geqslant\frac{1}{2}\dim(V\times V)\) ( \(\S\) 4, no 2, def. 2), and consequently is equal to \(\frac{1}{2}\dim(V\times V)\) (ibid., formula (4)). It follows that \(Q\tau(-Q)\) is neutral (ibid.).

This allows us to state the following definition:

DEFINITION 1. --- The set of types of nondegenerate quadratic forms of finite dimension over A, equipped with the addition defined by (2), is called the group of types of quadratic forms, or Witt group, of A.

Remarks. --- 1) Every nondegenerate quadratic form Q of finite dimension whose type is zero (that is, such that \(\theta(Q) = T_0\) with the notations above) is a neutral form. Indeed, there exist neutral forms N, N′ such that \(Q \tau N\) is equivalent to N′. This shows that Q is of even dimension, hence that there exists a neutral form \(N_1\) of the same dimension as Q. Since Q and \(N_1\) have the same type, it follows from prop. 1 that they are equivalent, hence that Q is neutral.

\begin{enumerate}
\def\labelenumi{\arabic{enumi})}
\setcounter{enumi}{1}
\tightlist
\item
  For every quadratic form Q of finite dimension over A, denote by \(\delta(Q)\) the class modulo 2 of the dimension of Q. We have
\end{enumerate}

\[
\delta (Q \tau Q ^ {\prime}) = \delta (Q) + \delta (Q ^ {\prime}).
\]

Since every neutral form N has even dimension, we have \(\delta(N) = 0\) ; the relation \(Q \sim Q'\) therefore implies \(\delta(Q) = \delta(Q')\) . Thus the restriction of \(\delta\) to the group \(\mathfrak{W}\) of types of quadratic forms on A is a homomorphism from \(\mathfrak{W}\) into the group \(\mathbf{Z}/(2)\) . This homomorphism is surjective when A has characteristic \(\neq 2\) , but is not so if A has characteristic 2, because a quadratic form of odd dimension is then degenerate since its associated bilinear form is alternating (cf.~§ 5).

\begin{enumerate}
\def\labelenumi{\arabic{enumi})}
\setcounter{enumi}{2}
\tightlist
\item
  Let \(a\) be an element \(\neq 0\) of A. If N is a neutral form, then so is \(aN\) . It follows that the relation \(Q \sim Q'\) implies \(aQ \sim aQ'\) . For every element T of the group \(\mathfrak{W}\) , we set
\end{enumerate}

\[
a. \mathrm{T} = \theta (a \mathrm{T}).\tag{3}
\]

We thus obtain an external law of composition between the group A* of nonzero elements of A and the group W. The following formulas follow immediately from the definition:

\[
a. (\mathrm{T} + \mathrm{T} ^ {\prime}) = a. \mathrm{T} + a. \mathrm{T} ^ {\prime}, \quad a b. \mathrm{T} = a. (b. \mathrm{T})\tag{4}
\]

\[
(a, b \text {   in   } A ^ {*}, T, T ^ {\prime} \text {   in   } \mathfrak {W}).
\]

On the other hand, if \(a\) , \(b\) and \(a + b\) are in \(\mathbf{A}^*\) one does not in general have \((a + b)\) . \(\mathrm{T} = a\) . \(\mathrm{T} + b\) . \(\mathrm{T}\) ( \(\mathrm{T} \in \mathfrak{W}\) ).

PROPOSITION 4. --- Let Q be a nondegenerate quadratic form on a finite-dimensional vector space E over A. Suppose A has characteristic \(\neq 2\) , and let \((x_{1},\ldots,x_{n})\) be an orthogonal basis of V. Let \(T_{1}\) denote the type of the quadratic form \(Q_{1}\) defined on the vector space A and such that \(Q_{1}(1)=1\) . The type of Q is then \(\sum_{i=1}^{n}\mathrm{Q}(x_{i}).T_{1}\) .

Indeed, the form Q is equivalent to

\[
(\mathrm{Q} (x _ {1}) \mathrm{Q} _ {1}) \tau \dots \tau (\mathrm{Q} (x _ {n}) \mathrm{Q} _ {1})
\]

COROLLARY. --- With the hypotheses and notations being those of prop. 3, the elements \(a\) . \(\mathrm{T}_{1}\) ( \(a \in \mathrm{A}^{*}\) ) form a set of generators of the group of types of quadratic forms on A.

Thus, determining the structure of the group of types of quadratic forms on A amounts to finding the Z-linear relations that exist among the elements of the form \(a.T_{1}\) . If \(b \in A^{*}\) , the form \(Q_{1}\) defined in Prop. 4 is clearly equivalent to \(b^{2}Q_{1}\) ; one therefore has \(a.T_{1} = ab^{2}.T_{1}\) , which shows that \(a.T_{1}\) depends only on the class of a modulo the subgroup \((A^{*})^{2}\) of squares of elements of \(A^{*}\) . Moreover, it follows from prop. 3 that one has \((-a).T_{1} = -a.T_{1}\) . However, there exist in general other Z-linear relations among the \(a.T_{1}\) besides those which are deduced from the relations we have just indicated.

PROPOSITION 5. --- Suppose that A is a maximal ordered field. Let Q be a nondegenerate quadratic form of finite dimension over A, and let (s, t) be its signature (§ 7, no. 2, def. 2). Then the type of Q is (s - t)·T₁, and the group W of types of quadratic forms over A is an infinite cyclic group generated by T₁.

Indeed, as \(\mathbf{A}^{*}/(\mathbf{A}^{*})^{2}\) has order 2 and that \((-1)\) . \(\mathrm{T}_{1} = -\mathrm{T}_{1}\) , \(\mathfrak{W}\) is generated by \(\mathrm{T}_{1}\) and is consequently monogenic. For all \(n > 0\) , \(n\) . \(\mathrm{T}_{1}\) is the type of positive nondegenerate quadratic forms of dimension \(n\) ; since these forms are not neutral, we have \(n\) . \(\mathrm{T}_{1} \neq 0\) , which shows that \(\mathfrak{W}\) is infinite. Finally, a form of signature \((s, t)\) is isomorphic, with the notations of Prop. 4, to the direct sum of s forms \(Q_{1}\) and of t forms \(-Q_{1}\) (§ 7, no. 2, th. 1); it follows that its type is \((s - t)\) . \(T_{1}\) .

\subsection{Ring of types of quadratic forms.}

We shall assume, in this section, that A is a field of characteristic ≠ 2.

Given two quadratic forms Q, Q' on vector spaces V, V' over A, we shall call the tensor product of Q and Q', and denote by Q ⊗ Q', the quadratic form on V ⊗ V' whose associated bilinear form is the tensor product (§ 1, no. 9, def. 11) of the bilinear forms associated with Q and Q'. One sees readily that Q ⊗ Q' satisfies the relation

\[
(Q \otimes Q ^ {\prime}) (x \otimes x ^ {\prime}) = Q (x) Q ^ {\prime} (x ^ {\prime}) \quad (x \in V, x ^ {\prime} \in V ^ {\prime}).\tag{5}
\]

If Q and Q' are nondegenerate and of finite dimension, then so is Q ⊗ Q' (§ 1, no. 9, prop. 9).

Let Q, Q′, Q″ be quadratic forms on the vector spaces V, V′, V″. Using the canonical isomorphism from V ⊗ V' onto V' ⊗ V (resp. from (V ⊗ V′) ⊗ V″ onto V ⊗ (V′ ⊗ V″), from (V × V′) ⊗ V″ onto (V ⊗ V″) × (V′ ⊗ V″)), we see immediately that Q ⊗ Q' is equivalent to Q' ⊗ Q (resp. (Q ⊗ Q′) ⊗ Q″ to Q ⊗ (Q′ ⊗ Q″), (Q τ Q′) ⊗ Q″ to (Q ⊗ Q″) τ (Q′ ⊗ Q″)).

Let Q and Q' be two non-degenerate quadratic forms of finite dimensions. If Q is neutral, then so is Q⊗Q'. Indeed, let V, V' be the spaces of definition of Q, Q', with dimensions 2n and n', and let W be a totally singular subspace of V of dimension n (§ 4, no. 2); then W⊗V' is a totally singular subspace, and its dimension is half that of V⊗V'; from this one deduces, as in prop. 3, that Q⊗Q' is neutral. Likewise, Q⊗Q' is neutral whenever Q' is neutral.

From this we deduce that, if Q, Q', Q₁, Q₁' are nondegenerate quadratic forms of finite dimension over A, and if one assumes that θ(Q₁) = θ(Q) and θ(Q₁') = θ(Q'), then one has θ(Q₁ ⊗ Q₁') = θ(Q ⊗ Q'). Indeed, it suffices to verify this in the case where \(Q_{1} = Q \tau N\) and \(Q_{1}' = Q' \tau N'\) , N and \(N'\) being neutral forms; in this case \(Q_{1} \otimes Q_{1}'\) is equivalent to

\[
(Q \otimes Q ^ {\prime}) \tau (Q \otimes N ^ {\prime} \tau Q ^ {\prime} \otimes N \tau N \otimes N ^ {\prime})
\]

and the second parenthesis denotes a neutral form; this proves our assertion.

Now let \(\mathfrak{W}\) be the group of types of quadratic forms over A. Let us define, on the set \(\mathfrak{W}\) a second law of composition, denoted multiplicatively, by the formula

\[
\mathrm{TT} ^ {\prime} = \theta (\mathrm{T} \otimes \mathrm{T} ^ {\prime})\tag{6}
\]

\[
(\mathbf {T}, \mathbf {T} ^ {\prime} \text {   in   } \mathfrak {W}).
\]

It follows immediately from what we have just seen that this composition law is commutative, associative, and distributive with respect to addition. It admits an identity element, namely the type \(T_{1}\) of the quadratic form \(Q_{1}\) defined on the vector space A and such that \(Q_{1}(1)=1\) : indeed, by (5), \(Q_{1}\otimes Q=Q\) for every quadratic form Q. The additive group W, equipped with the multiplication just defined, is therefore a commutative ring with identity element; it is called the ring of types of quadratic forms of A (or the Witt ring of A, when no confusion is to be feared).

Remarks. --- 1) It is clear that, if a is an element of A*, then

\[
a. (\mathrm{TT} ^ {\prime}) = (a. \mathrm{T}) \mathrm{T} ^ {\prime} = \mathrm{T} (a. \mathrm{T} ^ {\prime})\tag{7}
\]

for any elements T, T' of W. One will note moreover that one has a.T = TaT, denoting by Ta the type of the quadratic form aQ1 on A.

\begin{enumerate}
\def\labelenumi{\arabic{enumi})}
\setcounter{enumi}{1}
\tightlist
\item
  Since A has characteristic \(\neq 2\) , every element T of \(\mathfrak{W}\) takes the form \(\sum_{i=1}^{n} a_i \cdot T_1\) where \(a_i \in A^*\) (no. 2, prop. 4). We have
\end{enumerate}

\[
(\sum_ {i = 1} ^ {n} a _ {i}. \mathrm{T} _ {1}) (\sum_ {j = 1} ^ {q} b _ {j}. \mathrm{T} _ {1}) = \sum_ {i, j} a _ {i} b _ {j}. \mathrm{T} _ {1} \quad (a _ {i}, b _ {j} \text {   in   } \mathrm{A} ^ {*}).\tag{8}
\]

\begin{enumerate}
\def\labelenumi{\arabic{enumi})}
\setcounter{enumi}{2}
\tightlist
\item
  Suppose that A is a maximal ordered field. Then the ring W is isomorphic to Z (prop. 5), the integer corresponding to the type of a form of signature \((s, t)\) being s - t (ibid.). Since the tensor product of two forms Q, \(Q'\) of signatures \((s, t)\) , \((s', t')\) is a form of dimension \((s + t)(s' + t')\) , it follows, by an elementary calculation, that the signature of \(Q \otimes Q'\) is \((ss' + tt', st' + ts')\) .
\end{enumerate}

\section{Clifford algebras}

In this paragraph, we shall assume the ring A to be commutative. We shall denote by Q a quadratic form on the A-module E, and by \(\Phi\) the associated bilinear form (\(\S\) 3, no. 4).

\subsection{Definition and universal property of the Clifford algebra.}

DEFINITION 1. --- One calls the Clifford algebra of Q and denotes by C(Q) the quotient algebra of the tensor algebra T(E) of the module E by the two-sided ideal (denoted I(Q)) generated by the elements of the form \(x \otimes x - Q(x)\) . 1 ( \(x \in E\) ).

We will denote by \(\rho_{\mathbb{Q}}\) (or simply \(\rho\) when no confusion is to be feared) the mapping from E into C(Q) composed of the canonical mapping from E into T(E) and the canonical mapping \(\sigma\) of T(E) onto C(Q); the map \(\rho_{\mathbb{Q}}\) is said to be canonical. Note that C(Q) is generated by \(\rho_{\mathbb{Q}}(E)\) , and that, for \(x \in E\) , one has

\[
\rho (x) ^ {2} = \mathrm{Q} (x). 1;\tag{1}
\]

whence, replacing \(x\) by \(x + y\) ( \(x, y\) in E)

\[
\rho (x) \rho (y) + \rho (y) \rho (x) = \Phi (x, y). 1\tag{2}
\]

Example. --- If E admits a basis consisting of a single element \(e\) , T(E) is isomorphic to the polynomial algebra A{[}X{]}, and C(Q) is a quadratic extension of A, with basis \((1, u)\) , where \(u\) is the element \(u = \rho(e)\) and satisfies \(u^2 = Q(e)\) .

Denote by \(T^{h}\) the h-th tensor power \(\otimes E\) in \(\mathrm{T}(\mathrm{E})\) , and let \(T^{+}\) (resp. \(T^{-}\) ) be the sum of \(T^{h}\) for h even (resp. odd).

Since T(E) is the direct sum of \(T^{+}\) and \(T^{-}\) and since I(Q) is generated by elements of \(T^{+}\) , I(Q) is the direct sum of \(T^{+} \cap I(Q)\) and \(T^{-} \cap I(Q)\) , and C(Q) is the direct sum of the two submodules \(C^{+}(Q) = \sigma(T^{+})\) and \(C^{-}(Q) = \sigma(T^{-})\) (also denoted \(C^{+}\) and \(C^{-}\) ). The elements of \(C^{+}\) will be called even (resp. odd). We have the relations

\[
\mathrm{C} ^ {+} \mathrm{C} ^ {+} \subset \mathrm{C} ^ {+}, \quad \mathrm{C} ^ {+} \mathrm{C} ^ {-} \subset \mathrm{C} ^ {-}, \quad \mathrm{C} ^ {-} \mathrm{C} ^ {+} \subset \mathrm{C} ^ {-}, \quad \mathrm{C} ^ {-} \mathrm{C} ^ {-} \subset \mathrm{C} ^ {+}.\tag{3}
\]

In particular \(C^{+}\) is a subalgebra of \(C(Q)\) .

PROPOSITION 1. --- Let f be a linear map of E into an algebra D over A such that \(f(x)^{2} = \mathrm{Q}(x)\) . 1 for every \(x \in \mathbf{E}\) . There exists one and only one homomorphism \(\bar{f}\) of C(Q) into D such that \(f = \bar{f} \circ \rho_{\mathbb{Q}}\) .

The uniqueness of \(\bar{f}\) follows from the fact that C(Q) is generated by \(\rho_{\mathbb{Q}}(\mathrm{E})\) . Let \(h\) be the unique homomorphism of T(E) into D extending \(f\) ( \(h\) is defined by \(h(x_1 \otimes \cdots \otimes x_n) = f(x_1) \ldots f(x_n)\) ). We have

\[
h (x \otimes x - \mathrm{Q} (x). 1) = (f (x) ^ {2} - \mathrm{Q} (x)). 1 = 0,
\]

and consequently h vanishes on I(Q) and defines, by passage to the quotient, the homomorphism \(\bar{f}\) sought.

Prop. 1 expresses that C(Q) is the solution of a universal mapping problem (Ens., chap.~IV, § 3, no. 1).

Take in particular for D the opposite algebra of C(Q) and for \(f\) the map \(\rho\) ; prop. 1 implies that there exists one and only one anti-automorphism \(\beta\) of C(Q) whose restriction to \(\rho(E)\) is the identity; it is called the principal anti-automorphism of C(Q). It is clear that \(\beta^{2} = 1\) .

On the other hand, let Q' be a quadratic form on an A-module E' and f a linear map from E to E' such that Q'∘f = Q. We have \(\rho_{\mathbb{Q}'}(f(x))^{2} = \mathrm{Q}'(f(x)).1 = \mathrm{Q}(x).1\) , and consequently there exists one and only one homomorphism C(f) from C(Q) to C(Q') such that C(f)∘ρq = ρq'∘f.~If f is the identity, C(f) is the identity; if Q'' is a quadratic form on an A-module E'' and g a linear map from E' to E'' such that Q''∘g = Q', we have C(g∘f) = C(g)∘C(f). When E' is a submodule of E and f the canonical injection of E' into E (so that Q' is the restriction of Q to E'), we say that C(f) is the canonical homomorphism of C(Q') into C(Q).

Take in particular Q' = Q and for f the map \(x \to -x\) ; we see that there exists one and only one automorphism \(\alpha\) of C(Q) such that \(\alpha \circ \rho = -\rho\) ; it is called the principal automorphism of C(Q). It is clear that \(\alpha^{2} = 1\) , and that the restriction of \(\alpha\) to C⁺ (resp. C⁻) is the identity (resp. the map \(u \to -u\) ).

PROPOSITION 2. --- Let A' be a commutative ring, φ a homomorphism from A to A', Q' the quadratic form on E' = A' ⊗\_A E deduced from Q by extension of scalars (§ 3, n° 4, prop. 3). There exists one and only one isomorphism j from the algebra A' ⊗\_A C(Q) onto C(Q') such that j (1 ⊗ ρ\_Q(x)) = ρ\_Q′(1 ⊗ x) for all x ∈ E.

It suffices to show that the algebra \(C' = A' \otimes C(Q)\) and the map \(1 \otimes_{\rho_{Q}}\) of \(E'\) into \(C'\) form a solution of the same universal problem as \(C(Q')\) and \(\rho_{Q'}\) . Indeed, let \(D'\) be an algebra over \(A'\) and \(f'\) an A'-linear map from \(E'\) in \(D'\) such that \(f'(x')^{2} = Q'(x')\) . 1 for every \(x' \in E'\) . The map \(g : x \to f'(1 \otimes x)\) from E into \(D'\) (considered as an A-module via the homomorphism \(\varphi\) ) is A-linear and we have \(g(x)^{2} = Q'(1 \otimes x)\) . 1 = Q(x). 1 for all \(x \in E\) . There therefore exists one and only one A-homomorphism \(\bar{g}\) from C(Q) to \(D'\) such that \(\bar{g}(\rho_{Q}(x)) = f'(1 \otimes x)\) . Consequently there exists one and only one A'-homomorphism \(\overline{f'}\) from C' to \(D'\) such that \(\overline{f}'(1 \otimes \rho_{Q}(x)) = f'(1 \otimes x)\) for all \(x \in E\) ; by linearity it follows that \(\overline{f}'((1 \otimes \rho_{Q})(x')) = f'(x')\) for all \(x' \in E'\) . QED.

\subsection{Some operations in the tensor algebra.}

In this section we shall denote by \(e_x\) ( \(x \in E\) ) the linear mapping \(u \to x \otimes u\) of the tensor algebra T(E) into itself.

Lemma 1. — Let f be an element of the dual E* of E. There exists one and only one linear mapping \(i_{f}\) of T(E) into itself such that

\begin{enumerate}
\def\labelenumi{(\arabic{enumi})}
\setcounter{enumi}{3}
\tightlist
\item
\end{enumerate}

\[
i _ {f} (1) = 0\tag{5}
\]

\[
i _ {f} \circ e _ {x} + e _ {x} \circ i _ {f} = f (x). I
\]

for all \(x \in \mathbf{E}\)

(where I denotes the identity mapping). The mapping \(f \to i_f\) of E* into \(\mathfrak{U}(\mathrm{T}(\mathrm{E}))\) is linear. One has \(i_f(\mathrm{T}^n) \subset \mathrm{T}^{n-1}\) , \((i_f)^2 = 0\) , and \(i_f \circ i_g + i_g \circ i_f = 0\) for \(f\) , \(g\) in E*. The mapping \(i_f\) is zero on the subalgebra of

T(E) generated by the kernel of f. The ideal I(Q) is stable under \(i_{f}\) ; by passage to the quotient \(i_{f}\) therefore defines a linear mapping (still denoted \(i_{f}\) ) of C(Q) into itself.

Indeed, formula (5) can be written

\[
i _ {f} (x \otimes u) = - x \otimes i _ {f} (u) + f (x). u \quad (x \in \mathrm{E}, u \in \mathrm{T(E)}).\tag{6}
\]

Since (4) completely determines \(i_{f}\) on \(\mathbf{T}^{0}\) , and since (6) determines \(i_{f}\) on \(\mathbf{T}^{n}\) once its values on \(\mathbf{T}^{n-1}\) are known, the uniqueness of \(i_{f}\) is proved. On the other hand, for \(x\in\mathbf{E}\) and \(u\in\mathbf{T}^{n-1}\) , the right-hand side of (6) is bilinear on \(\mathbf{E}\times\mathbf{T}^{n-1}\) ; this proves the existence of \(i_{f}\) by induction on \(n\) (chap.~III, § 1, no. 2) and also proves, by induction on \(n\) , that \(i_{f}(\mathbf{T}^{n})\subset\mathbf{T}^{n-1}\) . If \(f=ag+bh\) ( \(a,b\) in A, \(g,h\) in E*) it is clear that \(ai_{g}+bi_{h}\) satisfies (4) and (5), hence is equal to \(i_{f}\) . We have

\[
(i _ {f}) ^ {2} \circ e _ {x} = - i _ {f} \circ e _ {x} \circ i _ {f} + f (x) i _ {f} = e _ {x} \circ (i _ {f}) ^ {2}
\]

and, since \((i_f)^2(1) = 0\) we deduce by induction on \(n\) that \((i_f)^2\) is zero on \(\mathbf{T}^n\) . Replacing \(f\) by \(f + g\) , the relation \((i_f)^2 = 0\) gives \(i_f \circ i_g + i_g \circ i_f = 0\) . An induction argument analogous to the preceding ones shows that \(i_f\) is zero on the subalgebra generated by the kernel of \(f\) . Finally, (6) implies that the set of elements \(u\) of I(Q) such that \(i_f(u) \in \mathrm{I}(Q)\) is a left ideal of T(E); moreover, if \(u = (x \otimes x - Q(x).1) \otimes v (x \in E, v \in T(E))\) , one has

\[
\begin{array}{r l} i _ {f} (u) & = f (x) x \otimes v - x \otimes i _ {f} (x \otimes v) - \mathrm{Q} (x) i _ {f} (v) \\ & = f (x) x \otimes v - f (x) x \otimes v + x \otimes x \otimes i _ {f} (v) - \mathrm{Q} (x) i _ {f} (v) \\ & = (x \otimes x - \mathrm{Q} (x). 1) \otimes i _ {f} (v); \end{array}
\]

therefore I(Q) is stable under \(i_{f}\) and the last assertion follows. QED.

Let F be a bilinear form on E; in the rest of this paragraph we shall denote by \(i_x^{\mathrm{F}}\) ( \(x \in \mathrm{E}\) ) the map \(i_f\) corresponding to the linear form \(f: y \to \mathrm{F}(x, y)\) on E.

Lemma 2. --- There exists a linear map \(\lambda_{\mathrm{F}}\) of T(E) into itself such that

\begin{enumerate}
\def\labelenumi{(\arabic{enumi})}
\setcounter{enumi}{6}
\tightlist
\item
\end{enumerate}

\[
\lambda_ {\mathrm{F}} (1) = 1\tag{8}
\]

\[
\lambda_ {\mathrm{F}} \circ e _ {x} = (e _ {x} + i _ {x} ^ {\mathrm{F}}) \circ \lambda_ {\mathrm{F}}\tag{\( (x \in E) \) .}
\]

Whatever the \(f \in \mathbf{E}^*\) , one has

\[
\lambda_ {\mathrm{F}} \circ i _ {f} = i _ {f} \circ \lambda_ {\mathrm{F}}.\tag{9}
\]

Indeed, formula (8) is equivalent to

\[
\lambda_ {\mathbb {F}} (x \otimes u) = x \otimes \lambda_ {\mathbb {F}} (u) + i _ {x} ^ {\mathrm{F}} (\lambda_ {\mathbb {F}} (u)) \quad (x \in \mathrm{E}, u \in \mathrm{T} (\mathrm{E})).\tag{10}
\]

Since (7) determines entirely \(\lambda_{\mathbb{F}}\) on \(\mathrm{T}^{0}\) , and since (10) determines \(\lambda_{\mathbb{F}}\) on \(\mathrm{T}^{n}\) once its values on \(\mathrm{T}^{n-1}\) are known, the uniqueness of \(\lambda_{\mathbb{F}}\) is proved. On the other hand, for \(x \in \mathrm{E}\) and \(u \in \mathrm{T}^{n-1}\) the right-hand side of (10) is bilinear on \(\mathrm{E} \times \mathrm{T}^{n-1}\) ; this proves the existence of \(\lambda_{\mathbb{F}}\) by induction on \(n\) . It remains to prove (9), which we shall do by induction; both sides of (9) vanish on \(\mathrm{T}^{0}\) ; suppose (9) holds on \(\sum_{h=0}^{n-1} \mathrm{T}^{h}\) ; we then have, for \(x \in \mathrm{E}\) and \(u \in \mathrm{T}^{n-1}\) :

\[
\begin{array}{r l} (\lambda_ {\mathbb {F}} \circ i _ {f}) (x \otimes u) & = (- \lambda_ {\mathbb {F}} \circ e _ {x} \circ i _ {f} + f (x) \lambda_ {\mathbb {F}}) (u) \\ & = - (e _ {x} + i _ {x} ^ {\mathrm{F}}) \circ \lambda_ {\mathbb {F}} \circ i _ {f} (u) + f (x) \lambda_ {\mathbb {F}} (u) \\ & = - (e _ {x} + i _ {x} ^ {\mathrm{F}}) \circ i _ {f} \circ \lambda_ {\mathbb {F}} (u) + f (x) \lambda_ {\mathbb {F}} (u) \\ & = (i _ {f} \circ e _ {x} \circ \lambda_ {\mathbb {F}} - f (x) \lambda_ {\mathbb {F}} + i _ {f} \circ i _ {x} ^ {\mathrm{F}} \circ \lambda_ {\mathbb {F}} + f (x) \lambda_ {\mathbb {F}}) (u) \\ & = (i _ {f} \circ (e _ {x} + i _ {x} ^ {\mathrm{F}}) \circ \lambda_ {\mathbb {F}}) (u) = (i _ {f} \circ \lambda_ {\mathbb {F}}) (x \otimes u), \end{array}
\]

whence our last assertion.

Lemma 3. --- Let F and G be two bilinear forms on E. Then we have \(\lambda_{\mathrm{F}} \circ \lambda_{\mathrm{G}} = \lambda_{\mathrm{F+G}}\) . For every bilinear form F on E, \(\lambda_{\mathrm{F}}\) is a bijection from T(E) onto itself.

Indeed, \(\lambda_{\mathbb{F}} \circ \lambda_{\mathbb{G}}\) possesses the characteristic properties (7) and (8) of \(\lambda_{\mathbb{F} + \mathbb{G}}: \text{one has } (\lambda_{\mathbb{F}} \circ \lambda_{\mathbb{G}})(1) = 1\) , and

\[
\begin{array}{r l} \lambda_ {\mathrm{F}} \circ \lambda_ {\mathrm{G}} \circ e _ {x} & = \lambda_ {\mathrm{F}} \circ (e _ {x} + i _ {x} ^ {\mathrm{G}}) \circ \lambda_ {\mathrm{G}} = (e _ {x} + i _ {x} ^ {\mathrm{G}} + i _ {x} ^ {\mathrm{F}}) \circ \lambda_ {\mathrm{F}} \circ \lambda_ {\mathrm{G}} \\ & = (e _ {x} + i _ {x} ^ {\mathrm{F+G}}) \circ \lambda_ {\mathrm{F}} \circ \lambda_ {\mathrm{G}}. \end{array}
\]

On the other hand, if F = 0, we have \(i_x^{\mathbb{F}} = 0\) for all \(x \in \mathbb{E}\) and consequently \(\lambda_{\mathbb{F}} = \mathbb{I}\) which implies that, for every F, we have \(\lambda_{\mathbb{P}} \circ \lambda_{-\mathbb{P}} = \lambda_{-\mathbb{F}} \circ \lambda_{\mathbb{F}} = \mathbb{I}\) .

\subsection{Basis of the Clifford algebra.}

PROPOSITION 3. --- Let Q and Q' be two quadratic forms and F a bilinear form on E such that Q'(x) = Q(x) + F(x, x) for all x ∈ E. The map λ\_F maps the ideal I(Q') onto the ideal I(Q), and defines an isomorphism (denoted \(\overline{\lambda_{\mathbf{F}}}\) ) from the A-module C(Q') onto the A-module C(Q).

As \(\lambda_{\mathbb{F}}\) is a bijection whose \(\lambda_{-\mathbb{F}}\) is the inverse bijection (Lemma 3), it suffices to prove the inclusion \(\lambda_{\mathbb{F}}(\mathrm{I}(\mathrm{Q}^{\prime})) \subset \mathrm{I}(\mathrm{Q})\) . Since \(\mathrm{I}(\mathrm{Q})\) is a left ideal stable under \(i_x^{\mathbb{F}}\) (Lemma 1), (8) shows that the set of \(u \in \mathrm{T}(\mathrm{E})\) such that \(\lambda_{\mathbb{F}}(u) \in \mathrm{I}(\mathrm{Q})\) is a left ideal. It therefore suffices to show that, for any \(u \in \mathrm{T}(\mathrm{E})\) and \(x \in \mathrm{E}\) , one has \(\lambda_{\mathbb{F}}(x \otimes x \otimes u - \mathrm{Q}'(x)u) \in \mathrm{I}(\mathrm{Q})\) . Now, by (8) and Lemma 1, we have

\[
\lambda_ {\mathrm{F}} \circ e _ {x} ^ {2} = (e _ {x} + i _ {x} ^ {\mathrm{F}}) ^ {2} \circ \lambda_ {\mathrm{F}} = (e _ {x} ^ {2} + \mathrm{F} (x, x)) \circ \lambda_ {\mathrm{F}},
\]

whence

\[
\begin{array}{r l} \lambda_ {\mathbf {F}} (x \otimes x \otimes u - \mathrm{Q} ^ {\prime} (x) u) & = (e _ {x} ^ {2} + \mathrm{F} (x, x) - \mathrm{Q} ^ {\prime} (x)) \circ \lambda_ {\mathbf {F}} (u) \\ & = (x \otimes x - \mathrm{Q} (x)) \otimes \lambda_ {\mathbf {F}} (u) \in \mathrm{I} (\mathrm{Q}). \end{array}
\]

Lemma 4. — If the quadratic form Q is zero, then C(Q) is none other than the exterior algebra of E.

Indeed, the exterior algebra of E is none other than the quotient of T(E) by the two-sided ideal J generated by the elements of the form \(a \otimes v \otimes a\) ( \(a \in E\) , \(v \in T(E)\) ) (chap.~III, § 5, no.~5 and no.~9). It is clear that I(Q) ⊂ J. It therefore suffices to show that \(a \otimes v \otimes a \in I(Q)\) . This is evident if \(v \in T^0\) . Suppose this assertion has been proved for \(n-1 \leqslant \nu \in \sum_{h=0}^{\infty} \mathrm{T}^{h}\) , and let \(x \in \mathrm{E}\) and \(u \in \mathrm{T}^{n-1}\) ; we have

\[
\begin{array}{c} \boldsymbol {a} \otimes \boldsymbol {x} \otimes \boldsymbol {u} \otimes \boldsymbol {a} = (\boldsymbol {a} + \boldsymbol {x}) \otimes (\boldsymbol {a} + \boldsymbol {x}) \otimes \boldsymbol {u} \otimes \boldsymbol {a} \\ - \boldsymbol {a} \otimes \boldsymbol {a} \otimes \boldsymbol {u} \otimes \boldsymbol {a} - \boldsymbol {x} \otimes \boldsymbol {a} \otimes \boldsymbol {u} \otimes \boldsymbol {a} - \boldsymbol {x} \otimes \boldsymbol {x} \otimes \boldsymbol {u} \otimes \boldsymbol {a} \end{array}
\]

and the four terms on the right-hand side belong to I(Q).

Suppose in particular that the module E admits a basis \((x_{i})_{i\in \mathbf{L}}\) , and let us totally order the index set L. We know (chap.~III, §5, no.~6) that the exterior algebra of E admits as a basis the family consisting of the elements \(x_{\mathrm{H}}\) where H ranges over the set of finite subsets of L and where \(x_{\mathrm{H}}\) is the element \(x_{h_1}\wedge \ldots \wedge x_{h_q}\) , \((h_1,\ldots ,h_q)\) denoting the strictly increasing sequence of elements of H.

On the other hand, consider the bilinear form F defined by \(\mathrm{F}(x_{i}, x_{j}) = -\Phi(x_{i}, x_{j})\) if i \textgreater{} j, \(\mathrm{F}(x_{i}, x_{j}) = 0\) if i \textless{} j and \(\mathrm{F}(x_{i}, x_{i}) = -\mathrm{Q}(x_{i})\) . It is clear that \(\mathrm{Q}(x) + \mathrm{F}(x, x) = 0\) ; with the notations of Prop. 3, it follows from this proposition and Lemma 4 that \(\bar{\lambda}_{\mathrm{F}}\) is an isomorphism of \(\wedge \mathrm{E}\) onto C(Q), which is therefore a free A-module. Let us prove, by induction on the number of elements of H, that we have

\[
\overline {{{{\lambda}}}} _ {\mathrm{F}} (x _ {\mathrm{H}}) = \rho (x _ {h _ {1}}) \dots \rho (x _ {h _ {q}})\tag{11}
\]

\((\mathrm{where}\ (h_{1},\ldots,h_{q})\) is the strictly increasing sequence of elements of H). This is evident if H has 0 or 1 element. Suppose (11) holds for sets having at most \(q-1\) elements. Consider a subset H of \(q\) elements, denote \(j\) its smallest element, and let us write \(\mathrm{H}=\{j\}\cup\mathrm{K}\) where K is a subset with \(q-1\) elements. By (8) and the induction hypothesis, we have

\[
\overline {{{\lambda}}} _ {\mathrm{F}} (x _ {\mathrm{H}}) = \overline {{{\lambda}}} _ {\mathrm{F}} (x _ {j} \wedge x _ {\mathrm{K}}) = \rho (x _ {j}) \overline {{{\lambda}}} _ {\mathrm{F}} (x _ {\mathrm{K}}) + i _ {x _ {j}} ^ {\mathrm{F}} (\overline {{{\lambda}}} _ {\mathrm{F}} (x _ {\mathrm{K}})) = x _ {\mathrm{H}} ^ {\prime} + i _ {x _ {j}} ^ {\mathrm{F}} (x _ {\mathrm{K}} ^ {\prime}),
\]

setting, for every finite subset J of L, \(x_{J}^{\prime}=\rho(x_{j_{1}})\ldots\rho(x_{j_{s}})\) , where \((j_{1},\ldots,j_{s})\) is the increasing sequence of elements of J. Now, for \(i\in\mathbf{K}\) , \(x_{i}\) belongs to the kernel of the linear form \(y\to\mathrm{F}(x_{j},y)\) , hence \(i_{x_{j}}^{\mathrm{F}}(x_{\mathrm{K}}^{\prime})=0\) (Lemma 1). This proves the desired result.

We have therefore proved the following theorem:

THEOREM 1. --- Suppose that the A-module E admits a basis \((x_{i})_{i\in L}\) , the index set L being equipped with a total order structure. For every finite subset H of L, set \(x_{\mathrm{H}} = \rho(x_{h_1})\rho(x_{h_2})\ldots \rho(x_{h_q})\) , where \((h_1,\ldots,h_q)\) is the strictly increasing sequence of the elements of H. Then the elements \(x_{\mathrm{H}}\) form a basis of the A-module C(Q).

COROLLARY 1. --- If E is a free module of dimension n, then C(Q) is a free module of dimension \(2^{n}\) ; moreover, if \(n > 0\) , C \(^{+}\) and C \(^{-}\) are free modules of dimension \(2^{n-1}\) .

This follows immediately from the properties of binomial coefficients.

COROLLARY 2. --- If E is a free module, the canonical mapping \(\rho\) from E into C(Q) and the mapping \(a \to a\) .1 of A into C(Q) are injective.

COROLLARY 3. --- Suppose that E is the direct sum of two free submodules \(E_{1}\) and \(E_{2}\) . Let \(Q_{i}\) be the restriction of Q to \(E_{i}\) and \(p_{i}\) the canonical map from C(Q \(_{i}\) ) into C(Q) (i = 1, 2). Then the

linear map \(p\) of \(\mathrm{C(Q_{1})\otimes C(Q_{2})}\) into \(\mathrm{C(Q)}\) deduced from the bilinear map \((a,b)\to p_{1}(a)p_{2}(b)\) of \(\mathrm{C(Q_{1})\times C(Q_{2})}\) in \(\mathrm{C(Q)}\) is a bijection.

It suffices indeed to consider the basis of E obtained by taking the union of a basis of \(E_{1}\) and a basis of \(E_{2}\) .

COROLLARY 4. --- With the hypotheses and notation of Cor. 3, suppose further that \(E_{1}\) and \(E_{2}\) are orthogonal, and transport to \(C(Q_{1})\otimes C(Q_{2})\) , by means of the bijection p, the algebra structure of \(C(Q)\) . If \(a_{i}\) and \(b_{i}\) are even or odd elements of \(C(Q_{i})\) (i=1, 2), we have \((a_{1}\otimes a_{2})(b_{1}\otimes b_{2})=\varepsilon(a_{1}b_{1})\otimes(a_{2}b_{2})\) , with \(\varepsilon=1\) unless \(a_{2}\) and \(b_{1}\) are odd, in which case \(\varepsilon=-1\) .

It suffices indeed to prove that \(p_2(a_2)p_1(b_1) = \varepsilon p_1(b_1)p_2(a_2)\) , and for that one may assume that \(p_2(a_2)\) (resp. \(p_1(b_1)\) ) is a product \(x_1 \ldots x_h\) (resp. \(y_1 \ldots y_k\) ) of elements of \(\rho_{\mathbb{Q}}(\mathbf{E}_2)\) (resp. \(\rho_{\mathbb{Q}}(\mathbf{E}_1)\) ). Since \(\mathbf{E}_1\) and \(\mathbf{E}_2\) are orthogonal, one has

\[
x _ {i} y _ {j} + y _ {j} x _ {i} = \Phi (x _ {i}, y _ {j}) = 0,
\]

whence

\[
x _ {1} \dots x _ {h} y _ {1} \dots y _ {k} = (- 1) ^ {h k} y _ {1} \dots y _ {k} x _ {1} \dots x _ {h}.
\]

The conclusions of Cor. 3 and 4 remain true if one omits the hypothesis that \(\mathbf{E}_1\) and \(\mathbf{E}_2\) are free modules (cf.~exerc. 1).

\subsection{Structure of the Clifford algebra.}

In this section, we will assume that A is a field, that E is a vector space of finite dimension m over A, and that the quadratic form Q is non-degenerate (which, according to th. 1 of § 5, no. 1, requires that m be even if A has characteristic 2). Since E is free, the canonical map ρ is injective (no. 3, cor. 2 of th. 1). We shall henceforth identify E with its image in C(Q).

THEOREM 2. --- Suppose that the dimension of E is an even number \(m = 2r\) and that Q be neutral ( \(\S 4, \text{n}^{\circ} 2\) ). Then the algebra C(Q) is separable (Chap.~VIII, \(\S 7, \text{n}^{\circ} 5\) , Def. 1) and is isomorphic to the algebra of endomorphisms of a vector space of dimension \(2^{r}\) on A. Moreover, if \(m > 0\) , C \(^{+}\) (Q) is separable and is the direct sum of two ideals isomorphic to the algebra of endomorphisms of a vector space of dimension \(2^{r-1}\) over A.

Indeed, since Q is neutral, one can decompose E as a direct sum of two totally singular subspaces N and P of dimension r (§ 4, no. 2, cor. 1 of prop. 2). Since the restriction of Q to N is zero, the subalgebra S of C(Q) generated by N is identified with the exterior algebra of N (no. 3, lemma 4). For \(n \in \mathbb{N}\) , we shall denote by \(e_n'\) the map \(t \to nt\) from S into itself.

Let \((n_{1},\ldots,n_{r})\) be a basis of N; we shall denote by \((p_{1},\ldots,p_{r})\) the basis of P such that \(\Phi(n_{i},p_{j})=\delta_{ij}\) ( \(\S 4\) , no. 2, prop. 2). For \(p\in P\) , we shall denote by \(p'\) the linear form \(n\to\Phi(n,p)\) on N, and \(i_{p}\) the endomorphism of S obtained by passing to the quotient from the endomorphism \(i_{p'}\) of T(N) associated with \(p'\) as was stated in Lemma 1 of No. 2. By (5), we have

\[
e _ {n} ^ {\prime} \circ i _ {p} + i _ {p} \circ e _ {n} ^ {\prime} = \Phi (n, p) \quad (n \in \mathrm{N}, p \in \mathrm{P}).\tag{12}
\]

Let us set, for \(x = n + p \in \mathbf{E}\) (with \(n \in \mathbf{N}\) and \(p \in \mathbf{P}\) ), \(s(x) = e_n' + i_p\) . It is clear that \(s\) is a linear map from \(\mathbf{E}\) into \(\mathfrak{L}(\mathbf{S})\) . Since we have

\[
s (x) ^ {2} = \left(e _ {n} ^ {\prime} + i _ {p}\right) ^ {2} = Q (n) + \Phi (n, p) = Q (x)
\]

By virtue of (12) and Lemma 1 (no. 2), s extends to a homomorphism (which we shall still denote by s) from C(Q) into ℒ(S) (no. 1, prop. 1). We shall show that this homomorphism is surjective, which, since C(Q) and ℒ(S) are both of dimension 2²ʳ, will imply that s is an isomorphism and will prove our first assertion.

Indeed, let I denote the interval \([1, r]\) . For every subset H of I, we set \(H' = I - H\) and we denote by \(n_{H}\) (resp. \(p_{H}\) ) the product of \(n_{i}\) (resp. \(p_{i}\) ) for \(i \in H\) , arranged in increasing order of indices. Recall that the \(n_{H}\) form a basis of S (no. 3, th. 1). Finally, for any two subsets H, K of I, set \(x_{H,K} = n_{H}p_{I}n_{K'}\) . We shall show that the elements \(s(x_{H,K})\) of \(s(C(Q))\) generate L(S). Now, if \(j \notin H\) , one has \(s(p_{j})(n_{H}) = i_{p_{j}}(n_{H}) = 0\) by lemma 1, since the \(n_{i}\) for \(i \in H\) belong to the kernel of the linear form \(n \to \Phi(n, p_{j})\) on N; on the other hand we have

\[
s \left(p _ {j}\right) \left(n _ {j} n _ {\mathrm{H}}\right) = \left(i _ {p _ {j}} \circ e _ {n _ {j}} ^ {\prime}\right) \left(n _ {\mathrm{H}}\right) = \Phi \left(p _ {j}, n _ {j}\right) n _ {\mathrm{H}} - n _ {j}. s \left(p _ {j}\right) \left(n _ {\mathrm{H}}\right) = n _ {\mathrm{H}}
\]

(by (12)). Since \(s\) is a homomorphism, we deduce, for any two subsets H, K of I, that \(s(p_{\mathrm{K}})(n_{\mathrm{H}}) = 0\) if \(\mathrm{K} \not\subset \mathrm{H}\) , and that \(s(p_{\mathbf{K}}) (n_{\mathbf{H}}) = \pm n_{\mathbf{H} - \mathbf{K}}\) if \(\mathrm{K} \subset \mathrm{H}\) . Since, for \(\mathrm{M} \subset \mathrm{I}\) and \(\mathrm{L} \subset \mathrm{I}\) , we have by definition \(s(n_{\mathbf{M}})(n_{\mathbf{L}}) = n_{\mathbf{M}}n_{\mathbf{L}}\) , and that \(n_{\mathbf{M}}n_{\mathbf{L}}\) is zero if \(\mathrm{M} \cap \mathrm{L} \neq \emptyset\) and is equal to \(\pm n_{\mathbf{M} \cup \mathbf{L}}\) in the contrary case, we conclude from the preceding that, for arbitrary subsets \(\mathrm{H}\) , \(\mathrm{K}\) , \(\mathrm{L}\) of \(\mathrm{I}\) , \(s(x_{\mathbf{H},\mathbf{K}})(n_{\mathbf{L}}) = s(n_{\mathbf{H}})s(p_{\mathbf{I}})s(n_{\mathbf{K}^{\prime}})(n_{\mathbf{L}})\) is zero if \(\mathrm{K} \neq \mathrm{L}\) and is equal to \(\pm n_{\mathbf{H}}\) if \(\mathrm{K} = \mathrm{L}\) . This shows that the \(s(x_{\mathbf{H},\mathbf{K}})\) generate \(\mathfrak{L}(\mathrm{S})\) and completes the proof of the first assertion.

To prove the second assertion, set \(S^{+} = S \cap C^{+}\) and \(S^{-} = S \cap C^{-}\) ; it is clear that \(S^{+}\) (resp. \(S^{-}\) ) is the subspace of S generated by the \(n_{\mathrm{H}}\) such that H has an even (resp. odd) number of elements, that S is the direct sum of \(S^{+}\) and \(S^{-}\) , and that \(s(C^{+})\) leaves \(S^{+}\) and \(S^{-}\) stable. Consequently \(s\) maps \(C^{+}\) into a subalgebra of \(\mathfrak{L}(S)\) , isomorphic to the product \(\mathfrak{L}(S^{+}) \times \mathfrak{L}(S^{-})\) ; the restriction of \(s\) to \(C^{+}\) is an isomorphism of \(C^{+}\) onto this subalgebra, since \(s\) is injective and \(C^{+}\) and \(\mathfrak{L}(S^{+}) \times \mathfrak{L}(S^{-})\) both have dimension \(2^{2r-1}\) (no. 2, cor. 1 of th. 1). QED.

COROLLARY. --- If m is even, but Q of arbitrary index, the algebra C(Q) is a central simple algebra of dimension \(2^{m}\) . Moreover, if m \textgreater{} 0, the subalgebra C⁺(Q) is separable, and its center Z has dimension 2 over A. When Z is a field, Z is a separable quadratic extension of A and C⁺(Q) is simple; otherwise, Z is a direct product of two fields isomorphic to A, and C⁺(Q) is then a direct product of two simple subalgebras of dimensions \(2^{m-2}\) .

Indeed, let A' be the algebraic closure of A, and Q' the quadratic form on \(E' = A' \otimes_{A} E\) obtained from Q by extension of scalars. We have seen that \(C(Q')\) is isomorphic to \(A' \otimes_{A} C(Q)\) (prop. 2), and it is clear that \(C^{+}(Q')\) is isomorphic to \(A' \otimes_{A} C^{+}(Q)\) . Since Q' is neutral (§ 4, no. 2, cor. 2 of prop. 3), the corollary is an immediate consequence of th. 2 and the permanence theorems of chap.~VIII, § 7.

Remarks. --- 1) Since the algebra C(Q) is simple, it has only one class of irreducible representations; these are called the spinor representations; when one fixes one's attention on one of these representations, say \(\tau\) , the elements of the space in which \(\tau\) takes place are called spinors. If Q is neutral, the restriction of \(\tau\) to \(C^{+}(Q)\) is, like that of \(s\) , sum of two inequivalent absolutely irreducible representations; the elements of the subspaces carrying these two representations are called semi-spinors. In the general case, if \(C^{+}(Q)\) is not simple, the restriction of \(\tau\) to \(C^{+}(Q)\) must, since it is faithful, contain subrepresentations belonging to each of the two classes of irreducible representations of \(C^{+}(Q)\) , and hence is a sum of two inequivalent absolutely irreducible representations, since this is the case after extending scalars to the algebraic closure A' of A. On the other hand, if \(C^{+}(Q)\) is simple, it has only one class of irreducible representations, and the restriction of \(\tau\) to \(C^{+}(Q)\) is therefore irreducible, since it decomposes under scalar extension to A' into two non-equivalent representations.

\begin{enumerate}
\def\labelenumi{\arabic{enumi})}
\setcounter{enumi}{1}
\tightlist
\item
  Suppose A has characteristic \(\neq 2\) , and let \((x_{1},\ldots,x_{m})\) ( \(m=2r\) ) be an orthogonal basis of E. Put
\end{enumerate}

\[
z = 2 ^ {r} x _ {1} \dots x _ {m} \in \mathrm{C(Q)};
\]

since \(x_{i}x_{j} + x_{j}x_{i} = 0\) for \(i \neq j\) , one has \(zx_{j} = -x_{j}z\) , which implies that \(z\) belongs to the center \(Z\) of \(C^{+}(Q)\) without belonging to \(A\) . We have

\[
z ^ {2} = 2 ^ {2 r} (- 1) ^ {r} \mathrm{Q} (x _ {1}) \dots \mathrm{Q} (x _ {m}) = (- 1) ^ {r} \mathrm{D}
\]

denoting by D the discriminant of \(\Phi\) with respect to the basis \((x_{j})\) (cf.~exerc. 9).

THEOREM 3. --- Suppose that the dimension of E is an odd number \(m = 2r + 1\) (so that A has characteristic \(\neq 2\) ).

\begin{enumerate}
\def\labelenumi{\alph{enumi})}
\item
  The algebra C \(^{+}\) (Q) is central simple. If Q has maximum index r, C \(^{+}\) (Q) is isomorphic to the algebra of endomorphisms of a vector space of dimension 2 \(^{r}\) over A.
\item
  The algebra C(Q) is separable. Its center Z has dimension 2, and C(Q) is isomorphic to Z⊗ \(_{A}\) C \(^{+}\) (Q), hence is simple or a direct product of two simple subalgebras.
\end{enumerate}

Indeed, let \(x_0\) be a non-isotropic vector of E, and F the orthogonal complement of \(x_0\) ; let us denote by \(Q_1\) the quadratic form \(y \to -Q(x_0)Q(y)\) on F; it is clear that \(Q_1\) is nondegenerate. Since \(x_0y = -yx_0\) (for \(y \in F\) ), we have \((x_0y)^2 = -Q(x_0)Q(y) = Q_1(y)\) , and consequently the map \(y \to x_0y\) of F into \(C^+(Q)\) extends to a homomorphism \(h\) of \(C(Q_1)\) in \(C^+(Q)\) (no. 1, prop. 1). But \(C(Q_1)\) is simple (th. 2) and has the same dimension \(2^{2r}\) as \(C^{+}(Q)\) ; this implies, since \(h(1) = 1\) , that \(h\) is an isomorphism. Moreover, if \(Q\) has index \(r\) , one can choose \(x_0\) in such a way that \(Q_1\) also has index \(r\) ( \(\S 4\) , no. 2, prop. 3), which proves \(a\) ).

Now let \((x_{1},\ldots ,x_{2r})\) be an orthogonal basis of F; set \(z = x_0x_1\dots x_{2r}\) . One immediately verifies that \(z\) commutes with \(x_{j}\) for \(j = 0,\dots,2r\) , hence belongs to the center of C(Q). Let Z be the subspace of C(Q) generated by 1 and \(z\) ; it is a subalgebra of the center of C(Q) and a quadratic extension of A, since \(z\) is odd and \(z^2\) is equal to the scalar \((-1)^r\mathrm{Q}(x_0)\ldots \mathrm{Q}(x_{2r})\) . Consider the homomorphism \(\theta\) of \(\mathbf{Z}\otimes_{\mathbf{A}}\mathbf{C}^{+}(\mathbf{Q})\) into C(Q) defined by \(\theta (u\otimes v) = uv\) . Since \(z\in \mathbb{C}^{-}\) and is invertible, the map \(u\to zu\) is an isomorphism of the module \(\mathbf{C}^+\) onto \(\mathbb{C}^{-}\) , which implies that \(\theta (\mathrm{Z}\otimes \mathrm{C}^{+})\) contains \(\mathbf{C}^{+}\) and \(\mathbb{C}^{-}\) , hence coincides with C(Q). Since \(\mathbf{Z}\otimes \mathbf{C}^{+}\) and C(Q) have the same dimension \(2^{2r + 1},\theta\) is an isomorphism; this proves \(b)\) taking into account the results of Chap.~VIII, § 7.

Note. --- The discriminant D of \(\Phi\) with respect to the basis \((x_{j})_{(j=0,\ldots,2r)}\) is equal to \(2^{2r+1}Q(x_{0})\ldots Q(x_{2r})\) . Consequently, Z is generated by 1 and by the odd element \(z' = 2^{r+1}z\) such that \(z'^{2} = (-1)^{r}2D\) . The algebra C(Q) is therefore simple if and only if \(2(-1)^{r}D\) is not a square in A.
% semantic-fragments: legacy-input-seam
\relax{}
\subsection{Clifford group.}

In this section, we assume that A is a field, that E has finite dimension m, and that Q is nondegenerate. We identify E with its canonical image in C(Q).

DEFINITION 2. --- We call the Clifford group of Q (resp. the special Clifford group of Q) the multiplicative group of invertible elements s of C(Q) (resp. C \(^{+}\) (Q)) such that sEs \(^{-1}\) = E.

In this section, we denote by G and G⁺ the Clifford group and the special Clifford group of Q. It is clear that G⁺ = G ∩ C⁺(Q).

THEOREM 4. --- Set, for \(s \in G\) and \(x \in E\) , \(\varphi(s). x = sxs^{-1}\) .

\begin{enumerate}
\def\labelenumi{\alph{enumi})}
\item
  The map \(\varphi\) is a homomorphism from G into the orthogonal group \(\mathbf{O}(\mathrm{Q})\) of Q and its kernel is the set of invertible elements of the center Z of C(Q).
\item
  The set \(E \cap G\) is the set of nonsingular vectors of E; for \(x \in E \cap G, -\varphi(x)\) is the symmetry with respect to the hyperplane orthogonal to x.
\item
  If dim(E) is even, we have \(\varphi(G) = \mathbf{O}(Q)\) , \(\varphi(G^{+})\) has index 2 in \(\mathbf{O}(Q)\) if \(E \neq \{0\}\) , and is equal to \(\mathbf{SO}(Q)\) if A has characteristic \(\neq 2\) .
\item
  If dim(E) is odd (which implies that A has characteristic ≠ 2), we have φ(G) = φ(G⁺) = SO(Q).
\end{enumerate}

Indeed, we have \(\mathrm{Q}(sxs^{-1}) = (sxs^{-1})^2 = sx^2s^{-1} = \mathrm{Q}(x)\) for \(s \in G\) and \(x \in E\) , which shows that \(\varphi(s) \in \mathbf{O}(Q)\) . For \(\varphi(s) = 1\) , it is necessary and sufficient that s commute with the elements of E, that is, belong to the center Z of C(Q). This proves a).

For an element \(x\) of E to belong to G, it must be invertible, that is, it must be a non-singular vector (since \(x^2 = Q(x)\) ). If this is so, we have \(x^{-1} = Q(x)^{-1}x\) , whence, for every \(y \in E\) ,

\[
x y x ^ {- 1} = \mathrm{Q} (x) ^ {- 1} x y x = \mathrm{Q} (x) ^ {- 1} x (\Phi (x, y) - x y) = - (y - \Phi (x, y) \mathrm{Q} (x) ^ {- 1} x)
\]

this proves b) (§ 6, no. 4).

Lemma 5. --- Every element s of G is of the form \(zs'\) , where z is an invertible element of Z and \(s'\) belongs to G ∩ C \(^{+}\) (Q) or to G ∩ C \(^{-}\) (Q); the subgroup G \(^{+}\) has index 2 in G when E ≠ \{0\}.

The second assertion follows evidently from the first, since the nonsingular vectors belong to \(\mathrm{G}\cap\mathrm{C}^{-}(\mathrm{Q})\) . Suppose first that dim(E) is even, and let \(s=t'+t''\) , with \(t'\in\mathrm{C}^{+}(\mathrm{Q})\) and \(t''\in\mathrm{C}^{-}(\mathrm{Q})\) ; by definition we have \(sx=(\varphi(s).x)s\) for all \(x\in E\) ; since \(t'x\) and \((\varphi(s).x)t'\) (resp. \(t''x\) and \((\varphi(s).x)t'')\) are odd (resp. even) elements, we have \(t'x=(\varphi(s).x)t'\) , whence \(s^{-1}t'x=xs^{-1}t'\) for all \(x\in E\) . From this we conclude that \(s^{-1}t'\in Z\) and since dim(E) is even, \(Z=A\) (no. 4, cor. of th. 2) , hence \(t'=as\) , where \(a\in A\) . If \(a\neq0\) , we therefore have \(s=a^{-1}t'\) and \(t'\in G\cap C^{+}(Q)\) ; if a=0, \(s=t''\in G\cap C^{-}(Q)\) and the lemma is proved in this case. If dim(E) is odd, A has characteristic \(\neq2\) , so for every \(s\in G\) , \(\varphi(s)\) is a product of symmetries with respect to nonsingular vectors \(x_{i} (i = 1,2,\ldots,h)\) ( \(\S 6, n^{o} 4\) , prop. 5); if we set \(s' = x_{1}x_{2}\ldots x_{h}\) , one has \(\varphi(s) = \varphi(s')\) , hence \(s = zs'\) , where \(z \in \mathbf{Z}\) , and \(s'\) belongs to \(C^{+}(Q)\) or to \(C^{-}(Q)\) according as \(h\) is even or odd.

Suppose dim(E) is even. Since every element \(u\) of \(\mathbf{O}(Q)\) extends in one and only one way to an automorphism \(\overline{u}\) of C(Q) (prop. 1), and since C(Q) is central simple (th. 2), \(\overline{u}\) is an inner automorphism (chap.~VIII, § 10, no. 1, th. 1). There therefore exists an element \(s\) of G such that \(\varphi(s) = u\) On the other hand, the center of C(Q) is contained in C \(^{+}\) , which implies that \(\varphi(G)/\varphi(G^{+})\) is isomorphic to G/G \(^{+}\) , hence that \(\varphi(G^{+})\) has index 2 in \(\varphi(G) = \mathbf{O}(Q)\) if E ≠ \{0\}. This proves the first two assertions of \(c\) ).

Finally, suppose that A has characteristic ≠ 2. Then every element u of O(Q) is a product of symmetries with respect to hyperplanes orthogonal to nonsingular vectors \(x_{i} (i = 1, \ldots, h) (\S 6, n^{o} 4, \text{prop. } 5)\) ; we therefore have \(u = (-1)^{h} \varphi(x_{1} \ldots x_{h})\) and det \((u) = (-1)^{h}\) . For u to belong to SO(Q), it is necessary and sufficient that h be even, which shows that \(\varphi(G^{+}) \supset \mathbf{SO}(Q)\) . Since SO(Q) has index 2 in O(Q) when \(E \neq \{0\}\) , one has \(\varphi(G^{+}) = \mathbf{SO}(Q)\) If E is of even dimension, which completes the proof of c). On the other hand, if the dimension of E is odd, \(\varphi(G)\) does not contain the orthogonal transformation \(x \to -x\) ; indeed, the latter extends to the principal automorphism \(\alpha\) of C(Q) (no. 1), and \(\alpha\) is not an inner automorphism since the center Z of C(Q) contains a nonzero element of C⁻(Q) (th. 3). We therefore have \(\varphi(G) \neq \mathbf{O}(Q)\) , and, since \(\varphi(G) \supset \varphi(G^{+}) \supset \mathbf{SO}(Q)\) and since SO(Q) is of index 2 in O(Q), we have \(\varphi(G) = \varphi(G^{+}) = \mathbf{SO}(Q)\) This proves d). QED.

The subgroup \(\varphi(G^{+})\) of \(\mathbf{O}(Q)\) , which is of index 2 if \(E \neq \{0\}\) , is called the group of rotations of E, and its elements are called rotations; it is denoted by \(\mathbf{O}^{+}(Q)\) Note that, if A is not of characteristic 2, we have \(\mathbf{O}^{+}(Q) = \mathbf{S}\mathbf{O}(Q)\) (cf.~exerc. 9).

PROPOSITION 4. — Let β be the principal antiautomorphism of C(Q) (no. 1). For every \(s \in G^{+}\) , \(\beta(s)s\) is a scalar. The map N : \(s \to \beta(s)s\) is a homomorphism of \(G^{+}\) in the multiplicative group A* of nonzero elements of A.

Indeed, for \(s \in G^{+}\) , one has \(sEs^{-1} = E\) , whence \(\beta(s)^{-1}E\beta(s) = E\) , which shows that \(\beta(s) \in G^{+}\) . Since \(sx = (\varphi(s).x)s\) for all \(x \in E\) , one has \(x\beta(s) = \beta(sx) = \beta(s)(\varphi(s).x)\) and consequently \(\beta(s)sx = \beta(s)(\varphi(s).x)s = x\beta(s)s\) , which implies that \(\beta(s)s\) belongs to the center of C(Q). Since, moreover, \(\beta(s)s\) belongs to C+(Q), \(\beta(s)s\) is a scalar (Ths. 2 and 3). Finally we have \(\beta(st)st = \beta(t)\beta(s)st = \beta(s)s\beta(t)t\) , that is, N(st) = N(s)N(t) for s, t in G+. QED.

The scalar \(N(s) = \beta(s)s (s \in G^{+})\) is called the spinorial norm of s. One denotes by \(G_{0}^{+}\) and one calls the reduced Clifford group the kernel of the homomorphism N. The image \(\varphi(G_{0}^{+})\) is denoted \(\mathbf{O}_{0}^{+}(Q)\) and is called the reduced orthogonal group of Q. Since the kernel of the restriction of \(\varphi\) to \(G^{+}\) is the set of even invertible elements of the center of C(Q) (Th. 4) and hence is identified with A* (Ths. 2 and 3), \(\varphi(G^{+})/\mathbf{O}_{0}^{+}(Q)\) is isomorphic to \(G^{+}/A^{*}G_{0}^{+}\) , hence also with \(N(G^{+})/N(A^{*})\) , and in particular is commutative. It is clear that \(N(A^{*})\) is the subgroup \((A^{*})^{2}\) of squares of elements of A*. If Q has index \textgreater0, there exist, for every \(a \in A^{*}\) , two elements x and y of E such that \(Q(x) = a\) and \(Q(y) = 1 (\S 4, n^{\circ} 2, \text{prop. } 4)\) ; since \(xy \in G^{+}\) and \(N(xy) = Q(x)Q(y) = a\) ; this shows that \(N(G^{+}) = A^{*}\) , hence that \(\varphi(G^{+})/\mathbf{O}_{0}^{+}(Q)\) is isomorphic to \(A^{*}/(A^{*})^{2}\) .

Exercises. --- ¶ 1) Prove Corollaries 3 and 4 of Theorem 1 of No. 3 when E₁ and E₂ are any two complementary submodules of E. (First establish Cor. 4: to do this, show that the tensor product C(Q₁) ⊗ C(Q₂) is endowed with an algebra structure by the sign convention made in the statement of Cor. 4, and that this algebra S solves the same universal mapping problem as C(Q); for this, consider, for every linear map f of E into an algebra D over A such that (f(x))² = Q(x)·1, the homomorphism \(\bar{f}_i\) of C(Qᵢ) into D such that \(\bar{f}_i(1) = 1\) , \(\bar{f}_i(\rho_{Q_i}(x_i)) = f(x_i)\) for \(x_i \in \mathrm{E}_i\) ( \(i = 1, 2\) ) and prove that there exists a homomorphism \(\bar{f}\) of the algebra S into D such that

\[
\bar {f} (z _ {1} \otimes z _ {2}) = \bar {f} _ {1} (z _ {1}) \bar {f} _ {2} (z _ {2})
\]

for \(z_{i}\in\mathrm{C}(\mathrm{Q}_{i})\) ( \(i=1,2\) ). To then prove Cor. 3, consider the quadratic form Q' which is the external direct sum of Q \(_{1}\) and Q \(_{2}\) (§ 3, No. 4), and note that one has Q'(x)=Q(x)+F(x,x), F being the bilinear form defined by F(x \(_{1}\) +x \(_{2}\) , y \(_{1}\) +y \(_{2}\) ) = -Φ(x \(_{1}\) , y \(_{2}\) ) for x \(_{i}\) , y \(_{i}\) in E \(_{i}\) ( \(i=1,2\) ); then use prop. 3.)

¶ 2) Suppose that E is the direct sum of two orthogonal submodules E \(_{1}\) , E \(_{2}\) and denote by Q \(_{i}\) the restriction of Q to E \(_{i}\) (i = 1, 2);

suppose moreover that there exists \(u \in \mathrm{C}^{+}(\mathrm{Q}_{2})\) such that \(u^{2} = a.1\) , \(a\) being invertible in A, and that \(u\rho_{\mathrm{Q}_{2}}(x_{2}) = -\rho_{\mathrm{Q}_{2}}(x_{2})\) for all \(x_{2} \in \mathrm{E}_{2}\) . Show that there exists an isomorphism \(\varphi\) of C(Q) onto the tensor product \(\mathrm{C}(a\mathrm{Q}_{1}) \otimes \mathrm{C}(\mathrm{Q}_{2})\) (as defined in Chap.~III, § 3, no. 1), having the property that for every \(x = x_{1} + x_{2}\) ( \(x_{i} \in \mathrm{E}_{i}, i = 1, 2\) ), one has

\[
\varphi (\rho_ {q} (x)) = \rho_ {a q _ {1}} (x _ {1}) \otimes u ^ {- 1} + 1 \otimes \rho_ {q _ {2}} (x _ {2}).
\]

(Prove the existence of the homomorphism \(\varphi\) as a consequence of the universal property of C(Q). On the other hand, there is a homomorphism \(g_{1}\) of C(aQ \(_{1}\) ) into C(Q) such that \(g_{1}(\rho_{aQ_{1}}(x_{1})) = h_{2}(u)\rho_{Q_{1}}(x_{1})\) , denoting by \(h_{2}\) the canonical homomorphism of C(Q \(_{2}\) ) into C(Q); then observe that \(g_{1}(z_{1})\) and \(h_{2}(z_{2})\) commute for \(z_{1} \in\) C(aQ \(_{1}\) ) and \(z_{2} \in\) C(Q \(_{2}\) ), and deduce from this the existence of an inverse homomorphism of \(\varphi\) .)

In which case does this result apply when A is a field of characteristic ≠ 2 (use Remark 2 following the cor. of th. 2)?

\begin{enumerate}
\def\labelenumi{\arabic{enumi})}
\setcounter{enumi}{2}
\tightlist
\item
  \begin{enumerate}
  \def\labelenumii{\alph{enumii})}
  \tightlist
  \item
    With the notations of no. 2, let \(x_{i}\) ( \(1 \leqslant i \leqslant n\) ) be elements of E such that \(f(x_{i}) = 0\) ; show that one has \(i_{j}(x_{1} \otimes x_{2} \otimes \cdots \otimes x_{n}) = 0\) . In particular, if F is a bilinear form on E such that F( \(x_{i}, x_{j}\) ) = 0 for \(i > j\) , one has \(i_{x_{1}}^{\mathrm{F}}(x_{2} \otimes x_{3} \otimes \cdots \otimes x_{n}) = 0\) .
  \end{enumerate}
\end{enumerate}

\begin{enumerate}
\def\labelenumi{\alph{enumi})}
\setcounter{enumi}{1}
\tightlist
\item
  With the notations of Prop. 3 of no. 3, let \(x_i\) ( \(1 \leqslant i \leqslant n\) ) be elements of E such that \(F(x_i, x_j) = 0\) for \(i > j\) . Show that we have \(\overline{\lambda}_{\mathrm{F}}(\rho_{\mathrm{Q}'}(x_1) \ldots \rho_{\mathrm{Q}'}(x_n)) = \rho_{\mathrm{Q}}(x_1) \ldots \rho_{\mathrm{Q}}(x_n)\) (using \(a\) ) and formula (10) of no. 2).
\item
  We assume that A is a field of characteristic \(\neq 2\) ; for every quadratic form Q on E, let \(\mu_{\mathrm{Q}}\) be the map \(\overline{\lambda}_{\mathrm{F}}\) of C(Q) onto \(\wedge\) E corresponding to \(F(x, y) = \frac{1}{2}\Phi(x, y)\) . Show that if the vectors \(x_i\) ( \(1 \leqslant i \leqslant n\) ) are pairwise orthogonal, we have
\end{enumerate}

\[
\mu_ {\mathrm{Q}} (x _ {1} x _ {2} \dots x _ {n}) = x _ {1} \wedge x _ {2} \wedge \dots \wedge x _ {n}.
\]

Deduce that if \(y_{i}\) ( \(1 \leqslant i \leqslant n\) ) are \(n\) arbitrary vectors, we have

\[
n! \mu_ {0} ^ {- 1} (y _ {1} \wedge y _ {2} \wedge \dots \wedge y _ {n}) = \sum_ {\sigma \in \mathfrak {S} _ {n}} \varepsilon_ {\sigma} y _ {\sigma (1)} y _ {\sigma (2)} \dots y _ {\sigma (n)}.
\]

\begin{enumerate}
\def\labelenumi{\arabic{enumi})}
\setcounter{enumi}{3}
\tightlist
\item
  Let \(C_{h}\) be the submodule of C(Q) generated by the products of k elements of \(\rho_{Q}(E)\) for \(0 \leqslant k \leqslant h\) . Show that the map
\end{enumerate}

\[
(x _ {1}, \dots , x _ {h}) \rightarrow \rho_ {\mathrm{Q}} (x _ {1}) \dots \rho_ {\mathrm{Q}} (x _ {h})
\]

defines by passing to quotients an alternating multilinear map from \(E^{h}\) into \(C_{h}/C_{h-1}\) . From this we deduce a linear map \(\pi_{h}\) of \(\bigwedge E\) into \(C_{h}/C_{h-1}\) ; show that \(\pi_{h}\) is an isomorphism when E is a free module. If f is an orthogonal transformation, we have \(C(f)\circ\pi_{h}=\pi_{h}\circ(\bigwedge f)\) . When A is a field of characteristic \(\neq2\) , show that \(\pi_{h}\) is the restriction of the map \(\mu_{Q}^{-1}\) defined in Exerc. 3 c).

\begin{enumerate}
\def\labelenumi{\arabic{enumi})}
\setcounter{enumi}{4}
\tightlist
\item
  With the notation of no. 2, consider the map \(i_{f}\) of \(\wedge\) from E into itself. Show that one has
\end{enumerate}

\[
i _ {f} \left(x _ {1} \wedge \dots \wedge x _ {h}\right) = \sum_ {i = 1} ^ {h} (- 1) ^ {i - 1} f \left(x _ {i}\right) \left(x _ {1} \wedge \dots \wedge x _ {i - 1} \wedge x _ {i + 1} \wedge \dots \wedge x _ {h}\right)
\]

and deduce that \(i_{f}\) is none other than the right interior product \(z \to z \sqcup f\) (chap.~III, § 8, no. 4, def. 2).

\begin{enumerate}
\def\labelenumi{\arabic{enumi})}
\setcounter{enumi}{5}
\item
  Show that, if E admits an orthogonal basis \((e_{1}, e_{2})\) for Q, and if we set \(\alpha_{i} = \mathrm{Q}(e_{i})\) For (i = 1, 2), the algebra C(Q) is isomorphic to the algebra of quaternions over A corresponding to the pair \((\alpha_{1}, \alpha_{2})\) .
\item
  Let A be a maximal ordered field, E a vector space of even dimension 2r over A, Q a nondegenerate quadratic form on E, positive or negative. Show that, if Q is positive, C(Q) is isomorphic to a matrix algebra over A when \(r(r-1)/2\) is even, to a matrix algebra over the quaternion algebra over A when \(r(r-1)/2\) is odd. If Q is negative, C(Q) is isomorphic to a matrix algebra over A when \(r(r+1)/2\) is even, to a matrix algebra over the quaternion algebra over A when \(r(r+1)/2\) is odd (use Exercises 2 and 6). What is the structure of C⁺(Q) in these different cases?
\item
  Let A be the ring \(\mathbf{Z}/(4)\) a left A-module \(\mathbf{Z}/(2)\) ; if we set \(Q(0)=0\) , \(Q(u)=1\) for the unique element \(u\neq0\) of E, Q is a quadratic form on E. Show that the A-modules C(Q) and \(\wedge\) They are not isomorphic (one will prove that C(Q) is isomorphic to the direct sum of two modules isomorphic to E).
\end{enumerate}

¶ 9) Suppose A is a field of characteristic 2, E a vector space of finite even dimension 2r, and Q a nondegenerate quadratic form on E.

\begin{enumerate}
\def\labelenumi{\alph{enumi})}
\tightlist
\item
  Let \((e_i)\) a symplectic basis (§ 5, no. 1) of E for the alternating bilinear form \(\Phi\) associated with Q. Show that the element
\end{enumerate}

\[
z = e _ {1} e _ {2} + e _ {3} e _ {4} + \dots + e _ {2 r - 1} e _ {2 r}
\]

of C \(^{+}\) (Q), together with the identity element, forms a basis of the center Z of C \(^{+}\) For Z to be a direct sum of two fields, it is necessary and sufficient that the element

\[
\Delta (Q) = Q \left(e _ {1}\right) Q \left(e _ {2}\right) + Q \left(e _ {3}\right) Q \left(e _ {4}\right) + \dots + Q \left(e _ {2 r - 1}\right) Q \left(e _ {2 r}\right)
\]

(called the pseudo-discriminant of Q with respect to the symplectic basis \((e_{i})\) ) is of the form \(\lambda^{2} + \lambda\) , with \(\lambda \in A\) .

\begin{enumerate}
\def\labelenumi{\alph{enumi})}
\setcounter{enumi}{1}
\tightlist
\item
  Let \(u\) a similarity for \(\Phi\) ( \(\S 6\) , no. 5) with multiplier \(\mu(u)\) , and let \(Q_{1}(x) = Q(u(x))\) . We set
\end{enumerate}

\[
u (e _ {2 i - 1}) = \sum_ {j = 1} ^ {r} a _ {i j} e _ {2 j - 1} + \sum_ {j = 1} ^ {r} b _ {i j} e _ {2 j},
\]

\[
u (e _ {2 i}) = \sum_ {j = 1} ^ {r} c _ {i j} e _ {2 j - 1} + \sum_ {j = 1} ^ {r} d _ {i j} e _ {2 j},
\]

and \(\mathrm{Q}(e_{2i - 1}) = \alpha_i, \mathrm{Q}(e_{2i}) = \beta_i (1 \leqslant i \leqslant r)\) . Show that we have

\[
\Delta (\mathrm{Q} _ {1}) = (\mu (u)) ^ {2} \Delta (\mathrm{Q}) + (\mathrm{D} (u)) ^ {2} + \mu (u) \mathrm{D} (u)
\]

where

\[
\mathbf {D} (u) = \sum_ {i, j} (\alpha_ {j} a _ {i j} c _ {i j} + \beta_ {j} b _ {i j} d _ {i j} + b _ {i j} c _ {i j})
\]

(Dickson invariant of u with respect to the basis \((e_{i})\) ). (Note that the element

\[
(\mu (u)) ^ {- 1} (u (e _ {1}) u (e _ {2}) + u (e _ {3}) u (e _ {4}) + \dots + u (e _ {2 r - 1}) u (e _ {2 r}))
\]

belongs to Z.) For \(u\) to be a similarity for Q (§ 4, exerc. 9) with multiplier \(\mu(u)\) , it is necessary and sufficient that D(u) = 0 or D(u) = \(\mu(u)\) ; the similarities for Q such that D(u) = 0 are called direct.

\begin{enumerate}
\def\labelenumi{\alph{enumi})}
\setcounter{enumi}{2}
\tightlist
\item
  Show that, if v is a similarity for \(\Phi\) , u a similarity for Q, one has
\end{enumerate}

\[
\mathrm{D} (u v) = \mu (v) \mathrm{D} (u) + \mu (u) \mathrm{D} (v)
\]

(consider the Dickson invariant of \(\varrho\) with respect to the symplectic basis formed by the \(u(e_{2i-1})\) and \((\mu(u))^{-1}u(e_{2i})\) for \(1 \leqslant i \leqslant r\) ). Deduce from this that the direct similarities for Q form a distinguished subgroup of index 2 in the group of similarities for Q.

\begin{enumerate}
\def\labelenumi{\alph{enumi})}
\setcounter{enumi}{3}
\item
  If \(u\) is the symmetry with respect to the hyperplane orthogonal to a nonsingular vector in E, show that \(\mathrm{D}(u)=1\) . Deduce from this that the group \(\varphi(\mathrm{G}^{+})\) (notations of no. 5) is the group \(\mathbf{SO}(\mathrm{Q})\) defined in exerc. 28 c) of § 6.
\item
  Let \(u \in \mathbf{O}(Q)\) , and suppose \(\Lambda\) algebraically closed. Show that, for \(u \in \mathbf{SO}(Q)\) , it is necessary and sufficient that the number of elementary divisors of the module \(E_u\) be even. (With the notations of exerc. 15 of § 5, first note that if \(p, q\) are two distinct irreducible factors of the minimal polynomial of \(u\) , the number of indecomposable submodules of which \(G(p, q)\) is the direct sum is equal to the number of indecomposable submodules of which \(G(q, p)\) is the direct sum. Note on the other hand that \(G(p, p) = \{0\}\) except for \(p = X - 1\) . Finally prove that one may restrict to the case where \(E_u\) is equal to \(G(p, p)\) (with \(p = X - 1\) ) and is indecomposable. There is then a symplectic basis ( \(e_i\) ) of E such that \(e_1, e_3, \ldots, e_{2k-1}\) form a basis of \((p(u))^{2r-k}(E)\) for \(1 \leqslant k \leqslant r\) ; show that \(e_1, e_3, \ldots, e_{2r-3}\) are singular vectors, and conclude that \(D(u) = 1\) .)
\end{enumerate}

¶ 10) Suppose A is a field, E a finite-dimensional vector space, Q a degenerate quadratic form on E; let M be a subspace complementary to E \(^{0}\) in E, let B be the (semisimple) Clifford algebra of the restriction of Q to M.

\begin{enumerate}
\def\labelenumi{\alph{enumi})}
\item
  We first assume that A has characteristic ≠ 2. Let L be the Clifford algebra of the restriction of Q to E⁰ (isomorphic to ∧ E⁰), and R₀ its radical (the ideal generated in L by E⁰, of codimension 1 in L); show that the radical R of C(Q) is obtained (up to isomorphism) by defining the algebra structure on \(B \otimes_{A} R_{0}\) as in cor. 4 of th. 1, that C(Q)/R is isomorphic to B, and that C(Q) is the direct sum of B and R.
\item
  Assume A has characteristic 2. Let F be the subspace of \(\mathbf{E}^0\) consisting of singular vectors \(x \in \mathbf{E}^0\) and let N be a complement of F with respect to \(\mathbf{E}^0\) . If \((a_i)_{1 \leqslant i \leqslant d}\) is a basis of N, and \(\mathrm{Q}(a_i) = \alpha_i\) the elements \(\alpha_i^{1/2}\) in an algebraic closure of A, are linearly independent over A. Let \((\alpha_i^{1/2})_{1 \leqslant i \leqslant e}\) be a 2-basis of the field \(\mathrm{A}_1 = \mathrm{A}(\alpha_1^{1/2}, \ldots, \alpha_d^{1/2})\) on A (Chap.~V, § 8, exerc. 1), and let \(h = \dim F\) . If \(\mathrm{B}_1\) is the central simple algebra \(\mathrm{B} \otimes_{A} \mathrm{A}_1\) , show that \(\mathrm{C(Q)}\) is isomorphic to the algebra \(\mathrm{B}_1 \otimes_{A_1} \mathrm{L}_1\) , where \(\mathrm{L}_1\) is the exterior algebra of a vector space of dimension \(h + d - e\) on \(\mathrm{A}_1\) . If \(\Re_1\) is the radical of \(\mathrm{L}_1\) (of codimension 1 (over \(\mathrm{A}_1\) ) in \(\mathrm{L}_1\) ), the radical \(\Re\) of \(\mathrm{C(Q)}\) is isomorphic to the algebra \(\mathrm{B}_1 \otimes_{A_1} \Re_1\) , \(\mathrm{C(Q)/R}\) is isomorphic to \(\mathrm{B}_1\) , and \(\mathrm{C(Q)}\) is a direct sum of \(\mathrm{B}_1\) and of \(\Re\) .
\item
  The hypotheses being those of \(b\) ; one assumes moreover that \(h=0\) , \(d=1\) , \(e=0\) (which implies that dim M is even). If Q \(_{0}\) is the restriction of Q to M, show that C \(^{+}\) (Q) is isomorphic to C(Q \(_{0}\) ).
\end{enumerate}

11) a) Under the hypotheses and notation of no. 5, show that N(G⁺) is the subgroup of A* generated by the products Q(x)Q(y), where x and y range over the set of non-isotropic vectors of E. (Reduce to the case A ≠ F₂; use prop. 5 and exerc. 28 c) of § 6, as well as exerc. 9 d) of § 9.) Case where Q has index ≥ 1.

\begin{enumerate}
\def\labelenumi{\alph{enumi})}
\setcounter{enumi}{1}
\item
  We assume in addition that \(\mathbf{A} \neq \mathbf{F}_2\) that E has dimension \(n \geqslant 2\) and that Q has index \(>0\) . Show that \(\mathbf{O}_0^+(\mathbf{Q})\) is the commutator subgroup of the group \(\mathbf{O}(\mathbf{Q})\) . (Reduce to the case \(n = 2\) using Exercises 17 c) and 28 e) of § 6.)
\item
  We keep the hypotheses of \(b\) , and we suppose in addition that A has characteristic \(\neq 2\) and that E is of even dimension. For the automorphism \(x \to -x\) of E belong to \(\mathbf{O}_0^+(\mathbf{Q})\) it is necessary and sufficient that the discriminant of Q (with respect to any basis of E) be a square.
\end{enumerate}

Let a be an invertible element of A; show that there exists one and only one isomorphism \(\theta_{a}\) of C \(^{+}\) (Q) onto C \(^{+}\) (aQ) such that

\[
\theta_ {a} (\rho (x) \rho (y)) = a ^ {- 1} \rho_ {1} (x) \rho_ {1} (y)
\]

for all \(x, y\) in E ( \(\rho\) and \(\rho_{1}\) denoting the canonical maps from E into C(Q) and C(aQ), respectively). Deduce that if A is a field of characteristic \(\neq 2\) , E a vector space of odd dimension and Q a non-degenerate quadratic form, C(Q) and C(aQ) are isomorphic (cf.~exerc. 7).

\begin{enumerate}
\def\labelenumi{\alph{enumi})}
\setcounter{enumi}{1}
\tightlist
\item
  Assume that A is a field, E a vector space of even dimension 2r \textgreater{} 0, Q a nondegenerate quadratic form. Let u be a similarity relative to Q; show that there exists one and only one A-automorphism \(\overline{u}\) of the algebra \(\mathrm{C}^{+}(\mathrm{Q})\) such that for \(1 \leqslant h \leqslant r\) and \(x_{i} \in E (1 \leqslant i \leqslant 2h)\) one has
\end{enumerate}

\[
\overline {{{u}}} (x _ {1} x _ {2} \dots x _ {2 h}) = \mu^ {- h} u (x _ {1}) u (x _ {2}) \dots u (x _ {2 h})
\]

where \(\mu\) denotes the multiplier of \(u\) . For \(\bar{u}\) to be an inner automorphism, it is necessary and sufficient that \(u\) be a direct similarity (§ 6, no. 5 and § 9, exerc. 9 b)). One assumes in addition \(r \geqslant 2\) ; then, for \(\bar{u}\) to be the identity, it is necessary and sufficient that \(u\) be a homothety.

Show that the automorphism \(\bar{u}\) of C \(^{+}\) (Q) is the restriction to C \(^{+}\) (Q) of an inner automorphism of C(Q).

¶ 13) Suppose that A is a field of characteristic ≠ 2, E a vector space of even dimension 2r. \textgreater{} 0, Q a nondegenerate quadratic form. We denote by \(E_{h}\) the inverse image of \(\Lambda E\) under the isomorphism \(\mu_{Q}\) of exerc. 3 c), so that \(E_{h}\) is the subspace of C(Q) generated by the products \(x_{1}x_{2}\ldots x_{h}\) , where the \(x_{i}\) ( \(1 \leqslant i \leqslant h\) ) are pairwise orthogonal.

\begin{enumerate}
\def\labelenumi{\alph{enumi})}
\tightlist
\item
  Let \(\alpha\) an element of A, \(s\) an element of E \(_2\) . Show that, if \((\alpha + s)^2 \in \mathbf{A}\) , then either \(s = 0\) , or else \(\alpha = 0\) and \(s = xy\) , where \(x\) and \(y\) are two orthogonal vectors. (Note that if \(t = \mu_Q(\alpha + s)\) , \(t \wedge t\)
\end{enumerate}

necessarily belongs to A + ∧ E, by expressing s using an orthogonal basis of E; deduce that one necessarily has t = β + x ∧ y with β ∈ A, x, y in E, using § 5, no. 1, cor. 2 of th. 1).

\begin{enumerate}
\def\labelenumi{\alph{enumi})}
\setcounter{enumi}{1}
\item
  Let \((x, y)\) , \((u, v)\) two pairs of orthogonal non-isotropic vectors in E, \(\mathrm{P}_{xy}\) and \(\mathrm{P}_{uv}\) the planes \(Ax + Ay\) , \(Au + Av\) respectively. For one to have \((xy)(uv) = -(uv)(xy)\) in C(Q), it is necessary and sufficient that \(\mathrm{P}_{xy} + \mathrm{P}_{uv}\) be a non-isotropic subspace of dimension 3, in which \(\mathrm{P}_{xy}\) and \(\mathrm{P}_{uv}\) are weakly orthogonal (§ 3, exerc. 11).
\item
  Let g be an A-automorphism of \(\mathrm{C}^{+}(\mathrm{Q})\) sending \(E_{2}\) into itself; show that there exists a similarity u relative to Q such that \(g = \bar{u}\) (exerc. 12 b)). (Let \((e_{i})_{1 \leqslant i \leqslant 2r}\) be an orthogonal basis of E; using a) and b), show that there exists an orthogonal basis \((x_{i})\) of E such that \(g(e_{1}e_{i}) = x_{1}x_{i}\) for \(2 \leqslant i \leqslant 2r\) .)
\end{enumerate}

\begin{enumerate}
\def\labelenumi{\arabic{enumi})}
\setcounter{enumi}{13}
\tightlist
\item
  We assume that A is a field of characteristic \(\neq 2\) E a vector space of dimension \(n\) on A, Q and Q₁ two non-degenerate quadratic forms on E; let \(\tilde{\Delta}\) (resp. \(\tilde{\Delta}_{1}\) ) be the class of the discriminant \(\Delta\) of Q (resp. of the discriminant \(\Delta_{1}\) of Q \(_{1}\) ) with respect to a basis of E, in the group A \(^{*}\) /(A \(^{*}\) )², a class that does not depend on the basis chosen in E. We assume \(n = 2\) .
\end{enumerate}

\begin{enumerate}
\def\labelenumi{\alph{enumi})}
\item
  Show that if \(\tilde{\Delta} = \tilde{\Delta}_{1}\) and if the Clifford algebras C(Q) and C(Q \(_{1}\) ) are isomorphic, then Q and Q \(_{1}\) are equivalent (consider on C(Q) the quadratic form \(z \to z\bar{z}\) , where \(z \to \bar{z}\) is the unique involutive antiautomorphism of C(Q) whose set of invariants is the center of C(Q) (cf.~exerc. 6 and chap.~VIII, § 11, exerc. 5 e)); apply Witt's theorem).
\item
  For C(Q) to be isomorphic to the matrix algebra \(\mathbf{M}_{2}(\mathrm{A})\) it is necessary and sufficient that, for at least one \(x \neq 0\) in E, there exists \(y\) , \(z\) in E such that Q(x) + Q(y)Q(z) = 0; if this is the case, this property holds for every \(x \neq 0\) in E.
\end{enumerate}

¶ 15) We keep the notations and general hypotheses of exerc. 14, but we assume n = 3.

\begin{enumerate}
\def\labelenumi{\alph{enumi})}
\item
  Show that \(\mathrm{C}^{+}(\mathrm{Q})\) is isomorphic to a quaternion algebra over A and that \(\beta\) is the antiautomorphism \(z\to \bar{z}\) of this algebra whose set of invariants is the center of \(\mathrm{C}^{+}(\mathrm{Q})\) ; if P is the subspace of \(\mathrm{C}^{+}(\mathrm{Q})\) consisting of pure quaternions (that is, those such that \(z = -\bar{z}\) ; chap.~VIII, § 11, exerc. 6), the restriction to P of the quadratic form \(z\to z\bar{z}\) is equivalent to \(\lambda Q\) , where \(\lambda \in A\) . Deduce that, in order for \(C^{+}(Q)\) to be a field, it is necessary and sufficient that \(Q\) be of index 0.
\item
  Show that if \(\tilde{\Delta} = \tilde{\Delta}_{1}\) and if the Clifford algebras C(Q) and C(Q \(_{1}\) ) are isomorphic, then Q and Q \(_{1}\) are equivalent. (First consider the case where - \(\Delta\) is a square in K, and show that in this case C \(^{+}\) (Q) and C \(^{+}\) (Q \(_{1}\) ) are isomorphic; then reason as in exerc. 14 a), using a) and exerc. 6 of chap.~VIII, § 11. In the general case, use exerc. 12 a)).
\item
  Show that the special Clifford group G \(^{+}\) (for the form Q) is identical to the group of invertible elements of C \(^{+}\) (Q). (If ( \(e_{1}\) , \(e_{2}\) , \(e_{3}\) ) is an orthogonal basis of E, and if \(j = e_{1}e_{2}e_{3}\) in C(Q), note that \(x \to xj\) is an isomorphism of vector spaces from E onto P).
\item
  Deduce from a) and c) that if Q has index 1, the group of rotations \(\mathbf{O}^{+}(\mathrm{Q})\) is isomorphic to the projective group \(\mathbf{PGL}_{2}(\mathrm{A})\) (chap.~II, \(2^{\mathrm{e}}\) éd., App. III, no. 6).
\end{enumerate}

We keep the general hypotheses and notations of exerc. 14, but we assume n = 4.

\begin{enumerate}
\def\labelenumi{\alph{enumi})}
\item
  Give an example where \(\tilde{\Delta} = \tilde{\Delta}_{1}\) and where \(C(Q)\) and \(C(Q_{1})\) are isomorphic, but where \(Q\) and \(Q_{1}\) are not equivalent (cf.~exerc. 7).
\item
  Let \((e_i)_{1\leqslant i\leqslant 4}\) an orthogonal basis of E for Q, \(Q_0\) the restriction of Q to the hyperplane H = Ae \(_1\) + Ae \(_2\) + Ae \(_3\) Show that, if Z is the center of C \(^+\) (Q), the algebra C \(^+\) (Q) is isomorphic to the tensor product Z ⊗ \(_A\) C \(^+\) (Q \(_0\) ). For every z ∈ C \(^+\) (Q), one has β(z)z ∈ Z ; for z to belong to the special Clifford group G \(^+\) it is necessary and sufficient that z be invertible and that β(z)z ∈ A.1. Deduce from this that the group O \(_0^+\) (Q) is isomorphic to the quotient by \{1,-1\} of the group of elements z ∈ Z ⊗ \(_A\) C \(^+\) (Q \(_0\) ) such that β(z)z = 1.
\item
  We assume that \(\Delta\) is not a square in A (which implies, by Witt's theorem, that the index of Q is 0 or 1). If \(Q_0'\) is the quadratic form obtained from \(Q_0\) by extension to \(A' = A(\sqrt{\Delta})\) of the field of scalars, deduce from \(b\) ) that \(\mathbf{O}_0^+(Q)\) is isomorphic to \(\mathbf{O}_0^+(Q_0')\) . In particular, if Q has index 1, \(\mathbf{O}_0^+(Q)\) is the commutator subgroup of \(\mathbf{O}(Q)\) and is isomorphic to \(\mathbf{PSL}_2(A')\) (cf.~exerc. 15 \(d\) ) and chap.~III, § 7, exerc. 8).
\item
  We assume that \(\Delta\) is a square in A (which implies, by Witt's theorem, that Q has index 0 or 2) and that \(Q(e_4) = 1\) . If one sets \(j = e_1e_2e_3\)  \(x \in E\) can be written in only one way \(x = \alpha e_4 + jz\) , where \(\alpha \in A\) and \(z\) is a pure quaternion (exerc. 15 a)) in \(L = C^+(Q_0)\) ; if we set \(\psi(x) = \alpha + z\) , \(\psi\) is a vector space isomorphism from E onto L. Let \(Z = Ac' + Ac''\) , where \(c'\) and \(c''\) are the two orthogonal idempotents in Z; every invertible element \(s \in C^+(Q)\) is written in a unique way \(s = uc' + vc''\) , where \(u\) and \(\nu\) belong to L; so that \(s\) belongs to the special Clifford group \(G^+\) , it is necessary and sufficient that \(u\bar{u} = v\bar{v}\) , and we then have \(\psi(sxs^{-1}) = u\psi(x)v^{-1}\) for all \(x \in E\) Deduce from this that the quotient of \(\mathbf{O}_0^+(Q)\) by its center (which is a group with 2 elements, cf.~exerc. 11 c)) is isomorphic to the product \(\mathbf{O}_0^+(Q_0) \times \mathbf{O}_0^+(Q_0)\) ; in particular, if Q has index 2, this quotient group is isomorphic to \(\mathbf{PSL}_2(A) \times \mathbf{PSL}_2(A)\) .
\end{enumerate}

\begin{enumerate}
\def\labelenumi{\arabic{enumi})}
\setcounter{enumi}{16}
\tightlist
\item
  Let K be a commutative field of characteristic ≠ 2, A the field \(\mathrm{K}(\mathrm{X}_{n})_{n\in\mathbb{N}}\) of rational fractions over K, with respect to a countable family of indeterminates (chap.~IV, § 3, n° 1). Let E be a vector space over A, having a countable basis \((e_{n})_{n\in\mathbb{N}}\) , and let Φ be a symmetric bilinear form on E, for which \((e_{n})\) is an orthogonal basis and such that \(\Phi(e_{n}, e_{n}) = \mathrm{X}_{n}\) for all \(n \in N\) . If one sets \(\mathrm{Q}(x) = \Phi(x, x)\) , show that the Clifford algebra C(Q) is a field (cf.~chap.~VIII, § 12, exerc. 14).
\end{enumerate}

\section{Angles}

Throughout this paragraph, A denotes a commutative field of characteristic \(\neq 2\) . Let E be a vector space of dimension 2 over A, and let \(\Phi\) be a nondegenerate symmetric bilinear form on E.

\subsection{Direct similarities in a plane.}

Recall (§ 6, no. 5) that a direct similarity of E is an automorphism u of the vector space E such that \(\Phi(u(x), u(y)) = (\det u)\Phi(x, y)\) for all \(x, y\) in E.

PROPOSITION 1. --- Let A(Φ) be the subalgebra of \(\mathfrak{L}_{\mathbf{A}}(\mathbf{E})\) generated by the direct similarities of E.

\begin{enumerate}
\def\labelenumi{\alph{enumi})}
\item
  Direct similarities are the invertible elements of A(Φ). The algebra A(Φ) is a commutative algebra of degree 2 over A. When E contains no nonzero isotropic vector, A(Φ) is a field, a quadratic extension of A; otherwise it is the direct product of two fields isomorphic to A.
\item
  Let \((e_{1}, e_{2})\) an orthogonal basis of E; set \(\alpha_{i} = \Phi(e_{i}, e_{i})\) (i = 1, 2) and \(\delta = -\alpha_{2}/\alpha_{1}\) Then the matrices of the elements of A( \(\Phi\) ) with respect to this basis are the matrices of the form \(\begin{pmatrix} a & \delta b \\ b & a \end{pmatrix}\) where \(a \in A\) , \(b \in A\) .
\item
  The space E is a free cyclic A(Φ)-module, generated by any non-isotropic vector.
\end{enumerate}

Indeed, let us introduce an auxiliary alternating bilinear form B ≠ 0 on E; then B is non-degenerate. There therefore exists an endomorphism w of E such that \(\Phi(x, y) = \mathrm{B}(w(x), y)\) for all x, y in E. For every invertible endomorphism u of E, we have

\[
\begin{array}{r l} \Phi (u (x), u (y)) & = \mathrm{B} (w u (x), u (y)) = \mathrm{B} (u u ^ {- 1} w u (x), u (y)) \\ & = (\det u) \mathrm{B} (u ^ {- 1} w u (x), y); \end{array}
\]

For u to be a direct similarity, it is necessary and sufficient that one have

\[
(\det u) \mathrm{B} (u ^ {- 1} w u (x), y) = (\det u) \Phi (x, y) = (\det u) \mathrm{B} (w (x), y)
\]

for all \(x, y\) in E; since det \(u \neq 0\) and since B is nondegenerate, this is equivalent to \(u^{-1}wu = w\) , or equivalently to \(uw = wu\) . Take for B the alternating bilinear form whose matrix \(S\) with respect to \((e_1, e_2)\) is \(\begin{pmatrix} 0 & 1 \\ -1 & 0 \end{pmatrix}\) by denoting \(R\) and \(W\) the matrices of \(\Phi\) and of \(w\) with respect to this basis, the relation \(\Phi(x, y) = \mathrm{B}(w(x), y)\) is written, by virtue of formula (47) of § 1, no. 10, \(R = ^t W.S\) ; making this explicit shows that we have \(W = \begin{pmatrix} 0 & \alpha_2 \\ -\alpha_1 & 0 \end{pmatrix}\) . If \(\begin{pmatrix} a & c \\ b & d \end{pmatrix}\) denotes the matrix of \(u\) with respect to \((e_1, e_2)\) , the relation \(uw = wu\) is therefore equivalent to the relations \(b\alpha_2 = -c\alpha_1\) , \(a\alpha_2 = d\alpha_2\) , \(a\alpha_1 = d\alpha_1\) that is, to \(a = d\) and \(c = \delta b\) ; this proves that the matrices of direct similarities are the invertible matrices of the form \(\begin{pmatrix} a & \delta b \\ b & a \end{pmatrix}(a, b\) in A).

Now the endomorphisms of E whose matrices are of the form \(\begin{pmatrix} a & \delta b \\ b & a \end{pmatrix}\) (a, b in A) form a vector subspace of dimension 2 of \(\mathfrak{L}_{\mathrm{A}}(\mathrm{E})\) generated by 1 and by the endomorphism w; since \(w^{2}\) is the homothety with ratio \(-\alpha_{1}\alpha_{2}\) , this subspace is the subalgebra \(\mathrm{A}(\Phi)\) of \(\mathfrak{L}_{\mathrm{A}}(\mathrm{E})\) generated by the direct similarities. The direct similarities are the invertible elements of \(\mathrm{A}(\Phi)\) , that is, those whose matrices satisfy \(a^{2}-\delta b^{2}\neq0\) . The fact that the algebra \(\mathrm{A}(\Phi)\) is commutative is evident. Let us apply to it the results of Chap.~II, §7, no.~7: if δ is not a square in A, that is, if no nonzero vector of E is isotropic, \(\mathrm{A}(\Phi)\) is a field; if, on the contrary, δ is a square in A, that is, if E contains nonzero isotropic vectors, \(\mathrm{A}(\Phi)\) is a direct sum of two fields isomorphic to A. This proves a) and b).

Finally, any non-isotropic vector of E can be taken as

first vector \(e_1\) of an orthogonal basis \((e_1, e_2)\) ( \(\S 6, n^0 1\) ); therefore its transforms \(u(e_1)\) by the elements \(u\) of A( \(\Phi\) ) are the vectors of the form \(ae_1 + be_2\) ( \(a, b\) in A), that is, all vectors of E, since all matrices \(\begin{pmatrix} a & \delta b \\ b & a \end{pmatrix}\) are matrices of elements of A( \(\Phi\) ). In other words, E is a cyclic A( \(\Phi\) )-module, generated by any non-isotropic vector. Moreover, we see that it is a A( \(\Phi\) -module monogenic free, since \(u(e_1) = ae_1 + be_2 = 0\) implies \(a = b = 0\) , hence \(u = 0\) . This proves \(c\) ).

Remarks. --- 1) Let \(\vartheta\) be the similarity with matrix \(\begin{pmatrix}0 & \delta \\ 1 & 0\end{pmatrix}\) with respect to \((e_1, e_2)\) ; the similarity \(w\) introduced in the proof of Prop. 1 is equal to \(-\alpha_1 v\) ; we have \(v^2 = \delta\) . The multiplier of the direct similarity \(u = a + b v\) ( \(a, b\) in A) is equal to the determinant of its matrix \(\begin{pmatrix}a & \delta b \\ b & a\end{pmatrix}\) that is, to \(a^2 - \delta b^2 = (a + b v)(a - b v) = u.\bar{u}\) , denoting by \(\bar{u}\) the conjugate of \(u\) in the algebra A(Φ) (Chap.~II, § 7, No. 7); in other words, the multiplier of \(u\) is the norm N( \(u\) ) of \(u\) in the algebra A(Φ) (ibid.). In particular, for a direct similarity \(u\) to be a rotation, it is necessary and sufficient that N( \(u\) ) = 1; for \(u\) to be a homothety, it is necessary and sufficient that \(u \in A^*\) .

\begin{enumerate}
\def\labelenumi{\arabic{enumi})}
\setcounter{enumi}{1}
\item
  The direct similarities \(u = a + bv\) ( \(a, b\) in A, \(b \neq 0\) ) whose square is a homothety are the homotheties and the scalar multiples \(bv\) of \(v\) , since \((a + bv)^{2} = (a^{2} + \delta b^{2}) + 2abv\) . These latter are none other than the automorphisms of the vector space E which transform every vector into an orthogonal vector; indeed the matrix of such an automorphism is necessarily of the form \(\begin{pmatrix} 0 & c \\ d & 0 \end{pmatrix} (c, d\) in A), and the condition of transforming the vector \(\lambda e_{1} + \mu e_{2}\) as an orthogonal vector is then written \(\lambda \mu(d\alpha_{2} + c\alpha_{1}) = 0\) .
\item
  One easily verifies that, for \(x\) , \(y\) in E and \(u \in A(\Phi)\) , one has \(\Phi(u(x), y) =: \Phi(x, \overline{u}(y))\) . Thus the adjoint endomorphism of a direct similarity \(u\) is the direct similarity \(\overline{u}\) conjugate of \(u\) in A( \(\Phi\) ).
\item
  Since every inverse similarity of E is the product of a direct similarity and the reflection with respect to \(Ae_{1}\) , the matrices of inverse similarities with respect to \((e_{1}, e_{2})\) are the matrices of the form \(\begin{pmatrix} a & -\delta b \\ b & -a \end{pmatrix}\) .
\end{enumerate}

We shall henceforth denote by S the group of similarities of E, by \(S^{+}\) the group of direct similarities, by H that of homotheties ≠ 0, and by \(O^{+}\) that of rotations. Recall that we have \(H \subset S^{+}\) (§ 6, no. 5).

Corollary 1. --- The group S \(^{+}\) of direct similarities is commutative. For any non-isotropic vectors x, y in E, there exists a unique direct similarity u such that y = u(x).

The first assertion follows from the fact that the algebra A(Φ) is commutative. Since E is a free monogenic A(Φ)-module generated by x (resp. y), there exists a unique element u (resp. u′) of A(Φ) such that \(y = u(x)\) (resp. \(x = u'(y)\) ); whence \(x = u(u'(x))\) , and \(uu'\) is the identity; this shows that u is invertible, and is therefore a direct similarity.

COROLLARY 2. --- The group O \(^{+}\) of rotations is commutative. For any vectors x, y in E such that Φ(x, x) = Φ(y, y) ≠ 0, there exists a unique rotation u such that y = u(x).

The first assertion follows from Cor. 1. This also shows that there exists a direct similarity \(u\) and exactly one such that \(y = u(x)\) ; since \(\Phi(u(x), u(x)) = \Phi(x, x)\) the multiplier of \(u\) is equal to 1, and \(u\) is therefore a rotation.

COROLLARY 3. --- The group S \(^{+}\) /H is commutative. It acts on the set of non-isotropic lines of E. Whatever the non-isotropic lines D, D′ of E may be, there exists one and only one element φ of S \(^{+}\) /H such that D' = φ(D).

This follows from Cor. 1 and from the fact that H leaves every line of E globally invariant.

PROPOSITION 2. --- The kernel of the canonical homomorphism from O \(^{+}\) into S \(^{+}\) /H is \{1, -1\}.

Indeed, 1 and -1 are the only elements of \(\mathbf{H} \cap \mathbf{O}^{+}\) .

PROPOSITION 3. --- The homomorphism \(u \to u/\overline{u} = u^{2}/\mathrm{N}(u)\) of \(\mathrm{S}^{+}\) into itself has H as its kernel, and defines, by passing to the quotient, an isomorphism from \(\mathrm{S}^{+}/\mathrm{H}\) onto \(\mathrm{O}^{+}\) .

Indeed, the relation \(u/\overline{u} = 1\) is equivalent to \(u = \overline{u}\) that is, to \(u \in A^* = H\) Thus H is the kernel of \(u \to u/\overline{u}\) . Since \(N(u/\overline{u}) = 1\) , \(u/\overline{u}\) is a rotation (Remark 1). It remains to show that every element \(v\) of \(O^+\) is of the form \(u/\overline{u}\) ( \(u \in S^+\) ). If \(1 + v\) is invertible, one can take \(u = 1 + v\) , because the relation \(N(v) = v\bar{v} = 1\) implies \(1 + v = v(1 + \bar{v})\) . Otherwise, one has \(N(1 + v) = (1 + v)(1 + \bar{v}) = 0\) that is, setting \(v = a + bw\) ( \(a \in A\) , \(b \in A\) , \(w^2 = \delta \in A\) ), \(2(1 + a) = 0\) , whence \(a = -1\) ; but the relations \(a = -1\) and \(N(v) = a^2 - \delta b^2 = 1\) imply \(b = 0\) , whence \(v = -1\) ; since \(\bar{w} = -w\) , it suffices, in this case, to take \(u = w\) .

When A(Φ) is a field, prop. 3 is a special case of Hilbert's theorem (chap.~V, § 11, no. 5, th. 3).

COROLLARY. — Let \(i: \mathrm{O}^{+} \to \mathrm{S}^{+}/\mathrm{H}\) and \(d: \mathrm{S}^{+}/\mathrm{H} \to \mathrm{O}^{+}\) denote the homomorphisms defined in props. 2 and 3, and write the abelian groups \(\mathrm{O}^{+}\) and \(\mathrm{S}^{+}/\mathrm{H}\) additively; we have \(d(i(\theta)) = 2\theta\) for \(\theta \in \mathrm{O}^{+}\) and \(i(d(\varphi)) = 2\varphi\) for \(\varphi \in \mathrm{S}^{+}/\mathrm{H}\) .

Indeed, if \(\bar{\nu}\) is a rotation, we have \(\bar{\nu} = \nu^{-1}\) , whence \(d(i(\nu)) = \nu/\bar{\nu} = \nu^{2}\) . On the other hand, if \(\varphi \in S^{+}/H\) , \(\varphi\) is the class mod. H of a direct similarity \(u\) , and \(d(\varphi) = u/\bar{u} = u^{2}/N(u)\) is congruent to \(u^{2}\) mod. H; whence the second formula.

\subsection{Plane trigonometry.}

In this section, we choose a generator w of the algebra A(Φ) such that \(w^{2} \in A\) . Such a generator is determined up to a homothety (no. 1, Remark 2), hence the element \(w^{2}\) of A, which we denote by δ, is determined modulo the multiplicative subgroup (A*)² of squares of nonzero elements of A.

Remark. — When -1 belongs to the class mod. (A*)² in question, one generally chooses w such that \(w^{2} = -1\) , which determines it up to sign. When this class contains 1 but not -1, one generally chooses w such that \(w^{2} = 1\) , which again determines it up to sign.

This being so, every element \(\nu\) of \(S^{+}\) can be written, in one and only one way, in the form

\[
\nu = c _ {w} (\nu) + s _ {w} (\nu) \varpi\tag{1}
\]

where \(c_w(v)\) , \(s_w(v)\) belong to A; the element \(s_w(v)/c_w(v)\) of the projective field \(\tilde{\mathbf{A}}\) (chap.~II, 2 \(^{e}\) ed., App. III, no. 5) is denoted \(t_w(v)\) ; it depends only on the class of \(v\) mod. H; thus \(t_w\) defines, by passage to the quotient, a map from S \(^{+}\) /H into the projective field \(\tilde{\mathbf{A}}\) , which we shall still denote by \(t_w\) by abuse of language. We shall often write \(c\) , \(s\) and \(t\) instead of \(c_w\) , \(s_w\) and \(t_w\) . We have \(c_{-w} = c_w\) , \(s_{-w} = -s_w\) and \(t_{-w} = -t_w\) .

PROPOSITION 4. — a) When \(w^2 = \delta\) is not a square in A (that is, when E contains no isotropic line), the map t from S⁺/H into \(\tilde{\mathbf{A}}\) is a bijection.

\begin{enumerate}
\def\labelenumi{\alph{enumi})}
\setcounter{enumi}{1}
\item
  When δ is the square of an element γ of A, the map t is a bijection from S⁺/H onto \(\tilde{\mathbf{A}}\) minus the elements 1/γ and -1/γ.
\item
  Writing \(\mathrm{S}^{+}/\mathrm{H}\) additively, we have , for \(\varphi, \varphi'\) in \(\mathrm{S}^{+}/\mathrm{H}\)
\end{enumerate}

\[
t (\varphi + \varphi^ {\prime}) = (t (\varphi) + t (\varphi^ {\prime})) / (1 + \delta t (\varphi) t (\varphi^ {\prime}))\tag{2}
\]

when \(t(\varphi)\) and \(t(\varphi')\) are finite and \(1 + t(\varphi)t(\varphi')\) is \(\neq 0\) .

Indeed, as \(S^{+}/H\) is a set of lines (with 0 removed) of \(A(\Phi)\) considered as a vector plane over A, t is injective. On the other hand, for an element \(a + bw\) ( \(a \in A\) , \(b \in A\) ) of \(A(\Phi)\) to be a direct similarity, it is necessary and sufficient that it be invertible, that is, that one have \(N(a + bw) = a^{2} - \delta b^{2} \neq 0\) , or again \((b/a)^{2} \neq 1/\delta\) ; this proves the surjectivity assertions in a) and b). Finally, the product of the similarities \(1 + t(\varphi)w\) and \(1 + t(\varphi')w\) is the similarity \(1 + \delta t(\varphi)t(\varphi') + (t(\varphi) + t(\varphi'))w\) , which proves c).

PROPOSITION 5. --- Let us write O \(^{+}\) additively. For every pair of elements θ, θ' of O \(^{+}\) we have

\begin{enumerate}
\def\labelenumi{(\arabic{enumi})}
\setcounter{enumi}{2}
\tightlist
\item
\end{enumerate}

\[
c (\theta) ^ {2} - \delta s (\theta) ^ {2} = 1\tag{4}
\]

\[
c (\theta + \theta^ {\prime}) = c (\theta) c (\theta^ {\prime}) + \delta s (\theta) s (\theta^ {\prime})\tag{5}
\]

\[
s (\theta + \theta^ {\prime}) = s (\theta) c (\theta^ {\prime}) + c (\theta) s (\theta^ {\prime}).
\]

The relation (3) indeed expresses that \(\mathrm{N}(c(\theta)+s(\theta)\omega)=1\) To prove (4) and (5), it suffices to compute, in \(\mathrm{A}(\Phi)\) the product of the rotations \(c(\theta)+s(\theta)\omega\) and \(c(\theta')+s(\theta')\omega\) .

PROPOSITION 6. — Let d be the isomorphism from S+/H onto O+ defined in Prop. 3. For every element φ of S+/H such that t = t(φ) is finite, one has

\[
s (d (\varphi)) = 2 t / (1 - \delta t ^ {2}), \quad c (d (\varphi)) = (1 + \delta t ^ {2}) / (1 - \delta t ^ {2}).\tag{6}
\]

Indeed \(\varphi\) is the class mod. H of the similarity \(1 + tw\) , and the rotation \(d(\varphi)\) is therefore \((1 + tw)^2/\mathrm{N}(1 + tw) = (1 + \delta t^2 + 2tw)/(1 - \delta t^2)\) (prop. 3), which proves (6).

COROLLARY. --- For every element θ of O⁺ such that \(t(\theta)\) is finite, we have

\[
s (2 \theta) = 2 t (\theta) / (1 - \delta t (\theta) ^ {2}), \quad c (2 \theta) = (1 + \delta t (\theta) ^ {2}) / (1 - \delta t (\theta) ^ {2}). \tag {7}
\]

For every element \(\varphi\) of \(S^{+}/H\) such that \(t(\varphi)\) is finite and \(1 + \delta t(\varphi)^{2} \neq 0\) , one has

\[
t (2 \varphi) = 2 t (\varphi) / (1 + \delta t (\varphi) ^ {2}).\tag{8}
\]

Indeed this follows immediately from prop. 6 and the corollary of prop. 3.

Remark. --- Formulas (6) remain true for \(t = \infty\) provided that the rational functions appearing on the right-hand side are replaced by their canonical extensions to the projective field. \(\tilde{\mathbf{A}}\) (chap.~II, 2 \(^{\text{e}}\) ed., App. III, no. 5); indeed, if \(t = \infty\) , \(\varphi\) is the class of \(w\) , and we have \(d(\varphi) = -1\) , \(s(d(\varphi)) = 0\) and \(c(d(\varphi)) = -1\) ; these are indeed the values taken by the canonical extensions of the right-hand sides for \(t = \infty\) The same holds for (7) when \(t(\theta) = \infty\) , and for (8) when \(t(\varphi) = \infty\) or that \(1 + \delta t(\varphi)^{2} = 0\) . Similarly, formula (2) remains true when only one of the elements \(t(\varphi)\) , \(t(\varphi')\) , \(t(\varphi)\) for example, is infinite, provided that its right-hand side is regarded as a rational function of \(t(\varphi)\) alone: indeed, the product of the similarities \(1 + t(\varphi')w\) and \(w\) is \(\delta t(\varphi') + w\) , while the value taken by the canonical extension of the right-hand side of (2) is \(1/\delta t(\varphi')\) . Finally, when \(t(\varphi)\) and \(t(\varphi')\) are finite and one has \(1 + \delta t(\varphi)t(\varphi') = 0\) , one has \(t(\varphi) + t(\varphi') = 0\) (otherwise \(t(\varphi)^{2}\) would be equal to \(1/\delta\) , which is impossible (prop. 4)); one may therefore agree that the value of the right-hand side of (2) is \(\infty\) , and this value is indeed that of the left-hand side. When \(t(\varphi)\) and \(t(\varphi')\) are both infinite, the right-hand side of (2) is not defined.

\subsection{Angles.}

We shall assume, in this section and the next, that A is an ordered field, hence of characteristic zero. Recall (Rectifications to Chap.~VI) that, if F is a vector space over A, the relation “ there exists λ \textgreater{} 0 such that y = λx ” is an equivalence relation between elements x, y of F - \{0\}, that every equivalence class for this relation is called an open half-line with origin 0, and that the union of an open half-line and \{0\} is called a closed half-line (or simply half-line) with origin 0; if D is a line and Δ a closed half-line contained in D, D is the union of Δ and -Δ, and contains no other closed half-line. We shall say that a half-line is isotropic if the line containing it is isotropic.

Recall also (ibid.) that, given a vector space of finite dimension n over A, an orientation on F is the choice of one of the two half-lines of the space \(\bigwedge^{n}F\) the n-vectors belonging to this half-line are said to be positive. A finite-dimensional vector space equipped with an orientation is said to be oriented.

The homotheties of E whose ratio is \textgreater0 obviously form a subgroup of index 2 in H; we will denote it by H \(^{+}\) . The canonical homomorphism i : O \(^{+}\) → S \(^{+}\) /H (cf.~prop. 2) is the composite of the canonical homomorphisms of S \(^{+}\) /H \(^{+}\) on S \(^{+}\) /H and of O \(^{+}\) in S \(^{+}\) /H \(^{+}\) ; since O \(^{+}\) ∩ H \(^{+}\) = \{1\} this latter homomorphism is injective.

PROPOSITION 7. --- Suppose that A is a maximal ordered field and that \(\Phi\) be a positive form (§ 7). Then the canonical homomorphisms of O \(^{+}\) into S \(^{+}\) /H \(^{+}\) and of O \(^{+}\) /\{1, -1\} into S \(^{+}\) /H are bijective, and S \(^{+}\) is isomorphic to O \(^{+}\) × H \(^{+}\) .

We have already seen that the homomorphisms in question are injective, and it suffices to show that the first one is surjective. Let \((e_{1}, e_{2})\) an orthonormal basis of E, and let w be the direct similarity such that \(w(e_{1}) = e_{2}\) (cor. 1 of prop. 1, no. 1); we then have \(w^{2} = -1\) (prop. 1, b)). Given any direct similarity transformation \(u = a + bw\) ( \(a \in \mathbf{A}, b \in \mathbf{A}\) ), we have \(\mathrm{N}(u) = a^2 + b^2 > 0\) and there exists one and only one rotation contained in the same half-line of \(\mathbf{A}(\Phi)\) as \(u\) , namely \((a^2 + b^2)^{1/2} u\) . QED.

COROLLARY. --- Given two half-lines D, D' with origin 0, there exists one and only one rotation v such that v(D) = D'.

The hypotheses indeed imply that E contains no isotropic lines. Our assertion then follows from cor. 1 of prop. 1 (no. 1).

We shall henceforth assume that A is a maximal ordered field, and that the form \(\Phi\) is positive. In the set of pairs \((\mathrm{D}_1, \mathrm{D}_2)\) of lines (resp. half-lines with origin 0) of E, the relation “there exists a direct similarity (resp. a rotation) u such that \(u(\mathrm{D}_1) = \mathrm{D}_1'\) and \(u(\mathrm{D}_2) = \mathrm{D}_2'\) ” is an equivalence relation between pairs \((\mathrm{D}_1, \mathrm{D}_2)\) and \((\mathrm{D}_1', \mathrm{D}_2')\) The equivalence class of the pair \((\mathrm{D}_1, \mathrm{D}_2)\) is called, by definition, the angle of the lines (resp. half-lines) \(\mathrm{D}_1, \mathrm{D}_2\) (taken in this order); we denote it \((\widehat{\mathrm{D}_1, \mathrm{D}_2})\) .

PROPOSITION 8. --- Suppose that A is a maximal ordered field, and that the form \(\Phi\) is positive. Let \(D_{1}, D_{2}, D_{1}^{\prime}, D_{2}^{\prime}\) be four lines (resp. half-lines) with origin 0 in E. For the angles \((\widehat{D_{1}, D_{2}})\) and \((\widehat{D_{1}^{\prime}, D_{2}^{\prime}})\) to be equal, it is necessary and sufficient that the angles \((\widehat{D_{1}, D_{1}^{\prime}})\) and \((\widehat{D_{2}, D_{2}^{\prime}})\) be equal.

Let us prove the necessity of the stated condition. Let \(u\) be a direct similarity (resp. a rotation) such that \(u(\mathrm{D}_{1}) = \mathrm{D}_{1}^{\prime}\) and \(u(\mathrm{D}_{2}) = \mathrm{D}_{2}^{\prime}\) . By cor. 3 of prop. 1, there exists a direct similarity (resp. by the corollary of prop. 7, a rotation) \(\nu\) such that \(\nu(\mathrm{D}_{1}) = \mathrm{D}_{2}\) . Since the group \(\mathrm{S}^{+}\) (resp. \(\mathrm{O}^{+}\) ) is commutative, we have \(\mathrm{D}_{2}^{\prime} = u(\nu(\mathrm{D}_{1})) = \nu(u(\mathrm{D}_{1})) = \nu(\mathrm{D}_{1}^{\prime})\) , and this shows that \((\widehat{\mathrm{D}_{1}, \mathrm{D}_{1}^{\prime})} = (\widehat{\mathrm{D}_{2}, \mathrm{D}_{2}^{\prime}})\) . Sufficiency follows from necessity by interchanging \(\mathrm{D}_{2}\) and \(\mathrm{D}_{1}^{\prime}\) .

It follows from prop. 8 that, to every angle \((\widehat{\mathrm{D}_1,\mathrm{D}_2})\) of lines (resp. half-lines) with origin 0 in E, there is canonically associated a well-determined element of \(\mathrm{S}^{+}/\mathrm{H}\) (resp. \(\mathrm{O}^{+}\) ), namely the class mod. H of the direct similarities \(\nu\) (resp. the rotation \(\nu\) ) such that \(u(D_1) = D_2\) for any representative ( \(D_1, D_2\) ) of the angle ( \(\widehat{D_1, D_2}\) ). We have thus defined a canonical bijection \(h\) (resp. \(h'\) ) from the set \(\mathfrak{A}_0\) of angles of lines (resp. \(\mathfrak{A}\) of angles of half-lines) onto \(S^+/H\) (resp. \(O^+\) ); in particular, for every \(\varphi \in \mathfrak{A}\) , one says that \(h(\varphi)\) is the rotation with angle \(\varphi\) . We shall transport to \(\mathfrak{A}_0\) and \(\mathfrak{A}\) , by means of \(h^{-1}\) and of \(h'^{-1}\) , the structures of commutative groups of \(S^+/H\) and of \(O^+\) , and we shall denote additively the groups \(\mathfrak{A}_0\) and \(\mathfrak{A}\) thus obtained. If one denotes by D, \(D', D''\) lines (resp. half-lines) with origin 0 in E, one has by definition

\[
\widehat {(\mathbf {D} , \mathbf {D} ^ {\prime \prime})} = \widehat {(\mathbf {D} , \mathbf {D} ^ {\prime})} + \widehat {(\mathbf {D} ^ {\prime} , \mathbf {D} ^ {\prime \prime})}\tag{9}
\]

(Chasles relation);

one deduces from this

\[
\widehat {(\mathrm{D} , \mathrm{D})} = 0, \quad \widehat {(\mathrm{D} , \mathrm{D} ^ {\prime})} = - (\widehat {\mathrm{D} ^ {\prime} , \mathrm{D}}).\tag{10}
\]

Remarks. --- 1) The set L of lines (resp. half-lines) with origin 0 in E is a homogeneous space of the abelian group S+/H (resp. O+) such that the identity element is the only operator leaving all elements of L invariant. One can therefore apply to L the formulas of chap.~II, 2 \(^{e}\) ed., App. II, n° 1; prop. 8 is thus a particular case of the "parallelogram rule", and formulas (9) and (10) are particular cases of formulas (2) (ibid.).

\begin{enumerate}
\def\labelenumi{\arabic{enumi})}
\setcounter{enumi}{1}
\tightlist
\item
  In the definition of the group of angles of lines, one can, instead of the group S+/H, use the group O+/\{−1, 1\} which is canonically isomorphic to it (prop. 7). The canonical homomorphism of O+ onto O+/\{−1, 1\} thus corresponds to a homomorphism of A onto A0, namely the one which, to the angle of the two half-lines Δ, Δ', associates the angle of the lines D, D' containing respectively Δ, Δ'. *In the case where the field A is the field of real numbers, the group A is thus a double covering of the group A0.*
\end{enumerate}

By prop. 8, all angles of lines (resp. half-lines) of the form \((\widehat{\mathbf{D}^{\prime},\mathbf{D}^{\prime\prime}})\) where \(\mathbf{D}'\) and \(\mathbf{D}''\) are orthogonal (resp.

of the form \((\widehat{\mathrm{D}}, -\mathrm{D})\) are equal: this indeed follows from Remark 2 of No. 1 (resp. is obvious). This angle of lines (resp. of half-lines) is called the right angle (resp. the straight angle); it is an element of order 2 of \(A_{0}\) (resp. of A).

PROPOSITION 9. --- Suppose that the field A is maximally ordered, and that \(\Phi\) is a positive form. For every integer \(n>1\) , the number of elements \(\theta\) of the group \(\mathfrak{A}_{0}\) of angles of lines (resp. \(\mathfrak{A}\) of angles of half-lines) such that \(n\theta=0\) is equal to \(n\) .

As \(\mathfrak{A}_{0}\) , \(\mathrm{S}^{+}/\mathrm{H}\) , \(\mathfrak{A}\) and \(\mathrm{O}^{+}\) are isomorphic (prop. 3), it suffices to carry out the proof for \(\mathrm{O}^{+}\) , that is, to show that there are exactly \(n\) rotations \(\varrho\) such that \(\varphi^{n}=1\) Since A is a maximal ordered field and that \(\mathrm{A}(\Phi)\) is a degree 2 extension field of A (prop. 1 a)), \(\mathrm{A}(\Phi)\) is an algebraically closed field (chap.~VI, § 2, no. 6, th. 3). Hence the roots \(n\) -th roots of unity in \(\mathrm{A}(\Phi)\) form a cyclic group of order \(n\) (chap. V, § 11, no. 1, th. 1). Since we have \(\mathrm{N}(u)=u\overline{u}\geqslant0\) for all \(u\in\mathrm{A}(\Phi)\) , the relation \(u^{n}=1\) implies that one has \(\mathrm{N}(u)=1\) , hence that \(u\) is a rotation (No. 1, Remark 1). This proves our assertion.

COROLLARY. --- The right angle (resp. straight angle) is the only element of order 2 in the group. \(\mathfrak{A}_{0}\) (resp. \(\mathfrak{A}\) ).

We will assume finally that the plane E is oriented.

Lemma 1. --- Let u be a direct similarity of E; all bivectors of the form \(x \wedge u(x)\) belong to the same closed half-line of \(\bigwedge^{2}\mathbf{E}\) .

The case where \(u\) is a homothety is trivial. Otherwise, we have \(x \wedge u(x) \neq 0\) for all \(x \neq 0\) ; let \(x, y\) be two vectors of \(E\) ( \(x \neq 0, y \neq 0\) ); there exists \(v \in S^+\) such that \(y = v(x)\) , whence \(y \wedge u(y) = v(x) \wedge uv(x) = v(x) \wedge v(u(x)) = (\det v)(x \wedge u(x))\) ; by taking an orthonormal basis of \(E\) , we see that det \(v\) is positive (prop. 1 b)); hence our assertion.

That being said, among the two generators \(w\) of \(\mathrm{A}(\Phi)\) such that \(w^{2} = -1\) there exists one and only one such that the bivector \(x \wedge w(x)\) be positive for all \(x\in\mathbf{E}\) . It is this generator that we will choose to define the functions \(c_{w}, s_{w}\) and \(t_{w}\) (no. 2). Let \(h\) and \(h'\) the canonical bijections defined above from the group \(\mathfrak{A}_{0}\) angles of lines on S+/H and of the group \(\mathfrak{A}\) of half-line angles on O+. The composed mappings \(t_{w}\circ h\) of \(\mathfrak{A}_{0}\) in the projective field \(\tilde{\mathbf{A}}\) , \(c_{w}\circ h'\) and \(s_{w}\circ h'\) of \(\mathfrak{A}\) in the field A are denoted respectively by tg, cos, and sin, and are called the tangent, cosine, and sine functions. The mapping \(\varphi\to1/\mathrm{tg}\varphi\) of \(\mathfrak{A}_{0}\) in \(\tilde{\mathbf{A}}\) is denoted cotg and is called the cotangent function. The functions cosine, sine, tangent, and cotangent are said to be the trigonometric functions. The composite mappings tg∘p and cotg∘p, where p denotes the canonical homomorphism from \(\mathfrak{A}\) on \(\mathfrak{A}_{0}\) (Remark 2 above) are still denoted tg and cotg by abuse of language.

Formulas (2), (8), (3), (4), (5), and (7) of §2 give, since here we have \(\delta = -1\)

\begin{enumerate}
\def\labelenumi{(\arabic{enumi})}
\setcounter{enumi}{10}
\tightlist
\item
\end{enumerate}

\[
\operatorname{tg} \left(\varphi + \varphi^ {\prime}\right) = (\operatorname{tg} \varphi + \operatorname{tg} \varphi^ {\prime}) / (1 - \operatorname{tg} \varphi \operatorname{tg} \varphi^ {\prime})\tag{12}
\]

\[
\operatorname{tg} (2 \varphi) = 2 \operatorname{tg} \varphi / (1 - \operatorname{tg} ^ {2} \varphi)
\]

for \(\varphi, \varphi'\) in \(\mathfrak{A}_0\) ;

\begin{enumerate}
\def\labelenumi{(\arabic{enumi})}
\setcounter{enumi}{12}
\tightlist
\item
\end{enumerate}

\[
\cos^ {2} \theta + \sin^ {2} \theta = 1\tag{14}
\]

\[
\cos (\theta + \theta^ {\prime}) = \cos \theta \cos \theta^ {\prime} - \sin \theta \sin \theta^ {\prime}\tag{15}
\]

\[
\sin (\theta + \theta^ {\prime}) = \sin \theta \cos \theta^ {\prime} + \cos \theta \sin \theta^ {\prime}\tag{16}
\]

\[
\sin (2 \theta) = 2 \operatorname{tg} \theta / (1 + \operatorname{tg} ^ {2} \theta),
\]

\[
\cos (2 \theta) = (1 - \mathrm{tg} ^ {2} \theta) / (1 + \mathrm{tg} ^ {2} \theta)
\]

for \(\theta, \theta'\) in \(\mathfrak{A}\) On the other hand, we have, by definition or as an easy consequence of the preceding formulas:

\begin{enumerate}
\def\labelenumi{(\arabic{enumi})}
\setcounter{enumi}{16}
\tightlist
\item
\end{enumerate}

\[
\operatorname{tg} \theta = \sin \theta / \cos \theta , \quad \cot \mathrm{g} \theta = \cos \theta / \sin \theta\tag{18}
\]

\[
1 + \mathrm{tg} ^ {2} \theta = 1 / \cos^ {2} \theta , \quad 1 + \cot \mathrm{g} ^ {2} \theta = 1 / \sin^ {2} \theta
\]

for θ ∈ 𝒜.

Given two nonzero vectors \(x, y\) of E, the angle between these two vectors (taken in this order) is called, and denoted \((\widehat{x}, y)\) the angle of the half-lines to which they belong. For every vector \(x\) from E we call the length of \(x\) is called, and denoted \(|x|\) , the element \(\Phi(x, x)^{1/2}\) of A.

PROPOSITION 10. --- Suppose that the field A is maximally ordered, that the plane E is oriented, and that the form \(\Phi\) is positive. For every pair of nonzero vectors \(x\) , \(y\) of E on \(a\)

\begin{enumerate}
\def\labelenumi{(\arabic{enumi})}
\setcounter{enumi}{18}
\tightlist
\item
\end{enumerate}

\[
\cos \widehat {(x , y)} = \Phi (x, y) / | x |. | y |\tag{20}
\]

\[
\widehat {\sin (x , y)}. e = (x \wedge y) / | x |. | y |,
\]

where e denotes the positive bivector such that \(\Phi_{(2)}(e, e) = 1\) .

Indeed, since the vectors \(x' = x / |x|\) and \(y' = y / |y|\) are both of length 1, there exists a rotation \(\nu\) and exactly one such that \(\nu(x') = y'\) (No. 1, cor. 2 of prop. 1). If we set \(\nu = a + bw\) ( \(a, b\) In A), we have by definition \(a = \cos(\widehat{x,y})\) and \(b = \sin(\widehat{x,y})\) . The relation \(y' = \nu(x') = ax' + bw(x')\) gives \(\Phi(x', y') = a\Phi(x', x') = a\) since \(x'\) and \(w(x')\) are orthogonal (Remark 2 of No. 1); this proves (19). On the other hand, this relation also gives \(x' \wedge y' = bx' \wedge w(x') = b.e\) by the definition of the extension \(\Phi_{(2)}\) of \(\Phi\) à \(\bigwedge^{2} E\) (§ 1, no. 9, formula (37)) and the choice of \(w\) ; this proves (20).

Remarks. --- 3) Given two lines (resp. rays) D, D' of an affine plane L attached to E, one calls the angle of D and D', denoted (D, D'), the angle made by their directions in E (resp. the rays with origin 0 of E corresponding to D and D') (chap.~II, 2). \(^{e}\) ed., App. II, nos. 1 and 3).

\begin{enumerate}
\def\labelenumi{\arabic{enumi})}
\setcounter{enumi}{3}
\tightlist
\item
  Let F be a vector space of arbitrary dimension over the maximal ordered field A, and \(\Psi\) a positive non-degenerate symmetric bilinear form on F. Given two vectors \(x\) , \(y\) linearly independent elements of F, let \(\mathbf{F}'\) the vector plane they span; we call the angle of \(x\) and \(y\) the angle of \(x\) and \(y\) considered as elements of the plane \(\mathbf{F}'\) ; it is denoted \((x,y)\) The cosine of this angle is, by virtue of (19), given by
\end{enumerate}

\[
\cos \widehat {(x , y)} = \Psi (x, y) / | x |. | y |\tag{21}
\]

\((\text{ou } |x| = \Psi(x, x)^{1/2}\) is also called the length of the vector \(x)\) and is therefore independent of the chosen orientation on \(F'\) ; the sine and the tangent of \((x, y)\) change sign if one changes the orientation of \(F'\) Given two nonzero proportional vectors \(x, y\) of \(F\) , set \((x, y) = 0\) by convention.
% semantic-fragments: legacy-input-seam
\relax{}
\subsection{Angular sectors.}

We shall first assume, without any further hypothesis, that E is an oriented plane over the ordered field A. We shall say that three half-lines \(D_{0}\) , \(D_{1}\) , \(D_{2}\) (with origin 0) of E form a direct sequence if, for \(x_{i} \in D_{i}\) , \(x_{i} \neq 0\) (i = 0, 1, 2), at least two of the bivectors \(x_{0} \wedge x_{1}\) , \(x_{1} \wedge x_{2}\) , \(x_{2} \wedge x_{0}\) are \textgreater0; in this case the sequences \(D_{1}\) , \(D_{2}\) , \(D_{0}\) and \(D_{2}\) , \(D_{0}\) , \(D_{1}\) are also direct. It is clear that three half-lines forming a direct sequence are distinct. Given two half-lines \(D_{1}\) , \(D_{2}\) of E, an open (resp. closed) angular sector with origin \(D_{1}\) and endpoint \(D_{2}\) is the set (or, by abuse of language, the union) of half-lines D such that the sequence \(D_{1}\) , D, \(D_{2}\) is direct (resp. such that \(D = D_{1}\) or \(D = D_{2}\) or that the sequence \(D_{1}\) , D, \(D_{2}\) is direct).

PROPOSITION 11. --- Let E be an oriented plane over an ordered field A, D₀ a ray of E, and G the set of rays of E distinct from D₀. The relation

« \(D_{1} = D_{2}\) , or the sequence \(D_{0}\) , \(D_{1}\) , \(D_{2}\) is direct »

between elements \(D_{1}\) , \(D_{2}\) of G is a total order relation in G.

Indeed the axioms of total order relations are trivially verified, except for transitivity. Let \(D_{1}\) , \(D_{2}\) , \(D_{3}\) be three half-lines such that the sequences, \(D_{0}\) , \(D_{1}\) , \(D_{2}\) and \(D_{0}\) , \(D_{2}\) , \(D_{3}\) are direct; we shall show that the sequence \(D_{0}\) , \(D_{1}\) , \(D_{3}\) is direct. To this end, take a vector \(x_{i} \neq 0\) in \(D_{i} (i = 0, 1, 2, 3)\) , choose a bivector e \textgreater{} 0 and set \(x_{i} \wedge x_{j} = a_{ij} e (a_{ij} \in A)\) . Writing \(e = x_{0} \wedge y (y \in E)\) , and taking \((x_{0}, y)\) as a basis of E, one readily verifies the relation

\[
a _ {0 1} a _ {2 3} + a _ {0 2} a _ {3 1} + a _ {0 3} a _ {1 2} = 0.\tag{22}
\]

This being so, if \(a_{01} \leqslant 0\) , one has \(a_{12} > 0\) and \(a_{20} > 0\) (since the first sequence is direct), then \(a_{23} > 0\) and \(a_{30} > 0\) (since \(a_{02} < 0\) and since the second sequence is direct), whence \(a_{13} > 0\) (by virtue of (22)); therefore the sequence \((\mathrm{D}_0, \mathrm{D}_1, \mathrm{D}_3)\) is direct in this case. Suppose now that \(a_{01} > 0\) . If \(a_{30} \leqslant 0\) , one has \(a_{02} > 0\) and \(a_{23} > 0\) (since the second sequence is direct), then \(a_{12} > 0\) (since \(a_{20} < 0\) and since the first sequence is direct), hence \(a_{13}>0\) (by (22)), and the sequence \((\mathrm{D}_{0}, \mathrm{D}_{1}, \mathrm{D}_{3})\) is direct. Finally, it is evidently the same if \(a_{01}>0\) and \(a_{30}>0\) . QED.

COROLLARY. --- Let \(D_{1}\) and \(D_{2}\) be two distinct rays of E. For every ray \(D_{0}\) of E such that the sequence \(D_{0}\) , \(D_{1}\) , \(D_{2}\) is direct, the set of half-lines D of E such that \(D_{1} < D < D_{2}\) (for the total order relation defined by \(D_{0}\) ) is equal to the open angular sector with origin \(D_{1}\) and endpoint \(D_{2}\) .

Indeed, given a half-line \(D_{3}\) we must show that the relations “the sequence \(D_{1}\) , \(D_{3}\) , \(D_{2}\) is direct” and “the sequences \(D_{0}\) , \(D_{1}\) , \(D_{3}\) and \(D_{0}\) , \(D_{3}\) , \(D_{2}\) are direct” are equivalent. To abbreviate, let us denote by (ijk) the relation “the sequence ( \(D_{i}\) , \(D_{j}\) , \(D_{k}\) ) is direct”. By prop. 11, the conjunction of (132) and (120) implies (130); likewise the conjunction of (201) and (213) implies (203); whence one half of our assertion. Conversely, suppose (012), (013), and (032); since the conjunction of (312) and (320) implies (310) (prop. 11), and since (310) and (013) are incompatible, (312) is false; this proves (132) and completes the proof.

By virtue of the preceding corollary, the open (resp. closed) angular sector with origin \(D_{1}\) and endpoint \(D_{2}\) is denoted \(\left[\begin{matrix}\mathrm{D}_{1},\mathrm{D}_{2}\end{matrix}\right]\) (resp. \(\left[\begin{matrix}\mathrm{D}_{1},\mathrm{D}_{2}\end{matrix}\right]\) ), it being understood that these are intervals for the order structure defined by any half-line \(D_{0}\) such that the sequence \(D_{0},D_{1},D_{2}\) is direct.

PROPOSITION 12. — Let A be a maximal ordered field, E an oriented plane over A, D₀ a ray of E, and G the set of rays of E distinct from D₀. The totally ordered sets A and G (prop. 11) are isomorphic.

Let, in fact, \((x, y)\) be a basis of E such that \(x \in -\mathrm{D}_{0}\) and the bivector \(x \wedge y\) be \(> 0\) . To each element \(t\) of A we associate the half-line \(f(t)\) to which the vector \((1 - t^{2})x + 2ty\) belongs. It is clear that \(f(\mathrm{A}) \subset \mathrm{G}\) . We show that \(f\) is strictly increasing. Indeed, for the sequence \(\mathrm{D}_{0}\) , \(f(t)\) , \(f(t')\) ( \(t, t'\) in A) to be direct, it is necessary and sufficient, by definition, that at least two of the elements

\[
- 2 t, \quad (1 - t ^ {2}) 2 t ^ {\prime} - (1 - t ^ {\prime 2}) 2 t, \quad 2 t ^ {\prime}
\]

be \textgreater0. Now the second is equal to \(2(t' - t)(1 + tt')\) . Therefore, if \(t < t'\) , one has either \(tt' \geqslant 0\) , hence \(t' > 0\) or -t \textgreater{} 0, or \(tt' < 0\) , hence -t \textgreater{} 0 and \(t' > 0\) ; in any case \(D_{0}\) , \(f(t)\) , \(f(t')\) is direct. Since A is totally ordered, f is an isomorphism of A onto the ordered set \(f(A)\) (Set Theory, Chap. III, § 1, no. 14, Prop. 13).

It remains to show that \(f\) is surjective. For this, consider the positive form \(\Phi\) on \(E\) such that \((x, y)\) be an orthonormal basis for \(\Phi\) . For every half-line \(D \in G\) , there exists one and only one angle \(\varphi\) such that \(2\varphi = (-D_0, D)\) ( \(n^o 1\) , Prop. 3); since \((-D_0, D)\) is not the straight angle, \(\varphi\) is not the right angle, and \(tg\varphi\) is therefore finite. Then, by virtue of formulas (16) ( \(n^o 3\) ), we have \(D = f(tg\varphi)\) . This completes the proof.

Exercises. --- 1) With the notations of no. 1, set Q(x) = Φ(x, x); the vector space E is canonically identified with C⁻(Q) (§ 9, no. 1). For every z ∈ C⁺(Q) and every x ∈ E, one has zx ∈ E; show that x ↦ zx is an element of A(Φ), and that z ↦ sz is an isomorphism from C⁺(Q) onto the algebra A(Φ).

\begin{enumerate}
\def\labelenumi{\arabic{enumi})}
\setcounter{enumi}{1}
\tightlist
\item
  The hypotheses and notation are those of no. 1 and of Exercise 1.
\end{enumerate}

\begin{enumerate}
\def\labelenumi{\alph{enumi})}
\item
  Let C be the set of \(x \in E\) such that \(\Phi(x, x) = 1\) (“unit circle”), and let \(\mathfrak{D}\) be the set of lines D whose intersection with C is nonempty (and consequently consists of two opposite elements of E). One calls a pointed line any pair \(\Delta = (D, z)\) formed by a line \(D \in \mathfrak{D}\) and one of the points \(z \in D \cap C\) . Show that if \(\Delta_1 = (D_1, z_1)\) , \(\Delta_2 = (D_2, z_2)\) are two pointed lines, there exists one and only one rotation \(u\) such that \(u(z_1) = z_2\) (and consequently \(u(D_1) = D_2\) ), which one expresses by writing \(u(\Delta_1) = \Delta_2\) . In the set of pairs \((\Delta_1, \Delta_2)\) of pointed lines, the relation “there exists a rotation \(u\) such that \(u(\Delta_1) = \Delta_1'\) and \(u(\Delta_2) = \Delta_2'\) ” is an equivalence relation. The set \(\mathfrak{A}_1\) of equivalence classes of pointed lines under this relation is called the set of angles of pointed lines, and the equivalence class to which a pair \((\Delta_1, \Delta_2)\) of pointed lines belongs is called the angle of this pair and is denoted \((\widehat{\Delta_1, \Delta_2})\) ; the relation \((\widehat{\Delta_1, \Delta_2}) = (\widehat{\Delta_1'}, \widehat{\Delta_2'})\) is equivalent to \((\widehat{\Delta_1, \Delta_1'}) = (\widehat{\Delta_2, \Delta_2'})\) and the rotation that maps \(\Delta_1\) to \(\Delta_2\) is called the rotation of angle \(\theta = (\Delta_1, \Delta_2)\) and denoted \(h_1(\theta)\) ; \(h_1\) is a bijection from \(\mathfrak{A}_1\) on \(O^+\) and one transports to \(\mathfrak{A}_1\) by means of \(h_1^{-1}\) the commutative group structure of \(O^+\) , writing the group \(\mathfrak{A}_1\) thus defined additively; one also calls a straight angle in \(\mathfrak{A}_1\) the angle of the pair formed by a pointed line \((D, z)\) and of the "opposite" pointed line ( \(D, -z\) ), which corresponds to the rotation \(x \to -x\) .
\item
  In the set \(\mathfrak{D}_{0} \supset \mathfrak{D}\) of non-isotropic lines, the notion of angle between lines is defined as in no. 3, the group \(\mathfrak{A}_{0}\) (using cor. 3 of prop. 1 of no. 1), the right angle in \(\mathfrak{A}_{0}\) and the canonical bijection \(h\) of \(\mathfrak{A}_{0}\) onto \(S^{+}/H\) . With the notations of the cor. of prop. 3, we set \(\bar{d} = h_{1}^{-1} \circ d \circ h\) and \(\bar{i} = h^{-1} \circ i \circ h_{1}\) so that \(\bar{i}\) is a homomorphism of \(\mathfrak{A}_{1}\) into \(\mathfrak{A}_{0}\) and \(\bar{d}\) a homomorphism of \(\mathfrak{A}_{0}\) into \(\mathfrak{A}_{1}\) ; \(\bar{d}\) is bijective, and the kernel of \(\bar{i}\) consists of 0 and the straight angle; moreover, we have \(\bar{d}(\bar{i}(\theta)) = 2\theta\) for \(\theta \in \mathfrak{A}_{1}\) and \(\bar{i}(\bar{d}(\varphi)) = 2\varphi\) for \(\varphi \in \mathfrak{A}_{0}\) . For \(\bar{i}\) is surjective (in other words, for \(\mathfrak{D}\) to be the set of all non-isotropic lines), it is necessary and sufficient that the field A be Pythagorean (chap. VI, § 2, exerc. 8 d)) and that there exist an orthonormal basis for \(\Phi\) ; this is equivalent to saying that \(\Phi(x, x)\) is a square for every \(x \in E\) .
\end{enumerate}

\begin{enumerate}
\def\labelenumi{\arabic{enumi})}
\setcounter{enumi}{2}
\tightlist
\item
  The hypotheses and notations are those of exerc. 2.
\end{enumerate}

\begin{enumerate}
\def\labelenumi{\alph{enumi})}
\tightlist
\item
  Let \(s\) be a symmetry with respect to a non-isotropic line D ( \(\S 6\) , no. 4). For every pointed line \(\Delta_{1} = (D_{1}, z_{1})\) , let
\end{enumerate}

\[
\Delta_ {2} = s (\Delta_ {1}) = (s (\mathrm{D} _ {1}), s (z _ {1})),
\]

and let \(\varphi = (\widehat{\mathrm{D}_1, \mathrm{D}})\) ; show that one has \((\widehat{\Delta_1, \Delta_2}) = \bar{d}(\varphi)\) .

\begin{enumerate}
\def\labelenumi{\alph{enumi})}
\setcounter{enumi}{1}
\item
  Show that every orthogonal transformation with determinant -1 is a symmetry \(s\) with respect to a non-isotropic line (cf. § 6, exerc. 15 e)); for every rotation \(u\) , one has \(sus^{-1} = u^{-1}\) .
\item
  If \(x, y\) are any two points of E such that \(\Phi(x, x) = \Phi(y, y) \neq 0\) , there exists one and only one symmetry with respect to a non-isotropic line, which transforms \(x\) to \(y\) .
\item
  Let \(s, s'\) be the symmetries with respect to two non-isotropic lines D, D', and let \(\varphi = (\widehat{\mathrm{D},\mathrm{D}'})\) ; for \(s's\) to be a rotation of angle \(\theta\) , it is necessary and sufficient that \(\bar{d}(\varphi) = \theta\) .
\item
  Show that the commutator subgroup of the orthogonal group \(\mathbf{O}(\mathrm{Q})\) is the image of \(\mathfrak{A}_{1}\) under the homomorphism \(\theta \to h_{1}(2\theta)\) ( \(\S 6\) , exerc. \(17a\) ); for this image to be equal to \(\mathrm{O}^{+}\) , it is necessary and sufficient that the homomorphism \(\bar{i}\) be surjective (exerc. \(2b\) ).
\end{enumerate}

\begin{enumerate}
\def\labelenumi{\arabic{enumi})}
\setcounter{enumi}{3}
\item
  The hypotheses and notation are those of exercise 2. Let \(a\) , \(b\) be two points of the circle C, \(\Delta_{a}\) , \(\Delta_{b}\) the pointed lines passing through \(a\) and \(b\) (and by 0) respectively, and let \(\theta = (\widehat{\Delta_{a}, \Delta_{b}})\) . For every \(x \in E\) , distinct from \(a\) and \(b\) , let D \(_{xa}\) (resp. D \(_{xb}\) ) be the affine line passing through \(a\) and \(x\) (resp. \(b\) and \(x\) ), and let D \(_{xa}'\) (resp. D \(_{xb}'\) ) be the direction of D \(_{xa}\) (resp. D \(_{xb}\) ); show that for \(x \in C\) ( \(x\) distinct from \(a\) and \(b\) ), it is necessary and sufficient that the angle \(\varphi = (\widehat{\mathrm{D}_{xa}', \mathrm{D}_{xb}'})\) be such that \(\bar{d}(\varphi) = \theta\) (using exercise 3). How does this result change when \(x = a\) or \(x = b\) (cf. § 6, exerc. 25) \(a\) )) ?
\item
  Using the notation of nos. 1 and 2 and of Exercise 2, we assume that \(\Phi\) has index 1, in other words that \(\delta\) is a square \(\gamma^{2}\) in A; we denote by D \(_{1}\) , D \(_{2}\) the isotropic lines of E, containing respectively the vectors \(e_{1}-\frac{1}{\gamma}e_{2}\) and \(e_{1}+\frac{1}{\gamma}e_{2}\) .
\end{enumerate}

\begin{enumerate}
\def\labelenumi{\alph{enumi})}
\item
  Show that the group of rotations \(O^{+}\) is isomorphic to the multiplicative group A* of the field A.
\item
  For every angle \(\theta \in \mathfrak{A}_1\) , set \(e_w(\theta) = c_w(h_1(\theta)) + \gamma s_w(h_1(\theta))\) . Show that \(\theta \to e_w(\theta)\) is an isomorphism of \(\mathfrak{A}_1\) onto the group \(A^*\) .
\item
  Let D and D′ be any two lines, and let \(\varphi = (\widehat{\mathbf{D},\mathbf{D}^{\prime}})\) . Show that the cross-ratio \(\left[ \begin{array}{cc}\mathrm{D}_{1} & \mathrm{D}_{2}\\ \mathrm{D}^{\prime} & \mathrm{D} \end{array} \right]\) (chap. II, 2nd ed., App. III, exerc. 5) is equal to \(e_w(\overline{d} (\varphi))\) ("Laguerre formula"). (Note that D1 and D2 are invariant under any direct similarity, and using exercise 4 c) of chap. II, 2nd ed., App. III, reduce to the case where D = Ae1.)
\end{enumerate}

\begin{enumerate}
\def\labelenumi{\arabic{enumi})}
\setcounter{enumi}{5}
\tightlist
\item
  We assume that A is an ordered field.
\end{enumerate}

\begin{enumerate}
\def\labelenumi{\alph{enumi})}
\item
  Let \(D_{1}\) , \(D_{2}\) be two non-isotropic rays with origin 0. Show that there exists a direct similarity u such that \(u(\mathrm{D}_{1}) = \mathrm{D}_{2}\) ; any other direct similarity having this property is of the form \(\varrho = su\) where s is a homothety with ratio \textgreater{} 0.
\item
  In the set of pairs \((\mathrm{D}_{1}, \mathrm{D}_{2})\) of half-lines that are not isotropic, the relation “there exists a direct similarity u such that \(u(\mathrm{D}_{1}) = \mathrm{D}_{1}'\) and \(u(\mathrm{D}_{2}) = \mathrm{D}_{2}'\) ” is an equivalence relation. The set A of equivalence classes of non-isotropic rays, under this relation, is called the set of angles of (non-isotropic) rays, and the equivalence class of a pair \((\mathrm{D}_{1}, \mathrm{D}_{2})\) of such half-lines is called the angle of this pair and denoted \((\widehat{\mathrm{D}_{1}, \mathrm{D}_{2}})\) ; if \(\theta = (\widehat{\mathrm{D}_{1}, \mathrm{D}_{2}})\) , one says that \(\theta\) is the angle of every direct similarity mapping \(D_{1}\) to \(D_{2}\) ; let \(h_{2}(\theta)\) be the class mod. \(H^{+}\) of these similarities, so that \(h_{2}\) is a bijection from the set A onto the group \(S^{+}/H^{+}\) ; one transports to A by means of \(h_{2}^{-1}\) the commutative group structure of \(S^{+}/H^{+}\) , writing the group A thus defined additively. Define a canonical injection \(\bar{j}\) of the group \(A_{1}\) of angles of pointed lines (exerc. 2) into the group A, such that \(h_{2} \circ \bar{j} \circ h_{2}^{-1}\) is the canonical injection j of \(O^{+}\) into \(S^{+}/H^{+}\) . For \(\bar{j}\) to be surjective, it is necessary and sufficient that the homomorphism \(\bar{i}\) of \(A_{1}\) into \(A_{0}\) (exerc. 2 b)) be surjective.
\item
  Show that in \(\mathfrak{A}\) the equation \(2\theta = 0\) has 2 solutions if \(\delta < 0\) and 4 solutions if \(\delta > 0\) . In the first case, the solution \(\varpi \neq 0\) of this equation is also called the straight angle.
\end{enumerate}

¶ 7) Under the hypotheses and notation of exercise 6, suppose the homomorphism \(\bar{j}\) bijective; we then define cos θ and sin θ for every θ ∈ A as in no. 3. Let T be the set of θ ∈ A such that sin θ ≥ 0.

\begin{enumerate}
\def\labelenumi{\alph{enumi})}
\item
  Show that for every \(\theta \in T\) , there exists one and only one angle \(\theta' \in T\) such that \(2\theta' = \theta\) ; we set \(\theta' = \theta/2\) .
\item
  Let L be the Z-module of formal linear combinations of the elements of T with coefficients in Z (chap. II, § 1, no. 8); we denote by \(\xi + \eta\) and \(\dot{\cdot} \xi\) the sum and the opposite in L. Let N be the submodule of L generated by the elements of L of the form \(\xi + \eta \dot{\cdot} (\xi + \eta)\) for all pairs \((\xi, \eta)\) of elements of T such that \(\xi + \eta \in T\) (sum taken in the group \(\mathfrak{A}\) ). Let \(f\) be the homomorphism from L to \(\mathfrak{A}\) which extends the canonical injection of T into \(\mathfrak{A}\) , and \(g\) be the endomorphism of L which extends the map \(\theta \to \theta/2\) from T into itself. We have \(f(N) = \{0\}\) and \(g(N) \subset N\) ; by passing to quotients, we deduce from \(f\) a homomorphism \(f\) from M = L/N into \(\mathfrak{A}\) , and from \(g\) an endomorphism \(g\) of M; we set \(g(\mu) = \mu/2\) and if \(g^m\)
\end{enumerate}

is the m-th iterate of g, \(g^{m}(\mu) = 2^{-m}\mu\) ; we have \(2^{m}(2^{-m}\mu) = \mu\) for all \(\mu \in M\) .

\begin{enumerate}
\def\labelenumi{\alph{enumi})}
\setcounter{enumi}{2}
\item
  Show that the restriction to T of the canonical map \(\psi\) from L onto M = L/N is injective, which allows one to identify T with a subset of M by means of \(\psi\) . Show that, if \(\lambda_1, \ldots, \lambda_m\) are elements \(\neq 0\) of T, the sum \(\lambda_1 + \lambda_2 + \cdots + \lambda_m\) cannot be 0 in M (consider the element \(2^{-m}(\lambda_1 + \cdots + \lambda_m)\) and reason by induction on m, noting that these elements belong to T).
\item
  Let \(M_{+}\) be the set of finite sums (in M) of elements of T; show that \(\mathbf{M}_{+}\cap(-\mathbf{M}_{+})=\left\{0\right\}\) and \(\mathbf{M}=\mathbf{M}_{+}\cup(-\mathbf{M}_{+})\) and consequently that \(M_{+}\) is the set of elements \(\geqslant0\) for a totally ordered group structure on M (note that for every \(\mu\in M_{+}\) , there exists an integer m such that \(2^{-m}\mu\in T\) ); this totally ordered group is said to be the group of measures of angles of half-lines. Show that the homomorphism f from M to A is surjective, and that its kernel is the set of integer multiples of \(2\varpi\) , where \(\varpi\) is the straight angle (exerc. 6 c)). Prove that T is identified with the interval \([0,\varpi]\) in M (prove by induction on m that if \(\mu,\lambda_{1},\ldots,\lambda_{m}\) belong to T and if we have \(\lambda_{1}+\cdots+\lambda_{m}\leqslant\mu\) , then \(\lambda_{1}+\cdots+\lambda_{m}\in T\) ). Show that in T (thus identified with an interval of M), the function \(\theta\to\cos\theta\) is strictly decreasing.
\item
  For the totally ordered group M to be Archimedean (chap. VI, § 1, exerc. 31), it is necessary and sufficient that the additive group of the field A be Archimedean. (To see that the condition is necessary, note that if sin θ is infinitely small with respect to the subfield Q of A (chap. VI, § 2, exerc. 1), then so is sin nθ for every integer n. To see that the condition is sufficient, note that if 0 ≤ θ ≤ π/4, then sin 2θ ≥ √2 sin θ).
\end{enumerate}

\begin{enumerate}
\def\labelenumi{\arabic{enumi})}
\setcounter{enumi}{7}
\tightlist
\item
  Let E be an oriented plane over an ordered field A. Let D', D'' be two distinct half-lines; let \(x'\) (resp. \(x''\) ) be a vector \(\neq 0\) in D′ (resp. D″); the angular sector (open or closed) with origin D' and endpoint D'' is said to be salient (resp. reentrant, flat) if \(x' \wedge x'' > 0\) (resp. \(x' \wedge x'' < 0\) , \(x' \wedge x'' = 0\) ).
\end{enumerate}

\begin{enumerate}
\def\labelenumi{\alph{enumi})}
\item
  For there to exist an automorphism of the vector space E mapping an open (resp. closed) angular sector \(\Sigma_{1}\) into an open (resp. closed) angular sector \(\Sigma_{2}\) , it is necessary and sufficient that \(\Sigma_{1}\) and \(\Sigma_{2}\) be both salient, or both reentrant, or both flat.
\item
  Show that the ordered set \(\left[D', D''\right]\) is isomorphic to the interval \(\left[0,1\right]\) of A (first consider the case of a salient sector and note that, in A, any two closed bounded intervals are isomorphic ordered sets).
\item
  With the notations and hypotheses of Exercise 7, define a canonical bijective map from the set T onto a flat angular sector, and show that this map is an isomorphism for the order structures.
\end{enumerate}

\begin{enumerate}
\def\labelenumi{\arabic{enumi})}
\setcounter{enumi}{8}
\item
  Let A be a Pythagorean ordered field, E a finite-dimensional vector space over A, Q a positive nondegenerate quadratic form on E, for which E admits an orthonormal basis (in other words, Q(x) is a square in A for every \(x \in E\) ). Show that the commutator subgroup \(\Omega(Q)\) of the orthogonal group \(\mathbf{O}(Q)\) is the rotation group \(\mathbf{O}^{+}(Q) = \mathbf{S}\mathbf{O}(Q)\) (use Exercise 3 e) of § 10 and Exercise 17 a) of § 6).
\item
  Let A be a maximal ordered field, E a finite-dimensional vector space over A, Q a positive nondegenerate quadratic form on E. For every orthogonal transformation \(u \in \mathbf{O}(Q)\) show that there exists a decomposition of E as a direct sum of pairwise orthogonal subspaces P, N, \(R_i\) ( \(1 \leqslant i \leqslant r\) ) having the following properties: \(1^\circ u(x) = x\) in P, \(u(x) = -x\) in N; \(2^\circ\) each of the \(R_i\) is of dimension 2, we have \(u(R_i) = R_i\) and the restriction of \(u\) to \(R_i\) is a rotation through the angle \(\theta_i\) distinct from 0 and from the straight angle. Moreover, for two decompositions of this nature, the subspaces P and N are the same, as is the sequence of elements \(\cos \theta_i\) up to order (cf. § 7, no. 3, cor. 2 of th. 2). Deduce that every rotation \(u \in \mathbf{O}^+(Q)\) is a commutator \(tst^{-1}s^{-1}\) , where \(s\) and \(t\) are in \(\mathbf{O}(Q)\) and \(s^2 = 1\) (cf. exerc. 3 d)).
\item
  Let A be a maximal ordered field, L the field of quaternions over A (relative to the pair \((-1,-1)\) ), E the subspace (of dimension 3) of L formed by the pure quaternions ( \(\S 9\) , exerc. 15 a)); every quaternion can therefore be written in a unique way \(s=\alpha.1+\nu\) , where \(\alpha\in A\) and \(\nu\in E\) ; we have \(\alpha^{2}-\nu^{2}=\rho.1\) , where \(\rho\in A\) and \(\rho\geqslant0\) ; we set \(\|\nu\|=\sqrt{\rho}\) .
\end{enumerate}

\begin{enumerate}
\def\labelenumi{\alph{enumi})}
\item
  Let \(\varphi(s)\) be the rotation \(x \to sxs^{-1}\) in E, for the positive nondegenerate quadratic form \(x \to \|x\|^2\) on E (§ 9, no. 5, th. 4 and exerc. 15). Show that if \(\nu \neq 0\) , the vectors of the line D ⊂ E containing \(\nu\) are invariant under \(\varphi(s)\) ; the restriction of \(\varphi(s)\) to the orthogonal plane P = D⁰ is a rotation of angle \(\theta\) such that (for a suitable orientation of P) one has \(\operatorname{tg} \frac{\theta}{2} = \| \nu \| / \alpha\) . (If (1, i, j, k) is the canonical basis of L over A, reduce to the case where \(\nu = \beta i\) , \(\beta \in A\) , and then compute \(sjs^{-1}\) ).
\item
  Show that every quaternion of norm 1 can be written \(tst^{-1}s^{-1}\) , where \(s \in L\) , \(t \in E\) (cf. exerc. 10), and that, if \(a\) , \(b\) are two pure quaternions of norm 1, there exists a quaternion \(s\) such that \(b = sas^{-1}\) .
\item
  For two quaternions \(s = \alpha + \nu\) , \(t = \beta + w\) , where \(\alpha, \beta\) are in A, \(\nu, w\) in E, to be permutable, it is necessary and sufficient that \(\nu = \lambda w\) , \(\lambda \in A\) ; for \(st = -ts\) , it is necessary and sufficient that \(\alpha = \beta = 0\) and that the vectors \(\nu\) and \(w\) in E be orthogonal (cf. § 3, exerc. 10).
\end{enumerate}

\begin{enumerate}
\def\labelenumi{\arabic{enumi})}
\setcounter{enumi}{11}
\tightlist
\item
  Let A be a maximal ordered field, and let L be an n-dimensional Euclidean space over A, whose metric form \(\Phi\) is positive nondegenerate; we denote by \(d(x, y)\) the distance \(\sqrt{\Phi(x - y, x - y)}\) between two points of L ( \(\S 7\) , no. 1, Remark). For every point \(c \in L\) and every element \(\rho > 0\) of A, one calls sphere with center c and radius \(\rho\) the set of \(x \in L\) such that \(d(x, c) = \rho\) ; the spheres are therefore nondegenerate affine quadrics in L, admitting a center ( \(\S 6\) , Exercise 25).
\end{enumerate}

\begin{enumerate}
\def\labelenumi{\alph{enumi})}
\item
  Show that at every point x of a sphere S with center c, the tangent hyperplane to S at x (§ 6, Exercise 25) is perpendicular (§ 6, Exercise 22) to the line passing through c and x.
\item
  Let S be a sphere with center c and radius \(\rho\) , let a be a point of L, and let D be a line passing through a and meeting S in two distinct points \(x_{1}\) , \(x_{2}\) (resp. tangent to S at a point x). Show that one has
\end{enumerate}

\[
\Phi (x _ {1} - a, x _ {2} - a) = (d (a, c)) ^ {2} - \rho^ {2} \quad (\text { resp. } (d (x, a)) ^ {2} = (d (a, c)) ^ {2} - \rho^ {2}).
\]

The element \((d(a, c))^{2} - \rho^{2}\) of A is called the power of a with respect to S.

\begin{enumerate}
\def\labelenumi{\alph{enumi})}
\setcounter{enumi}{2}
\item
  Let \(S_{1}\) , \(S_{2}\) be two spheres not having the same center; show that the set of points of L whose powers with respect to \(S_{1}\) and \(S_{2}\) are equal, is a hyperplane perpendicular to the line passing through the centers of the two spheres, and containing their intersection \(S_{1} \cap S_{2}\) ; this hyperplane is said to be the radical hyperplane of \(S_{1}\) and \(S_{2}\) .
\item
  Let \(S_{1}\) , \(S_{2}\) be two spheres with respective centers \(c_{1}\) , \(c_{2}\) and respective radii \(\rho_{1}\) , \(\rho_{2}\) . Show that the following properties are equivalent:
\end{enumerate}

α) The intersection \(S_{1} \cap S_{2}\) is nonempty and, for every point of this intersection, the tangent hyperplanes to \(S_{1}\) and \(S_{2}\) at this point are perpendicular.

\(\beta_{1}\) The power of \(c_{1}\) with respect to \(S_{2}\) is \(\rho_{1}^{2}\) .

\(\beta_{2})\) The power of \(c_{2}\) with respect to \(\mathbf{S}_{1}\) is \(\rho_{2}^{2}\) .

\(\gamma_{1})\) The radical hyperplane of \(\mathbf{S}_{1}\) and \(\mathbf{S}_{2}\) is the polar hyperplane (§ 6, exerc. 25) of \(c_{1}\) with respect to \(\mathbf{S}_{2}\) .

\(\gamma_{2})\) The radical hyperplane of \(\mathbf{S}_{1}\) and \(\mathbf{S}_{2}\) is the polar hyperplane of \(c_{2}\) with respect to \(\mathbf{S}_{1}\) .

When these properties hold, the spheres \(S_{1}\) , \(S_{2}\) are said to be orthogonal.

\begin{enumerate}
\def\labelenumi{\alph{enumi})}
\setcounter{enumi}{4}
\tightlist
\item
  Let \(S_{1}\) , \(S_{2}\) be two orthogonal spheres, \(c_{1}\) , \(c_{2}\) their centers. Show that if \(\varpi_{1}\) , \(\varpi_{2}\) are the powers of a point x with respect to \(S_{1}\) , \(S_{2}\) respectively, we have \(\varpi_{1} + \varpi_{2} = 2\Phi(x - c_{1}, x - c_{2})\) . Converse.
\end{enumerate}

\begin{enumerate}
\def\labelenumi{\arabic{enumi})}
\setcounter{enumi}{12}
\tightlist
\item
  The hypotheses and notation are the same as in Exercise 12. Given a point \(c \in \mathrm{L}\) and an element \(\alpha \neq 0\) of \(\mathbf{A}\) , we call inversion with pole \(c\) and of power \(\alpha\) the involutive permutation \(u\) of the set \(\mathrm{L} - \{c\}\) which is such that, for all \(x \in \mathrm{L} - \{c\}\) , \(u(x)\) belongs to the line passing through \(c\) and \(x\) and satisfies the relation \(\Phi(x - c, u(x) - c) = \alpha\) . By abuse of language, we say that \(u\) is an inversion in \(\mathrm{L}\) .
\end{enumerate}

\begin{enumerate}
\def\labelenumi{\alph{enumi})}
\item
  If \(u\) , \(\nu\) are two inversions with the same pole \(c\) and powers \(\alpha\) , \(\beta\) , \(uv^{-1}\) is the restriction to L - \(\{c'\}\) of the homothety (chap. II, 2nd ed., App. II, exerc. 6) with center \(c\) and ratio \(\alpha\beta^{-1}\) .
\item
  Let S be a sphere (exerc. 12) containing c. Show that the image of S - \{c\} under an inversion with pole c is a hyperplane perpendicular to the line joining c to the center of S (by an abuse of language, this hyperplane is said to be the image of S under the inversion in question). Converse.
\item
  Let S be a sphere not containing c. Show that, if \(\varpi\) is the power of c with respect to S (exerc. 12 b)), the image of S under an inversion with pole c and power \(\alpha\) is the image of S under a homothety with center c and ratio \(\alpha/\varpi\) . If n = 2 and if, for every \(x \in S\) , we denote by T (resp. T') the tangent to S (resp. \(u(S)\) ) at the point x (resp. \(u(x)\) ), by D the line passing through c, x and \(u(x)\) , show that one has \((\widehat{\mathrm{D},\mathrm{T}}) = (\widehat{\mathrm{T}^{\prime},\mathrm{D}})\) .
\item
  Let \(S_{1}\) , \(S_{2}\) be two orthogonal spheres (exerc. 12 d)), \(S_{1}'\) , \(S_{2}'\) their images under an inversion with pole c. If c does not belong to \(S_{1}\) nor to \(S_{2}\) , show that \(S_{1}'\) and \(S_{2}'\) are orthogonal spheres. If \(c \in S_{1}\) and \(c \notin S_{2}\) , \(S_{1}'\) is a hyperplane containing the center of \(S_{2}'\) . If \(c \in S_{1} \cap S_{2}\) , \(S_{1}'\) and \(S_{2}'\) are perpendicular hyperplanes (§ 6, exerc. 22). Converse.
\item
  Let \(u\) be an inversion with pole \(c\) and of power \(\alpha = \rho^2 > 0\) and C the sphere with center \(c\) and radius \(\rho\) . If \(x_1, x_2\) are two distinct points lying on a line passing through \(c\) , and distinct from \(c\) , the following properties are equivalent: \(\alpha\) \(x_1\) and \(x_2\) are transformed into each other by \(u\) ; \(\beta\) \(x_1\) and \(x_2\) are conjugate with respect to C (§ 6, exerc. 25); \(\gamma\) every sphere containing \(x_1\) and \(x_2\) is orthogonal to C. One also says that \(u\) is the inversion with sphere C.
\end{enumerate}

¶ 14) The hypotheses and notations being the same as in exerc. 12, one takes an origin 0 in L. Let E \(_{1}\) be the vector space direct sum of L and a space Af \(_{1}\) of dimension 1; we denote by Q \(_{1}\) the quadratic form on E \(_{1}\) such that for x ∈ L and η ∈ A, we have

\[
\mathrm{Q} _ {1} (x + \eta f _ {1}) = \mathrm{Q} (x) + \eta^ {2},
\]

form which is positive and nondegenerate; we denote by C the sphere with center 0 and radius 1 in E \(_{1}\) (for Q \(_{1}\) ).

In the Euclidean space E \(_{1}\) let s be the inversion with pole -f \(_{1}\) and power 2 (exerc. 13); its restriction s \(_{0}\) to L maps L onto C -\{-f \(_{1}\) \}; s \(_{0}\) (resp. s \(_{0}^{-1}\) ) is called, by abuse of language, the stereographic projection of L onto C (resp. of C onto L) with viewpoint -f \(_{1}\) . For every inversion u in L, with pole c, s \(_{0}\) us \(_{0}^{-1}\) is an involutive permutation of the complement in C of the set \{s \(_{0}\) (c), -f \(_{1}\) \}; we extend it to an involutive permutation u' of C by setting u'(s \(_{0}\) (c)) = -f \(_{1}\) , u'(-f \(_{1}\) ) = s \(_{0}\) (c), and one says that u' is an inversion in C. Likewise, for every symmetry v in L with respect to a hyperplane, s \(_{0}\) vs \(_{0}^{-1}\) is an involutive permutation of the complement in C of the set \{-f \(_{1}\) \}; one extends it to an involutive permutation v' of C by setting v'(-f \(_{1}\) ) = -f \(_{1}\) and we say that v' is a symmetry in C. The group of permutations of C generated by inversions and symmetries is called the conformal group of C (or of L, by abuse of language).

\begin{enumerate}
\def\labelenumi{\alph{enumi})}
\item
  Show that the conformal group of C is generated by the symmetries \(v'\) and the inversions \(u'\) corresponding to the inversions \(u\) in L, of power \(>0\) . (Use exerc. 13 a) and note that in L every translation, as well as the homothety \(x \to -x\) , are products of symmetries with respect to hyperplanes).
\item
  Let u be an inversion in L of power \textgreater{} 0; show that the corresponding inversion \(u'\) in C is the restriction to C of a well-determined transformation \(u_{1}'\) which is either an inversion of power \textgreater{} 0 in \(E_{1}\) whose sphere (exerc. 13 c)) is orthogonal to C, or a symmetry with respect to a hyperplane of \(E_{1}\) passing through 0 (consider in \(E_{1}\) the inversion \(u_{1}\) with the same pole and the same power as u). Formulate the corresponding proposition for the symmetry \(v'\) in C corresponding to a symmetry v in L with respect to a hyperplane.
\item
  In the vector space \(E_{2} = A \times E_{1}\) , consider the quadratic form \(Q_{2}\) such that, for \(\zeta \in A\) , \(y \in E_{1}\) one has
\end{enumerate}

\[
\mathrm{Q} _ {2} ((\zeta , y)) = \zeta^ {2} - \mathrm{Q} _ {1} (y),
\]

a form which is nondegenerate and of signature \((1, n + 1)\) . Identify \(E_{1}\) with its canonical image in projective space \(\mathbf{P}(E_{2})\) (chap. II, 2nd ed., App. III, no. 4). Show (with the notations of b)) that \(u'\) is also the restriction to C of a projective linear map \(\overline{u}''\) obtained by passing to quotients from a transformation \(u''\) of the orthogonal group \(\mathbf{O}(Q_{2})\) , which is a symmetry with respect to a non-isotropic hyperplane in \(E_{2}\) . Formulate the corresponding proposition for \(\varphi'\) . Deduce from this that the conformal group of L is isomorphic to the quotient of the group \(\mathbf{O}(Q_{2})\) by its center (use prop. 5 and exerc. 17 c) of § 6). Conclude from this that every transformation of the conformal group is a product of at most \(n + 2\) transformations which are inversions or symmetries in L (cf. § 6, exerc. 15 e)).

\begin{enumerate}
\def\labelenumi{\alph{enumi})}
\setcounter{enumi}{3}
\tightlist
\item
  Let \(\Sigma\) be the set whose elements are the spheres and the hyperplanes in the affine space L. Deduce from \(b\) ) that there exists a bijection of \(\Sigma\) onto the complement, in \(\mathbf{P}(\mathrm{E}_2)\) , of the set of \(x \in \mathrm{E}_1\) such that \(\mathrm{Q}_1(x) \leqslant 1\) , so that to two orthogonal spheres there correspond two points conjugate with respect to C.
\end{enumerate}

\begin{enumerate}
\def\labelenumi{\arabic{enumi})}
\setcounter{enumi}{14}
\item
  Generalize the definitions and results of Exercises 12 to 14 to the case where A is a Pythagorean field and where there exists an orthonormal basis for \(\Phi\) .
\item
  Let A be a commutative field, V a vector space over A of odd dimension \(2r + 1\) , F the product space A × V, \(\Psi\) a nondegenerate alternating form on F; in the projective space P(F), of dimension \(2r + 1\) , we say that the set \(C_{0}\) of lines that are the canonical images of the totally isotropic planes of F (for \(\Psi\) ) is the (projective) linear complex associated with \(\Psi\) .
\end{enumerate}

Assume henceforth that A is maximally ordered; let \(\Phi\) be a nondegenerate positive symmetric form on V. Let D be the line orthogonal to V (for \(\Psi\) ) in F, which is contained in V, and let H be the hyperplane orthogonal to D (for \(\Phi\) ) in V; in F, H is a non-isotropic subspace for \(\Psi\) , whose orthogonal for \(\Psi\) is therefore a plane P supplementary to H and containing D. In the affine space \(E = \{1\} \times V \subset F\) , one also says that the set C of intersections with E of the totally isotropic planes of F (for \(\Psi\) ) not contained in V, is the (affine) linear complex associated with \(\Psi\) ; the line \(\Delta = P \cap E\) is called the axis of the linear complex C (for the Euclidean space structure defined on E by the metric form \(\Phi\) ).

\begin{enumerate}
\def\labelenumi{\alph{enumi})}
\item
  Show that, in E, every translation equal to a direction vector of \(\Delta\) (chap. II, \(2^{\text{e}}\) ed., App. II, no. 3) leaves C invariant (cf. § 4, exerc. 6).
\item
  Let \((e_i)_{0\leqslant i\leqslant 2r}\) be an orthonormal basis of V for \(\Phi\) , such that \(e_0\in \mathrm{D}\) and that one has \(\Psi (e_{2i - 1},e_{2i}) = \rho_{i} > 0\) for \(1\leqslant i\leqslant r,\Psi (e_j,e_k) = 0\) for the pairs of indices that are not of the form \((e_{2i - 1},e_{2i})\) or \((e_{2i},e_{2i - 1})\)
\end{enumerate}

(§ 7, no. 3, prop. 6). One takes in the affine space E an origin \(a \in \Delta\) , and one sets \(\Psi(a, e_0) = \rho_0\) . Let \(x\) be a vector belonging to the plane generated by \(e_{2i-1}\) and \(e_{2i}\) , and let \(y\) be the vector of this same plane such that \(\Phi(x, y) = 0\) , \(\Phi(y, y) = 1\) and \(x \wedge y = \lambda e_{2i-1} \wedge e_{2i}\) with \(\lambda > 0\) . Let \(E_i\) be the affine linear variety of dimension 3 spanned by the points \(a\) , \(a + e_0\) , \(a + e_{2i-1}\) , \(a + e_{2i}\) in E, and let \(R_x\) be the intersection of \(E_i\) and of the affine hyperplane (in E) generated by the lines of C containing the point \(a + x\) . Show that the direction of the lines orthogonal (for \(\Phi\) ) to the plane \(R_x\) in \(E_i\) is a line \(L_x\) of the plane \(Ae_0 + Ay\) , such that if \(\theta = (\widehat{D, L_x})\) one has

\[
\operatorname{tg} \theta = \frac {\rho_ {i}}{\rho_ {0}} \sqrt {\Phi (x , x)}
\]

when the plane \(Ae_{0} + Ay\) is oriented so that \(e_{0} \wedge y\) is positive.

¶ 17) a) Let A be a commutative field, J : \(\xi \to \overline{\xi}\) an involutive automorphism of A, E a vector space of dimension 2 over A, \(\Phi\) a nondegenerate Hermitian sesquilinear form (not alternating) on E, \((e_1, e_2)\) an orthogonal basis of E for \(\Phi\) ( \(\S\) 6, no. 1, th. 1), such that

\[
\Phi (\xi_ {1} e _ {1} + \xi_ {2} e _ {2}, \eta_ {1} e _ {1} + \eta_ {2} e _ {2}) = \alpha \xi_ {1} \bar {\eta} _ {1} + \beta \xi_ {2} \bar {\eta} _ {2}
\]

(α and β belonging to the field K of invariants of J); we set γ = β/α. We identify the point \(\xi_{1}e_{1} + \xi_{2}e_{2} \in E\) with the element \(\xi_{1} + \xi_{2}\rho\) of the ring B defined by the conditions \(\rho^{2} = -\gamma\) , \(\rho\xi = \bar{\xi}\rho\) for \(\xi \in A\) (§ 3, exerc. 4 a)). For every \(x = \xi_{1} + \xi_{2}\rho \in B\) , set \(\tilde{x} = \bar{\xi}_{1} - \xi_{2}\rho\) and N(x) = x \(\tilde{x}\) = \(\tilde{x}\) x, so that \(x \to \tilde{x}\) is an involutive antiautomorphism of the algebra B (over K) and that one has \(\Phi(x, x) = \alpha N(x)\) . Show that every similarity for Φ whose determinant is equal to the multiplier (also called a direct similarity) can be written in one and only one way \(x \to xy\) , where y is a non-isotropic vector of E, and that its multiplier is N(y).

\begin{enumerate}
\def\labelenumi{\alph{enumi})}
\setcounter{enumi}{1}
\item
  Suppose first that J is the identity, hence K = A. If A has characteristic ≠ 2, recover in this way the results of no. 1. Develop the corresponding theory when A has characteristic 2 (cf. § 4, exerc. 14; one will distinguish two cases, according as γ is or is not a square in A).
\item
  We assume \(J \neq 1\) , so that A is a separable quadratic extension of K. Then \(\Phi\) necessarily satisfies condition (T) (§ 4, no. 2 and exerc. 1); if \(\Phi\) has index 0, B is a reflexive field with center K (chap. VIII, § 11, exerc. 4), and if \(\Phi\) has index 1, B is isomorphic to \(\mathbf{M}_2(\mathrm{K})\) .
\end{enumerate}

\(d)\) We assume \(\mathrm{J} \neq 1\) and A has characteristic \(\neq 2\) ; we then have \(\mathrm{A} = \mathrm{K}(\theta)\) with \(\theta^2 = -\delta \in \mathrm{K}\) . If S is the group of direct similarities for \(\Phi\) , H the group of homotheties in E, with ratio \(\neq 0\) and belonging to K (a group isomorphic to \(\mathrm{K}^*\) ), show that the group S/H is isomorphic to the group of rotations \(\mathrm{O}^+(\mathrm{Q})\) , where Q is a nondegenerate quadratic form on a vector space F of dimension 3 over K, such that there exists an orthogonal basis \((f_1, f_2, f_3)\) of F for which one has

\[
\mathrm{Q} \left(\zeta_ {1} f _ {1} + \zeta_ {2} f _ {2} + \zeta_ {3} f _ {3}\right) = \gamma \zeta_ {1} ^ {2} + \delta \zeta_ {2} ^ {2} + \gamma \delta \zeta_ {3} ^ {2}
\]

(cf. § 9, exerc. 15).

¶ 18) a) Let A be a commutative field, \(\xi \to \bar{\xi}\) an involutive automorphism of A, E a vector space of even dimension \(2m\) over A, \(\Phi\) a nondegenerate Hermitian form of index 0 on E, satisfying condition (T), \(\Delta\) the discriminant of \(\Phi\) with respect to a basis of E. Let M(Φ) be the group of multipliers of similarities for \(\Phi\) (§ 4, exerc. 8). Show that M(Φ) is a subgroup of the multiplicative group of elements of A of the form \(\alpha\bar{\alpha} - (-1)^{m}\beta\bar{\beta}\Delta\) . (Argue by induction on \(m\) , using exerc. 17, as well as the following two remarks: \(1^{\circ}\) if \(u\) is a similarity with multiplier \(\mu\) , \(x\) a vector of E, \(y = u(x)\) and \(z = u(y)\) , there exists a unitary transformation \(\nu\) such that \(\nu(y) = y\) and \(\nu(z) = \mu x\) ; \(2^{\circ}\) if \(\alpha\) , \(\beta\) , \(\lambda\) are three elements \(\neq 0\) of A such that there exist \(a\) , \(b\) , \(c\) , \(d\) in A satisfying the conditions \(\lambda = a\bar{a} + \alpha c\bar{c}\) , \(\lambda = b\bar{b} + \beta dd\bar{d}\) , then there exist \(s\) , \(t\) in A satisfying the condition \(\lambda = s\bar{s} - \alpha\beta t\bar{t}\) ).

\begin{enumerate}
\def\labelenumi{\alph{enumi})}
\setcounter{enumi}{1}
\tightlist
\item
  Let K be a maximal ordered field, \(\mathrm{K}_{1} = \mathrm{K}((t_{1}))\) the field of formal power series in one indeterminate \(t_{1}\) , with coefficients in K (chap. IV, § 5, no. 7), \(\mathrm{A} = \mathrm{K}_{1}((t_{2}))\) the field of formal power series in a second indeterminate \(t_{2}\) , with coefficients in \(K_{1}\) . Let E be a vector space of dimension 6 over A, Q a nondegenerate quadratic form on E, such that there exists an orthogonal basis \((e_{i})\) for which one has
\end{enumerate}

\[
\mathrm{Q} (\sum_ {i = 1} ^ {6} \xi_ {i} e _ {i}) = \xi_ {1} ^ {2} + \xi_ {2} ^ {2} + \xi_ {3} ^ {2} + \xi_ {4} ^ {2} + t _ {1} \xi_ {5} ^ {2} + t _ {2} \xi_ {6} ^ {2}.
\]

Show that there exists no similarity for Q with multiplier \(t_{1}t_{2}\) .

\backmatter
\chapter*{Historical Note}
\phantomsection
\addcontentsline{toc}{chapter}{Historical Note}

(N.-B. --- The Roman numerals refer to the bibliography placed at the end of this note.)

The theory of quadratic forms, in its modern aspect, goes back hardly beyond the second half of the eighteenth \(^{e}\) century, and, as we shall see, it developed above all to meet the needs of Arithmetic, Analysis, and Mechanics. But the fundamental notions of this theory in reality made their appearance from the beginnings of "Euclidean" geometry, of which they form the framework. For this reason, one cannot trace its history without speaking, at least in a summary way, of the development of "elementary geometry" since antiquity. Of course, we can only concern ourselves with the evolution of a few general ideas, and the reader should not expect to find here precise information on the history of any particular theorem, concerning which it will suffice for us to refer to the specialized historical or didactic works (*). It goes without saying also that, when we speak below of the various possible interpretations of one and the same theorem in various algebraic or geometric languages, we in no way mean to say that these "translations" have at all times been as familiar as they are today; quite the contrary, it is the principal goal of this Note to show how, very gradually, mathematicians became aware of these affinities between questions often of very different aspect; we shall also have to show how, in doing so, they were led to introduce some coherence into the mass of geometric theorems bequeathed by the ancients, and finally to try to delimit exactly what should be understood by "geometry."

Apart from the Babylonians' discovery of the formula for solving the quadratic equation ((I), pp. 183-189), it is thus under their geometric disguise that the birth of the principal concepts of the theory of quadratic forms must be noted. These present themselves first as squares of distances (in the plane or in three-dimensional space), and the corresponding notion of "orthogonality" is introduced by means of the right angle, defined by Euclid as half of a straight angle (Elements, Book I, Def. 10); the notions of distance and right angle being related by the Pythagorean theorem, the keystone of the Euclidean edifice (*). The idea of angle seems to have been introduced very early in Greek mathematics (which probably received it from the Babylonians, well versed in the use of angles through their long astronomical experience). It is known that in the classical period, only angles less than two right angles are defined (Euclid's "definition" is moreover as vague and unusable as the one he gives for the straight line or the plane); the notion of orientation is not brought out, although Euclid uses (without axiom or definition) the fact that a line divides the plane into two regions, which he carefully distinguishes when necessary (**). At this stage, the idea of the group of plane rotations thus emerges only in a very imperfect manner, through the addition (introduced, also without explanation, by Euclid) of non-oriented angles of half-lines, which is defined, in principle, only when the sum is at most equal to two right angles (***). As for trigonometry, it is disdained by geometers and left to surveyors and astronomers; it is the latter (Aristarchus, Hipparchus, and especially Ptolemy (V)) who establish the fundamental relations between sides and angles of a right triangle (plane or spherical) and draw up the first tables (these are tables giving the chord of the arc cut off by an angle \(\theta < \pi\) on a circle of radius \(r\) , in other words the number \(2r \sin \frac{\theta}{2}\) ; the introduction of the sine, which is more convenient to handle, is due to the Hindu mathematicians of the Middle Ages); in the computation of these tables, the addition formula for arcs, unknown at that time, is replaced by the equivalent use of Ptolemy's theorem (possibly going back to Hipparchus) on quadrilaterals inscribed in a circle (cf. Esp. vect. top., chap. V, § 1, exerc. 5). It should also be noted that Euclid and Heron give propositions equivalent to the formula

\[
a ^ {2} = b ^ {2} + c ^ {2} - 2 b c \cos A
\]

between sides and angles of an arbitrary plane triangle; but one can hardly see in this a first appearance of the notion of the bilinear form associated with a metric form, for lack of the idea of a vector calculus that would emerge only in the nineteenth \(^{e}\) century.

Displacements (or motions, the distinction between the two notions not being clear in antiquity---nor even much later) were known to Euclid; but, for reasons unknown to us, he seems to show a marked reluctance to make use of them (for example in the "cases of equality of triangles," where one has the impression that he employs the notion of displacement only for want of having been able to formulate an appropriate axiom ((III bis), t. I, p. 225-227 and 249)); however, it is to the notion of displacement (rotation about an axis) that he has recourse for the definition of cones of revolution and spheres (Elements, Book XI, defs. 14 and 18), as does Archimedes for that of quadrics of revolution. But the general idea of transformation, applied to all of space, is almost foreign to mathematical thought before the end of the eighteenth \(^{e}\) century (*);

and before the seventeenth \(^{e}\) century, there is likewise no trace of the notion of composition of motions, nor a fortiori of composition of displacements. This does not mean, of course, that the Greeks were not particularly sensitive to the “regularities” and “symmetries” of figures, which we now connect with the notion of the group of displacements; their theory of regular polygons and even more that of regular polyhedra — one of the most remarkable chapters of all their mathematics — is there to prove the contrary (*).

Finally, the last of the essential contributions of Greek mathematics, in the domain that concerns us, is the theory of conics (as for quadrics, the Greeks knew only certain quadrics of revolution, and did not carry their study very far, except for the sphere). It is interesting to note here that, although the Greeks never had the idea of the fundamental principle of analytic geometry (essentially for lack of a manageable algebra), they commonly used, for the study of particular "figures," the "ordinates" with respect to two (or even more than two) axes in the plane (closely related to the figure, which is one of the fundamental points where their method differs from that of Fermat and Descartes, whose axes are fixed independently of the figure under consideration). In particular, the first examples of conics (other than the circle) that appear in connection with the problem of the duplication of the cube are the curves given by the equations \(y^{2} = ax\) , \(y = bx^{2}\) , \(xy = c\) (Menaechmus, a student of Eudoxus, mid- \(1v^{e}\) century) (**); and it is the equation of conics (usually with respect to two oblique axes formed by a diameter and the tangent at one of its points of intersection with the curve) that is most often used in the study of problems relating to these curves (whereas "focal" properties play only a very subdued role, contrary to what school traditions dating back only to the \(x1x^{e}\) century might lead one to believe). From this vast theory, what we must above all retain here is the notion of conjugate diameters (already known to Archimedes), and the property that now serves as the definition of the polar of a point, given by Apollonius (IV) when the point is exterior to the conic (the polar thus being for him the line joining the points of contact of the tangents drawn from this point); from our point of view, these are two examples of "orthogonality" with respect to a quadratic form distinct from the metric form, but of course the connection between these notions and the classical notion of perpendiculars could not possibly have been conceived at that time.

There is scarcely any other progress to report before Descartes and Fermat; but from the very beginnings of analytic geometry, the algebraic theory of quadratic forms begins to emerge from its geometric shell: Fermat knows that a second-degree equation in the plane represents a conic ((VII a), p. 100-102) and sketches analogous ideas on quadrics (VII b). With the development of analytic geometry in 2 and 3 dimensions during the eighteenth \(^{e}\) century there appear (especially in connection with conics and quadrics) two of the central problems of the theory: the reduction of a quadratic form to a sum of squares and the search for its "axes" with respect to the metric form. For conics, these two problems are too elementary to give rise to important algebraic progress; for an arbitrary number of variables, the first is solved by Lagrange in 1759, in connection with the maxima of functions of several variables (IX a). But this problem is almost immediately overshadowed by that of the search for axes, even before the invariance of the rank had been formulated (*); as for the law of inertia, it is discovered only around 1850 by Jacobi (XIX b), who proves it by the same reasoning as is used today, and Sylvester (XX), who merely states it as almost evident (**).

The problem of reducing a quadric to its axes already presents substantially greater algebraic difficulties than the analogous problem for conics; and Euler, who is the first to tackle it, is not in a position to prove the reality of the eigenvalues, which he admits after a sketch of a justification lacking probative value ((VIII a), p. 379-392) (**). If this point is correctly established around 1800, one must wait for Cauchy to prove the corresponding theorem for forms in an arbitrary number \(n\) of variables (XIV b). It is also Cauchy who, around the same time, proves that the characteristic equation giving the eigenvalues is invariant under any change of rectangular axes ((XIV a), p. 252) (****); but for \(n = 2\) or \(n = 3\) , this invariance was intuitively "evident" because of the geometric interpretation of the eigenvalues by means of the axes of the corresponding conic or quadric. Moreover, in the course of research on this subject, the elementary symmetric functions of the eigenvalues also presented themselves in a natural way (with various geometric interpretations, relating in particular to Apollonius's theorems on conjugate diameters), and in particular the discriminant, which (known for a long time for n = 2 in connection with the theory of the second-degree equation) appears for the first time for n = 3 in Euler ((VIII a) p. 382); the latter encounters it in connection with the classification of quadrics (in expressing the condition for a quadric to have no point at infinity) and does not mention its invariance under changes of rectangular axes. But a little later, with the beginnings of the arithmetic theory of quadratic forms with integer coefficients, Lagrange notes (for n = 2) a particular case of invariance of the discriminant under a linear but non-orthogonal change of variables ((IX b), p. 699), and Gauss establishes, for n = 3, the "covariance" of the discriminant under any linear transformation ((XI a), p. 301-302) (*). Once the general formula for the multiplication of determinants had been proved by Cauchy and Binet, the extension of Gauss's formula to an arbitrary number of variables was immediate; it is this that, around 1845, was to give the first impetus to the general theory of invariants.

To the two notions which, among the Greeks, took the place of the theory of displacements --- namely the notion of motion and that of "symmetry" of a figure --- there is added a third in the seventeenth \(^{e}\) and eighteenth \(^{e}\) centuries with the problem of the change of rectangular axes, which is substantially equivalent to this theory. Euler devotes several works to this question, striving above all to obtain manageable parametric representations for the formulas of change of axes. We know what use Mechanics was to make of the three angles that he introduces for this purpose for \(n = 3\) ((VIII a), p. 371-378). But he does not stop there; he considers in 1770 the general problem of orthogonal transformations for arbitrary \(n\) , observes that one arrives here at the goal by introducing \(n(n - 1)/2\) angles as parameters, and finally, for \(n = 3\) and \(n = 4\) , gives rational representations for rotations (in terms, respectively, of 4 homogeneous parameters and of 8 homogeneous parameters linked by a relation), which are none other than those obtained later by means of the theory of quaternions (cf. § 9, exerc. 15 and 16), and whose origin he does not indicate (VIII c) (**).

On the other hand, Euler also indicates how to translate analytically the search for the “symmetries” of plane figures, and it is in this connection that he is led to prove, in substance, that a plane displacement is a rotation, or a translation, or a translation followed by a symmetry ((VIII a), pp. 197-199). The rise of Mechanics at that time moreover leads to the general study of displacements; but at first only “infinitely small” displacements tangent to continuous motions are in question: these are apparently the only ones that occur in the researches of Torricelli, Roberval, and Descartes on the composition of motions and the instantaneous center of rotation for plane motions (cf. Hist. Note of Book IV, chaps. I-II-III). The latter is defined in a general way by Johann Bernoulli; d'Alembert in 1749, Euler the following year, extend this notion by proving the existence of an instantaneous axis of rotation for motions leaving a fixed point. The analogous theorem for finite displacements is stated only in 1775 by Euler (VIII d), in a memoir in which he discovers at the same time that the determinant of a rotation is equal to 1; the following year, he proves the existence of a fixed point for plane similarities (VIII e). But it will be necessary to wait for the works of Chasles, starting in 1830 (XV a), to finally have a coherent theory of finite and infinitely small displacements.
% semantic-fragments: legacy-input-seam
\relax{}
\begin{center}
\(\ast\qquad\ast\)
\end{center}

We thus arrive at what may be called the golden age of geometry, which falls roughly between the publication dates of Monge's Descriptive Geometry (1795) (X) and F. Klein's "Erlangen Program" (1872) (XXV b). The essential advances we owe to this sudden revival of geometry are the following:

\begin{enumerate}
\def\labelenumi{\Alph{enumi})}
\item
  The notion of an element at infinity (point, line, or plane), introduced by Desargues in the seventeenth century (VI), but which hardly manifests itself in the eighteenth century only as an abuse of language, is rehabilitated and systematically used by Poncelet (XII), who thus makes projective space the general framework for all geometric phenomena.
\item
  At the same time, with Monge, and above all Poncelet, the transition to complex projective geometry takes place. The notion of imaginary point, sporadically used during the eighteenth century, is exploited here (concurrently with that of the point at infinity) to yield statements independent of the "cases of figure" of real affine geometry. If, at first, the justifications offered in support of these innovations remain highly awkward (especially on the part of the adherents of the "synthetic" school of geometry, where the use of coordinates came to be regarded as a stain), one cannot fail to recognize therein, under the name of "principle of contingent relations" in Monge, or of "principle of continuity" in Poncelet, the first germ of the idea of "specialization" in modern algebraic geometry (*).
\end{enumerate}

One of the first results arising from these conceptions is the observation that, in complex projective space, all nondegenerate conics (resp. quadrics) are of the same nature; this leads Poncelet to the discovery of the "isotropic" elements: "Circles placed arbitrarily in a plane," he says, "are therefore not entirely independent of one another, as one might at first believe; they have ideally two imaginary points in common at infinity" ((XII), p.~48). Further on, he likewise introduces the "umbilic," the imaginary conic at infinity common to all spheres ((XII), p.~370); and if he does not speak particularly of the isotropic generators of the sphere, he at least explicitly emphasizes the existence of rectilinear generators, real or imaginary, for all quadrics (ibid., p.~371) (*); notions of which his successors (notably Plücker and Chasles), even more than he himself, make great use, in particular in the study of the "focal" properties of conics and quadrics.

\begin{enumerate}
\def\labelenumi{\Alph{enumi})}
\setcounter{enumi}{2}
\item
  The notions of point transformation and composition of transformations are likewise formulated in a general way and introduced systematically as means of proof. Apart from displacements and projections, only a few particular transformations had been known up to that time: certain planar projective transformations, of the type \(x' = a/x\) , \(y' = y/x\) used by La Hire and Newton, the “affinity” \(x' = ax\) , \(y' = by\) from Clairaut and Euler, and finally certain particular quadratic transformations, in Newton again, Maclaurin, and Braikenridge. Monge, in his Descriptive Geometry, shows all the use that can be made of plane projections in three-dimensional geometry. In Poncelet, one of the systematic methods of proof, employed to excess, consists in reducing by projection the properties of conics to those of the circle (a method already applied on occasion by Desargues and Pascal); and in order to pass similarly from a quadric to a sphere, he invents the first example of a projective transformation in space, the "homology" ((XII), p.~357); finally, it is also he who introduces the first examples of birational transformations of a curve into itself. In 1827, Möbius ((XIII a), p.~217) (and independently Chasles in 1830 (XV b)) define the most general projective linear transformations; at the same period there appear inversion (cf.~§ 10, exerc. 13) and other types of quadratic transformations, whose study will inaugurate the theory of birational transformations, which will develop in the second half of the nineteenth century.
\item
  The notion of duality appears in full light and is consciously linked to the theory of bilinear forms. The theory of poles and polars with respect to conics, which since Apollonius had made progress only in the work of Desargues and La Hire, is extended to quadrics by Monge, who, together with his students, perceives the possibility of transforming known theorems into new results by this means (*). But it is again to Poncelet that the credit belongs for having erected these remarks into a general method in his theory of transformations "by reciprocal polars," and for having made of it a particularly effective tool of discovery. A little later, notably with Gergonne, Plücker, Möbius, and Chasles, the general notion of duality detaches itself from the connection with quadratic forms, still too narrow in Poncelet. In particular, Möbius, in examining the various possibilities of duality in 3-dimensional space (defined by a bilinear form), discovers in 1833 the duality with respect to an alternating bilinear form (XIII b) (**), especially studied in the nineteenth century, in the form of the theory of “linear complexes” (cf.~§ 10, exerc. 16) and developed in connection with the “geometry of lines” and the “Plücker coordinates” introduced by Cayley, Grassmann, and Plücker around 1860.
\item
  From the very beginnings of projective geometry, the intensive study of the properties of classical geometry in their relations with projective space had quickly led to dividing them into "projective properties" and "metric properties"; and it is doubtless no exaggeration to see in this separation one of the clearest manifestations, at that time, of what was to become the modern notion of structure. But Poncelet, who first introduced this distinction and this terminology, was already aware of what connects these two types of properties; and, addressing in his Treatise the problems concerning angles, whose properties "do not seem to form part of those we have called projective\ldots, they nevertheless follow in so simple a manner," he said, "from the principles that form the basis {[}of this work{]}\ldots, which I do not believe that any other geometric theory can lead to in a manner at once more direct and simpler. One will not be at all surprised by this, if one considers that the projective properties of figures are necessarily the most general of those that can belong to them; so that they must include, as simple corollaries, all the other particular properties or relations of extension" ((XII), p.~248). To tell the truth, after this declaration, one is somewhat surprised to see him approach questions of angles in a very roundabout way, by attaching them to the focal properties of conics, instead of bringing the cyclic points directly into play; and in fact, it was only 30 years later that Laguerre (still a student at the École Polytechnique) gave the expression of an angle between lines by means of the cross-ratio of these lines and the isotropic lines with the same origin (cf.~§ 10, exerc. 5) ((XXI), vol. II, p.~13). Finally, with Cayley (XVIII d) there is clearly expressed the fundamental idea that the "metric" properties of a plane figure are none other than the "projective" properties of the figure augmented by the cyclic points --- a decisive milestone toward the "Erlangen program".
\item
  Hyperbolic non-Euclidean geometry, which came into being around 1830, at first remains somewhat apart from the movement whose main lines we are tracing. Arising from concerns of an essentially logical order touching on the foundations of classical geometry, this new geometry is presented by its inventors (*) in the same axiomatic and "synthetic" form as Euclid's geometry, and without any connection to projective geometry (whose introduction following the classical model even seemed excluded a priori, since the notion of a unique parallel disappears in this geometry); it is doubtless for this reason that for a long time it attracted little interest from the French, German, and English schools of projective geometry. Thus, when Cayley, in the fundamental memoir cited above (XVIII d), had the idea of replacing the cyclic points (considered as a conic "degenerate tangentially") by an arbitrary conic (which he called the "absolute"), he in no way thought of connecting this idea to the geometry of Lobachevsky-Bolyai, although he indicates how his conception leads to new expressions for the "distance" between two points, and although he mentions its links with spherical geometry. The situation changes around 1870, when the non-Euclidean geometries, following the diffusion of the works of Lobachevsky, and the publication of the works of Gauss and of Riemann's inaugural lecture, came to the forefront of mathematical actuality. Following the path traced by Riemann, Beltrami, without knowing Cayley's work, rediscovered in 1868 the expressions for distance given by the latter, but in a quite different context, by considering the interior of a circle as an image of a surface of constant curvature, in which geodesics are represented by straight lines (XXIV); it is Klein who, two years later, made (independently of Beltrami) the synthesis of these various points of view, which he completed by the discovery of elliptic non-Euclidean space (XXV a) (**).
\item
  In the second half of the period we are considering here, a period of critical reflection sets in, during which the partisans of "synthetic" geometry, not content with having banished coordinates from their proofs, claim to do without real numbers even in the axioms of geometry. The principal representative of this school is von Staudt, who essentially succeeded in performing this tour de force (XXIII), much admired in his time and even well into the twentieth century; and if today one no longer attributes the same importance to ideas of this order, whose possibilities for fruitful application have proved rather meager, one must nevertheless recognize that the efforts of von Staudt and his disciples contributed to clarifying ideas on the role of real or complex "scalars" in classical geometry, and thereby to introducing the modern conception of geometries over an arbitrary base field.
\end{enumerate}

Around 1860, "synthetic" geometry was at its zenith, but the end of its reign was fast approaching. Having remained heavy and inelegant throughout the eighteenth century, analytic geometry, in the hands of Lamé, Bobillier, Cauchy, Plücker, and Möbius, finally acquired the elegance and concision that would enable it to compete on equal terms with its rival. Above all, from about 1850 onward, the ideas of group and invariant, at last formulated precisely, gradually took center stage, and it became apparent that the theorems of classical geometry are nothing other than the expression of identical relations among invariants or covariants of the similarity group (*), just as those of projective geometry express the identities (or "syzygies") among covariants of the projective group. This is the thesis masterfully expounded by F. Klein in the celebrated "Erlangen program" (XXV b), where he advocates abandoning the sterile controversies between the "synthetic" and "analytic" tendencies; if, he says, the accusation leveled against the latter of granting a privileged role to an arbitrary system of axes "was only too often justified with regard to the defective manner in which the coordinate method was formerly used, it collapses when it comes to a rational application of this method.\ldots{} The domain of spatial intuition is not forbidden to the analytic method.\ldots{}", and he emphasizes that “one must not underestimate the advantage that a well-suited formalism brings to subsequent research, in that it anticipates thought, so to speak” ((XXV b), pp. 488–490).

This leads to a rational and "structural" classification of the theorems of "geometry" according to the group from which they derive: the linear group for projective geometry, the orthogonal group for metric questions, and the symplectic group for the geometry of the "linear complex." But under this pitiless clarity, classical geometry—with the exceptions of algebraic geometry and differential geometry (**), now constituted as autonomous sciences—suddenly withers and loses all its brilliance. Already the generalization of methods based on the use of transformations had rendered the formation of new theorems somewhat mechanical: "Today," said Chasles in 1837 in his \emph{Aperçu historique}, "anyone may come forward, take any known truth, and subject it to the various general principles of transformation; he will extract from it other truths, different or more general; and these will be susceptible to similar operations; so that one may multiply, almost to infinity, the number of new truths deduced from the first.\ldots{} Anyone who wishes may therefore, in the current state of science, generalize and create in Geometry; genius is no longer indispensable for adding a stone to the edifice" ((XV b), p.~268-269). But the situation becomes much clearer with the progress of invariant theory, which finally succeeds (at least for the “classical” groups) in formulating general methods that make it possible in principle to write out all algebraic covariants and all their “syzygies” in a purely automatic manner; a victory which, at the same time, marks the death, as a field of research, of classical invariant theory itself, and of “elementary” geometry, which has practically become a mere dictionary for it. To be sure, nothing allows one to foresee a priori, among the infinity of “theorems” that can thus be unrolled at will, which ones will have, in an appropriate geometric language, a simplicity and elegance comparable to the classical results, and there remains a restricted domain where numerous amateurs continue to exercise themselves happily (geometry of the triangle, of the tetrahedron, of algebraic curves and surfaces of low degree, etc.). But for the professional mathematician, the mine is exhausted, since there are no longer any problems of structure there, capable of having repercussions on other parts of mathematics; and this chapter of group theory and invariant theory may be considered closed until further notice. (*)

\[
\ast \ast
\]

Thus, after the Erlangen program, Euclidean and non-Euclidean geometries, from a purely algebraic point of view, became simple languages, more or less convenient, for expressing the results of the theory of bilinear forms, whose progress goes hand in hand with that of the theory of invariants (*). Everything concerning the notion of the rank of a bilinear form and the relations between these forms and linear transformations is definitively clarified by the work of Frobenius (XXVII a). It is also to Frobenius that we owe the canonical expression of an alternating form on a free Z-module (§ 5, no. 1, th. 1) (XXVII b); however, skew-symmetric determinants had already appeared in Pfaff, at the beginning of the century, in connection with the reduction of differential forms to a normal form; Jacobi, who in 1827 takes up this problem again (XIX a), knows that a skew-symmetric determinant of odd order is zero, and it is he who forms the expression of the Pfaffian and shows that it is a factor of the skew-symmetric determinant of even order; but he had not perceived that the latter is the square of the Pfaffian, and this point was established only by Cayley in 1849 (XVIII b). The notion of the symmetric bilinear form associated with a quadratic form is the most elementary case of the process of "polarization," one of the fundamental tools of the theory of invariants. Under the name of "scalar product," this notion would enjoy immense fortune, first with the popularizers of "vector calculus," then, from the twentieth century, thanks to the unforeseen generalization brought to it by the theory of Hilbert space (see the historical note in Book V). It is also this latter theory that will bring to light the notion of the adjoint of an operator (which previously had scarcely manifested itself except in the theory of linear differential equations, and, in tensor calculus, through the dance of co- and contravariant indices under the baton of the metric tensor); it is this theory, finally, that will give full prominence to the notion of Hermitian form, introduced first by Hermite in 1853 in connection with arithmetic investigations ((XXII), p.~237), but which remained somewhat on the margins of the main mathematical currents until around 1925 and the applications of complex Hilbert spaces to quantum theories.

The study of the orthogonal group and the group of similarities --- clearly conceived and treated as such since the middle of the nineteenth century, and having become the heart of the theory of quadratic forms — as well as of the other "classical" groups (the linear group, the symplectic group, and the unitary group) — takes on, on the other hand, an increasingly important role. We can only mention here the essential part played by these groups in the theory of Lie groups and differential geometry on the one hand, and in the arithmetic theory of quadratic forms (see, for example, (XXXIII) and (XXXV)) on the other (**); to this circumstance, as well as to the extension of the concept of duality to the most diverse questions, is due the fact that there is scarcely any modern mathematical theory in which bilinear forms do not intervene in one way or another. We must in any case note that it was the study of the rotation group (in three dimensions) that led Hamilton to the discovery of quaternions (XVII); this discovery was generalized by W. Clifford, who in 1876 introduced the algebras bearing his name and proved that they are tensor products of quaternion algebras, or of quaternion algebras and a quadratic extension (XXVIII). Rediscovered four years later by Lipschitz (XXIX), who used them to give a parametric representation of orthogonal transformations in n variables (generalizing those that Cayley had obtained for n = 3 (XVIII a) and n = 4 (XVIII c) via the theory of quaternions (cf. § 9, exercises 15 and 16)), these algebras, and the notion of "spinor" derived from them (see (XXXII) and (XXXIV)), were also to enjoy great popularity in the modern era by virtue of their use in quantum theories.

Finally, we have only a word left to say about the evolution of ideas that led to the almost total abandonment of any restriction on the ring of scalars in the theory of sesquilinear forms — a tendency common to all modern algebra, but one that perhaps manifested itself here earlier than elsewhere. We have already noted the fruitful introduction of geometry over the field of complex numbers (which, moreover, throughout the nineteenth century, was not without a perpetual and sometimes perilous confusion between this geometry and real geometry); here the clarity comes above all from the axiomatic studies of the end of the nineteenth century on the foundations of geometry (XXX). In the course of these investigations, Hilbert and his followers, in particular, while examining the relations among the various axioms, were led to construct suitable counterexamples, where the “base field” (commutative or not) possessed more or less pathological properties, and they thus accustomed mathematicians to “geometries” of an entirely new type. From the analytic point of view, Galois had already considered linear transformations where coefficients and variables took their values in a finite prime field ((XVI), p.~27); in developing these ideas, Jordan (XXVI) is naturally led to envisage the classical groups over these fields, groups whose intervention manifests itself in varied domains of mathematics. Dickson, around 1900, extended Jordan's investigations to all finite fields, and more recently it has been realized that a large part of the Jordan–Dickson theory extends to the case of an absolutely arbitrary “base field”; this is due essentially to the general properties of isotropic vectors and to Witt's theorem, which, trivial in the classical cases, were established for an arbitrary base field only in 1936 (XXXI) (*).

But in thus pushing the study of sesquilinear forms toward an ever greater “abstraction,” it proved extremely suggestive to retain as is the terminology which, in the case of 2- and 3-dimensional spaces, came from classical geometry, and to extend it to the n-dimensional case and even to infinite-dimensional spaces. Outmoded as an autonomous and living science, classical geometry was thereby transfigured into a universal language of contemporary mathematics, of incomparable flexibility and convenience.

\section*{Bibliography}
\phantomsection
\addcontentsline{toc}{section}{Bibliography}

\begin{enumerate}
\def\labelenumi{(\Roman{enumi})}
\item
  O. NEUGEBAUER, Lectures on the History of Ancient Mathematical Sciences, Vol. I: Pre-Greek Mathematics, Berlin (Springer), 1934.
\item
  J. TROPFKE, History of Elementary Mathematics, vols. IV–VI, Berlin–Leipzig (de Gruyter), 1923–24.
\item
  Euclid's Elements, 5 vols., ed. J. L. Heiberg, Leipzig (Teubner), 1883–88.
\end{enumerate}

(III bis) T. L. HEATH, The thirteen books of Euclid's Elements\ldots, 3 vols., Cambridge, 1908.

\begin{enumerate}
\def\labelenumi{(\Roman{enumi})}
\setcounter{enumi}{3}
\item
  T. L. HEATH, Apollonius of Perga, Treatise on conic sections, Cambridge (Univ. Press), 1896.
\item
  Ptolemaei Cl. Opera, ed. J. L. Heiberg, 2 vols., Leipzig (Teubner), 1898–1903.
\item
  G. DESARGUES, WORKS\ldots vol. I, Paris (Leiber), 1864: Brouillon project d'une atteinte aux événements des rencontres d'un cône avec un plan, pp. 103-230.
\item
  P. FERMAT, Œuvres, vol. I, Paris (Gauthier-Villars), 1891: a) Ad locos planos et solidos Isagoge, pp. 91-110 (French translation, ibid., vol. III, pp. 84-101); b) Isagoge ad locos ad superficiem, pp. 111-117 (French translation, ibid., vol. III, pp. 102-108).
\item
  L. EULER: a) Introductio in Analysin Infinitorum (Opera Omnia (1), vol. IX, Zürich-Leipzig-Berlin (O. Füssli and B. G. Teubner), 1945); b) Theoria motus corporum solidorum seu rigidorum (Opera Omnia (2), vol. III, Zürich-Leipzig-Berlin (O. Füssli and B. G. Teubner), 1948); c) Problema algebraicum ob affections prorsus singulares memorabile (Opera Omnia, (1), vol. VI, Leipzig-Berlin (Teubner), 1921, pp. 287-315); d) Formulae generales pro translatione quacunque corporum rigidorum, Novi Comm. Acad. Sc. imp. Petrop., vol. XX (1776), pp. 189-207; e) De centro similitudinis. (Opera Omnia (1), vol. XXVI, Zürich (O. Füssli), 1956, pp. 276-285).
\item
  J. L. LAGRANGE, Œuvres, Paris (Gauthier-Villars), 1867-1892: a) Researches on the method of maxima and minima, vol. I, pp. 3–20; b) Researches in arithmetic, vol. III, pp. 695–795.
\item
  G. MONGE, Descriptive Geometry, Paris, 1798.
\item
  C. F. GAUSS, Works: a) Disquisitiones arithmeticae, vol. I, Göttingen, 1870; b) Mutations of Space, vol. VIII, Göttingen, 1900, pp. 357–362.
\item
  J.-V. PONCELET, Treatise on the Projective Properties of Figures, vol. I, 2nd ed., Paris (Gauthier-Villars), 1865.
\item
  A. F. Möbius, Collected Works, Leipzig (Hirzel), 1885–87: a) The Barycentric Calculus, vol. I, pp. 1–388; b) On a Special Kind of Dual Relations between Figures in Space, vol. I, pp. 489–515 (= Crelle's Journal, vol. X, 1833); c) On a New Treatment of Analytic Spherics, vol. II, pp. 1–54.
\item
  A.-L. CAUCHY: a) Lessons on the applications of infinitesimal calculus to geometry (Complete Works, (2), vol. V, Paris (Gauthier-Villars), 1903); b) On the equation by means of which one determines the secular inequalities of the planets (Complete Works, (2), vol. IX, Paris (Gauthier-Villars), 1891, pp. 174–195).
\item
  M. CHASLES: a) Note on the general properties of the system of two bodies, Bull. de Férussac, vol. XIV (1830), pp. 321–326; b) Historical survey of the origin and development of methods in geometry, Brussels, 1837.
\item
  E. GALOIS, Mathematical Works, Paris (Gauthier-Villars), 1897.
\item
  W. R. HAMILTON, Lectures on Quaternions, Dublin, 1853.
\item
  A. CAYLEY, Collected Mathematical Papers, Cambridge, 1889-1898: a) On certain results relating to quaternions, vol. I, p.~123-126 (= Phil. Mag., 1845); b) Sur les déterminants gauches, vol. I, p.~410-413 (= J. de Crelle, vol. XXXVIII (1848)); c) Recherches ultérieures sur les déterminants gauches, vol. II, p.~202-215 (= J. de Crelle, vol. L (1855)); d) A sixth memoir on quantics, vol. II, p.~561-592 (= Phil. Trans., 1859).
\item
  C. G. J. JACOBI, Gesammelte Werke, Berlin (G. Reimer), 1881-1891: a) On the Pfaffian method\ldots vol. IV, pp. 17–29; b) On a fundamental algebraic theorem and its applications, vol. III, pp. 593–598.
\item
  J. J. SYLVESTER, Collected Mathematical Papers, vol. I, Cambridge, 1904: A demonstration of the theorem that every homogeneous quadratic polynomial is reducible by real orthogonal substitution to the form of a sum of positive and negative squares, pp. 378–381 (= Phil. Mag., 1852).
\item
  E. LAGUERRE, Works, vol. II, Paris (Gauthier-Villars), 1905.
\item
  C. HERMITE, Works, vol. I, Paris (Gauthier-Villars), 1905: On the theory of quadratic forms, pp. 200–263 (= J. de Crelle, vol. XLVII (1854)).
\item
  K. G. V. von STÄUDT, Contributions to the Geometry of Position, Nuremberg, 1856.
\item
  E. BELTRAMI: a) Essay on the interpretation of non-Euclidean geometry, Giorn. di Mat., vol. VI (1868), pp. 284–312; b) Fundamental theory of spaces of constant curvature, Ann. di Mat. (2), vol. II (1868–69), pp. 232–255.
\item
  F. KLEIN, Collected Mathematical Works, vol. I, Berlin (Springer), 1921: a) On the so-called non-Euclidean geometry, p.~254-305 (= Math. Ann., vol. IV (1871)); b) Comparative considerations on recent geometrical research, p.~460-497 (= Math. Ann., vol. XLIII (1893)).
\item
  C. JORDAN, Treatise on substitutions and algebraic equations, Paris (Gauthier-Villars), 1870.
\item
  G. FROBENIUS; a) On linear substitutions and bilinear forms, J. de Crelle, vol. LXXXIV (1878), p.~1-63; b) Theory of linear forms with integer coefficients, J. de Crelle, vol. LXXXVI (1879), p.~146-208.
\item
  W. K. CLIFFORD, Mathematical Papers, London (Macmillan), 1882: a) On the classification of geometric algebras, p.~397-401; b) Applications of Grassmann's extensive algebras, p.~266-276 (= Amer. Journ. of Math., vol. I (1878)).
\item
  R. LIPSCHITZ, Investigations on the sums of squares, Bonn, 1886.
\item
  D. HILBERT, Foundations of Geometry, Leipzig (Teubner), 1899.
\item
  E. WITT, Theory of quadratic forms in arbitrary fields, J. de Crelle, vol. CLXXVI (1937), p.~31-44.
\item
  E. CARTAN, Lectures on the theory of spinors, Actual. Sci. et Industr., nos. 643 and 701, Paris (Hermann), 1938.
\item
  C. L. SIEGEL, Symplectic Geometry, Amer. Journ. of Math., vol. LXV (1943), p.~1-86.
\item
  C. CHEVALLEY, The algebraic theory of spinors, New York (Columbia Univ. Press), 1954.
\item
  M. EICHLER, Quadratic forms and orthogonal groups, Berlin-Göttingen-Heidelberg (Springer), 1952.
\end{enumerate}

\chapter*{Index of Notations}
\phantomsection
\addcontentsline{toc}{chapter}{Index of Notations}

\(b^{J}, b^{J'} : 1, 2.\)

The reference numbers indicate successively the section and the number (or, exceptionally, the exercise).

\(F^{J}\) (an antiautomorphism of the ring B of scalars of the right module F): 1, 2.

\(N^{0}, M^{0}\) (N, M submodules): 1, 3.

\(u^{*}\) (u homomorphism): 1, 8.

\(M(x), \mathbf{x} (x \text{ element of a free module}) : 1, 10.\)

\(M(u)\) (u homomorphism from a free module to a free module): 1, 10.

\(^{t}M, \dot{M}^{J} (M \text { matrix } ): 1, 10.\)

\(\mathrm{D}_{\Phi}(x_{1},\ldots ,x_{n}),\mathrm{D}_{\Phi}(\mathrm{S}):2.\)

\(\overline{\alpha}\) ( \(\alpha\) scalar): 3.

\(\operatorname{Pf}(R):5,2.\)

\(\mathbf{Sp}(\Phi), \mathbf{Sp}(2m, \mathrm{A}), \mathbf{Sp}_{2m}(\mathrm{A}) : 5, 3.\)

\(\mathbf{U}(\Phi), \mathbf{S}\mathbf{U}(\Phi) (\Phi \text{ Hermitian form}) : 6, 2.\)

O(Q), SO(Q) (Q quadratic form): 6, 2.

\(\mathbf{U}(n,\mathrm{A}),\mathbf{S}\mathbf{U}(n,\mathrm{A}),\mathbf{O}(n,\mathrm{A}),\mathbf{SO}(n,\mathrm{A}):6,2.\)

\(Q^{\top}Q^{\prime}, Q \sim Q^{\prime}(Q, Q^{\prime}\text{ quadratic forms}): 8, 1.\)

θ(Q) (Q quadratic form): 8, 1.

\(\mathbf{T} + \mathbf{T}', a. \mathbf{T} (\mathbf{T}, \mathbf{T}'\) types of quadratic forms): 8, 2.

\(\delta (\mathbb{Q})\) (Q quadratic form): 8, 2.

\(\mathbf{Q} \otimes \mathbf{Q}'\) (Q, \(\mathbf{Q}'\) quadratic forms): 8, 3.

\(\mathrm{TT}^{\prime}\) (T, \(\mathrm{T}^{\prime}\) types of quadratic forms): 8, 3.

\(C(Q), I(Q), \rho_{Q}, \rho, C^{+}(Q), C^{-}(Q), C^{+}, C^{-} (Q \text{ quadratic form}): 9, 1.\)

\(\alpha, \beta, \mathrm{G}(f): 9, 1.\)

\(e_x, i_f, i_x^{\mathrm{F}}, \lambda_{\mathrm{F}}: 9, 2.\)

\(\lambda_{\mathbf{F}}:9,3.\)

A(Φ) (Φ symmetric bilinear form on a vector space of dimension 2): 10, 1.

\(\overline{u} (u\) direct similarity) \(i,d:10,1\) .

\((\widehat{\mathbf{D}_1,\mathbf{D}_2})\) \((\mathbf{D}_1,\mathbf{D}_2\) lines or half-lines): 10, 3.

\[
\mathfrak {A}, \mathfrak {A} _ {0}, h, h ^ {\prime}: 1 0, 3.
\]

\(|x|, \widehat{(x,y)}(x,y\text{ vectors}):10,3.\)

cos θ, sin θ, tan θ, cot θ: 10, 3.

\(\left.\right)D_{1}, D_{2}\left\{\left(D_{1}, D_{2}\right)\left(D_{1}, D_{2} \text{ half-lines in an oriented plane}\right): 10, 4.\right.\)

\chapter*{Terminological Index}
\phantomsection
\addcontentsline{toc}{chapter}{Terminological Index}

The reference numbers indicate successively the paragraph and the number (or, exceptionally, the exercise).

Adjoint (right, left) of a homomorphism: 1, 8.

Adjoint of a semilinear map: 7, 3.

Clifford algebra of a quadratic form: 9, 1.

Alternating (matrix): 3, 1.

Angle of two half-lines: 10, 3, and exercise 6.

Angle between two rays in an affine plane: 10, 3.

Angle between two lines: 10, 3 and exercise 2.

Angle of two lines in an affine plane: 10, 3.

Angle between two directed lines: 10, exercise 2.

Angle between two vectors in a plane: 10, 3.

Angle between two vectors in a space of dimension \(\geqslant 2:10,3\) .

Right angle: 10, 3 and exercise 2.

Angle of a rotation: 10, 3.

Angle of a direct similarity: 10, exercise 6.

Straight angle: 10, 3 and exercise 6.

Ring of types of quadratic forms: 8, 3.

Witt ring: 8, 3.

Antiautomorphism of a ring: 1, 2.

Principal antiautomorphism of a Clifford algebra: 9, 1.

Antihermitian (form): 3, 1.

Antihermitian (matrix): 3, 1.

Antisymmetric (matrix): 3, 1.

Bilinear map: 1, 1.

Degenerate bilinear map (on the right, on the left): 1, 1.

Non-degenerate bilinear map: 1, 1.

Nondegenerate bilinear map associated with a bilinear map: 1, 3.

Bilinear map obtained by extension of scalars from a bilinear map: 1, 4.

Canonical map of a module into its Clifford algebra: 9, 1.

Canonical mapping from the set of Hermitian endomorphisms to the set of Hermitian forms: 7, 3.

Canonical mapping of the set \(\mathfrak{A}\) on \(\mathrm{S}^{+}/\mathrm{H}:10,3\) .

Canonical mapping of the set \(\mathfrak{A}\) on \(O^{+}:10,3\) .

Right (left) linear map associated with a bilinear map: 1, 1.

Right (left) semilinear map associated with a sesquilinear form: 1, 6.

Right (left) sesquilinear map for an antiautomorphism: 1, 2. Degenerate (right, left) sesquilinear map: 1, 2. Nondegenerate sesquilinear map: 1, 2. Nondegenerate sesquilinear map associated with a sesquilinear map: 1, 3. Sesquilinear map obtained by extension of scalars from a sesquilinear map: 1, 4. Associated (linear map) with a bilinear map: 1, 1. Associated (semilinear map) with a sesquilinear form: 1, 6. Associated (bilinear form) with a quadratic form: 3, 4. Orthogonal automorphism: 6, 2. Orthogonal automorphism in n variables: 6, 2. Principal automorphism of a Clifford algebra: 9, 1. Symplectic automorphism: 5, 3. Symplectic automorphism in 2m variables: 5, 3. Unitary automorphism: 6, 2. Unitary automorphism in n variables: 6, 2. Axis of an affine linear complex: 10, exerc. 16. Orthogonal basis: 6, 1. Orthonormal basis: 6, 1. Symplectic basis: 5, 1. Bilinear (map): 1, 1. Bilinear (form): 1, 6. Bimodule: 1, 1. Canonical: see Canonical map and Canonical homomorphism. Center of a quadric: 6, exerc. 25 and 10, exerc. 12. Unit circle: 10, exerc. 2. Chasles (relation of): 10, 3. Clifford (algebra of): 9, 1. Clifford (group of): 9, 5. Linear complex: 10, exerc. 16. Condition (C): 6, 1. Condition (C'): 6, 1. Condition (T): 4, 2. Isotropic cone: 6, exerc. 23. Conformal (group): 10, exerc. 14. Affine conic: 6, exerc. 25. Projective conic: 6, exerc. 23. Conjugate (affine linear varieties): 6, exerc. 25. Conjugate (projective linear varieties): 6, exerc. 23. Cosine of an angle: 10, 3. Cotangent of an angle: 10, 3. Witt decomposition: 4, 2. Degenerate (bilinear map) on the right (left): 1, 1. Degenerate (sesquilinear map) on the right (left): 1, 2. Degenerate (affine conic): 6, exerc. 25. Degenerate (projective conic): 6, exerc. 23. Degenerate (affine quadric): 6, exerc. 25. Degenerate (projective quadric): 6, exerc. 23. Half-line: 10, 3. Closed half-line: 10, 3. Isotropic half-line: 10, 3. Open half-line: 10, 3. Displacement: 6, 6. Dimension of a quadratic form: 8, 1. Direct (similarity): 6, 5. Discriminant of a sesquilinear form with respect to a system of elements: 2. Distance: 7, 1. Right (angle): 10, 3. Pointed line: 10, exerc. 2.

Odd element of a Clifford algebra: 9, 1. Isotropic element: 4, 1. Even element of a Clifford algebra: 9, 1. Singular element: 4, 1. Elements orthogonal with respect to a sesquilinear map: 1, 3. Elements orthogonal with respect to a quadratic form: 3, 4. Hermitian endomorphism: 7, 3. Normal endomorphism: 7, 3. ε-Hermitian (form): 3, 1. ε-Hermitian (matrix): 3, 1. Equivalent (quadratic forms): 3, 4. Equivalent (sesquilinear forms): 1, 6. Definition space of a quadratic form: 8, 1. Euclidean space: 6, 6. Hermitian space: 6, 6. Oriented space: 10, 3. Euclidean (space): 6, 6. Extension of a sesquilinear form to a tensor power: 1, 9. Extension of a sesquilinear form to an exterior power: 1, 9. External (direct sum) of quadratic forms: 3, 4. External (direct sum) of quadratic modules: 3, 4. Weakly orthogonal (subspaces): 3, exerc. 11. Closed (angular sector): 10, 4. Closed (half-line): 10, 3. Cosine function: 10, 3. Cotangent function: 10, 3. Sine function: 10, 3. Tangent function: 10, 3. Trigonometric function: 10, 3. Anti-Hermitian form: 3, 1. Bilinear form: 1, 6. Bilinear form associated with a quadratic form: 3, 4. Hermitian bilinear form: 3, 1. Hermitian form: 3, 1. Negative (positive) Hermitian form: 7, 1. Inverse form of a bilinear (sesquilinear) form: 1, 7. Metric form: 6, 6. Neutral form: 4, 2 and 8, 1. Quadratic form: 3, 4. Negative (positive) quadratic form: 7, 1. Non-degenerate quadratic form: 3, 4. Quadratic form obtained by extension of scalars from a quadratic form: 3, 4. Sesquilinear form: 1, 6. Equivalent quadratic forms: 3, 4. Equivalent sesquilinear forms: 1, 6. Gram-Schmidt (orthogonalization process of): 6, 1. Clifford group: 9, 5. Reduced Clifford group: 9, 5. Special Clifford group: 9, 5. Rotation group: 9, 5. Group of types of quadratic forms: 8, 2. Witt group: 8, 2. Orthogonal group associated with Q: 6, 2. Orthogonal group in n variables: 6, 2. Special orthogonal group: 6, 2. Special unitary group: 6, 2. Symplectic group associated with Φ: 5, 3. Symplectic group in 2m variables: 5, 3. Unitary group associated with Φ: 6, 2.

\begin{Shaded}
\begin{Highlighting}[]
\NormalTok{Groupe unitaire à n variables : 6, 2.}
\NormalTok{Hermitien (endomorphisme) : 7, 3.}
\NormalTok{Hermitien (espace) : 6, 6.}
\NormalTok{Hermitienne (forme) : 3, 1.}
\NormalTok{Hermitienne (matrice) : 3, 1.}
\NormalTok{Homomorphisme canonique de l\textquotesingle{}algèbre de Clifford d\textquotesingle{}un sous{-}module de E dans l\textquotesingle{}algèbre de Clifford de E : 9, 1.}
\NormalTok{Homomorphisme métrique : 4, 3.}
\NormalTok{Hyperplan radical de deux sphères : 10, exerc. 12.}
\NormalTok{Image réciproque d\textquotesingle{}une application bilinéaire : 1, 1.}
\NormalTok{Image réciproque d\textquotesingle{}une application sesquilinéaire : 1, 2.}
\NormalTok{Impair (élément) dans une algèbre de Clifford : 9, 1.}
\NormalTok{Indice d\textquotesingle{}une forme quadratique (sesquilinéaire) : 4, 2.}
\NormalTok{Invariant de Dickson : 9, exerc. 9.}
\NormalTok{Inverse (forme) : 1, 7.}
\NormalTok{Inverse (similitude) : 6, 5.}
\NormalTok{Inversion : 10, exerc. 13.}
\NormalTok{Inversion de sphère C : 10, exerc. 13.}
\NormalTok{Involution dans un groupe linéaire : 6, 3.}
\NormalTok{Isotrope (demi{-}droite) : 10, 3.}
\NormalTok{Isotrope (élément) : 4, 1.}
\NormalTok{Isotrope (sous{-}module) : 4, 1.}
\NormalTok{Isotrope (variété linéaire) : 6, 6.}
\NormalTok{Laguerre (formule de) : 10, exerc. 5.}
\NormalTok{Loi d\textquotesingle{}inertie : 7, 2.}
\NormalTok{Longueur d\textquotesingle{}un vecteur : 10, 3.}
\NormalTok{Matrice alternée : 3, 1.}
\NormalTok{Matrice antihermitienne : 3, 1.}
\NormalTok{Matrice antisymétrique : 3, 1.}
\NormalTok{Matrice d\textquotesingle{}une application satisfaisant à (G), (D), (GD) ou (DG) : 1, 10.}
\NormalTok{Matrice d\textquotesingle{}une application bilinéaire : 1, 1.}
\NormalTok{Matrice d\textquotesingle{}une application sesquilinéaire : 1, 2.}
\NormalTok{Matrice ε{-}hermitienne : 3, 1.}
\NormalTok{Matrice hermitienne : 3, 1.}
\NormalTok{Matrice normale : 7, exerc. 17.}
\NormalTok{Matrice orthogonale : 6, 2.}
\NormalTok{Matrice symétrique : 3, 1.}
\NormalTok{Matrice symplectique : 5, 3.}
\NormalTok{Matrice unitaire : 6, 2.}
\NormalTok{Métrique (forme) : 6, 6.}
\NormalTok{Métrique (homomorphisme) : 4, 3.}
\NormalTok{Module quadratique : 3, 4.}
\NormalTok{Multiplicateur d\textquotesingle{}une similitude : 6, 5 et 6.}
\NormalTok{Négatif (n{-}vecteur) : 10, 3.}
\NormalTok{Négative (forme hermitienne) : 7, 1.}
\NormalTok{Négative (forme quadratique) : 7, 1.}
\NormalTok{Neutre (forme quadratique) : 8, 1.}
\NormalTok{Neutre (forme sesquilinéaire) : 4, 2.}
\NormalTok{Non dégénérée (application bilinéaire) : 1, 1.}
\NormalTok{Non dégénérée (application sesquilinéaire) : 1, 2.}
\NormalTok{Non dégénérée (forme quadratique) : 3, 4.}
\NormalTok{Normal (endomorphisme) : 7, 3.}
\NormalTok{Norme spinorielle : 9, 5.}
\NormalTok{Noyau d\textquotesingle{}un module quadratique : 3, 4.}
\NormalTok{n{-}vecteur négatif (positif) : 10, 3.}
\NormalTok{Orientation d\textquotesingle{}un espace : 10, 3.}
\NormalTok{Orienté (espace) : 10, 3.}
\NormalTok{Orthogonal (automorphisme) : 6, 2.}
\NormalTok{Orthogonal (groupe) : 6, 2.}
\NormalTok{Orthogonal (groupe spécial) : 6, 2.}
\end{Highlighting}
\end{Shaded}

Orthogonal (projector): 6, 3. Orthogonal (submodule) to a submodule: 1, 3. Orthogonal (vector) to a linear manifold: 6, 6. Orthogonal (basis): 6, 1. Orthogonal (matrix): 6, 2. Orthogonal (projection) onto a linear manifold: 6, 6. Orthogonal (transformation): 6, 2. Orthogonal (manifold) to a linear manifold: 6, 6. Orthogonal (parts): 1, 3 and 3, 4. Orthogonal (linear manifolds): 6, 6. Orthogonalization (process of) of Gram-Schmidt: 6, 1. Orthogonal (elements): 1, 3 and 3, 4. Orthonormal (basis): 6, 1. Open (angular sector): 10, 4. Open (half-line): 10, 3. Even (element) in a Clifford algebra: 9, 1. Orthogonal parts: 1, 3 and 3, 4. Perpendicular (linear manifolds): 6, exerc. 22. Pfaffian of an alternating matrix: 5, 2. Flat (angle): 10, 3. Point of view of a stereographic projection: 10, exerc. 14. Polar of a linear manifold with respect to a quadric: 6, exerc. 23 and 25. Pole of a hyperplane with respect to a quadric: 6, exerc. 23 and 25. Pole of an inversion: 10, exerc. 13. Positive (n-vector): 10, 3. Positive (Hermitian form): 7, 1. Positive (quadratic form): 7, 1. Principal (antiautomorphism): 9, 1. Principal (automorphism): 9, 1. Tensor product of quadratic forms: 8, 3. Tensor product of sesquilinear forms: 1, 9. Orthogonal projector: 6, 3. Orthogonal projection onto a linear manifold: 6, 6. Stereographic projection: 10, exerc. 14. Pseudo-discriminant: 9, exerc. 9. Power of a point with respect to a sphere: 10, exerc. 12. Power of an inversion: 10, exerc. 13. Pythagoras (theorem of): 6, 6. Quadratic (form): 3, 4. Quadratic (module): 3, 4. Affine quadric: 6, exerc. 25. Projective quadric: 6, exerc. 23. Rank of a bilinear (sesquilinear) form: 1, 6. Radius of a sphere: 10, exerc. 12. Reduced (Clifford group): 9, 5. Chasles relation: 10, 3. Spinor representations: 9, 4. Rotation: 9, 5. Rotation of angle φ: 10, 3 and exerc. 2. Rotations (group of): 9, 5. Closed angular sector: 10, 4. Open angular sector: 10, 4. Flat angular sector: 10, exerc. 8. Reentrant angular sector: 10, exerc. 8. Salient angular sector: 10, exerc. 8. Semi-spinor: 9, 4. Sesquilinear (map): 1, 2. Sesquilinear (form): 1, 6. Signature of a Hermitian form: 7, 2. Similarity: 6, 5 and 6.

Direct similarity: 6, 5 and 9, exercise 9.

Inverse similarity: 6, 5.

Singular (element): 4, 1.

Singular (submodule): 4, 1.

Sine of an angle: 10, 3.

Direct sum of bilinear (sesquilinear) maps: 1, 3.

External direct sum of quadratic forms (quadratic modules): 3, 4.

Isotropic submodule: 4, 1.

Orthogonal submodule to a submodule: 1, 3.

Singular submodule: 4, 1.

Totally isotropic submodule: 4, 1.

Totally singular submodule: 4, 1.

Sphere: 10, exercise 12.

Orthogonal spheres: 10, exercise 12.

Spinor: 9, 4.

Direct sequence of half-lines: 10, 4.

Symmetry with respect to a hyperplane: 6, 4.

Symmetry with respect to a vector subspace: 6, 3.

Symmetry with respect to an affine linear manifold: 6, 6.

Symmetric (matrix): 3, 1.

Symplectic (automorphism): 5, 3.

Symplectic (base): 5, 1.

Symplectic (group): 5, 3.

Symplectic (matrix): 5, 3.

Symplectic (transformation): 5, 3.

Tangent of an angle: 10, 3.

Tangent (linear variety) to a quadric: 6, exercises 23 and 25.

Pythagorean theorem: 6, 6.

Witt's theorem: 4, 3.

Totally isotropic (submodule): 4, 1.

Totally orthogonal (submodule) to a submodule: 1, 3.

Totally orthogonal (variety) to a linear variety: 6, 6.

Totally singular (submodule): 4, 1.

Orthogonal transformation: 6, 2.

Symplectic transformation: 5, 3.

Unitary transformation: 6, 2.

Trigonometric (function): 10, 3.

Type of a quadratic form: 8, 1.

Unitary (automorphism): 6, 2.

Unitary (group): 6, 2.

Unitary (special group): 6, 2.

Unitary (matrix): 6, 2.

Unitary (transformation): 6, 2.

Isotropic linear manifold: 6, 6.

Linear manifold orthogonal to a linear manifold: 6, 6.

Totally isotropic linear manifold: 6, 6.

Orthogonal linear manifolds: 6, 6.

Vector orthogonal to a linear manifold: 6, 6.

Witt (ring of): 8, 3.

Witt (decomposition of): 4, 2.

Witt (group of): 8, 2.

Witt (theorem of): 4, 3.

\chapter*{\texorpdfstring{Definitions of Chapter \Roman{chapter}}{Definitions of Chapter IX}}
\phantomsection
\addcontentsline{toc}{chapter}{\texorpdfstring{Definitions of Chapter \Roman{chapter}}{Definitions of Chapter IX}}

\section*{Sesquilinear forms:}

Let A be a ring, \(\xi \to \overline{\xi}\) an involutive antiautomorphism of A, that is, a bijection of A onto itself such that \(\overline{(\xi + \eta)} = \overline{\xi} + \overline{\eta}\) , \(\overline{(\xi\eta)} = \overline{\eta} . \overline{\xi}\) , \(\overline{\overline{\xi}} = \xi\) for all \(\xi, \eta\) in A. A sesquilinear form \(\Phi\) on a left A-module E is a map from E × E into A such that

\[
\begin{array}{c c} \Phi (x + x ^ {\prime}, y) = \Phi (x, y) + \Phi (x ^ {\prime}, y), & \Phi (x, y + y ^ {\prime}) = \Phi (x, y) + \Phi (x, y ^ {\prime}) \\ \Phi (\alpha x, y) = \alpha \Phi (x, y), & \Phi (x, \alpha y) = \Phi (x, y) \bar {\alpha} \end{array}
\]

for all \(\alpha\in\mathbf{A},x,x^{\prime},y,y^{\prime}\) in E.

Let \(\varepsilon\) be an element of the center of A. One says that a sesquilinear form \(\Phi\) on E is \(\varepsilon\)-hermitian if \(\Phi(y, x) = \varepsilon \overline{\Phi(x, y)}\) for all \(x \in E\) , \(y \in E\) ; if \(\varepsilon = 1\) (resp. \(\varepsilon = -1\) ) one says that \(\Phi\) is hermitian (resp. anti-hermitian).

When the antiautomorphism \(\xi\to\overline{\xi}\) is the identity (which implies that the ring A is commutative), the corresponding sesquilinear forms are the bilinear forms. One then says "symmetric" instead of "hermitian", and "antisymmetric" instead of "antihermitian" (cf.~Chap.~III). A bilinear form \(\Phi\) such that \(\Phi(x,x)=0\) is said to be alternating; it is then antisymmetric, and the converse is true when A is a field of characteristic \(\neq2\) (cf.~Chap.~III).

One says that a form \(\varepsilon\)-hermitian satisfies condition (T) if for all \(x \in E\) , there exists \(\lambda \in A\) such that \(\Phi(x, x) = \lambda + \varepsilon \overline{\lambda}\) . This condition is always satisfied when \(\varepsilon = 1\) and \(A\) is a field of characteristic \(\neq 2\) , or when \(\Phi\) is alternating.

\section*{Quadratic forms:}

Let A be a commutative ring. A quadratic form Q on an A-module E is a map from E into A such that \(\mathrm{Q}(\alpha x) = \alpha^{2}\mathrm{Q}(x)\) for \(\alpha \in A, x \in E\) , and that the map

\[
(x, y) \rightarrow \Phi (x, y) = Q (x + y) - Q (x) - Q (y)
\]

is a bilinear form (necessarily symmetric), said to be associated with the quadratic form Q. One has \(\Phi(x, x) = 2Q(x)\) ; conversely, for every bilinear form \(\Psi\) on E, \(x \to \Psi(x, x)\) is a quadratic form on E;

when A is a field of characteristic ≠ 2, quadratic forms and symmetric bilinear forms on E therefore correspond bijectively.

\section*{Orthogonal elements:}

Let \(\Phi\) be a form \(\varepsilon\)-hermitian on a left A-module E. One says that two elements \(x, y\) of E are orthogonal (for \(\Phi\) ) if \(\Phi(x, y) = 0\) ; this relation is symmetric in \(x\) and \(y\) . For every submodule M of E, the set of \(x \in E\) which are orthogonal to all elements of M is a submodule denoted \(\mathbf{M}^{\circ}\) and called the orthogonal of the submodule M. One says that \(\Phi\) is non-degenerate if \(\mathbf{E}^{\circ} = \{0\}\) . When E is a finite-dimensional vector space and \(\Phi\) is non-degenerate, one has codim \(\mathbf{M}^{\circ} = \dim \mathbf{M}\) and \(\mathbf{M}^{\circ\circ} = \mathbf{M}\) for every subspace M of E; moreover, for every pair of subspaces M, N of E, one has

\[
(\mathbf {M} + \mathbf {N}) ^ {0} = \mathbf {M} ^ {0} \cap \mathbf {N} ^ {0}, \quad (\mathbf {M} \cap \mathbf {N}) ^ {0} = \mathbf {M} ^ {0} + \mathbf {N} ^ {0}.
\]

Suppose the ring A is commutative, and let Q be a quadratic form on the A-module E, \(\Phi\) the bilinear form associated with Q. One says that two elements of E are orthogonal (for Q) if they are orthogonal for \(\Phi\) ; the orthogonal of a submodule M (for Q) is the orthogonal M \(^{o}\) of M for \(\Phi\) . One says that Q is non-degenerate if \(\Phi\) is non-degenerate.

\section*{Isotropic elements and singular elements:}

Let \(\Phi\) be a form \(\varepsilon\)-hermitian on a left A-module E. An element \(x \in E\) is said to be isotropic if \(\Phi(x, x) = 0\) ; a submodule M of E is said to be isotropic if \(\mathbf{M} \cap \mathbf{M}^{\circ} \neq \{0\}\) (in other words, if the restriction of \(\Phi\) to M is degenerate); M is said to be totally isotropic if \(\mathbf{M} \subset \mathbf{M}^{\circ}\) (in other words, if the restriction of \(\Phi\) to M is zero). For every submodule M of E, \(\mathbf{M} \cap \mathbf{M}^{\circ}\) is totally isotropic. If E is a finite-dimensional vector space and if \(\Phi\) is nondegenerate, the following three conditions are equivalent for a subspace M of E : \(1^{\circ}\) M is nonisotropic ; \(2^{\circ}\) M \(^{\circ}\) is nonisotropic ; \(3^{\circ}\) E is the direct sum of M and M \(^{\circ}\) . Suppose the ring A is commutative, and let Q be a quadratic form on an A-module E, \(\Phi\) the bilinear form associated with Q. An element \(x \in E\) is said to be singular if \(\mathrm{Q}(x) = 0\) ; a submodule M of E is said to be singular (resp. totally singular) if \(\mathbf{M} \cap \mathbf{M}^{\circ}\) contains a singular element \(\neq 0\) (resp. if the restriction of Q to M is zero). Every singular element (resp. singular submodule, totally singular submodule) is isotropic (resp. isotropic, totally isotropic); the converse is true when A is a field of characteristic \(\neq 2\) .

\end{document}
